\documentclass[UTF8, 12pt, fontset=fandol, oneside]{ctexart}
\usepackage{researchnotes}

\graphicspath{{math/数分习题/}}

\title{第四章\quad 一元函数积分学}
\author{}
\date{}

\begin{document}

\maketitle

\noindent\textbf{导读}\quad §4.2 针对数学院系学生，其余四节适合本书各类读者.

本章讨论如下几方面内容：积分与极限，可积性，积分值的估计，积分不等式与定积分的几类典型问题，几个著名的不等式，反常积分.

\section*{§4.1\quad 积分与极限}

\subsection*{一、利用积分求极限}

\noindent\textbf{要点}\quad 定积分是积分和的极限，因此求某个表达式的极限，若能将表达式写成某可积函数的积分和（或 Darboux 和），那么极限就等于此函数的积分.

\noindent\textbf{例 4.1.1}\quad 求
\[
\lim_{n\to\infty}
\frac{\left[1^\alpha+3^\alpha+\cdots+(2n+1)^\alpha\right]^{\beta+1}}
{\left[2^\beta+4^\beta+\cdots+(2n)^\beta\right]^{\alpha+1}}
\quad(\alpha,\beta\ne-1).
\]

\noindent\textbf{解}\quad
\[
\begin{aligned}
&\frac{\left[1^\alpha+3^\alpha+\cdots+(2n+1)^\alpha\right]^{\beta+1}}
{\left[2^\beta+4^\beta+\cdots+(2n)^\beta\right]^{\alpha+1}}\\
={}&2^{\alpha-\beta}
\frac{\left\{\dfrac2n\left[\left(\dfrac1n\right)^\alpha+
\left(\dfrac3n\right)^\alpha+\cdots+
\left(\dfrac{2n+1}{n}\right)^\alpha\right]\right\}^{\beta+1}}
{\left\{\dfrac2n\left[\left(\dfrac2n\right)^\beta+
\left(\dfrac4n\right)^\beta+\cdots+
\left(\dfrac{2n}{n}\right)^\beta\right]\right\}^{\alpha+1}}\\
={}&2^{\alpha-\beta}
\frac{\left[\displaystyle\sum_{i=1}^n
\left(\frac{2i-1}{n}\right)^\alpha\frac2n+
\left(\frac{2n+1}{n}\right)^\alpha\frac2n\right]^{\beta+1}}
{\left[\displaystyle\sum_{i=1}^n
\left(\frac{2i}{n}\right)^\beta\frac2n\right]^{\alpha+1}}\\
\longrightarrow{}&2^{\alpha-\beta}
\frac{\left(\displaystyle\int_0^2 t^\alpha\,dt\right)^{\beta+1}}
{\left(\displaystyle\int_0^2 t^\beta\,dt\right)^{\alpha+1}}
=2^{\alpha-\beta}\frac{(\beta+1)^{\alpha+1}}{(\alpha+1)^{\beta+1}}
\quad(n\to\infty).
\end{aligned}
\]
这里把
\[
\sum_{i=1}^n\left(\frac{2i-1}{n}\right)^\alpha\frac2n
\quad\text{与}\quad
\sum_{i=1}^n\left(\frac{2i}{n}\right)^\beta\frac2n
\]
分别看成 $f(t)=t^\alpha$ 与 $g(t)=t^\beta$ 在 $[0,2]$ 上的积分和，其分划是将 $[0,2]$ $n$ 等分，$\xi_i$ 分别取小区间的中点与右端点.

\noindent\textbf{练习}\quad 设 $f(x)$ 在 $[0,1]$ 上连续，求证：
\[
\lim_{n\to\infty}\frac1n\sum_{k=1}^{n-1}(-1)^{k+1}f\left(\frac kn\right)=0.
\tag{1}
\]
（南开大学）

\noindent\textbf{提示}\quad 记
\[
a_n=\frac1n\sum_{k=1}^n(-1)^{k+1}f\left(\frac kn\right).
\tag{2}
\]
式 (2) 比 (1) 多一项 $\dfrac1n f\left(\dfrac nn\right)$，但
\[
\lim_{n\to\infty}\frac1n f\left(\frac nn\right)=0,
\]
故“原问题”等价于证明
\[
\lim_{n\to\infty}a_n=0.
\]

\noindent\textbf{证 I}\quad
\[
a_n=
\begin{cases}
\dfrac1n\displaystyle\sum_{k=1}^n f\left(\dfrac kn\right)
-\dfrac2n\displaystyle\sum_{k=1}^m f\left(\dfrac{2k}{n}\right),
&\text{当 }n=2m\text{ 时},\\[1.2em]
\dfrac1n\displaystyle\sum_{k=1}^n f\left(\dfrac kn\right)
-\dfrac2n\displaystyle\sum_{k=1}^m f\left(\dfrac{2k}{n}\right),
&\text{当 }n=2m+1\text{ 时}.
\end{cases}
\]
（注意意思是
\[
b_1-b_2+b_3-b_4=(b_1+b_2+b_3+b_4)-2(b_2+b_4),
\]
\[
b_1-b_2+b_3-b_4+b_5=(b_1+b_2+b_3+b_4+b_5)-2(b_2+b_4).
\]
）
注意到
\[
\frac1n\sum_{k=1}^n f\left(\frac kn\right)\to\int_0^1f(x)\,dx\stackrel{\text{记}}{=}I
\quad(n\to\infty).
\]
\[
\frac2n\sum_{k=1}^m f\left(\frac{2k}{n}\right)
\]
是：先将 $[0,1]$ 进行 $2m=n$ 等分，然后从左至右，每两个小区间合并成为一个新的小区间，于是新的小区间组成 $[0,1]$ 的一个新分划，新小区间长度为 $\dfrac1m=\dfrac2n$，$\xi_k$ 取（第 $k$ 个新小区间的）右端点 $\dfrac{2k}{n}$，则
\[
\frac2n\sum_{k=1}^m f\left(\frac{2k}{n}\right)
\]
是 $\displaystyle\int_0^1f(x)\,dx$ 的一个积分和，故
\[
\frac2n\sum_{k=1}^m f\left(\frac{2k}{n}\right)\to\int_0^1f(x)\,dx=I
\quad(n\to\infty).
\]
至此有 $\displaystyle\lim_{n\to\infty}a_n=I-I=0$.

\noindent\textbf{证 II}\quad 在 $[0,1]$ 上，$f(x)$ 连续，可积. 根据可积的充要条件：有
\[
\lim_{n\to\infty}\sum_{i=1}^n\omega_i\Delta x_i=0
\]
（其中 $\omega_i$ 是 $f$ 在第 $i$ 个小区间上的振幅）.

在式 (2) 中将 $a_n$ 相邻正、负两项合并，则
\[
\left|\frac1n f\left(\frac kn\right)-\frac1n f\left(\frac{k+1}{n}\right)\right|
\le\omega_k\frac1n,
\]
所以
\[
\lim_{n\to\infty}a_{2m}
\le \lim_{m\to\infty}\sum_{k=1}^m
\left|\frac1n f\left(\frac kn\right)-\frac1n f\left(\frac{k+1}{n}\right)\right|
\le \lim_{m\to\infty}\frac12\sum_{k=1}^m\omega_k\frac1m=0,
\]
\[
\begin{aligned}
\lim_{n\to\infty}a_{2m+1}
&\le \lim_{m\to\infty}\sum_{k=1}^m
\left|\frac1n f\left(\frac kn\right)-\frac1n f\left(\frac{k+1}{n}\right)\right|
\left|+\frac1n f\left(\frac nn\right)\right|\\
&\le \lim_{m\to\infty}\frac12\sum_{k=1}^{m-1}\omega_k\frac1m
+\frac1n|f(1)|=0.
\end{aligned}
\]
于是，$\displaystyle\lim_{n\to\infty}a_n=0$（见例 1.2.13）. 式 (1) 获证.

作为练习，试利用一致连续，通过 $\varepsilon$--$\delta,\varepsilon$--$N$ 方法证明上述结论.

按定积分的定义，不一定要将区间 $n$ 等分，只要最长的小区间长度趋于零即可. 下面看一个非 $n$ 等分的例子.

\noindent\textbf{例 4.1.2}\quad 求极限
\[
\lim_{n\to\infty}\left(b^{1/n}-1\right)
\sum_{i=0}^{n-1} b^{1/n}\sin b^{(2i+1)/(2n)}
\quad(b>1).
\]

\noindent\textbf{解}\quad 原式
\[
=\lim_{n\to\infty}\sum_{i=0}^{n-1}
\left(\sin b^{(2i+1)/(2n)}\right)
\left(b^{(i+1)/n}-b^{i/n}\right).
\]
这里的和式，可看成函数 $\sin x$ 在 $[1,b]$ 上按分划
\[
1=b^{0/n}<b^{1/n}<b^{2/n}<\cdots<b^{n/n}=b
\]
所作的积分和. 其中 $\Delta x_i=b^{(i+1)/n}-b^{i/n}$ 为小区间
$[b^{i/n},b^{(i+1)/n}]$ 的长度，最大区间长度
\[
\lambda:0\le\lambda=\max_{1\le i\le n}\Delta x_i\le b\left(b^{1/n}-1\right)\to0,
\]
$\xi_i=b^{(2i+1)/(2n)}\in[b^{i/n},b^{(i+1)/n}]$ 为小区间两端点的比例中项，因此
\[
\text{原极限}=\int_1^b\sin x\,dx=\cos1-\cos b.
\]

\noindent\textbf{例 4.1.3}\quad 证明：
\[
\lim_{n\to\infty}\prod_{i=0}^{n-1}\left(2+\cos i\frac\pi n\right)^{\pi/n}
=\left(\frac{\sqrt3+2}{2}\right)^\pi.
\]

\noindent\textbf{提示}\quad 先取对数. 积分可用分部积分法及变量替换 $\pi-x=t$ 求解.（“博士数学论坛”（现为“博士数学家园”）对此题有详细讨论.）

\subsection*{二、积分的极限}

\noindent\textbf{要点}\quad 当极限的表达式里含有定积分时，我们把这种极限称为积分的极限. 对这种极限，以前讨论的各种方法原则上都是适用的，所不同的，这里需要充分运用积分的各种特性和运算法则，有时也可转化为某函数的积分和或 Darboux 和的极限，从而转化为新的定积分.

\noindent\textbf{例 4.1.4}\quad 求极限：
\[
1)\ \lim_{n\to\infty}\int_0^{\pi/2}\sin^n x\,dx;\quad
2)\ \lim_{n\to\infty}\int_0^1\frac{x^n}{1+\sqrt{x}}\,dx.
\]

\noindent\textbf{解}\quad 1）$\forall\varepsilon>0$（不妨设 $0<\varepsilon<\dfrac\pi2$），
\[
\begin{aligned}
0\le\int_0^{\pi/2}\sin^n x\,dx
&=\int_0^{\pi/2-\varepsilon/2}\sin^n x\,dx
+\int_{\pi/2-\varepsilon/2}^{\pi/2}\sin^n x\,dx\\
&\le\int_0^{\pi/2-\varepsilon/2}
\sin^n\left(\frac\pi2-\frac\varepsilon2\right)\,dx
+\int_{\pi/2-\varepsilon/2}^{\pi/2}1\,dx\\
&\le\left(\frac\pi2-\frac\varepsilon2\right)
\sin^n\left(\frac\pi2-\frac\varepsilon2\right)+\frac\varepsilon2.
\end{aligned}
\]
因 $0<\sin\left(\dfrac\pi2-\dfrac\varepsilon2\right)<1$，所以
\[
\left(\frac\pi2-\frac\varepsilon2\right)
\sin^n\left(\frac\pi2-\frac\varepsilon2\right)\to0\quad(n\to\infty).
\]
故 $\exists N>0$，当 $n>N$ 时，
\[
\left(\frac\pi2-\frac\varepsilon2\right)
\sin^n\left(\frac\pi2-\frac\varepsilon2\right)<\frac\varepsilon2,
\]
从而上式 $<\dfrac\varepsilon2+\dfrac\varepsilon2=\varepsilon$.

原极限为 $0$.

2）因
\[
0\le\int_0^1\frac{x^n}{1+\sqrt{x}}\,dx
\le\int_0^1x^n\,dx=\frac1{n+1}\to0\quad(n\to\infty).
\]
所以
\[
\lim_{n\to\infty}\int_0^1\frac{x^n}{1+\sqrt{x}}\,dx=0.
\]

\noindent\textbf{注}\quad 该例十分典型，特别第 1）小题，将区间分为两段：一段上函数有界，可将区间长度取得任意小，然后固定分点；另一段区间长度有限，函数一致趋向零，因而两段上的积分都任意小，整个积分趋向零.

\noindent\textbf{练习}\quad 求证：
\[
1)\ \lim_{n\to\infty}\int_a^{\pi/2}(1-\sin x)^n\,dx=0,\quad \text{其中 }a\in(0,1);
\]
\[
2)\ \lim_{n\to\infty}\int_0^{\pi/2}(1-\sin x)^n\,dx=0;\quad(\text{北京航空航天大学})
\]
\[
3)\ \lim_{n\to\infty}\int_0^1 e^{x^n}\,dx=1.\quad(\text{武汉大学})
\]

\noindent\textbf{提示}\quad 3）也可看成求证：
\[
\lim_{n\to\infty}\int_0^1(e^{x^n}-1)\,dx=0.
\]

\noindent\textbf{例 4.1.5}\quad 设 $f(x)$ 在 $[0,1]$ 上连续，试证：
\[
\lim_{h\to0^+}\int_0^1\frac{h}{h^2+x^2}f(x)\,dx
=\frac\pi2 f(0).
\]

\noindent\textbf{证 I}\quad
\[
\int_0^1\frac{h}{h^2+x^2}f(x)\,dx
=\int_0^{h^{1/4}}\frac{hf(x)}{h^2+x^2}\,dx
+\int_{h^{1/4}}^1\frac{hf(x)}{h^2+x^2}\,dx=I_1+I_2,
\]
其中
\[
\begin{aligned}
I_1&=\int_0^{h^{1/4}}\frac{hf(x)}{h^2+x^2}\,dx
=f(\xi)\int_0^{h^{1/4}}\frac{h}{h^2+x^2}\,dx\quad(0\le\xi\le h^{1/4})\\
&=f(\xi)\arctan\frac{x}{h}\bigg|_0^{h^{1/4}}
=f(\xi)\arctan\frac1{h^{3/4}}\to f(0)\frac\pi2
\quad(h\to0^+),
\end{aligned}
\]
\[
\begin{aligned}
|I_2|
&=\left|\int_{h^{1/4}}^1\frac{h}{h^2+x^2}f(x)\,dx\right|
\le M\int_{h^{1/4}}^1\frac{h}{h^2+x^2}\,dx\quad(|f(x)|\le M)\\
&=M\left(\arctan\frac1h-\arctan\frac1{h^{3/4}}\right)\to0
\quad(h\to0^+).
\end{aligned}
\]

\noindent\textbf{证 II}\quad （拟合法）因
\[
\lim_{h\to0^+}\int_0^1\frac{h}{h^2+x^2}\,dx=\frac\pi2,
\]
故极限值可改写为
\[
\frac\pi2f(0)=\lim_{h\to0^+}\int_0^1\frac{h}{h^2+x^2}f(0)\,dx.
\]
问题归结为证明：
\[
\lim_{h\to0^+}\int_0^1
\frac{h}{h^2+x^2}\left[f(x)-f(0)\right]\,dx=0.
\]
但是
\[
\int_0^1\frac{h}{h^2+x^2}[f(x)-f(0)]\,dx
=\left(\int_0^\delta+\int_\delta^1\right)
\frac{h}{h^2+x^2}[f(x)-f(0)]\,dx.
\]
因为 $f(x)$ 在 $x=0$ 处连续，所以 $\forall\varepsilon>0$，当 $\delta>0$ 充分小时，在 $[0,\delta]$ 上，
$|f(x)-f(0)|<\dfrac\varepsilon\pi$，从而
\[
\left|\int_0^\delta\frac{h}{h^2+x^2}[f(x)-f(0)]\,dx\right|
\le\int_0^\delta\frac{h|f(x)-f(0)|}{h^2+x^2}\,dx
\le\frac\varepsilon\pi\int_0^\delta\frac{h}{h^2+x^2}\,dx
=\frac\varepsilon\pi\arctan\frac\delta h\le\frac\varepsilon\pi\cdot\frac\pi2=\frac\varepsilon2.
\]
再将 $\delta$ 固定，这时第二个积分
\[
\left|\int_\delta^1\frac{h}{h^2+x^2}[f(x)-f(0)]\,dx\right|
\le h\int_\delta^1\frac1{x^2}|f(x)-f(0)|\,dx=h\cdot M_0.
\]
故当 $0<h<\dfrac{\varepsilon}{2M_0}$ 时，
\[
\left|\int_0^1\frac{h}{h^2+x^2}[f(x)-f(0)]\,dx\right|
<\frac\varepsilon2+\frac\varepsilon2=\varepsilon.
\]
证毕.

\noindent\textbf{证 III}\quad
\[
\int_0^1\frac{hf(x)}{h^2+x^2}\,dx
=\int_0^1\frac{h(f(x)-f(0))}{h^2+x^2}\,dx
+f(0)\int_0^1\frac{h}{h^2+x^2}\,dx=I_1+I_2.
\]
然后证明 $I_1\to0$，$I_2\to\dfrac\pi2 f(0)$.

\noindent\textbf{练习 1}\quad 设 $f$ 为连续函数，证明：
\[
\lim_{n\to\infty}\frac2\pi\int_0^1\frac{n}{n^2x^2+1}f(x)\,dx=f(0).
\]
（武汉大学，复旦大学）

\noindent\textbf{练习 2}\quad 设 $f(x)$ 在 $[a,b]$ 上可积，在 $x=b$ 处左连续，求证：
\[
\lim_{n\to\infty}\frac{n+1}{(b-a)^{n+1}}\int_a^b(x-a)^n f(x)\,dx=f(b).
\]
（浙江大学）

\noindent\textbf{提示}\quad （拟合法）由
\[
\frac{n+1}{(b-a)^{n+1}}\int_a^b(x-a)^n\,dx=1,
\]
知
\[
f(b)=\frac{n+1}{(b-a)^{n+1}}\int_a^b(x-a)^n f(b)\,dx.
\]
于是
\[
\frac{n+1}{(b-a)^{n+1}}\int_a^b(x-a)^n f(x)\,dx-f(b)
=\frac{n+1}{(b-a)^{n+1}}\int_a^b(x-a)^n[f(x)-f(b)]\,dx.
\]
再提示：因 $f(x)$ 在 $x=b$ 处左连续，故 $\forall\varepsilon>0(\varepsilon<1)$，$\exists\delta>0$，当 $b-\delta<x<b$ 时，有
\[
|f(x)-f(b)|<\frac\varepsilon2.
\tag{1}
\]
因 $f(x)$ 可积，必有界：$\exists M>\dfrac12$ 使得
\[
|f(x)|\le M\quad(x\in[a,b]).
\tag{2}
\]
固定 $\delta(<b-a)$，则
\[
\lim_{n\to\infty}\frac{(b-a-\delta)^{n+1}}{(b-a)^{n+1}}=0.
\]
这说明：存在充分大的 $n$ 使得
\[
0<2M\left(\frac{b-a-\delta}{b-a}\right)^{n+1}\le\frac\varepsilon2<1.
\tag{3}
\]
因而
\[
\begin{aligned}
&\left|\frac{n+1}{(b-a)^{n+1}}\int_a^b(x-a)^n(f(x)-f(b))\,dx\right|\\
\le{}&\frac{n+1}{(b-a)^{n+1}}
\left(\int_a^{b-\delta}+\int_{b-\delta}^b\right)
(x-a)^n|f(x)-f(b)|\,dx\\
\le{}&\frac{2M}{(b-a)^{n+1}}\int_a^{b-\delta}d(x-a)^{n+1}
+\frac\varepsilon2\cdot\frac1{(b-a)^{n+1}}
\int_{b-\delta}^b d(x-a)^{n+1}\\
={}&2M\left(\frac{b-a-\delta}{b-a}\right)^{n+1}
+\frac\varepsilon2\left[1-\left(\frac{b-a-\delta}{b-a}\right)^{n+1}\right]
<\frac\varepsilon2+\frac\varepsilon2=\varepsilon.
\end{aligned}
\]

下两例使用两边夹法则及其推广形式.

\noindent\textbf{例 4.1.6}\quad 设 $f(x)$ 严减，在 $[0,1]$ 上连续，$f(0)=1,f(1)=0$. 试证明：$\forall\delta\in(0,1)$，有
\[
1)\ \lim_{n\to\infty}\frac{\displaystyle\int_\delta^1(f(x))^n\,dx}
{\displaystyle\int_0^\delta(f(x))^n\,dx}=0;
\quad
2)\ \lim_{n\to\infty}\frac{\displaystyle\int_0^\delta(f(x))^{n+1}\,dx}
{\displaystyle\int_0^1(f(x))^n\,dx}=1.
\]

\noindent\textbf{证}\quad 1）（利用两边夹法则.）因 $f(x)$ 严减，
\[
0<f(\delta)<f(\delta/2),\qquad
\left(\frac{f(\delta)}{f(\delta/2)}\right)^n\to0\quad(n\to\infty).
\]
故对任意固定的 $\delta\in(0,1)$，有
\[
\begin{aligned}
0\le\frac{\displaystyle\int_\delta^1(f(x))^n\,dx}
{\displaystyle\int_0^\delta(f(x))^n\,dx}
&\le
\frac{\displaystyle\int_\delta^1(f(x))^n\,dx}
{\displaystyle\int_0^{\delta/2}(f(x))^n\,dx}\\
&\le
\frac{\displaystyle\int_\delta^1(f(\delta))^n\,dx}
{\displaystyle\int_0^{\delta/2}(f(\delta/2))^n\,dx}
\le
\left(\frac{f(\delta)}{f(\delta/2)}\right)^n\frac{1-\delta}{\delta/2}\to0.
\end{aligned}
\]

2）（利用两边夹法则的推广形式（参看例 1.3.11）.）因 $f(x)$ 严减，$f(0)=1,f(1)=0$，知
\[
0<f(x)<1\quad(x\in(0,1)).
\]
根据连续性，$\forall\varepsilon>0$，$\exists\delta_1:0<\delta_1<\delta$，使得
\[
f(x)>1-\varepsilon\quad(\forall x\in[0,\delta_1]).
\]
于是
\[
1=\frac{\displaystyle\int_0^1(f(x))^n\,dx}
{\displaystyle\int_0^1(f(x))^n\,dx}
\ge
\frac{\displaystyle\int_0^1(f(x))^{n+1}\,dx}
{\displaystyle\int_0^1(f(x))^n\,dx}
\ge
\frac{\displaystyle\int_0^\delta(f(x))^{n+1}\,dx}
{\displaystyle\int_0^1(f(x))^n\,dx}
\]
\[
\ge
\frac{\displaystyle\int_0^{\delta_1}(f(x))^{n+1}\,dx}
{\displaystyle\int_0^1(f(x))^n\,dx}
\ge
(1-\varepsilon)
\frac{\displaystyle\int_0^{\delta_1}(f(x))^n\,dx}
{\displaystyle\int_0^{\delta_1}(f(x))^n\,dx+
\displaystyle\int_{\delta_1}^1(f(x))^n\,dx}
\to(1-\varepsilon)\quad(n\to\infty).
\]
由 $\varepsilon>0$ 的任意性知
\[
\lim_{n\to\infty}
\frac{\displaystyle\int_0^\delta(f(x))^{n+1}\,dx}
{\displaystyle\int_0^1(f(x))^n\,dx}=1.
\]

\noindent\textbf{例 4.1.7}\quad 设 $f(x)\ge0,g(x)>0$，两函数在 $[a,b]$ 上连续，求证：
\[
\lim_{n\to\infty}
\left[\int_a^b(f(x))^n g(x)\,dx\right]^{1/n}
=\max_{a\le x\le b}f(x).
\]
（南京大学）

\noindent\textbf{解}\quad 因 $f(x)$ 在 $[a,b]$ 上连续，所以在 $[a,b]$ 上达到最大值. 即存在 $x_0\in[a,b]$，使得
\[
f(x_0)=\max_{a\le x\le b}f(x)\stackrel{\text{记}}{=}M.
\]
于是
\[
\left[\int_a^b(f(x))^n g(x)\,dx\right]^{1/n}
\le \left(\int_a^b M^n g(x)\,dx\right)^{1/n}
=M\left(\int_a^b g(x)\,dx\right)^{1/n}\to M.
\]
又由于 $f$ 在 $x_0$ 处连续，$\forall\varepsilon>0$，$\exists[\alpha,\beta]\subset[a,b]$，使得
\[
f(x)>M-\varepsilon\quad(x\in[\alpha,\beta]),
\]
故
\[
\left[\int_a^b(f(x))^n g(x)\,dx\right]^{1/n}
\ge \left(\int_\alpha^\beta(f(x))^n g(x)\,dx\right)^{1/n}
\ge (M-\varepsilon)\left(\int_\alpha^\beta g(x)\,dx\right)^{1/n}
\to M-\varepsilon.
\]
由 $\varepsilon>0$ 的任意性，知
\[
\lim_{n\to\infty}
\left[\int_a^b(f(x))^n g(x)\,dx\right]^{1/n}
=M=\max_{a\le x\le b}f(x).
\]

下两例将问题转化为积分和的极限.

\noindent\textbf{例 4.1.8}\quad 设 $f(x)$ 在 $[a,b]$ 上可导，$f'(x)$ 在 $[a,b]$ 上可积. $\forall n\in\mathbb{N}$，记
\[
A_n=\sum_{i=1}^n f\left(a+i\frac{b-a}{n}\right)\frac{b-a}{n}
-\int_a^b f(x)\,dx,
\]
试证：
\[
\lim_{n\to\infty}nA_n=\frac{b-a}{2}(f(b)-f(a)).
\]

\noindent\textbf{解 I}\quad 转化为 $f'(x)$ 的积分和.

$1^\circ$ 令 $x_i=a+i\dfrac{b-a}{n}$，则
\begin{equation}
\begin{aligned}
nA_n
&=n\left(\sum_{i=1}^n f(x_i)\frac{b-a}{n}
-\sum_{i=1}^n\int_{x_{i-1}}^{x_i}f(x)\,dx\right)\\
&=n\sum_{i=1}^n\int_{x_{i-1}}^{x_i}(f(x_i)-f(x))\,dx\\
&=n\sum_{i=1}^n\int_{x_{i-1}}^{x_i}f'(\eta_i)(x_i-x)\,dx
\quad(\eta_i\in(x_{i-1},x_i)).
\end{aligned}
\tag{1}
\end{equation}

$2^\circ$ 因 $(x_i-x)$ 不变号，导函数有介值性，因此应用积分第一中值定理，不难知道：$\exists\xi_i\in[x_{i-1},x_i]$ 使得
\[
\int_{x_{i-1}}^{x_i}f'(\eta_i)(x_i-x)\,dx
=f'(\xi_i)\int_{x_{i-1}}^{x_i}(x_i-x)\,dx.
\]
于是式 (1) 成为
\[
\begin{aligned}
nA_n
&=n\sum_{i=1}^n f'(\xi_i)\int_{x_{i-1}}^{x_i}(x_i-x)\,dx
=\frac n2\sum_{i=1}^n f'(\xi_i)(x_i-x_{i-1})^2\\
&=\frac n2\cdot\frac{b-a}{n}\sum_{i=1}^n f'(\xi_i)(x_i-x_{i-1})
\to \frac{b-a}{2}\int_a^b f'(x)\,dx\\
&=\frac{b-a}{2}(f(b)-f(a))\quad(n\to\infty).
\end{aligned}
\]

\noindent\textbf{证 II}\quad 利用两边夹法则及 $f'$ 的 Darboux 和
\[
\lim_{\lambda\to0}\sum_{i=1}^n m_i\Delta x_i
=\lim_{\lambda\to0}\sum_{i=1}^n M_i\Delta x_i
=\int_a^b f'(x)\,dx.
\]
令 $x_i=a+i\dfrac{b-a}{n}$，将 $[a,b]$ $n$ 等分，作分划，如上有式 (1). 记
\[
m_i=\inf_{x_{i-1}\le x\le x_i}f'(x),\qquad
M_i=\sup_{x_{i-1}\le x\le x_i}f'(x)\quad(i=1,2,\cdots,n),
\]
则
\[
m_i(x_i-x)\le f'(\eta_i)(x_i-x)\le M_i(x_i-x),
\]
\[
m_i\int_{x_{i-1}}^{x_i}(x_i-x)\,dx
\le\int_{x_{i-1}}^{x_i}f'(\eta_i)(x_i-x)\,dx
\le M_i\int_{x_{i-1}}^{x_i}(x_i-x)\,dx,
\]
\[
\frac{m_i}{2}(x_i-x_{i-1})^2
\le\int_{x_{i-1}}^{x_i}f'(\eta_i)(x_i-x)\,dx
\le\frac{M_i}{2}(x_i-x_{i-1})^2.
\]
注意 $x_i-x_{i-1}=\dfrac{b-a}{n}$，上式代入式 (1) 得
\[
\frac12(b-a)\sum_{i=1}^n m_i\frac{b-a}{n}
\le nA_n\le
\frac{b-a}{2}\sum_{i=1}^n M_i\frac{b-a}{n}.
\]
其中 $\sum_{i=1}^n m_i\dfrac{b-a}{n}$，$\sum_{i=1}^n M_i\dfrac{b-a}{n}$ 为 $f'(x)$ 的 Darboux 和. 令 $n\to\infty$，取极限知
\[
\lim_{n\to\infty}nA_n
=\frac{b-a}{2}\int_a^b f'(x)\,dx
=\frac{b-a}{2}[f(b)-f(a)].
\]

\noindent\textbf{例 4.1.9}\quad 设 $f(x)$ 在 $[a,b]$ 上可积，$g(x)\ge0$，$g$ 是以 $T>0$ 为周期的函数，在 $[0,T]$ 上可积，试证：
\[
\lim_{n\to\infty}\int_a^b f(x)g(nx)\,dx
=\frac1T\int_0^T g(x)\,dx\int_a^b f(x)\,dx.
\]

\noindent\textbf{证}\quad 因 $g(x)$ 以 $T$ 为周期，因此 $g(nx)$ 以 $\dfrac Tn$ 为周期，当 $n$ 充分大时，$[a,b]$ 含有 $g(nx)$ 的多个周期. 为了把区间变成 $\dfrac Tn$ 的整倍数，取足够大的正整数 $m$，使得
\[
[A,B]=[-mT,mT]\supset[a,b].
\]
这时 $[A,B]$ 相当 $g(nx)$ 的 $2mn$ 个周期. 令
\[
F(x)=
\begin{cases}
f(x), & x\in[a,b],\\
0, & x\in[A,B]\setminus[a,b],
\end{cases}
\]
于是 $F(x)$ 在 $[A,B]$ 上可积，且
\[
I_n\equiv\int_a^b f(x)g(nx)\,dx
=\int_A^B F(x)g(nx)\,dx.
\]
将 $[A,B]$ $2mn$ 等分，作分划 $A=x_0<x_1<\cdots<x_{2mn}=B$. 每个小区间恰是 $g(nx)$ 的一个周期，小区间长度等于 $\dfrac Tn$. 于是
\[
I_n=\int_A^B F(x)g(nx)\,dx
=\sum_{i=1}^{2mn}\int_{x_{i-1}}^{x_i}F(x)g(nx)\,dx.
\]
注意到 $g(x)\ge0$，应用第一积分中值定理，得
\[
I_n=\sum_{i=1}^{2mn}c_i\int_{x_{i-1}}^{x_i}g(nx)\,dx,
\]
其中
\[
c_i:m_i\equiv\inf_{x_{i-1}\le x\le x_i}F(x)
\le c_i\le
M_i\equiv\sup_{x_{i-1}\le x\le x_i}F(x).
\]
因 $[x_{i-1},x_i]$ 是 $g(nx)$ 的一个周期，令 $nx=t$，则
\[
\int_{x_{i-1}}^{x_i}g(nx)\,dx
=\int_0^{T/n}g(nx)\,dx
=\frac1n\int_0^T g(t)\,dt.
\]
代入上式，得
\[
I_n=\frac1T\int_0^Tg(x)\,dx\cdot
\sum_{i=1}^{2mn}c_i\frac Tn.
\]
注意：
\[
\sum_{i=1}^{2mn}m_i\frac Tn
\le \sum_{i=1}^{2mn}c_i\frac Tn
\le \sum_{i=1}^{2mn}M_i\frac Tn.
\]
其左、右两端分别为 $F(x)$ 在 $[A,B]$ 上的 Darboux 和. 故
\[
\lim_{n\to\infty}I_n
=\frac1T\int_0^Tg(x)\,dx\cdot\int_A^B F(x)\,dx
=\frac1T\int_0^Tg(x)\,dx\cdot\int_a^b f(x)\,dx.
\]

\noindent\textbf{例 4.1.10}\quad 证明 Riemann 引理：若 $f(x)$ 在 $[a,b]$ 上可积，$g(x)$ 以 $T$ 为周期，在 $[0,T]$ 上可积，则
\[
\lim_{n\to\infty}\int_a^b f(x)g(nx)\,dx
=\frac1T\int_0^T g(x)\,dx\int_a^b f(x)\,dx.
\]

\noindent\textbf{注}\quad 与上例不同之处在于去掉了条件 $g(x)\ge0$，不能用上例的方法直接证明本例. 但欲证的极限式关于 $g$ 有可加性，利用这一点，我们可以通过取函数正部和负部的方法化为非负的情况.

\noindent\textbf{证}\quad 定义
\[
g^+(x)=
\begin{cases}
g(x),& g(x)\ge0,\\
0,& g(x)<0,
\end{cases}
\qquad
g^-(x)=
\begin{cases}
-g(x),& g(x)\le0,\\
0,& g(x)>0.
\end{cases}
\]
$g^+(x),g^-(x)$ 分别称为 $g(x)$ 的正部、负部. $g(x)$ 以 $T$ 为周期，在 $[0,T]$ 上可积，故 $g^+(x),g^-(x)$ 也如此，且 $g(x)=g^+(x)-g^-(x)$. 因而
\[
\begin{aligned}
\lim_{n\to\infty}\int_a^b f(x)g(nx)\,dx
&=\lim_{n\to\infty}\int_a^b f(x)[g^+(nx)-g^-(nx)]\,dx\\
&=\lim_{n\to\infty}\int_a^b f(x)g^+(nx)\,dx
-\lim_{n\to\infty}\int_a^b f(x)g^-(nx)\,dx\\
&=\frac1T\int_0^T g^+(x)\,dx\int_a^b f(x)\,dx
-\frac1T\int_0^T g^-(x)\,dx\int_a^b f(x)\,dx\\
&=\frac1T\int_0^T[g^+(x)-g^-(x)]\,dx\cdot\int_a^b f(x)\,dx\\
&=\frac1T\int_0^T g(x)\,dx\cdot\int_a^b f(x)\,dx.
\end{aligned}
\]

\noindent\textbf{练习}\quad 设函数 $f(x)$ 在区间 $[0,1]$ 上有一阶连续导数，且 $f(0)=f(1)$，$g(x)$ 是周期为 $1$ 的连续函数. 记
\[
a_n=\int_0^1 f(x)g(nx)\,dx,
\]
求证：$\displaystyle\lim_{n\to\infty}na_n=0$.（北京大学）

\noindent\textbf{注意}\quad 根据 Riemann 引理（例 4.1.10），知 $\displaystyle\lim_{n\to\infty}a_n=0$，但该题不仅要证明 $a_n$ 是无穷小量，还要证明比 $\dfrac1n$ 高阶.

\noindent\textbf{提示}\quad 由题设可知：
\[
G(x)=\int_0^x g(t)\,dt
\]
以 $1$ 为周期，且 $G'(x)=g(x)$ 连续.

再提示：
\[
\begin{aligned}
na_n
&=\int_0^1 f(x)g(nx)\,d(nx)
=\int_0^1 f(x)\,dG(nx)\\
&=f(x)G(nx)\bigg|_0^1-\int_0^1 f'(x)G(nx)\,dx
=-\int_0^1 f'(x)G(nx)\,dx.
\end{aligned}
\]
再利用 Riemann 引理，当 $n\to\infty$ 时，
\[
\int_0^1 f'(x)G(nx)\,dx
\to \int_0^1G(x)\,dx\cdot\int_0^1f'(x)\,dx
=\int_0^1G(x)\cdot[f(1)-f(0)]\,dx=0.
\]
本章 §4.5 的例 4.5.32 将 Riemann 引理推广到反常积分.

前面我们利用积分和来求极限. 不仅如此，有时我们还可借助可积的充要条件来求极限.

\noindent\textbf{例 4.1.11}\quad 设 $f(x)$ 在 $[a,b]$ 上可积，在 $[a,b]$ 外保持有界. 试证：函数 $f(x)$ 具有积分连续性，即
\[
\lim_{h\to0}\int_a^b|f(x+h)-f(x)|\,dx=0.
\]

\noindent\textbf{分析}\quad 问题在于证明：$\forall\varepsilon>0$，$\exists\delta>0$，当 $|h|<\delta$ 时，有
\[
\int_a^b|f(x+h)-f(x)|\,dx<\varepsilon.
\]
为了利用可积性，将 $[a,b]$ 作一分划，例如令 $x_i=a+i\dfrac{b-a}{n}$，将其 $n$ 等分，这时小区间长度为 $\dfrac{b-a}{n}$，
\[
\int_a^b|f(x+h)-f(x)|\,dx
=\sum_{i=1}^n\int_{x_{i-1}}^{x_i}|f(x+h)-f(x)|\,dx.
\]
若令 $|h|<\dfrac{b-a}{n}$，则点 $(x+h)$ 要么位于 $x$ 所在的（第 $i$ 个）小区间上，从而
$|f(x+h)-f(x)|<\omega_i$（$\omega_i$ 为第 $i$ 个小区间上 $f$ 的振幅）；要么 $(x+h)$ 在相邻的一个区间上，例如，当 $h>0$ 时有
\[
|f(x+h)-f(x)|\le |f(x+h)-f(x_i)|+|f(x_i)-f(x)|\le\omega_{i+1}+\omega_i;
\]
当 $h<0$ 时有
\[
|f(x+h)-f(x)|\le |f(x+h)-f(x_{i-1})|+|f(x_{i-1})-f(x)|
\le\omega_{i-1}+\omega_i.
\]
总之有 $|f(x+h)-f(x)|\le\omega_{i-1}+\omega_i+\omega_{i+1}$. 对第 $1,n$ 两个小区间只要将 $\omega_0,\omega_{n+1}$ 看成 $2M$ 即可（$M$ 表示 $f$ 的界：$|f(x)|\le M(\forall x)$）. 于是
\[
\sum_{i=1}^n\int_{x_{i-1}}^{x_i}|f(x+h)-f(x)|\,dx
\le 3\sum_{i=1}^n\omega_i\Delta x_i+2M\frac{b-a}{n}
\quad\left(\Delta x_i=\frac{b-a}{n}\right).
\]
由于 $f$ 在 $[a,b]$ 上可积，$\forall\varepsilon>0$，当初只要把 $n$ 取得足够大，使得
\[
\sum_{i=1}^n\omega_i\Delta x_i<\frac\varepsilon6,\qquad
2M\frac{b-a}{n}<\frac\varepsilon2.
\]
取 $\delta<\dfrac{b-a}{n}$，则当 $|h|<\delta$ 时，便恒有
\[
0\le\int_a^b|f(x+h)-f(x)|\,dx
\le3\sum_{i=1}^n\omega_i\Delta x_i+2M\frac{b-a}{n}<\varepsilon.
\]

在求积分的极限时，也经常用到 L'Hospital 法则. 如

\noindent\textbf{例 4.1.12}\quad 设 $f(x)$ 在 $[A,B]$ 上连续，$A<a<b<B$. 试证：
\[
\lim_{h\to0}\int_a^b\frac{f(x+h)-f(x)}{h}\,dx=f(b)-f(a).
\]

\noindent\textbf{证}\quad
\[
\begin{aligned}
\lim_{h\to0}\int_a^b\frac{f(x+h)-f(x)}{h}\,dx
&=\lim_{h\to0}\frac1h\left[\int_a^b f(x+h)\,dx-\int_a^b f(x)\,dx\right]\\
&=\lim_{h\to0}\frac1h\left[\int_{a+h}^{b+h}f(t)\,dt-\int_a^b f(x)\,dx\right].
\end{aligned}
\]
用 $\dfrac00$ 型 L'Hospital 法则，上式
\[
=\lim_{h\to0}[f(b+h)-f(a+h)]=f(b)-f(a).
\]
证毕.

\noindent\textbf{练习}\quad 已知 $f(x)$ 在 $x=0$ 处连续可导，且 $f(0)=0,f'(0)=5$，求极限：
\[
\lim_{x\to0}\frac1x\int_0^1 f(xt)\,dt.
\]
（浙江大学）

\noindent\textbf{解 I}\quad
\[
\begin{aligned}
\lim_{x\to0}\frac1x\int_0^1 f(xt)\,dt
&=\lim_{x\to0}\frac{\displaystyle\int_0^x f(u)\,du}{x^2}
\stackrel{\text{L'Hospital 法则}}{=}
\lim_{x\to0}\frac{f(x)}{2x}\\
&=\frac12\lim_{x\to0}\frac{f(x)-f(0)}x
=\frac12f'(0)=\frac52.
\end{aligned}
\]

\noindent\textbf{解 II}\quad 对 $f(x)$ 应用 Taylor 公式，在 $x=0$ 处展开：
\[
f(xt)=f(0)+f'(0)xt+o(xt)=5xt+o(xt).
\]
因此
\[
\lim_{x\to0}\frac1x\int_0^1f(xt)\,dt
=\lim_{x\to0}\frac1x\int_0^1[5xt+o(xt)]\,dt=\frac52.
\]

\noindent\textbf{注}\quad 1）为什么说
\[
\lim_{x\to0}\frac1x\int_0^1o(xt)\,dt=0?
\]
因为 $f$ 连续，$f(u)=5u+o(u)$，因此余项 $o(u)=f(u)-5u$ 也连续，故
\[
\left(\int_0^x o(u)\,du\right)'=o(x).
\]
于是
\[
\lim_{x\to0}\frac1x\int_0^1o(xt)\,dt
=\lim_{x\to0}\frac1{x^2}\int_0^x o(u)\,du
\stackrel{\text{L'Hospital 法则}}{=}
\lim_{x\to0}\frac{o(x)}{2x}=0.
\]

2）也可用 $\varepsilon$--$\delta$ 方法证明，因
\[
\lim_{u\to0^+}\frac{o(u)}u=0,
\]
对于 $0<u\le x$，$\forall\varepsilon>0$，$\exists\delta>0$，当 $0<x<\delta$ 时，有
\[
\left|\frac{o(u)}u\right|<\varepsilon.
\]
因此
\[
\frac1x\int_0^1o(xt)\,dt
=\left|\frac1{x^2}\int_0^x o(u)\,du\right|
\le\frac1{x^2}\int_0^x\left|u\frac{o(u)}u\right|\,du
<\frac\varepsilon{x^2}\int_0^x u\,du=\frac\varepsilon2.
\]

下例中的序列 $\{x_n\}$ 是通过积分等式定义的，我们来求它的极限.

\noindent\textbf{例 4.1.13}\quad 设 $f(x)$ 在 $[a,b]$ 上非负、连续、严格递增. 由积分中值定理，$\forall n\in\mathbb{N}$，$\exists x_n\in[a,b]$，使得
\[
f^n(x_n)=\frac1{b-a}\int_a^b f^n(x)\,dx.
\]
求极限 $\displaystyle\lim_{n\to\infty}x_n$（注意这里 $f^n$ 是 $f$ 的 $n$ 次幂）.

\noindent\textbf{分析}\quad 首先，通过变换可把 $[a,b]$ 上的问题化为 $[0,1]$ 上类似的问题. 令
\[
x=a+t(b-a),\quad x_n=a+t_n(b-a),\quad F(t)=f[a+t(b-a)],
\]
则 $F(t)\ge0$ 严增 $(t\in[0,1])$，且
\[
F^n(t_n)=\int_0^1F^n(t)\,dt.
\]
只要求出了 $\displaystyle\lim_{n\to\infty}t_n$，则
\[
\lim_{n\to\infty}x_n=a+(b-a)\lim_{n\to\infty}t_n.
\]
为了猜测极限 $\displaystyle\lim_{n\to\infty}t_n$ 的值，考虑 $F(t)=t$ 的情况. 此时
\[
F^n(t_n)=t_n^n=\int_0^1t^n\,dt=\frac1{n+1},
\]
\[
t_n=\sqrt[n]{\frac1{n+1}}\to1\quad(n\to\infty).
\]
因此，我们希望能证明，在一般情况下仍有 $\displaystyle\lim_{n\to\infty}t_n=1$. 由于 $t_n\in[0,1]$，为此只要证明：$\forall\varepsilon>0$，当 $n$ 充分大时，有 $1-\varepsilon<t_n$. 注意 $F(t)\ge0$ 严增，这等价于
\[
F^n(1-\varepsilon)<F^n(t_n)=\int_0^1F^n(t)\,dt.
\tag{1}
\]
此式解不出 $n$，我们设法将式 (1) 右端进行化简和缩小，只要使缩小后的量大于式 (1) 左端的 $F^n(1-\varepsilon)$，则式 (1) 自然成立. 任取 $\xi\in(0,1)$，有
\[
\int_0^1F^n(t)\,dt\ge\int_\xi^1F^n(t)\,dt
\ge F^n(\xi)\cdot(1-\xi).
\]
故要式 (1) 成立，只要使
\[
F^n(\xi)\cdot(1-\xi)>F^n(1-\varepsilon),
\tag{2}
\]
亦即
\[
\left[\frac{F(1-\varepsilon)}{F(\xi)}\right]^n<1-\xi.
\tag{3}
\]
因 $0\le F(x)$ 严增，取 $\xi=1-\dfrac\varepsilon2>1-\varepsilon$，则
\[
0<\frac{F(1-\varepsilon)}{F(\xi)}<1,\qquad
\lim_{n\to\infty}\left[\frac{F(1-\varepsilon)}{F(\xi)}\right]^n=0.
\]
故 $n$ 充分大时式 (3) 成立. 故有式 (1) 成立，等价地有 $1-\varepsilon<t_n\le1$，这就证明了
\[
\lim_{n\to\infty}t_n=1,
\]
从而
\[
\lim_{n\to\infty}x_n=b.
\]

\subsection*{单元练习 4.1}

\noindent\textbf{4.1.1}\quad 设 $f(x)$ 在 $[0,1]$ 上连续，且 $f(x)>0$，求极限
\[
\lim_{n\to\infty}
\sqrt[n]{f\left(\frac1n\right)f\left(\frac2n\right)\cdots
f\left(\frac{n-1}{n}\right)f(1)}.
\]
（上海科技大学）
\[
\left(e^{\int_0^1\ln f(x)\,dx}\right)
\]

\noindent\textbf{4.1.2}\quad 考虑积分 $\displaystyle\int_0^1(1-x)^n\,dx$，证明：
\[
C_n^0-\frac12C_n^1+\frac13C_n^2-\cdots+\frac{(-1)^n}{n+1}C_n^n
=\frac1{n+1}.
\]

\noindent\textbf{提示}\quad 在
\[
(1-x)^n=\sum_{k=0}^n(-1)^kC_n^kx^k
\]
两端同时积分.

\noindent\textbf{4.1.3}\quad 设 $f(x)$ 在 $[0,1]$ 上可微，且对任意 $x\in(0,1)$ 有 $|f'(x)|\le M$. 求证：对任意正整数 $n$，有
\[
\left|\int_0^1f(x)\,dx-\frac1n\sum_{i=1}^n f\left(\frac in\right)\right|
\le\frac Mn,
\]
其中 $M$ 是一个与 $x$ 无关的常数.（南开大学）

\noindent\textbf{提示}\quad 将 $[0,1]$ $n$ 等分，每段上应用微分中值定理.
\[
\text{左端}
=\left|\frac1n\sum_{i=1}^nf(\xi_i)-\frac1n\sum_{i=1}^nf\left(\frac in\right)\right|
\stackrel{\text{微分中值定理}}{\le}
\frac1n\left|\sum_{i=1}^nf'(\eta_i)\left(\xi_i-\frac in\right)\right|
\le\text{右端}.
\]

\noindent\textbf{4.1.4}\quad 若 $f(x)$ 在 $[a,b]$ 上可积，$g(x)$ 是以 $T$ 为周期的函数，且在 $[0,T]$ 上可积. 试证：
\[
\lim_{\lambda\to+\infty}\int_a^b f(x)g(\lambda x)\,dx
=\frac1T\int_0^Tg(x)\,dx\int_a^bf(x)\,dx.
\]

\noindent\textbf{提示}\quad 参看例 4.1.9 和例 4.1.10.

\noindent\textbf{4.1.5}\quad 设 $s(x)=4[x]-2[2x]+1$，其中 $[x]$ 代表数 $x$ 的整数部分（即不超过 $x$ 的整数之最大值），$n$ 代表自然数，$f(x)$ 在 $[0,1]$ 上可积. 证明
\[
\lim_{n\to\infty}\int_0^1f(x)s(nx)\,dx=0.
\]
（兰州大学）

\noindent\textbf{提示}\quad $s(x)$ 周期为 $1$，且 $\displaystyle\int_0^1s(x)\,dx=0$，用 Riemann 引理.

\noindent\textbf{4.1.6}\quad 设 $f_0(x)$ 在 $[0,1]$ 上可积，$f_0(x)>0$，
\[
f_n(x)=\sqrt{\int_0^x f_{n-1}(t)\,dt},\quad n=1,2,\cdots.
\]
1）若记 $f_1(x)$ 在 $[0,1]$ 上最大值为 $M$，试证：
\[
f_n(x)\le M^{1/2^{n-1}}a_nx^{1-1/2^{n-1}},\quad n=1,2,\cdots,
\]
其中
\[
a_n=\prod_{k=1}^{n-2}
\left(\frac{2^k}{2^{k+1}-1}\right)^{1/2^{n-1-k}}.
\]

2）若记 $f_1(x)$ 在 $[\delta,1]$ 上的最小值为 $m>0$（$0<\delta<1$ 任意给定），试证：
\[
m^{1/2^{n-1}}a_n(x-\delta)^{1-1/2^{n-1}}
\le f_n(x)\le
M^{1/2^{n-1}}a_nx^{1-1/2^{n-1}},\quad n=1,2,\cdots.
\]

3）已知
\[
\lim_{n\to\infty}\sum_{k=1}^{n-2}
\frac1{2^{n-1-k}}\left[\ln\left(1+\frac1{2^{k+1}-1}\right)\right]=0
\]
（见例 5.1.57 及其练习），试证：$\displaystyle\lim_{n\to\infty}f_n(x)=\dfrac x2$.

\noindent\textbf{证}\quad 1）因 $f_0(x)$ 在 $[0,1]$ 上可积，故 $f_1(x)$ 在 $[0,1]$ 上连续，有界，即 $\exists M$ 使得
\[
0\le f_1(x)\le M,\quad \forall x\in[0,1].
\]
得
\[
f_3(x)=\sqrt{\int_0^x f_2(t)\,dt}
\le\sqrt{\int_0^x M^{1/2}t^{1/2}\,dt}
=M^{1/2^2}\left(\frac{2}{2^2-1}\right)^{1/2}x^{1-1/2^2},
\]
\[
f_4(x)=\sqrt{\int_0^x f_3(t)\,dt}
\le M^{1/2^3}
\left(\frac{2}{2^2-1}\right)^{1/2^2}
\left(\frac{2^2}{2^3-1}\right)^{1/2}
x^{1-1/2^3},
\]
\[
\cdots\cdots
\]
\[
f_n(x)=\sqrt{\int_0^x f_{n-1}(t)\,dt}
\le M^{1/2^{n-1}}a_nx^{1-1/2^{n-1}},
\tag{1}
\]
其中
\[
a_n=\prod_{k=1}^{n-2}
\left(\frac{2^k}{2^{k+1}-1}\right)^{1/2^{n-1-k}}.
\tag{2}
\]

2）记 $f_1(x)$ 在 $[\delta,1]$ 上的最小值为 $m>0$（$0<\delta<1$ 任意给定），定义
\[
g(x)=
\begin{cases}
m, & x\in[\delta,1],\\
0, & x\in[0,\delta).
\end{cases}
\]
于是 $\forall x\in[0,1]$，恒有 $f_1(x)\ge g(x)$，可知
\[
f_2(x)=\sqrt{\int_0^x f_1(t)\,dt}
\ge\sqrt{\int_\delta^xg(t)\,dt}
=m^{1-1/2}(x-\delta)^{1-1/2},
\]
\[
\cdots\cdots
\]
\[
f_n(x)\ge m^{1/2^{n-1}}a_n(x-\delta)^{1-1/2^{n-1}}.
\tag{3}
\]
将 (1), (3) 联立，可得
\[
m^{1/2^{n-1}}a_n(x-\delta)^{1-1/2^{n-1}}
\le f_n(x)\le
M^{1/2^{n-1}}a_nx^{1-1/2^{n-1}}.
\tag{4}
\]

3）注意到
\[
\lim_{n\to\infty}M^{1/2^{n-1}}
=\lim_{n\to\infty}m^{1/2^{n-1}}=1,
\quad
\lim_{n\to\infty}(x-\delta)^{1-1/2^{n-1}}=x-\delta,\quad
\lim_{n\to\infty}x^{1-1/2^{n-1}}=x.
\]
下面证明：
\[
\lim_{n\to\infty}a_n
\stackrel{\text{记}}{=}
\lim_{n\to\infty}\prod_{k=1}^{n-2}
\left(\frac{2^k}{2^{k+1}-1}\right)^{1/2^{n-1-k}}
=\frac12.
\tag{5}
\]
事实上，
\[
\lim_{n\to\infty}\left[
\prod_{k=1}^{n-2}
\left(\frac{2^k}{2^{k+1}-1}\right)^{1/2^{n-1-k}}
\right]=\frac12
\]
\[
\Longleftrightarrow
\lim_{n\to\infty}\sum_{k=1}^{n-2}
\frac1{2^{n-1-k}}\ln\frac{2^k}{2^{k+1}-1}
=\ln\frac12
\]
\[
\Longleftrightarrow
\lim_{n\to\infty}\sum_{k=1}^{n-2}
\frac1{2^{n-1-k}}\left(\ln\frac{2^k}{2^{k+1}-1}-\ln\frac12\right)=0
\]
\[
\Longleftrightarrow
\lim_{n\to\infty}\sum_{k=1}^{n-2}
\frac1{2^{n-1-k}}\left[
\ln\left(1+\frac1{2^{k+1}-1}\right)\right]=0.
\]
所以式 (5) 成立，故
\[
\lim_{n\to\infty}f_n(x)=\frac x2.
\]
再在式 (4) 里令 $n\to\infty$，取极限可得
\[
\frac12(x-\delta)\le\liminf f_n(x)\le\limsup f_n(x)\le\frac12x.
\]
然后，利用 $\delta>0$ 的任意性，得 $\displaystyle\lim_{n\to\infty}f_n(x)=\dfrac x2$. 证毕.

\noindent\textbf{注}\quad 本题是道经典题目，通常是应用 Toeplitz 定理来证明的. 现行教材并未提及 Toeplitz 定理，因此这里希望回避 Toeplitz 定理直接证明. 此外（为了化解难度）将一问改成了三小题.

\noindent\textbf{4.1.7}\quad 设 $f(x),g(x)$ 在 $[a,b]$ 上连续，$f(x)>0,g(x)>0$，求
\[
\lim_{p\to\infty}\left(\int_a^b g(x)f^p(x)\,dx\right)^{1/p^2}.
\]

\noindent\textbf{提示}\quad 参看例 4.1.7.

\noindent\textbf{4.1.8}\quad 设 $f(x)$ 在 $[a,b]$ 上二次可微，且 $f''(x)$ 在 $[a,b]$ 上可积，记
\[
B_n=\int_a^b f(x)\,dx-\frac{b-a}{n}
\sum_{i=1}^n f\left[a+(2i-1)\frac{b-a}{2n}\right],
\]
试证：
\[
\lim_{n\to\infty}n^2B_n
=\frac{(b-a)^2}{24}[f'(b)-f'(a)].
\]

\noindent\textbf{提示}\quad 跟例 4.1.8 不同的是：将区间 $[a,b]$ $n$ 等分，则小区间 $[x_{i-1},x_i]$ 的端点为
\[
x_i=a+i\frac{b-a}{n},
\]
中点为
\[
\eta_i=a+(2i-1)\frac{b-a}{2n}.
\]
题中的
\[
B_n=\int_a^b f(x)\,dx-\frac{b-a}{n}\sum_{i=1}^n f(\eta_i).
\tag{1}
\]
积分
\[
\int_a^b f(x)\,dx=\sum_{i=1}^n\int_{x_{i-1}}^{x_i}f(x)\,dx.
\tag{2}
\]
将 $f(x)$ 在 $\eta_i$ 点按 Taylor 公式展开：
\[
f(x)=f(\eta_i)+f'(\eta_i)(x-\eta_i)
+f''(\xi_i)\frac{(x-\eta_i)^2}{2}.
\tag{3}
\]
再提示
\[
\begin{aligned}
n^2B_n
&\stackrel{(1),(2)}{=}
n^2\left(\sum_{i=1}^n\int_{x_{i-1}}^{x_i}f(x)\,dx
-\sum_{i=1}^n\int_{x_{i-1}}^{x_i}f(\eta_i)\,dx\right)\\
&=n^2\sum_{i=1}^n\int_{x_{i-1}}^{x_i}[f(x)-f(\eta_i)]\,dx\\
&\stackrel{(3)}{=}
n^2\sum_{i=1}^n\int_{x_{i-1}}^{x_i}
\left[f'(\eta_i)(x-\eta_i)+f''(\xi_i)\frac{(x-\eta_i)^2}{2}\right]\,dx,
\end{aligned}
\]
其中第一个积分
\[
\int_{x_{i-1}}^{x_i}f'(\eta_i)(x-\eta_i)\,dx
=f'(\eta_i)\int_{x_{i-1}}^{x_i}(x-\eta_i)\,dx=0,
\]
第二个积分
\[
\begin{aligned}
\int_{x_{i-1}}^{x_i}f''(\xi_i)\frac{(x-\eta_i)^2}{2}\,dx
&=f''(\xi_i)\int_{x_{i-1}}^{x_i}\frac{(x-\eta_i)^2}{2}\,dx\\
&=\frac12f''(\xi_i)\cdot\frac13(x-\eta_i)^3\bigg|_{x_{i-1}}^{x_i}
=\frac{f''(\xi_i)}{24n^3}(b-a)^3.
\end{aligned}
\]
所以
\[
n^2B_n=\frac{(b-a)^2}{24}
\sum_{i=1}^n f''(\xi_i)\frac{b-a}{n}
\xrightarrow[n\to\infty]{}
\frac{(b-a)^2}{24}\int_a^b f''(x)\,dx
=\frac{(b-a)^2}{24}[f'(b)-f'(a)].
\]

\noindent\textbf{4.1.9}\quad 设
\[
A_n=\frac1{n+1}+\frac1{n+2}+\cdots+\frac1{2n},\qquad
B_n=\frac2{2n+1}+\frac2{2n+3}+\cdots+\frac2{4n-1}.
\]
试证：
\[
\lim_{n\to\infty}n(\ln2-A_n)=\frac14,\qquad
\lim_{n\to\infty}n^2(\ln2-B_n)=\frac1{32}.
\]

\noindent\textbf{提示}\quad 两小题分别是例 4.1.8 和习题 4.1.8 的特例. 现在函数是
\[
f(x)=\frac1{1+x},
\]
区间是 $[a,b]=[0,1]$，
\[
\ln2=\int_0^1\frac{dx}{1+x}.
\]
再提示
\[
\lim_{n\to\infty}n(\ln2-A_n)
=\lim_{n\to\infty}n\left(\int_0^1\frac{dx}{1+x}
-\sum_{i=1}^n\frac1{1+\frac in}\cdot\frac1n\right)
\stackrel{\text{例 4.1.8}}{=}
-\frac12\frac1{1+x}\bigg|_0^1=\frac14,
\]
\[
\lim_{n\to\infty}n^2(\ln2-B_n)
=\lim_{n\to\infty}n^2\left(\int_0^1\frac{dx}{1+x}
-\sum_{i=1}^n\frac1{1+\frac{2i-1}{2n}}\cdot\frac1n\right)
\]
\[
\stackrel{\text{习题 4.1.8}}{=}
\frac1{24}\frac{-1}{(1+x)^2}\bigg|_0^1
=\frac1{32}.
\]

\noindent\textbf{4.1.10}\quad 设 $f(x)$ 在 $[a,b]$ 上可积，记
\[
f_{in}=f\left(a+i\frac{b-a}{n}\right).
\]
试利用不等式
\[
|\ln(1+x)-x|\le2x^2\quad\left(\text{当 }|x|<\frac12\text{ 时}\right)
\]
证明：
\[
\lim_{n\to\infty}
\left(1+f_{1n}\frac{b-a}{n}\right)
\left(1+f_{2n}\frac{b-a}{n}\right)\cdots
\left(1+f_{nn}\frac{b-a}{n}\right)
=e^{\int_a^b f(x)\,dx}.
\]

\noindent\textbf{提示}\quad 记
\[
x_{in}=a+i\frac{b-a}{n}\quad(i=0,1,2,\cdots,n)
\]
（原式取对数，将连乘变为连加）.

\noindent\textbf{证}\quad 原式两端同时取对数，知原式等价于
\[
\lim_{n\to\infty}\ln\prod_{i=1}^n
\left(1+f(x_{in})\frac{b-a}{n}\right)
=\int_a^b f(x)\,dx
=\lim_{n\to\infty}\sum_{i=1}^n f(x_{in})\frac{b-a}{n}.
\tag{1}
\]
注意到
\[
\ln\prod_{i=1}^n\left(1+f(x_{in})\frac{b-a}{n}\right)
=\sum_{i=1}^n\ln\left(1+f(x_{in})\frac{b-a}{n}\right).
\]
知原式等价于
\[
\lim_{n\to\infty}
\left[\sum_{i=1}^n
\ln\left(1+f(x_{in})\frac{b-a}{n}\right)
-f(x_{in})\frac{b-a}{n}\right]=0.
\tag{2}
\]
因 $f(x)$ 在 $[a,b]$ 上可积，故 $f(x)$ 有界：$\exists M>0$，使得 $|f(x)|\le M$. 因此当 $n>2M(b-a)$ 时，有
\[
\left|f(x_{in})\frac{b-a}{n}\right|<\frac12.
\]
故在已知的不等式 $|\ln(1+x)-x|\le2x^2\left(|x|<\dfrac12\right)$ 中，可令
\[
x=f(x_{in})\frac{b-a}{n},
\]
得
\[
\begin{aligned}
&\left|\sum_{i=1}^n\left[
\ln\left(1+f(x_{in})\frac{b-a}{n}\right)
-f(x_{in})\frac{b-a}{n}\right]\right|\\
\le{}&2\sum_{i=1}^n\left(|f(x_{in})|\frac{b-a}{n}\right)^2
=\frac{2(b-a)}{n}\sum_{i=1}^n(f(x_{in}))^2\frac{b-a}{n}
\longrightarrow0\cdot\int_a^b f^2(x)\,dx=0.
\end{aligned}
\]
（因 $f$ 可积 $\Rightarrow f^2$ 亦可积）式 (2) 获证，式 (1) 成立. 如此，
\[
\lim_{n\to\infty}\prod_{i=1}^n
\left(1+f(x_{in})\frac{b-a}{n}\right)
=\lim_{n\to\infty}
e^{\ln\prod_{i=1}^n(1+f(x_{in})\frac{b-a}{n})}
=e^{\lim_{n\to\infty}\sum_{i=1}^n f(x_{in})\frac{b-a}{n}}
=e^{\int_a^b f(x)\,dx}.
\]

\noindent\textbf{4.1.11}\quad 设 $f(x)$ 是在 $[-1,1]$ 上可积，在 $x=0$ 处连续的函数. 记
\[
\varphi_n(x)=
\begin{cases}
(1-x)^n, & 0\le x\le1,\\
e^{nx}, & -1\le x\le0,
\end{cases}
\]
证明：
\[
\lim_{n\to\infty}\frac n2\int_{-1}^1f(x)\varphi_n(x)\,dx=f(0).
\]
（浙江大学）

\noindent\textbf{提示}\quad 参看例 4.1.5（用拟合法）.

\section*{§4.2\quad 定积分的可积性}

\noindent\textbf{导读}\quad 本节内容是数学院系该课程难点之一，难度较大，其他院系学生可从略. 数学院系学生以正文为主，习题部分可作机动.

本节主要讨论如何证明一个函数在给定区间上是否可积的问题. 因为关于这一部分的知识各书上写法不尽相同，为了与读者取得共同语言，首先我们对基本内容作一概述.

\noindent\textbf{Riemann 积分的定义}\quad 函数 $f(x)$ 在 $[a,b]$ 上可积是指：当区间 $[a,b]$ 的分划无限加细时，积分和
\[
\sum_{i=1}^n f(\xi_i)\Delta x_i
\]
有确定的极限. 详细地说：若对 $[a,b]$ 的任一分划
\[
T:a=x_0<x_1<\cdots<x_n=b
\]
及任意 $\xi_i\in[x_{i-1},x_i]\ (i=1,2,\cdots,n)$ 作积分和 $\sum_{i=1}^n f(\xi_i)\Delta x_i$，当
\[
\lambda=\max_{1\le i\le n}\Delta x_i
=\max_{1\le i\le n}(x_i-x_{i-1})\to0
\]
时，极限
\[
I=\lim_{\lambda\to0}\sum_{i=1}^n f(\xi_i)\Delta x_i
\tag{A}
\]
存在，则称 $f(x)$ 在 $[a,b]$ 上可积. 其中式 (A) 用 $\varepsilon$--$\delta$ 表述，即 $\forall\varepsilon>0$，$\exists\delta>0$，当 $\lambda<\delta$ 时，有
\[
\left|I-\sum_{i=1}^n f(\xi_i)\Delta x_i\right|<\varepsilon.
\tag{B}
\]
直接应用定义判断可积性要克服两重困难：一是分划 $T$ 的任意性，二是 $\xi_i\in[x_{i-1},x_i]$ 选取的任意性. 为了使问题简化，引入 Darboux 和的概念. 记
\[
m_i=\inf_{x_{i-1}\le x\le x_i}f(x),\qquad
M_i=\sup_{x_{i-1}\le x\le x_i}f(x),
\]
则
\[
\overline S\equiv\overline S(T)\equiv\sum_{i=1}^nM_i\Delta x_i
\quad\text{与}\quad
\underline S\equiv\underline S(T)\equiv\sum_{i=1}^nm_i\Delta x_i
\]
分别称为 $f$ 的 Darboux 上和与 Darboux 下和.

关于 Darboux 和有如下五条重要性质：

（1）对于任意分划 $T$ 及 $\xi_i$ 的任意选法，恒有
\[
\underline S(T)\le\sum_{i=1}^n f(\xi_i)\Delta x_i\le\overline S(T).
\]

（2）当分划加细时，Darboux 下和不会减小，Darboux 上和不会增大，即
\[
\underline S(T)\nearrow,\qquad \overline S(T)\searrow\quad(\text{当细分时}).
\]

（3）对任意两分划 $T$ 与 $T'$，恒有 $\underline S(T)\le\overline S(T')$.

（4）$I^0\equiv\inf_T\overline S(T)$ 称为 $f(x)$ 在 $[a,b]$ 的上积分，$I_0\equiv\sup_T\underline S(T)$ 称为 $f(x)$ 在 $[a,b]$ 的下积分，上、下积分与 Darboux 和有关系：
\[
\underline S\le I_0\le I^0\le\overline S.
\]

（5）
\[
\lim_{\lambda\to0}\underline S(T)=I_0,\qquad
\lim_{\lambda\to0}\overline S(T)=I^0.
\]
在此基础上，我们获得如下可积充要条件.

\noindent\textbf{定理 1}\quad $f(x)$ 在 $[a,b]$ 上可积的充要条件是
\[
\lim_{\lambda\to0}\sum_{i=1}^n\omega_i\Delta x_i=0,
\tag{C}
\]
其中 $\omega_i=M_i-m_i$ 为 $f(x)$ 在 $[x_{i-1},x_i]$ 上的振幅. 式 (C) 用 $\varepsilon$--$\delta$ 表述，即 $\forall\varepsilon>0$，$\exists\delta>0$，当 $\lambda<\delta$ 时，
\[
0\le\sum\omega_i\Delta x_i<\varepsilon.
\]

由此，利用 Darboux 和的性质（4）、（5），直接可看出有

\noindent\textbf{定理 2}\quad $f(x)$ 在 $[a,b]$ 上可积的充要条件是 $I_0=I^0$.

下面讨论如何证明可积性问题.

\subsection*{一、直接用定义证明可积性}

\noindent\textbf{要点}\quad （1）直接应用式 (B) 找 $\delta$；（2）反证法：用式 (B) 的反面形式，找出矛盾.

\noindent\textbf{例 4.2.1}\quad 设 $f(x),F(x)$ 在 $[a,b]$ 上连续，且当 $a<x<b$ 时 $F'(x)=f(x)$. 试用定义直接证明 $f(x)$ 在 $[a,b]$ 上可积，且积分值可通过如下的 Newton--Leibniz 公式计算：
\[
\int_a^b f(x)\,dx=F(b)-F(a).
\]

\noindent\textbf{证}\quad $\forall\varepsilon>0$，要式 (B) 成立，即要
\[
\left|\sum_{i=1}^n f(\xi_i)\Delta x_i-(F(b)-F(a))\right|<\varepsilon.
\tag{1}
\]
对任意分划 $T:a=x_0<x_1<x_2<\cdots<x_n=b$，将 $F(b)-F(a)$ 改写为
\[
\begin{aligned}
F(b)-F(a)
&=\sum_{i=1}^n(F(x_i)-F(x_{i-1}))\quad(\text{应用微分中值定理})\\
&=\sum_{i=1}^nF'(\eta_i)(x_i-x_{i-1})
=\sum_{i=1}^nf(\eta_i)\Delta x_i,\quad
\eta_i\in[x_{i-1},x_i].
\end{aligned}
\]
于是
\[
\text{式 (1) 左端}
=\left|\sum_{i=1}^n(f(\xi_i)-f(\eta_i))\Delta x_i\right|
\le\sum_{i=1}^n|f(\xi_i)-f(\eta_i)|\Delta x_i
\quad(\xi_i,\eta_i\in[x_{i-1},x_i]).
\]
因此，问题只要证明
\[
\sum_{i=1}^n|f(\xi_i)-f(\eta_i)|\Delta x_i<\varepsilon.
\tag{2}
\]
但由于 $f$ 在 $[a,b]$ 上连续，所以一致连续：$\forall\varepsilon>0$，$\exists\delta>0$，当 $x',x''\in[a,b]$，$|x''-x'|<\delta$ 时，有
\[
|f(x'')-f(x')|<\frac{\varepsilon}{b-a}.
\]
因此，当 $\lambda<\delta$ 时，有 $|\xi_i-\eta_i|\le\Delta x_i\le\lambda<\delta$，故
\[
\sum_{i=1}^n|f(\xi_i)-f(\eta_i)|\Delta x_i
\le\sum_{i=1}^n\frac{\varepsilon}{b-a}\Delta x_i
=\frac{\varepsilon}{b-a}\sum_{i=1}^n\Delta x_i=\varepsilon.
\]
问题获证.

\noindent\textbf{例 4.2.2}\quad 证明：$f(x)$ 在 $[a,b]$ 上可积的充要条件是：对于任何一个使得 $\lambda_k\to0$ 的分划序列 $\{T_k\}$，所作的积分和
\[
\sum_{i=1}^{n_k}f(\xi_i)\Delta x_i
\]
其极限
\[
\lim_{k\to\infty}\sum_{i=1}^{n_k}f(\xi_i)\Delta x_i
\]
恒存在并且相同（不妨记为 $I$）.

\noindent\textbf{证}\quad 必要性明显，只证充分性.

若 $f(x)$ 在 $[a,b]$ 上不可积，利用式 (B) 的反面形式，即：$\exists\varepsilon_0>0$，$\forall\dfrac1k>0$，$\exists$ 分划 $T_k$ 及 $\xi_i^{(k)}\in[x_{i-1},x_i]$，虽然对应的 $\lambda_k<\dfrac1k$，但
\[
\left|I-\sum_{i=1}^{n_k}f(\xi_i^{(k)})\Delta x_i\right|\ge\varepsilon_0.
\]
如此，我们得到一个分划序列 $\{T_k\}$，虽然 $\lambda_k\to0$，但是
\[
\lim_{k\to\infty}\sum_{i=1}^{n_k}f(\xi_i^{(k)})\Delta x_i\ne I,
\]
与已知条件矛盾.

\subsection*{二、利用定理证明可积性}

\noindent\textbf{a.\quad 利用定理 2 证明可积性}

\noindent\textbf{要点}\quad 要证明 $f(x)$ 在 $[a,b]$ 上可积，按定理 2，只要证明 $f(x)$ 在 $[a,b]$ 的上、下积分相等，即 $I_0=I^0$.

\noindent\textbf{例 4.2.3（定理 1'）}\quad $f(x)$ 在 $[a,b]$ 上可积的充要条件是：$\forall\varepsilon>0$，$\exists$ 分划 $T$ 使得
\[
\sum_{i=1}^n\omega_i\Delta x_i<\varepsilon.
\]

\noindent\textbf{证}\quad 根据定理 1，必要性明显，只要证明充分性.

已知 $\forall\varepsilon>0$，$\exists$ 分划 $T$ 使得
\[
0<\overline S(T)-\underline S(T)=\sum_{i=1}^n\omega_i\Delta x_i<\varepsilon.
\]
由 Darboux 和性质（4）：
\[
\underline S(T)\le I_0\le I^0\le\overline S(T),
\]
故
\[
0\le I^0-I_0\le\overline S(T)-\underline S(T)<\varepsilon.
\]
由 $\varepsilon>0$ 的任意性，知 $I^0=I_0$，所以 $f(x)$ 在 $[a,b]$ 上可积.（请注意定理 1 和本题（定理 1'）的差别.）

\noindent\textbf{b.\quad 利用定理 1 与定理 1' 证明可积性}

\noindent\textbf{要点}\quad 利用定理 1 与定理 1'（例 4.2.3）证明可积性，关键在于证明当 $\lambda$ 充分小时 $\sum\omega_i\Delta x_i$ 能小于任意事先指定的正数 $\varepsilon$：$\sum\omega_i\Delta x_i<\varepsilon$.

\noindent\textbf{方法 I}\quad 若 $\sum_{i=1}^n\omega_i$ 有界，可以利用
\[
\sum\omega_i\Delta x_i\le\lambda\sum\omega_i.
\]
（如证明单调函数的可积性.）

\noindent\textbf{方法 II}\quad 证明 $\omega_i<\varepsilon\ (i=1,2,\cdots,n)$，从而
\[
\sum\omega_i\Delta x_i<\varepsilon\sum\Delta x_i<\varepsilon(b-a).
\]
（如证明连续函数的可积性.）

\noindent\textbf{方法 III}\quad 利用
\[
\sum\omega_i\Delta x_i=\sum'\omega_i\Delta x_i+\sum''\omega_i\Delta x_i,
\]
$\sum'$ 中，
\[
\omega_i<\frac{\varepsilon}{b-a};
\]
$\sum''$ 中，
\[
\sum''\Delta x_i<\frac{\varepsilon}{2\Omega},
\]
其中
\[
\Omega=\sup_{x\in[a,b]}f(x)-\inf_{x\in[a,b]}f(x)
\]
是 $f$ 的全振幅.

\noindent\textbf{方法 IV}\quad 利用 $\omega_i^f\le\omega_i^g$（其中 $\omega_i^f$ 与 $\omega_i^g$ 分别表示函数 $f$ 与 $g$ 在第 $i$ 个小区间上的振幅），从 $g$ 的可积性推出 $f$ 的可积性（例如：$f(x)$ 在 $[a,b]$ 上可积，用此法可证
 $|f(x)|$ 在 $[a,b]$ 上亦可积）.

下面是应用这些方法的例题.

\noindent\textbf{例 4.2.4}\quad 设 $f(x)$ 是 $[a,b]$ 上的有界变差函数（意指 $f(x)$ 在 $[a,b]$ 上的全变差），
\[
M=\sup_T\left\{\sum_{i=1}^n|f(x_i)-f(x_{i-1})|\right\}<+\infty,
\]
试证 $f(x)$ 在 $[a,b]$ 上可积.

\noindent\textbf{提示}\quad 证明对任意分划 $T$ 而言，有 $\sum_{i=1}^n\omega_i\le M$，从而应用方法 I.

\noindent\textbf{例 4.2.5}\quad 设 $f(x)$ 在 $[a,b]$ 上每一点处的极限存在并皆为零，试证：$f(x)$ 在 $[a,b]$ 上可积，且
\[
\int_a^b f(x)\,dx=0.
\]

\noindent\textbf{证}\quad 设 $x_0\in[a,b]$ 为任意一点，因 $\lim_{x\to x_0}f(x)=0$，$\forall\varepsilon_1>0$，$\exists\delta_{x_0}>0$，当 $x\in(x_0-\delta_{x_0},x_0+\delta_{x_0})$ 时，有
\[
|f(x)|<\varepsilon_1\quad(x\ne x_0).
\tag{1}
\]
如此，$\{(x_0-\delta_{x_0},x_0+\delta_{x_0})\}$ 组成了区间 $[a,b]$ 的一个开覆盖. 由有限覆盖定理，其中存在有限子覆盖 $\{(x_i-\delta_{x_i},x_i+\delta_{x_i})\}_{i=1}^k$. 至此，我们证明了：除有限个点 $x_1,x_2,\cdots,x_k$ 之外，恒有
\[
|f(x)|<\varepsilon_1\quad(x\ne x_1,x_2,\cdots,x_k).
\tag{2}
\]
至此，容易用方法 III 及定理 $1'$ 证明 $f$ 在 $[a,b]$ 上可积. 事实上，$\forall\varepsilon>0$，令 $\varepsilon_1=\dfrac{\varepsilon}{4(b-a)}$，如上有式 (2) 成立. 取 $M>\max\{|f(x_1)|,|f(x_2)|,\cdots,|f(x_k)|,\varepsilon_1\}$，作一分划 $T$，使得含 $x_1,x_2,\cdots,x_k$ 的各小区间之总长 $\sum'\Delta x_i<\dfrac{\varepsilon}{4M}$，则
\[
\sum\omega_i\Delta x_i
=\sum'\omega_i\Delta x_i+\sum''\omega_i\Delta x_i
\le 2M\frac{\varepsilon}{4M}+2\varepsilon_1\cdot(b-a)=\varepsilon.
\]
（其中 $\sum'$ 表示含 $x_1,x_2,\cdots,x_k$ 的各小区间对应项之和；$\sum''$ 是其余各项之和.）$f$ 可积性获证. 既然可积，点 $\xi_i$ 不论怎样选，积分和的极限都相同. 如上，每次只要 $\xi_i$ 选得与式 (2) 中的 $x_1,x_2,\cdots,x_k$ 不同，易证积分和的极限为零.

另外，利用“每个 $\varepsilon>0$，最多只有有限个点使得 $|f(x)|$ 不小于 $\varepsilon$”，我们也可直接证明
\[
\lim_{\lambda\to0}\sum_{i=1}^n f(\xi_i)\Delta x_i=0.
\]

下面是应用方法 IV 的例子.

\noindent\textbf{例 4.2.6}\quad 设 $f(x)$ 在 $[a,b]$ 上可微，试证 $f'(x)$ 在 $[a,b]$ 上可积的充要条件是：存在可积函数 $g(x)$ 使得
\[
f(x)=f(a)+\int_a^x g(t)\,dt.
\tag{1}
\]

\noindent\textbf{证}\quad 必要性只要令 $g(x)=f'(x)$ 即可得到，这里只证充分性（我们证明：对任一分划而言，在小区间上 $f'(x)$ 的振幅小于或等于 $g$ 的振幅：$\omega_i^{f'}\le\omega_i^g$）.

设 $T:a=x_0<x_1<\cdots<x_n=b$ 是 $[a,b]$ 的任一分划，记
\[
m_i^g=\inf_{x_{i-1}\le x\le x_i}g(x),\quad
M_i^g=\sup_{x_{i-1}\le x\le x_i}g(x),
\]
则
\[
\omega_i^g=M_i^g-m_i^g.
\]
设 $x\in[x_{i-1},x_i]$ 为任意一点，$x+\Delta x\in[x_{i-1},x_i]$，则由 (1)，
\[
\frac{\Delta f}{\Delta x}
=\frac{f(x+\Delta x)-f(x)}{\Delta x}
=\frac1{\Delta x}\int_x^{x+\Delta x}g(x)\,dx.
\]
注意到 $m_i^g\le g(x)\le M_i^g$，所以 $m_i^g\le\dfrac{\Delta f}{\Delta x}\le M_i^g$. 令 $\Delta x\to0$ 取极限，得
\[
m_i^g\le f'(x)\le M_i^g,\quad \forall x\in[x_{i-1},x_i].
\]
因此，$f'(x)$ 在 $[x_{i-1},x_i]$ 上的振幅
\[
\omega_i^{f'}=
\sup_{x_{i-1}\le x\le x_i}f'(x)
-\inf_{x_{i-1}\le x\le x_i}f'(x)
\le M_i^g-m_i^g=\omega_i^g,
\]
故
\[
0\le\sum\omega_i^{f'}\Delta x_i\le\sum\omega_i^g\Delta x_i.
\]
因 $g$ 在 $[a,b]$ 上可积，$\lim_{\lambda\to0}\sum\omega_i^g\Delta x_i=0$，故 $\lim_{\lambda\to0}\sum\omega_i^{f'}\Delta x_i=0$，因此 $f'(x)$ 在 $[a,b]$ 上可积.

\noindent\textbf{注}\quad 关于充分性，下面的证法是错误的：将式 (1) 对 $x$ 求导，得 $f'(x)=g(x)$，由 $g(x)$ 可积知 $f'(x)$ 可积. 错误是 $g(x)$ 未必连续，从而 $\left(\int_a^x g(t)\,dt\right)'=g(x)$ 未必成立.

上述方法以方法 III 应用最广，下面我们用它来证 Riemann 函数的可积性及定理 3.

\noindent\textbf{例 4.2.7}\quad 证明 Riemann 函数 $R(x)$ 在 $[0,1]$ 上可积：
\[
R(x)=
\begin{cases}
\dfrac1q, & x=\dfrac pq\ (p,q\text{ 为正整数，}\dfrac pq\text{ 为既约真分数}),\\
0, & x\text{ 为无理数}.
\end{cases}
\]

\noindent\textbf{证}\quad $\forall\varepsilon>0$，$[0,1]$ 上使得 $R(x)\ge\dfrac{\varepsilon}{2}$ 的点最多只有有限个. 事实上，要 $\dfrac1q\ge\dfrac{\varepsilon}{2}$，即要 $q\le\dfrac2\varepsilon$，因此使 $R(x)>\dfrac{\varepsilon}{2}$ 的点 $x=\dfrac pq\left(0\le p\le q<\dfrac2\varepsilon\right)$ 最多只有有限个，记为 $k$ 个. 作分划之后，令
\[
\sum\omega_i\Delta x_i=\sum'\omega_i\Delta x_i+\sum''\omega_i\Delta x_i,
\]
其中 $\sum'$ 表示含有上述那种例外点的小区间对应项的和，$\sum''\omega_i\Delta x_i$ 是其余各项的和，$\sum''\omega_i\Delta x_i$ 中每项 $\omega_i<\dfrac{\varepsilon}{2}$，因此
\[
\sum''\omega_i\Delta x_i<\frac{\varepsilon}{2}\sum''\Delta x_i\le\frac{\varepsilon}{2}\cdot1=\frac{\varepsilon}{2}.
\]
在 $\sum'$ 中，因这种项最多只有 $2k$ 个，故 $\sum'\omega_i\Delta x_i\le\sum'\Delta x_i\le2k\lambda$. 要 $2k\lambda<\dfrac{\varepsilon}{2}$，只要取 $\delta=\dfrac{\varepsilon}{4k}$，则 $\lambda<\delta$ 时，
\[
\sum\omega_i\Delta x_i=\sum'\omega_i\Delta x_i+\sum''\omega_i\Delta x_i<\frac{\varepsilon}{2}+\frac{\varepsilon}{2}=\varepsilon.
\]

\noindent\textbf{例 4.2.8（定理 3）}\quad $f(x)$ 在 $[a,b]$ 上可积的充要条件是 $\forall\varepsilon>0$，$\forall\sigma>0$，存在分划 $T$，使得振幅 $\omega_i\ge\varepsilon$ 的那些小区间 $[x_{i-1},x_i]$ 的长度之和
\[
\sum_{\omega_i\ge\varepsilon}\Delta x_i<\sigma.
\]
（通俗地说，即是振幅不能任意小的那些小区间之总长可以任意小.）

\noindent\textbf{证}\quad $1^\circ$ 必要性. 目的在于证明：$\forall\varepsilon>0$，$\forall\sigma>0$，有分划 $T$，使得 $\sum_{\omega_i\ge\varepsilon}\Delta x_i<\sigma$. 因为
\[
\varepsilon\sum_{\omega_i\ge\varepsilon}\Delta x_i
\le\sum_{\omega_i\ge\varepsilon}\omega_i\Delta x_i
\le\sum_{i=1}^n\omega_i\Delta x_i,
\]
可见只要 $\sum_{i=1}^n\omega_i\Delta x_i<\varepsilon\sigma$，与上式联立即得 $\varepsilon\sum_{\omega_i\ge\varepsilon}\Delta x_i<\varepsilon\sigma$，从而 $\sum_{\omega_i\ge\varepsilon}\Delta x_i<\sigma$.

因为已知 $f(x)$ 在 $[a,b]$ 上可积，所以，$\forall\varepsilon>0$，$\forall\sigma>0$，存在 $T$（分划）使得
\[
\sum_{i=1}^n\omega_i\Delta x_i<\varepsilon\sigma.
\]
这样必要性获证.

$2^\circ$ 充分性. 要证 $f(x)$ 在 $[a,b]$ 上可积，即：$\forall\varepsilon>0$，要找分划 $T$，使得
\[
\sum_{i=1}^n\omega_i\Delta x_i<\varepsilon.
\tag{1}
\]
但已知：$\forall\varepsilon_1>0$，$\forall\sigma>0$，$\exists T$ 使得 $\sum_{\omega_i\ge\varepsilon_1}\Delta x_i<\sigma$，从而
\[
\begin{aligned}
\sum_{i=1}^n\omega_i\Delta x_i
&=\sum_{\omega_i\ge\varepsilon_1}\omega_i\Delta x_i
+\sum_{\omega_i<\varepsilon_1}\omega_i\Delta x_i\\
&\le \Omega\sum_{\omega_i\ge\varepsilon_1}\Delta x_i
+\varepsilon_1\sum_{\omega_i<\varepsilon_1}\Delta x_i\\
&\le\Omega\cdot\sigma+\varepsilon_1(b-a),
\end{aligned}
\tag{2}
\]
其中
\[
\Omega=M-m=\sup_{a\le x\le b}f(x)-\inf_{a\le x\le b}f(x)
\]
为 $f(x)$ 在 $[a,b]$ 上的全振幅（因此 $\omega_i\le\Omega$）.

可见要式 (1) 成立，只要事先取 $\varepsilon_1=\dfrac{\varepsilon}{2(b-a)}$，$\sigma=\dfrac{\varepsilon}{2\Omega}$，则式 (2) 表明
\[
\sum_{i=1}^n\omega_i\Delta x_i\le\Omega\cdot\sigma+\varepsilon_1(b-a)<\varepsilon.
\]
故这时的 $T$ 即为所求.

\noindent\textbf{c.\quad 利用定理 3 证明可积性}

\noindent\textbf{要点}\quad 要证一函数 $f(x)$ 在 $[a,b]$ 上可积，此法的关键在于，对任意的 $\varepsilon>0$ 和 $\sigma>0$，找一分划 $T$，使得
\[
\sum_{\omega_i\ge\varepsilon}\Delta x_i<\sigma.
\]

\noindent\textbf{例 4.2.9}\quad 设 $y=f(u)$ 在 $[A,B]$ 上连续，$u=\varphi(x)$ 在 $[a,b]$ 上可积；当 $x\in[a,b]$ 时，$A\le\varphi(x)\le B$. 试证 $F(x)=f(\varphi(x))$ 在 $[a,b]$ 上可积.

\noindent\textbf{证}\quad 因为 $f(u)$ 在 $[A,B]$ 上连续，所以在 $[A,B]$ 上一致连续：$\forall\varepsilon>0$，$\exists\delta>0$，当 $u',u''\in[A,B]$，$|u'-u''|<\delta$ 时，有
\[
|f(u')-f(u'')|<\frac{\varepsilon}{2}.
\tag{1}
\]
因此作分划以后，在 $[x_{i-1},x_i]$ 上，若 $\varphi(x)$ 的振幅 $\omega_i^\varphi<\delta$，则 $F(x)=f(\varphi(x))$ 的振幅 $\omega_i^F<\varepsilon$.（事实上，这时 $\forall x',x''\in[x_{i-1},x_i]$，记 $u'=\varphi(x')$，$u''=\varphi(x'')$，则
\[
|u'-u''|=|\varphi(x')-\varphi(x'')|\le\omega_i^\varphi<\delta,
\]
从而
\[
|F(x')-F(x'')|=|f(u')-f(u'')|<\frac{\varepsilon}{2},
\]
故
\[
\omega_i^F=\sup_{x_{i-1}\le x',x''\le x_i}|F(x')-F(x'')|<\frac{\varepsilon}{2}<\varepsilon.)
\]
由此可见，在 $[x_{i-1},x_i]$ 上若 $\omega_i^F\ge\varepsilon$，必有 $\omega_i^\varphi\ge\delta$. 故
\[
\sum_{\omega_i^F\ge\varepsilon}\Delta x_i
\le
\sum_{\omega_i^\varphi\ge\delta}\Delta x_i.
\tag{2}
\]
如此，$\forall\varepsilon>0$，$\forall\sigma>0$，首先按式 (1) 找出 $\delta>0$. 再由 $\varphi(x)$ 在 $[a,b]$ 上可积，利用定理 3 的必要性，对 $\delta>0$ 与 $\sigma>0$，存在分划 $T$ 使得 $\sum_{\omega_i^\varphi\ge\delta}\Delta x_i<\sigma$. 于是由式 (2) 知
\[
\sum_{\omega_i^F\ge\varepsilon}\Delta x_i\le\sum_{\omega_i^\varphi\ge\delta}\Delta x_i<\sigma.
\]
根据定理 3 的充分性，这便证明了 $F(x)$ 在 $[a,b]$ 上可积.

最后，我们讨论可积性与连续的关系.

\noindent\textbf{例 4.2.10}\quad 证明：若 $f(x)$ 在 $[a,b]$ 上的不连续点可以用有限个总长任意小的区间所覆盖，则 $f(x)$ 在 $[a,b]$ 上可积.

\noindent\textbf{提示}\quad 利用定理 3，或定理 $1'$ 容易证明.

\noindent\textbf{例 4.2.11}\quad 若 $f(x)$ 在 $[a,b]$ 上可积，则 $f(x)$ 的连续点在 $[a,b]$ 上处处稠密.

\noindent\textbf{证}\quad 问题归于证明 $f(x)$ 在 $[a,b]$ 上至少有一个连续点. 事实上，若能如此，则 $\forall[\alpha,\beta]\subset[a,b]$，因为 $f(x)$ 在 $[\alpha,\beta]$ 上可积，故 $f(x)$ 在 $[\alpha,\beta]$ 上有连续点，这就证明了连续点处处稠密.

为了证明 $f(x)$ 在 $[a,b]$ 上至少有一个连续点，我们采用区间套的办法：

因为 $f(x)$ 在 $[a,b]$ 上可积，所以 $\lim_{\lambda\to0}\sum\omega_i\Delta x_i=0$，对 $\varepsilon_1=\dfrac12$，$\exists$ 分划 $T_1$，使得
\[
\sum\omega_i\Delta x_i<\varepsilon_1(b-a).
\tag{1}
\]
如此，至少存在一个小区间 $[x_{i-1},x_i]$，使得其上 $f(x)$ 的振幅 $\omega_i<\varepsilon$. 因为不然的话，$\sum\omega_i\Delta x_i\ge\varepsilon_1\sum\Delta x_i\ge\varepsilon_1(b-a)$，与式 (1) 矛盾. 将此小区间适当收缩，总可以使得它的长度 $x_i-x_{i-1}<\dfrac12(b-a)$，使它的两端点在 $(a,b)$ 内. 记缩小后的区间为 $[a_1,b_1]$，则 $a<a_1<b_1<b$，$b_1-a_1<\dfrac12(b-a)$，$f(x)$ 在 $[a_1,b_1]$ 上的振幅
\[
\omega^f[a_1,b_1]<\varepsilon_1=\frac12.
\]
用 $[a_1,b_1]$ 取代上面的 $[a,b]$，类似推理可知对 $\varepsilon_2=\dfrac1{2^2}$，存在 $[a_2,b_2]\subset[a_1,b_1]$，$a_1<a_2<b_2<b_1$，$b_2-a_2<\dfrac12(b_1-a_1)<\dfrac14(b-a)$，$f(x)$ 在 $[a_2,b_2]$ 上的振幅
\[
\omega^f[a_2,b_2]<\varepsilon_2=\frac1{2^2}.
\]
如此无限做下去，我们可得一区间套：
\[
[a,b]\supset[a_1,b_1]\supset\cdots\supset[a_n,b_n]\supset\cdots,
\]
\[
0<b_n-a_n<\frac1{2^n}(b-a)\to0\quad(n\to\infty),
\]
且 $a_n<a_{n+1}<b_{n+1}<b_n$，$f(x)$ 在 $[a_n,b_n]$ 上的振幅
\[
\omega^f[a_n,b_n]<\varepsilon_n=\frac1{2^n}.
\]
根据区间套定理，$\exists\xi\in[a_n,b_n]\ (n=1,2,\cdots)$. 因为 $a_n$ 严↗ $\xi$，$b_n$ 严↘ $\xi$，所以
\[
a_n<\xi<b_n\quad(n=1,2,\cdots).
\]
现在容易看出，$f(x)$ 在点 $\xi$ 连续. 事实上，$\forall\varepsilon>0$，可取 $n$ 使得 $\dfrac1{2^n}<\varepsilon$. 从而令 $\delta=\min\{b_n-\xi,\xi-a_n\}$，则当 $|x-\xi|<\delta$ 时，$x\in[a_n,b_n]$，从而
\[
|f(x)-f(\xi)|<\omega^f[a_n,b_n]<\frac1{2^n}<\varepsilon.
\]
即表明 $f(x)$ 在 $\xi$ 处连续.

\noindent\textbf{☆例 4.2.12}\quad 证明：若 $f(x)\ge0$ 在 $[a,b]$ 上有定义并且可积，则等式
\[
\int_a^b f(x)\,dx=0
\]
成立的充要条件是 $f(x)$ 在连续点上恒为零.

\noindent\textbf{提示}\quad （必要性）若 $x_0\in[a,b]$ 为 $f$ 的连续点，$f(x_0)>0$，则 $\exists\delta>0$ 使得
\[
\int_a^b f(x)\,dx\ge\int_{x_0-\delta}^{x_0+\delta}f(x)\,dx>0.
\]
（充分性）可利用积分定义与上例结果得到.

\noindent\textbf{注}\quad 由该例可得如下明显事实：

1）连续函数 $f(x)\ge0$，若在 $[a,b]$ 上积分为 0，则 $f(x)\equiv0$（$\forall x\in[a,b]$）.

2）假设 $f(x)$ 和 $g(x)$ 在 $[a,b]$ 上连续，$f(x)\ge g(x)$.

若 $\exists x_0\in[a,b]$ 使得 $f(x_0)>g(x_0)$，则 $\int_a^b f(x)\,dx>\int_a^b g(x)\,dx$；

反之，若 $\int_a^b f(x)\,dx=\int_a^b g(x)\,dx$，则 $f(x)\equiv g(x)$.

（以上两点常被应用，默认为基本常识.）

\noindent\textbf{*例 4.2.13}\quad 设 $\{f_n(x)\}$ 是 $(-\infty,+\infty)$ 上的连续函数序列，一致有界. 若对任意区间 $[a,b]\subset(-\infty,+\infty)$，恒有 $\lim_{n\to\infty}\int_a^b|f_n(x)|\,dx=0$，则在 $(-\infty,+\infty)$ 上，对内闭绝对可积的函数 $h(x)$，恒有
\[
\lim_{n\to\infty}\int_a^b f_n(x)h(x)\,dx=0.
\]
（北京大学）

（注\quad $|f_n(x)|$ 一致有界，意指：$\exists M>0$，使得：$\forall f_n$ 和 $\forall x\in(-\infty,+\infty)$，恒有 $|f_n(x)|\le M$.）

\noindent\textbf{证}\quad 对 $[a,b]$ 任作分划：$a=a_0<a_1<\cdots<a_n=b$. $M_i^h,m_i^h$ 分别表示 $h(x)$ 在 $[a_{i-1},a_i]$ 上的上、下确界，$M^h\stackrel{\text{记}}{=}\max_{1\le i\le n}\{|M_i^h|,|m_i^h|\}$，则
\[
\begin{aligned}
\left|\int_a^b f_n(x)h(x)\,dx\right|
&\le \sum_{i=0}^{n-1}\int_{a_i}^{a_{i+1}}|f_n(x)|\,|h(x)-m_i^h|\,dx\\
&\quad+\sum_{i=0}^{n-1}\int_{a_i}^{a_{i+1}}|f_n(x)m_i^h|\,dx\\
&\le M\sum_{i=0}^{n-1}\omega_i^h(a_{i+1}-a_i)
+M^h\int_a^b|f_n(x)|\,dx,
\end{aligned}
\tag{1}
\]
其中 $\omega_i^h\stackrel{\text{记}}{=}M_i^h-m_i^h$. $\forall\varepsilon>0$，$h(x)$ 可积，（利用可积的充要条件）$\exists N_1>0$，当 $n>N_1$ 时，
\[
M\sum_{i=0}^{n-1}\omega_i^h(a_{i+1}-a_i)<\frac{\varepsilon}{2}.
\]
因 $\lim_{n\to\infty}\int_a^b|f_n(x)|\,dx=0$，$\exists N_2>0$，当 $n>N_2$ 时，$\int_a^b|f_n(x)|\,dx<\dfrac{\varepsilon}{2M^h}$.

故当 $n>\max\{N_1,N_2\}$ 时，由式 (1) 得
\[
\left|\int_a^b f_n(x)h(x)\,dx\right|\le\frac{\varepsilon}{2}+\frac{\varepsilon}{2}=\varepsilon.
\]

\noindent\textbf{*例 4.2.14}\quad 若函数 $f(x)$ 在区间 $[a,b]$ 上有界，给出并证明 $f(x)$ 在 $[a,b]$ 上（Riemann）积分和的极限 $\lim_{\lambda(\Delta)\to0}\sum f(\xi_i)(x_i-x_{i-1})$ 收敛的 Cauchy 准则.（北京大学）

\noindent\textbf{提示}\quad $\lim_{\lambda(\Delta)\to0}\sum f(\xi_i)(x_i-x_{i-1})$ 收敛的充要条件是：$\forall\varepsilon>0$，$\exists\delta>0$，对于 $[a,b]$ 的任意两个分划 $T_1,T_2$，以及 $\xi_i$ 的任意两组取法，只要它们最大、最小区间的长度 $\lambda_1,\lambda_2<\delta$，则它们对应的积分和之差
\[
\|S(\xi_1,T_1)-S(\xi_2,T_2)\|<\varepsilon.
\tag{1}
\]
证必要性：可利用三角不等式；证充分性：可利用本节“定理一”“可积的充要条件”.

\noindent\textbf{再提示}\quad （必要性）当 $\lambda\to0$ 时，$S(\xi,T)\stackrel{\text{记}}{=}\sum f(\xi_i)\Delta x_i\to I$. 则 $\forall\varepsilon>0$，$\exists\delta>0$，当 $\lambda<\delta$ 时，有
\[
|S(\xi,T)-I|<\frac{\varepsilon}{2}.
\]
现设 $S(\xi_1,T_1)$ 和 $S(\xi_2,T_2)$ 是任意两不同的积分和；对应的 $\lambda_1,\lambda_2<\delta$，那么由
\[
|S(\xi_1,T_1)-I|<\frac{\varepsilon}{2}\quad\text{和}\quad |S(\xi_2,T_2)-I|<\frac{\varepsilon}{2},
\]
可得
\[
|S(\xi_1,T_1)-S(\xi_2,T_2)|
\le |S(\xi_1,T_1)-I|+|I-S(\xi_2,T_2)|<\frac{\varepsilon}{2}+\frac{\varepsilon}{2}=\varepsilon.
\]

（充分性）若条件 (1) 成立，亦即：$\forall\varepsilon>0$，$\exists\delta>0$，当 $\lambda_1,\lambda_2<\delta$ 时恒有
\[
|S(\xi_1,T_1)-S(\xi_2,T_2)|<\frac{\varepsilon}{2},
\tag{2}
\]
那么对应任意分划 $T:a=x_0<x_1<\cdots<x_n=b$，当 $\lambda=\max(x_i-x_{i-1})<\delta$ 时，有
\[
\left|\sum_{i=1}^n[f(\xi_{1i})\Delta x_i-\sum_{i=1}^n f(\xi_{2i})\Delta x_i]\right|
=\left|\sum_{i=1}^n[f(\xi_{1i})-f(\xi_{2i})]\Delta x_i\right|<\frac{\varepsilon}{2}.
\]
取上确界
\[
\omega_i=\sup\{|f(\xi_{1i})-f(\xi_{2i})|\mid \xi_{1i},\xi_{2i}\in[x_{i-1},x_i]\},
\]
则由上式知
\[
\sum_{i=1}^n\omega_i\Delta x_i\le\frac{\varepsilon}{2}<\varepsilon.
\]
根据定理 1 的充要条件，知：$f(x)$ 在 $[a,b]$ 上可积，$\lim_{\lambda(\Delta)\to0}\sum f(\xi_i)\Delta x_i$ 收敛.

\noindent\textbf{练习}\quad 设 $f(x)$ 在 $[0,a]$ 上二次可导，$f(0)=f'(0)=f'(a)=0$，$f(a)=1$，且 $\forall x\in[0,a]$ 有 $f''(x)\le1$；
\[
g(x)=
\begin{cases}
x, & 0\le x\le\dfrac a2,\\
a-x, & \dfrac a2<x\le a.
\end{cases}
\]
试证：

1）$[0,a]$ 上，$f'(x)\le g(x)$；

2）$\exists x_0\in[0,a]$，使得 $f'(x_0)<g(x_0)$；

3）$a>2$.（南开大学）

\noindent\textbf{证}\quad 1）$x\in\left[0,\dfrac a2\right]$（应用 Lagrange 定理）：$f'(x)=f'(0)+f''(\xi)x\le x=g(x)$（$0<\xi<x$）.

同理，$x\in\left[\dfrac a2,a\right]$：$f'(x)=f'(a)+f''(\eta)(a-x)\le(a-x)=g(x)$（$x<\eta<a$）.

2）若结论不成立，则 $f'(x)\equiv g(x)$，推得 $f''\left(\dfrac a2\right)$ 不存在，矛盾.

3）根据 1）和 2），$f'(x)$ 和 $g(x)$ 符合例 4.2.12 注 2）的条件，因此
\[
1=f(x)\bigg|_0^a=\int_0^a f'(x)\,dx<\int_0^a g(x)\,dx
=\frac{x^2}{2}\bigg|_0^{a/2}-\frac{(a-x)^2}{2}\bigg|_{x=a/2}^a
=\frac{a^2}{4}.
\]
因 $a$ 为正，由此即得：$a>2$.（证毕.）

\noindent\textbf{※附：积分定义的简化}

\noindent\textbf{引言}\quad 众所周知，两三角形，边、角分别相等，共有 6 个等式，但判断两三角形全等，一般只需 3 个即可（工作量减少一半）. 定积分的定义有两个任意性：“分划 $T$ 任意”“$\xi_i$ 的选法任意”. 能否只用一个任意性？回答是肯定的.

\noindent\textbf{积分定义简化之一：}

\[
\text{取消分划的任意性（只用“等分法”），保留 $\xi_i$ 选法的任意性.}
\]

\noindent\textbf{※例 4.2.15（积分的等价定义（一））}\quad 如下三条等价：

（i）（按原定义）$f(x)$ 在 $[a,b]$ 上可积，且 $\int_a^b f(x)\,dx=I$；（含两个任意性）

（ii）（按等分定义）将 $[a,b]$ $n$ 等分，$d_i=a+i\dfrac{b-a}{n}\ (i=1,2,\cdots,n)$，任取 $\xi_i\in[d_{i-1},d_i]$，作积分和：$\sum_{i=1}^n f(\xi_i)\dfrac{b-a}{n}$，则
\[
\lim_{n\to\infty}\sum_{i=1}^n f(\xi_i)\frac{b-a}{n}\stackrel{\text{存在}}{=}I
\]
（其中 $I$ 的大小与 $\xi_i$ 的选取无关）（只含一个任意性）；

（iii）（可积的充要条件）将 $[a,b]$ $n$ 等分，用 $M_i$ 和 $m_i$ 分别表示 $f(x)$ 在第 $i$ 个小区间上的“上、下确界”，则当 $n\to\infty$ 时，相应的上、下 Darboux 和：$\sum_{i=1}^n M_i\dfrac{b-a}{n}$，$\sum_{i=1}^n m_i\dfrac{b-a}{n}$ 有有限的极限，且都等于 $I$. 即
\[
\lim_{n\to\infty}\sum_{i=1}^nM_i\frac{b-a}{n}
=\lim_{n\to\infty}\sum_{i=1}^n m_i\frac{b-a}{n}=I.
\]
记 $\omega_i=M_i-m_i$，则有 $\lim_{n\to\infty}\sum_{i=1}^n\omega_i\dfrac{b-a}{n}=\lim_{n\to\infty}\sum_{i=1}^n(M_i-m_i)\dfrac{b-a}{n}=0$.

\noindent\textbf{证}\quad （证 (i)$\Rightarrow$(ii)$\Rightarrow$(iii)$\Rightarrow$(i)）

$1^\circ$（(i)$\Rightarrow$(ii)）无须证明，因对任意分划成立，等分当然也成立.

$2^\circ$（(ii)$\Rightarrow$(iii)）已知 $\lim_{n\to\infty}\sum_{i=1}^n f(\xi_i)\dfrac{b-a}{n}\stackrel{\text{存在}}{=}I$，因此，$\forall\varepsilon>0$，$\exists N>0$，当 $n>N$ 时，有
\[
\left|\sum_{i=1}^n f(\xi_i)\frac{b-a}{n}-I\right|<\frac{\varepsilon}{2}.
\]
此时 $f(x)$ 在 $[a,b]$ 上有界（不可能无界，否则 $|f(\xi_i)|$ 能任意取大，极限不可能存在），因此小区间上的上、下确界 $M_i$ 和 $m_i$ 都应是有界数. 利用确界定义，在第 $i$ 个小区间中存在 $\xi_i$，使得 $0\le M_i-f(\xi_i)<\dfrac{\varepsilon}{2(b-a)}\ (i=1,2,\cdots,n)$. 于是
\[
\left|\sum_{i=1}^n M_i\frac{b-a}{n}-I\right|
\le \frac{b-a}{n}\left|\sum_{i=1}^n(M_i-f(\xi_i))\right|
+\left|\sum_{i=1}^n f(\xi_i)\frac{b-a}{n}-I\right|
<\frac{\varepsilon}{2}+\frac{\varepsilon}{2}=\varepsilon,
\]
表明
\[
\lim_{n\to\infty}\sum_{i=1}^n M_i\frac{b-a}{n}=I.
\]
同理可证
\[
\lim_{n\to\infty}\sum_{i=1}^n m_i\frac{b-a}{n}=I.
\]
进而有
\[
\lim_{n\to\infty}\sum_{i=1}^n\omega_i\frac{b-a}{n}
=\lim_{n\to\infty}\sum_{i=1}^n(M_i-m_i)\frac{b-a}{n}=0.
\]

$3^\circ$（(iii)$\Rightarrow$(i)）要证明 $f(x)$ 在 $[a,b]$ 上可积.（利用本节定理 1）只需证明：任意分划
\[
T:a=x_0<x_1<\cdots<x_k=b.
\tag{1}
\]
当 $\lambda=\max_{i=1,\cdots,k}|x_i-x_{i-1}|\to0$ 时，有
\[
\lim_{\lambda\to0}\sum_{i=1}^k\omega_i(x_i-x_{i-1})=0,
\]
其中 $\omega_i$ 是函数 $f(x)$ 在 $[x_{i-1},x_i]$ 上的振幅，$\omega_i=M_i-m_i$（这里的 $M_i,m_i$ 分别表示 $f(x)$ 在 $[x_{i-1},x_i]$ 上的上、下确界）.

现在条件 (iii) 成立，说明 $f(x)$ 有界，即 $\exists M>0$，使得 $|f(x)|\le M\ (\forall x\in[a,b])$. 因此
\[
0\le\omega_i=M_i-m_i\le |M_i|+|m_i|\le2M.
\tag{2}
\]
在 $[a,b]$ 内部插入 $n-1$ 个“等分点”，将 $[a,b]$ 进行 $n$ 等分：
\[
T_{n\text{等分}}:a=d_0<d_1<d_2<\cdots<d_{n-1}<b,\quad d_{i+1}-d_i=\frac{b-a}{n}.
\]
根据条件 (iii)：$\forall\varepsilon>0$，$\exists N>0$，当 $n>N$ 时，
\[
\sum_{i=1}^n\omega_i'\frac{b-a}{n}<\frac{\varepsilon}{2}
\]
（$\omega_i'$ 表示 $f$ 在 $[d_{i-1},d_i]$ 上的振幅）.

现将 $n$ 固定，对 $[a,b]$ 作任意分划 $T$（充分细分），使得
\[
\lambda_T<\min\left\{\frac{b-a}{n},\frac{\varepsilon}{4(n-1)M}\right\}.
\]
于是 $T$ 的每个小区间内部最多只含 $T_{n\text{等分}}$ 的一个等分点（$d_i$）（否则 $\lambda_T\ge\dfrac{b-a}{n}$，矛盾）. 如此，$T$ 的小区间分为两类：第一类是内部含有 $T_{n\text{等分}}$ 的某一等分点的小区间；第二类是内部不含等分点的小区间，整个小区间落入某个 $[d_{i-1},d_i]$.

因为 $T_{n\text{等分}}$ 只有 $n-1$ 个分点，$T$ 的第一类小区间最多只有 $n-1$ 个；其余（第二类）小区间分别属于 $T_{n\text{等分}}$ 的某个 $[d_{i-1},d_i]$，将同属 $[d_{i-1},d_i]$ 的小区间归为第 $i$ 组. 因此第 $i$ 组小区间上的振幅 $\omega_i\le\omega_i'$（因为振幅随范围增大而增大），因此
\[
\sum_{\text{第 }i\text{组}}\omega_i(x_i-x_{i-1})
\le\sum_{\text{第 }i\text{组}}\omega_i'(x_i-x_{i-1})
=\omega_i'\sum_{\text{第 }i\text{组}}(x_i-x_{i-1})
\le\omega_i'(d_i-d_{i-1})=\omega_i'\frac{b-a}{n}.
\]
故
\[
\sum_{\text{第二类}}\omega_i(x_i-x_{i-1})
=\sum_{i=1}^n\left[\sum_{\text{第 }i\text{组}}\omega_i(x_i-x_{i-1})\right]
\le\sum_{i=1}^n\omega_i'\frac{b-a}{n}<\frac{\varepsilon}{2}.
\]
另外，（因 $f$ 有界，$\omega_i\le2M$）第一类小区间 $\omega_i$ 对应和：
\[
\sum_{\text{第一类}}\omega_i(x_i-x_{i-1})
\le2M(n-1)\lambda<2M(n-1)\frac{\varepsilon}{4(n-1)M}=\frac{\varepsilon}{2}.
\]
因此，对分划 $T$ 有
\[
\sum_{i=1}^n\omega_i(x_i-x_{i-1})
=\sum_{\text{第一类}}\omega_i(x_i-x_{i-1})
+\sum_{\text{第二类}}\omega_i(x_i-x_{i-1})
<\frac{\varepsilon}{2}+\frac{\varepsilon}{2}=\varepsilon.
\]
根据本节定理 1，$f(x)$ 在 $[a,b]$ 上可积. 如此，
\[
\int_a^b f(x)\,dx=\lim_{\lambda\to0}\sum_{i=1}^n f(\xi_i)\Delta x_i
=\lim_{n\to\infty}\sum_{i=1}^n f(\xi_i)\frac{b-a}{n}=I.
\]

\noindent\textbf{积分定义简化之二：}

\[
\text{取消 $\xi_i$ 选法的任意性（$\xi_i$ 只取端点），保留分划的任意性.}
\]

\noindent\textbf{※例 4.2.16（积分的等价定义（二））}\quad 函数 $f(x)$ 在 $[a,b]$ 上积分存在且等于 $I$ 的充要条件是：在 $[a,b]$ 上，任取分划
\[
T:a=x_0<x_1<\cdots<x_n=b,
\tag{1}
\]
作积分和
\[
\sum_{i=1}^n f(\xi_i)(x_i-x_{i-1}),
\tag{2}
\]
这里 $\xi_i$ 只需选端点（$\xi_i=x_i$ 或 $x_{i-1}$）. 记 $\lambda=\max_{i=1,\cdots,n}|x_i-x_{i-1}|$，则极限 $\lim_{\lambda\to0}\sum_{i=1}^n f(\xi_i)(x_i-x_{i-1})\stackrel{\text{存在}}{=}I$.

更精细地说，即 $\forall\varepsilon>0$，$\exists\delta>0$，当 $\lambda<\delta$ 时，有
\[
\left|\sum_{i=1}^n f(\xi_i)(x_i-x_{i-1})-I\right|<\varepsilon.
\tag{3}
\]

\noindent\textbf{注}\quad 可见新定义简化许多，必要性明显成立，下面只证充分性.

\noindent\textbf{分析}\quad （充分性）要证明 $f(x)$ 在 $[a,b]$ 上可积且积分等于 $I$.

根据积分的等价定义（一），只需证明：$\forall\varepsilon>0$，$\exists N>0$，当 $n>N$ 时，将 $[a,b]$ $n$ 等分，作分划：
\[
T_{n\text{等分}}:a=d_0<d_1<d_2<\cdots<d_{n-1}<d_n=b,\quad d_{i+1}-d_i=\frac{b-a}{n},
\]
则有
\[
\left|\sum_{i=1}^n f(\xi_i)\frac{b-a}{n}-I\right|<\varepsilon\quad(\forall\xi_i\in[d_i,d_{i+1}]).
\tag{4}
\]
（如果将等价定义（一）的条件简称为“等分 + 任选”，定义（二）的条件简称为“任分 + $\xi_i$ 取端点”，那么现在的任务就是由后者导出前者.）

已知：$\forall\varepsilon>0$，能找到 $\delta>0$，当 $\lambda<\delta$ 时式 (3) 成立. 因此只要取 $N>\dfrac{b-a}{\delta}$，则当 $n>N$ 时，就有
\[
\lambda_{\text{等分}}=\max_{i=1,\cdots,n}|d_{i+1}-d_i|=\frac{b-a}{n}<\frac{b-a}{N}<\delta.
\]
假若其中的 $\xi_i$ 刚好都取在端点（即 $\xi_i=d_i$ 或 $d_{i+1}$），那么此时式 (3) 直接导出式 (4)，问题自明. 但 $\xi_i\in[d_i,d_{i+1}]$ 是任取的，$\xi_i$ 很可能不在端点上.

如果 $\xi_i$ 不是端点（$d_i<\xi_i<d_{i+1}$），也有办法：可将 $\xi_i$ 作为新分点，将 $[d_i,d_{i+1}]$ 分为左、右两半：$[d_i,\xi_i]$ 和 $[\xi_i,d_{i+1}]$（凡不是端点的 $\xi_i$，统统都这样处理），如此得一新分划，记作 $T_{\text{新}}$；

i）对 $T_{\text{新}}$ 而言，每个 $\xi_i$ 都取端点；

ii）增加分点意味着分划更细密，故 $\lambda_{\text{新}}<\lambda_{\text{等分}}=\dfrac{b-a}{n}<\delta$；

iii）注意，当一小区间一剖为二时，两个小区间之和就是原来的小区间，且
\[
f(\xi_i)(\xi_i-d_i)+f(\xi_i)(d_{i+1}-\xi_i)=f(\xi_i)(d_{i+1}-d_i).
\]
故 $T_{n\text{等分}}$ 的积分和里，若用 $\sum'$ 表示 $\xi_i$ 原本选于端点的那些项之和，$\sum''$ 表示原本 $\xi_i$ 不是端点（需要一剖为二的）的那些项之和，那么式 (4) 中的积分和：
\[
\sum_{i=1}^n f(\xi_i)\frac{b-a}{n}
=\sum' f(\xi_i)\frac{b-a}{n}
+\sum''[f(\xi_i)(\xi_i-d_i)+f(\xi_i)(d_{i+1}-\xi_i)].
\tag{5}
\]
右端是新分划的积分和（$\lambda_{\text{新}}<\delta$，且 $\xi_i$ 都在端点），符合式 (3) 的要求. 因此
\[
\left|\left[\sum' f(\xi_i)\frac{b-a}{n}
+\sum'' f(\xi_i)(\xi_i-d_i)+f(\xi_i)(d_{i+1}-\xi_i)\right]-I\right|<\varepsilon.
\tag{6}
\]
式 (5) 左端的积分和 $\sum_{i=1}^n f(\xi_i)\dfrac{b-a}{n}$ 是分划 $T_{n\text{等分}}$（将 $[a,b]$ $n$ 等分）的积分和，其中 $\xi_i\in[d_i,d_{i+1}]$ 是任意的，由式 (5) 和 (6) 知
\[
\left|\sum_{i=1}^n f(\xi_i)\frac{b-a}{n}-I\right|<\varepsilon.
\tag{7}
\]
（此式说明在“等分 + 任选”条件下右式 (7) 成立.）

根据等价定义（一）的充分性，知 $f(x)$ 在 $[a,b]$ 上可积，且积分等于 $I$. 证毕.

\noindent\textbf{注}\quad 1）这里是用等价定义（一）的充分性来证明的，实际上也可直接用积分定义来证明（留作练习）.

2）原积分定义有两个任意性，判断可积性比较困难，采用新定义，确实方便很多；但是，假若已知函数可积，那么两个任意性反而会给我们带来很多好处. 巧妙地作分划和选 $\xi_i$，常可得到极好的结果（教材里有不少例子）.

另外，本节定理 1（可积 $\Leftrightarrow\lim_{\lambda\to0}\sum_{i=1}^n\omega_i(x_i-x_{i-1})=0$）可化解 $\xi_i$ 带来的困难，也许这正是教材不提“简化”的原因. 值得注意的是：例 4.2.15 里的等价形式 (iii)，实际上两个任意性都已被简化.

3）这里介绍“简化”的主要目的是扩大眼界，拓展思路，讨论方法. 由于教材并未涉及“简化”的问题，建议读者（尤其是考生）谨慎引用，不然说你“擅自降低难度，不得也罢！”

\noindent\textbf{单元练习 4.2}

\noindent\textbf{4.2.1}\quad 设函数 $f(u)$ 在区间 $[A,B]$ 上连续，$g(x)$ 在 $[a,b]$ 上可积，当 $x\in[a,b]$ 时，$A\le g(x)\le B$. 试用各种不同的方法证明 $f(g(x))$ 在 $[a,b]$ 上可积.

\noindent\textbf{证}\quad $1^\circ$ 因 $f(u)$ 连续，故在 $[A,B]$ 上有界，即 $\exists M>0$ 使得
\[
|f(u)|<M\quad(\forall u\in[A,B]).
\tag{1}
\]
$2^\circ$ Cantor 定理，它还在 $[A,B]$ 上一致连续：$\forall\varepsilon>0$，$\exists\delta>0$，使得 $\forall u_1,u_2\in[A,B]$，当 $|u_1-u_2|<\delta$ 时，有
\[
|f(u_1)-f(u_2)|<\frac{\varepsilon}{2(b-a)}.
\tag{2}
\]
$3^\circ$ $u=g(x)$ 在 $[a,b]$ 上可积，利用例 4.2.8，知：当 $[a,b]$ 的分划充分细密时，使得 $\omega_i^g\ge\delta$ 的小区间的总长能任意小（小于事先任意指定的正数），例如小于 $\dfrac{\varepsilon}{4M}$，即
\[
\sum_{\omega_i^g\ge\delta}\Delta x_i<\frac{\varepsilon}{4M}.
\tag{3}
\]
其余的 $\omega_i^g<\delta$，从而由式 (2) 知
\[
\omega_i^{f\circ g}<\frac{\varepsilon}{2(b-a)}.
\tag{4}
\]
其中 $\omega_i^{f\circ g}$ 表示 $f(g(x))$ 在第 $i$ 个小区间上的振幅.

$4^\circ$ 故
\[
\begin{aligned}
\sum\omega_i^{f\circ g}\Delta x_i
&=\sum_{\omega_i^g<\delta}\omega_i^{f\circ g}\Delta x_i
+\sum_{\omega_i^g\ge\delta}\omega_i^{f\circ g}\Delta x_i\\
&\le \frac{\varepsilon}{2(b-a)}\sum_{\omega_i^g<\delta}\Delta x_i
+2M\sum_{\omega_i^g\ge\delta}\Delta x_i\\
&\le \frac{\varepsilon}{2(b-a)}(b-a)+2M\frac{\varepsilon}{4M}
=\frac{\varepsilon}{2}+\frac{\varepsilon}{2}=\varepsilon.
\end{aligned}
\]
此式表明 $f(g(x))$ 在 $[a,b]$ 上可积.（证毕.）

\noindent\textbf{4.2.2}\quad 试用多种方法证明 $f(x)$ 在 $[0,1]$ 上可积，设
\[
1)\ f(x)=\operatorname{sgn}\left(\sin\frac{\pi}{x}\right);\qquad
2)\ f(x)=
\begin{cases}
\dfrac1x-\left[\dfrac1x\right], & x\ne0\text{ 时},\\
0, & x=0\text{ 时}.
\end{cases}
\]

\noindent\textbf{提示}\quad 利用例 4.2.10.

\noindent\textbf{再提示}\quad 两小题的函数都只有一串间断点 $\left\{\dfrac1n\right\}_{n=1}^\infty$，$x=0$ 是该数列唯一聚点. $\forall\varepsilon>0$，$\exists N>0$，当 $n>N$ 时，有 $\left|\dfrac1n-0\right|<\dfrac{\varepsilon}{4}$，即：区间 $\left(-\dfrac{\varepsilon}{4},\dfrac{\varepsilon}{4}\right)$ 包含了聚点 0 以及该数列中的无穷多项 $\left\{\dfrac1n\right\}_{n>N}$. 剩下的 $N$ 项可分别用长度小于 $\dfrac{\varepsilon}{2N}$ 的小区间覆盖，因此全部间断点可用 $N+1$ 个总长小于
\[
\frac{\varepsilon}{2}+N\frac{\varepsilon}{2N}=\varepsilon
\]
的小区间覆盖. 利用例 4.2.10，函数在 $[0,1]$ 上可积.

\noindent\textbf{4.2.3}\quad 若 $f(x),g(x)$ 在 $[a,b]$ 上可积，试证 $\max\{f(x),g(x)\}$ 及 $\min\{f(x),g(x)\}$ 在 $[a,b]$ 上亦可积.

\noindent\textbf{提示}\quad 记 $\Phi(x)=\max\{f(x),g(x)\}$，不难证明：任一小区间上，
\[
\omega_i^\Phi\le\max\{\omega_i^f,\omega_i^g\}.
\tag{1}
\]
因 $f(x)$ 和 $g(x)$ 可积，（根据定理 3）$\forall\varepsilon>0$ 和 $\forall\sigma>0$，当分划充分细密时，有
\[
\sum_{\omega_i^f\ge\varepsilon}\Delta x_i<\frac{\sigma}{2}
\quad\text{和}\quad
\sum_{\omega_i^g\ge\varepsilon}\Delta x_i<\frac{\sigma}{2}.
\]
根据下面的注释，可得
\[
\sum_{\omega_i^\Phi\ge\varepsilon}\Delta x_i<\frac{\sigma}{2}+\frac{\sigma}{2}=\sigma.
\]
利用定理 3 知，$\Phi(x)=\max\{f(x),g(x)\}$ 在 $[a,b]$ 上可积.

类似可证 $\min\{f(x),g(x)\}$ 在 $[a,b]$ 上可积.

\noindent\textbf{注释}\quad $\max\{\omega_i^f,\omega_i^g\}$ 表示 $\omega_i^f,\omega_i^g$ 中的大者，以 $\omega_i^f\le\omega_i^g$ 为例进行说明. 此时 $\max\{\omega_i^f,\omega_i^g\}=\omega_i^g$，如式 (1)：$\omega_i^\Phi\le\max\{\omega_i^f,\omega_i^g\}=\omega_i^g$. 可见若 $\omega_i^\Phi\ge\varepsilon$，则 $\omega_i^g\ge\varepsilon$. 也就是说：满足 $\omega_i^\Phi\ge\varepsilon$ 的小区间之集合，是满足 $\omega_i^g\ge\varepsilon$ 的小区间集合里的子集.（若 $\omega_i^g\le\omega_i^f$，情况类似.）可见
\[
\sum_{\omega_i^\Phi\ge\varepsilon}\Delta x_i
\le
\sum_{\omega_i^g\ge\varepsilon}\Delta x_i
+
\sum_{\omega_i^f\ge\varepsilon}\Delta x_i.
\]

\noindent\textbf{4.2.4}\quad 试用定理 3 重新证明 Riemann 函数在 $[0,1]$ 上可积.

\noindent\textbf{提示}\quad 重复例 4.2.7 中的相应论述，可知：$\forall\varepsilon>0$，使得（Riemann 函数）$R(x)\ge\varepsilon$ 的点只有有限个. 因此作分划时，含此种点的小区间最多只有有限个. 当分划足够细密时，此类（$\omega_i^{R(x)}\ge\varepsilon$）小区间之总长能任意小，小于事先任意指定的 $\sigma>0$，即 $\sum_{\omega_i^{R(x)}\ge\varepsilon}\Delta x_i<\sigma$. 利用例 4.2.8，知 $R(x)$ 在该区间上可积.

\noindent\textbf{4.2.5}\quad 设 $f(x)$ 在 $[a,b]$ 上可积，$[f(x)]$ 表示对 $f(x)$ 的值取整数部分，试问 $[f(x)]$ 在 $[a,b]$ 上是否一定可积.

\noindent\textbf{提示}\quad 跟题 4.2.3 类似，$f(x)$ 的取整数部分是 $[f(x)]$，取整数后，同一区间上的振幅只会比未取整数时小，不会变大，即 $\omega_i^{[f(x)]}\le\omega_i^{f(x)}$. 故 $\forall\varepsilon>0,\forall\sigma>0$，
\[
\sum_{\omega_i^{f(x)}\ge\varepsilon}\Delta x_i<\sigma
\Rightarrow
\sum_{\omega_i^{[f(x)]}\ge\varepsilon}\Delta x_i
\le
\sum_{\omega_i^{f(x)}\ge\varepsilon}\Delta x_i<\sigma
\Rightarrow [f(x)]\text{ 可积.}
\]

\noindent\textbf{4.2.6}\quad 设 $f(x)$ 在 $[a,b]$ 上可积，试证：对于 $[a,b]$ 上任一可积函数 $g(x)$，恒有 $\int_a^b f(x)g(x)\,dx=0$，则函数 $f(x)$ 在连续点上恒为零.

\noindent\textbf{4.2.7}\quad 设在 $[-1,1]$ 上的连续函数 $f(x)$ 满足如下条件：对 $[-1,1]$ 上的任意偶连续函数 $g(x)$，积分 $\int_{-1}^1 f(x)g(x)\,dx=0$. 试证 $f(x)$ 是 $[-1,1]$ 上的奇函数.（武汉大学）

\noindent\textbf{提示}\quad 只需证明 $\int_{-1}^1[f(x)+f(-x)]^2\,dx=0$，并注意到 $g(x)=f(x)+f(-x)$ 是偶函数.

\section*{§4.3\quad 有关积分的几类典型问题}

\subsection*{一、积分值估计}

有些函数虽然可积，但原函数不能用初等函数的有限形式表达，或说这种函数的积分“积不出”，无法应用 Newton-Leibniz 公式计算，只能用其他方法对积分值进行估计或近似计算. 另一种情况是，被积函数没有明确给出，只知道它的结构或某些性质，希望对积分值给出某种估计.

\noindent\textbf{* a.\quad 利用 Darboux 和估计积分值}

\noindent\textbf{要点}\quad 若 $\underline S=\sum_{i=1}^n m_i\Delta x_i$，$\overline S=\sum_{i=1}^n M_i\Delta x_i$ 分别表示积分 $I=\int_a^b f(x)\,dx$ 的下、上 Darboux 和，那么积分存在时，有估计
\[
\underline S\le\int_a^b f(x)\,dx\le\overline S.
\]

\noindent\textbf{*例 4.3.1}\quad 求 $A,B$，使得 $A\le\int_0^1\sqrt{1+x^4}\,dx\le B$，要求 $B-A\le0.1$.（首都师范大学）

\noindent\textbf{注}\quad Чебышев 定理告诉我们，二项微分式的不定积分 $\int x^\alpha(1+x^\beta)^\gamma\,dx$ 除以下三种情况可以积分出来之外，其余情况积不出：

（i）$\gamma=$ 整数；（ii）$\gamma\ne$ 整数，但 $\dfrac{\alpha+1}{\beta}=$ 整数；（iii）$\gamma\ne$ 整数，$\dfrac{\alpha+1}{\beta}\ne$ 整数，但 $\dfrac{\alpha+1}{\beta}+\gamma=$ 整数. 现在 $\alpha=0,\beta=4,\gamma=\dfrac12$，不属于这三种情况，因此该积分积不出；但 $\sqrt{1+x^4}$ 连续，积分有意义. 下面我们用 Darboux 和对积分进行估计.

\noindent\textbf{解}\quad 将区间 $[0,1]$ $n$ 等分，利用 $\sqrt{1+x^4}$ 的单调性，每个小区间上端点到达上、下确界，因此
\[
\underline S_n=\frac1n\sum_{i=1}^n\sqrt{1+\left(\frac{i-1}{n}\right)^4}
\le\int_0^1\sqrt{1+x^4}\,dx
\le\frac1n\sum_{i=1}^n\sqrt{1+\left(\frac in\right)^4}
=\overline S_n,
\]
这时 $\overline S_n-\underline S_n=\dfrac1n(\sqrt2-1)$. 要使 $\overline S_n-\underline S_n<0.1$，只要取 $n=5$，于是
\[
A=\frac15\sum_{i=1}^5\sqrt{1+\left(\frac{i-1}{5}\right)^4}=1.053,\quad
B=\frac15\sum_{i=1}^5\sqrt{1+\left(\frac i5\right)^4}=1.135.
\]

\noindent\textbf{*例 4.3.2}\quad 若 $f(x)$ 在 $[a,b]$ 上可积，$f(x)>0$，试证：
\[
\int_a^b f(x)\,dx>0.
\]
（该结论很重要，经常用到.）

\noindent\textbf{证 I}\quad $f(x)>0$，$\sum_{i=1}^n f(\xi_i)\Delta x_i>0$，因而由可积性知
\[
\int_a^b f(x)\,dx=\lim_{\lambda\to0}\sum_{i=1}^n f(\xi_i)\Delta x_i\ge0.
\]
（现用反证法证明这里的等号不可能发生.）设 $\int_a^b f(x)\,dx=0$，则对于 $f(x)$ 的上 Darboux 和，有 $\lim_{\lambda\to0}\sum_{i=1}^n M_i\Delta x_i=0$. 从而 $\forall\varepsilon_1>0$，$\exists$ 分划 $T$ 使得
\[
\sum_{i=1}^nM_i\Delta x_i<\varepsilon_1(b-a).
\]
由此至少存在一个 $M_i<\varepsilon_1$，否则每个 $M_i\ge\varepsilon_1$，应有 $\sum_{i=1}^nM_i\Delta x_i\ge\varepsilon_1\sum\Delta x_i=\varepsilon_1(b-a)$，与上式矛盾.

将 $M_i<\varepsilon_1$ 的这个小区间记为 $[a_1,b_1]$，于是 $f(x)$ 在 $[a_1,b_1]$ 上可积. 把 $[a_1,b_1]$ 取作上面的 $[a,b]$，重复上述推理，可得到 $[a_2,b_2]\subset[a_1,b_1]$，使得
\[
\sup_{a_2\le x\le b_2} f(x)\le\varepsilon_2.
\]
如此无限进行下去，可以得到一串区间
\[
[a,b]\supset[a_1,b_1]\supset[a_2,b_2]\supset\cdots\supset[a_n,b_n]\supset[a_{n+1},b_{n+1}]\supset\cdots,
\]
使得每个 $[a_n,b_n]$ 上
\[
\sup_{a_n\le x\le b_n} f(x)\le\varepsilon_n\quad(n=1,2,\cdots).
\]
根据区间套定理，$\exists\xi\in[a_n,b_n]\ (n=1,2,\cdots)$，从而
\[
0\le f(\xi)<\varepsilon_n\quad(n=1,2,\cdots).
\]
令 $\varepsilon_n\to0$，可知 $f(\xi)=0$. 与已知条件 $f(x)>0$ 矛盾.

\noindent\textbf{证 II}\quad 利用例 4.2.11 的结论即得.

\noindent\textbf{b.\quad 利用变形求估计及积分估计的应用}

\noindent\textbf{要点}\quad 若 $f(x)$ 在 $[a,b]$ 上可积，一般来说我们可以通过各种变形来对积分 $\int_a^b f(x)\,dx$ 之值进行估计. 例如用变量替换、分部积分、中值公式、Taylor 公式等，使积分变成易于估计的形式. 另外，被积函数放大、缩小，区间放大和缩小也是获得估计的重要方法.

\noindent\textbf{利用变量替换进行变形}

\noindent\textbf{例 4.3.3}\quad 证明 $\int_0^{\sqrt{2\pi}}\sin x^2\,dx>0$.（国外赛题）

\noindent\textbf{证}\quad 令 $x^2=y$，作变换
\[
I=\int_0^{\sqrt{2\pi}}\sin x^2\,dx
=\frac12\int_0^{2\pi}\frac{\sin y}{\sqrt y}\,dy
=\frac12\int_0^\pi\frac{\sin y}{\sqrt y}\,dy
+\frac12\int_\pi^{2\pi}\frac{\sin y}{\sqrt y}\,dy
=I_1+I_2.
\]
在 $I_2$ 中作变换 $z=y-\pi$，
\[
I_2=\frac12\int_\pi^{2\pi}\frac{\sin y}{\sqrt y}\,dy
=-\frac12\int_0^\pi\frac{\sin z}{\sqrt{z+\pi}}\,dz
=-\frac12\int_0^\pi\frac{\sin y}{\sqrt{y+\pi}}\,dy.
\]
于是
\[
I=\frac12\int_0^\pi\left(\frac{\sin y}{\sqrt y}-\frac{\sin y}{\sqrt{y+\pi}}\right)\,dy
=\frac12\int_0^\pi\sin y\cdot\left(\frac1{\sqrt y}-\frac1{\sqrt{y+\pi}}\right)\,dy.
\]
在 $(0,\pi)$ 内，被积函数 $\sin y\cdot\left(\dfrac1{\sqrt y}-\dfrac1{\sqrt{y+\pi}}\right)>0$.

个别点不影响积分值. 由例 4.3.2 知 $I>0$.

\noindent\textbf{利用分部积分进行变形}

若 $f(x),g(x)$ 在 $[a,b]$ 上有连续导数，$f(x)g(x)\big|_a^b=0$，则
\[
\int_a^b f'(x)g(x)\,dx
=f(x)g(x)\bigg|_a^b-\int_a^b f(x)g'(x)\,dx
=-\int_a^b f(x)g'(x)\,dx.
\]
若 $f(x),g(x)$ 在 $[a,b]$ 上有 $2n$ 阶连续导数，且 $f^{(i)}(a)=f^{(i)}(b)=g^{(i)}(a)=g^{(i)}(b)=0\ (i=0,1,2,\cdots,n-1)$，则反复利用分部积分法可得
\[
\int_a^b f^{(2n)}(x)g(x)\,dx=\int_a^b f(x)g^{(2n)}(x)\,dx.
\tag{A}
\]
此式给出了一个重要的变形.

\noindent\textbf{例 4.3.4}\quad 设 $f(x)$ 在 $[a,b]$ 上有 $2n$ 阶连续导数，$|f^{(2n)}(x)|\le M$，$f^{(i)}(a)=f^{(i)}(b)=0\ (i=0,1,2,\cdots,n-1)$. 试证：
\[
\left|\int_a^b f(x)\,dx\right|
\le\frac{(n!)^2M}{(2n)!(2n+1)!}(b-a)^{2n+1}.
\]

\noindent\textbf{证}\quad 令 $g(x)=(x-a)^n(b-x)^n$，则 $g^{(i)}(a)=g^{(i)}(b)=0\ (i=0,1,2,\cdots,n-1)$，且 $g^{(2n)}(x)=(-1)^n(2n)!$. 因此，利用上面刚介绍的公式 (A)，有
\[
\int_a^b f(x)\,dx
=\frac{(-1)^n}{(2n)!}\int_a^b f(x)(-1)^n(2n)!\,dx
=\frac{(-1)^n}{(2n)!}\int_a^b f(x)g^{(2n)}(x)\,dx
\]
\[
=\frac{(-1)^n}{(2n)!}\int_a^b f^{(2n)}(x)g(x)\,dx;
\]
所以
\[
\left|\int_a^b f(x)\,dx\right|
\le\frac1{(2n)!}\int_a^b|f^{(2n)}(x)|(x-a)^n(b-x)^n\,dx
\le\frac{M}{(2n)!}\int_a^b(x-a)^n(b-x)^n\,dx.
\]
令 $x=a+t(b-a)$，
\[
\int_a^b(x-a)^n(b-x)^n\,dx
=(b-a)^{2n+1}\int_0^1t^n(1-t)^n\,dt,
\]
令 $t=\sin^2\theta$，
\[
\int_0^1t^n(1-t)^n\,dt
=2\int_0^{\pi/2}\sin^{2n+1}\theta\cos^{2n+1}\theta\,d\theta
\quad(\text{利用 Wallis 公式})
\]
\[
=2\cdot\frac{(2n)!!(2n)!!}{[2(2n+1)]!!}
=2\cdot\frac{2^n n!\cdot2^n n!}{2^{2n+1}(2n+1)!}
=\frac{(n!)^2}{(2n+1)!}.
\]
（或者
\[
\int_0^1t^n(1-t)^n\,dt
=B(n+1,n+1)
=\frac{\Gamma(n+1)\Gamma(n+1)}{\Gamma(2n+2)}
=\frac{(n!)^2}{(2n+1)!}.
\]
）
总之，
\[
\left|\int_a^b f(x)\,dx\right|
\le\frac{(n!)^2M}{(2n)!(2n+1)!}(b-a)^{2n+1}.
\]

\noindent\textbf{利用缩放被积函数或积分区间进行变形}

\noindent\textbf{☆例 4.3.5}\quad $f(x)$ 在 $[0,1]$ 上有连续二阶导数，$f(0)=f(1)=0$，$f(x)\ne0$（当 $x\in(0,1)$ 时），试证：
\[
\int_0^1\left|\frac{f''(x)}{f(x)}\right|\,dx\ge4.
\tag{1}
\]
（四川大学）

\noindent\textbf{证}\quad 因 $(0,1)$ 内 $f(x)\ne0$，故 $f(x)$ 在 $(0,1)$ 内恒正或恒负（否则由介值性，必有零点在 $(0,1)$ 内，与 $f(x)\ne0$ 矛盾）. 不妨设 $f(x)>0$（$<0$ 的情况类似可证）. 因 $f(x)$ 在 $[0,1]$ 上连续，故存在 $c\in[0,1]$，使得 $f(c)=\max\limits_{0\le x\le1}f(x)$. 于是 $\forall a,b:0<a<b<1$，有
\[
\int_0^1\left|\frac{f''(x)}{f(x)}\right|\,dx
\ge\int_a^b\left|\frac{f''(x)}{f(c)}\right|\,dx
=\frac1{f(c)}\int_a^b|f''(x)|\,dx
\ge\frac1{f(c)}\left|\int_a^b f''(x)\,dx\right|
=\frac1{f(c)}|f'(b)-f'(a)|.
\tag{2}
\]
下面我们来恰当地选取 $a,b$，得到所需的估计. 注意到 $f(0)=f(1)=0$，应用 Lagrange 公式，$\exists\xi\in(0,c)$，使得
\[
f'(\xi)=\frac{f(c)-f(0)}{c-0}=\frac{f(c)}c;
\tag{3}
\]
$\exists\eta\in(c,1)$，使得
\[
f'(\eta)=\frac{f(1)-f(c)}{1-c}=\frac{f(c)}{c-1}.
\tag{4}
\]
在式 (2) 中令 $a=\xi,b=\eta$，将 (3) 和 (4) 代入，注意到
\[
\sqrt{c(1-c)}\le\frac{c+(1-c)}2=\frac12,
\]
得
\[
\int_0^1\left|\frac{f''(x)}{f(x)}\right|\,dx
\ge\frac1{f(c)}|f'(\eta)-f'(\xi)|
=\frac1{c(1-c)}\ge4.
\]
式 (1) 获证.

\noindent\textbf{练习 1}\quad 设函数 $f(x)$ 在 $[2,4]$ 上有连续的一阶导数，且 $f(2)=f(4)=0$. 试证：
\[
\max_{2\le x\le4}f'(x)\ge\int_2^4f(x)\,dx.
\]
（华中师范大学）

\noindent\textbf{提示}\quad
\[
\int_2^4f(x)\,dx=\int_2^3f(x)\,dx+\int_3^4f(x)\,dx=I+J.
\]
下面研究 $I$（$J$ 类似）.
\[
I=\int_2^3[f(x)-f(2)]\,dx
\stackrel{\text{Lagrange 定理}}{=}\int_2^3 f'(\xi_1)(x-2)\,dx
\le\frac12\max_{2\le x\le4}f'(x).
\]

\noindent\textbf{☆练习 2}\quad 已知 $f(x)$ 在 $[0,1]$ 上有直到二阶的连续导数，试证：
\[
\max_{x\in[0,1]}|f'(x)|
\le |f(1)-f(0)|+\int_0^1|f''(x)|\,dx.
\tag{1}
\]
（华中科技大学）

\noindent\textbf{证}\quad $1^\circ$ 因 $f'(x)$ 连续，$\exists\eta\in[0,1]$，使得
\[
|f'(\eta)|=\max_{x\in[0,1]}|f'(x)|.
\tag{2}
\]
$2^\circ$ 应用 Lagrange 定理，$\exists\xi\in(0,1)$，使得
\[
f(1)-f(0)=f'(\xi)(1-0)=f'(\xi).
\tag{3}
\]
$3^\circ$ 注意 $\forall(\alpha,\beta)\subset[0,1]$，恒有
\[
\int_0^1|f''(x)|\,dx
\ge\int_\alpha^\beta|f''(x)|\,dx
\ge\left|\int_\alpha^\beta f''(x)\,dx\right|
=|f'(\beta)-f'(\alpha)|
\ge\bigl||f'(\beta)|-|f'(\alpha)|\bigr|.
\]
再令 $\alpha=\xi,\beta=\eta$，利用式 (2), (3)，得
\[
\int_0^1|f''(x)|\,dx
\ge |f'(\eta)|-|f'(\xi)|
\ge \max_{x\in[0,1]}|f'(x)|-|f(1)-f(0)|.
\]
移项即得式 (1). 证毕.

\noindent\textbf{*练习 3}\quad 已知 $f(x)$ 在 $[0,1]$ 上有二阶连续导数，且 $f(1)=f(0)=0$. 试证：
\[
\int_0^1|f'(x)|\,dx
\le\int_0^1|f''(x)|\,dx+2\int_0^1|f(x)|\,dx.
\tag{1}
\]

\noindent\textbf{证}\quad $1^\circ$
\[
\int_0^1|f'(x)|\,dx
\stackrel{\text{积分中值定理}}{=}|f'(\xi)|
\le\max_{x\in[0,1]}|f'(x)|
\stackrel{f'\text{连续}}{=}|f'(x_0)|.
\tag{2}
\]
$2^\circ$
\[
2\int_0^1|f(x)|\,dx
\stackrel{\text{积分中值定理}}{=}2|f(\eta)|
=2|f(\eta)-f(\zeta)|
\stackrel{\text{Lagrange 定理}}{=}2|f'(\alpha)(\eta-\zeta)|
\ge |f'(\alpha)|.
\tag{3}
\]
其中：若 $\eta\in[0,\frac12]$，取 $\zeta=1$；若 $\eta\in(\frac12,1]$，取 $\zeta=0$.

$3^\circ$ 因 $x_0,\alpha\in[0,1]$，故
\[
\int_0^1|f''(x)|\,dx
\ge\left|\int_\alpha^{x_0}f''(x)\,dx\right|
=|f'(x_0)-f'(\alpha)|.
\tag{4}
\]
于是
\[
\text{式 (1) 左端}
=\int_0^1|f'(x)|\,dx
\stackrel{\text{式 (2)}}{\le}|f'(x_0)|
\le |f'(x_0)-f'(\alpha)|+|f'(\alpha)|
\stackrel{\text{式 (4),(3)}}{\le}
\int_0^1|f''(x)|\,dx+2\int_0^1|f(x)|\,dx
=\text{式 (1) 右端}.
\]
式 (1) 获证.

\noindent\textbf{小结}\quad 要证不等式：“$A\le B$”，常采用缩放法，即：将小者 $A$（简化）适当放大；将大者 $B$（简化）适当缩小，使得
\[
A\le(\text{简化放大后的小者 }A_{\text{新}})
\le(\text{简化缩小后的大者 }B_{\text{新}})\le B.
\]

\noindent\textbf{*练习 4}\quad 已知 $f(x)$ 在 $[0,1]$ 上有直到二阶的连续导数. 求证：
\[
\int_0^1|f'(x)|\,dx
\le\int_0^1|f''(x)|\,dx+8\int_0^1|f(x)|\,dx.
\tag{1}
\]
（华中科技大学）

\noindent\textbf{证}\quad $\forall\alpha\in(0,\frac14)$，$\forall\beta\in(\frac34,1)$. 根据 Lagrange 定理，$\exists\xi\in(\alpha,\beta)$，使得
\[
|f'(\xi)|
=\left|\frac{f(\beta)-f(\alpha)}{\beta-\alpha}\right|
\le2|f(\beta)-f(\alpha)|
\le2|f(\alpha)|+2|f(\beta)|.
\tag{2}
\]
利用积分的 Newton-Leibniz 公式，$\forall x\in[0,1]$，有
\[
|f'(x)|=\left|\int_\xi^x f''(t)\,dt+f'(\xi)\right|
\le\int_\xi^x|f''(t)|\,dt+|f'(\xi)|.
\]
再用 (2) 代入，得
\[
|f'(x)|\le\int_0^1|f''(t)|\,dt+2|f(\alpha)|+2|f(\beta)|.
\tag{3}
\]
注意到：此式里 $x,\alpha,\beta$ 已相互独立，在式 (3) 两端同时取积分 $\int_0^{1/4}(\cdot)\,d\alpha$；再取积分 $\int_{3/4}^1(\cdot)\,d\beta$，可得
\[
\frac1{16}|f'(x)|
\le\frac1{16}\int_0^1|f''(t)|\,dt
+\frac12\int_0^{1/4}|f(\alpha)|\,d\alpha
+\frac12\int_{3/4}^1|f(\beta)|\,d\beta
\le\frac1{16}\int_0^1|f''(t)|\,dt+\frac12\int_0^1|f(t)|\,dt.
\]
故
\[
|f'(x)|\le\int_0^1|f''(t)|\,dt+8\int_0^1|f(t)|\,dt.
\]
再在 $[0,1]$ 上对 $x$ 求积分，即得式 (1):
\[
\int_0^1|f'(x)|\,dx
\le\int_0^1|f''(x)|\,dx+8\int_0^1|f(x)|\,dx.
\]

\noindent\textbf{说明}\quad 作为名校考研压轴题，张卫发现并推荐了他和网友的解法，此处采用，并略作改进. 跟上题相比，本题减少了条件：$f(1)=f(0)=0$，但也付出了一定的代价，因为若将式 (1) 看作对积分 $\int_0^1|f'(x)|\,dx$ 的估计，则新结论相对较弱（系数原为 2，现为 8）.

\noindent\textbf{利用微分中值公式或 Taylor 公式对被积函数进行变形}

\noindent\textbf{☆例 4.3.6}\quad 设 $f(x)$ 在 $[a,b]$ 上二次连续可微，$f(\frac{a+b}{2})=0$，试证：
\[
\left|\int_a^b f(x)\,dx\right|\le\frac{M(b-a)^3}{24},
\]
其中 $M=\sup\limits_{a\le x\le b}|f''(x)|$.（中山大学）

\noindent\textbf{证}\quad 将 $f(x)$ 在 $x=\frac{a+b}{2}$ 处用 Taylor 公式展开，注意到 $f(\frac{a+b}{2})=0$，有
\[
f(x)=f'\left(\frac{a+b}{2}\right)\left(x-\frac{a+b}{2}\right)
+\frac1{2!}f''(\xi)\left(x-\frac{a+b}{2}\right)^2.
\]
右端第一项在 $[a,b]$ 上的积分为零. 故
\[
\left|\int_a^b f(x)\,dx\right|
\le\frac1{2!}\int_a^b|f''(\xi)|\left(x-\frac{a+b}{2}\right)^2\,dx
\le\frac16M\left(x-\frac{a+b}{2}\right)^3\bigg|_a^b
=\frac{M}{24}(b-a)^3.
\]

\noindent\textbf{练习 1}\quad 已知 $f(x)$ 在 $[a,b]$ 上二阶可导，$f(a)=f(b)=0$. 试证：$\exists\xi\in(a,b)$，使得
\[
\int_a^b f(x)\,dx=\frac1{12}f''(\xi)(a-b)^3.
\]
（南京大学）

\noindent\textbf{证 I}\quad 令
\[
F(x)=\int_a^x f(t)\,dt-\frac12(x-a)f(x),\quad
G(x)=-\frac1{12}(x-a)^3,
\]
在 $[a,b]$ 上连续两次应用 Cauchy 中值定理，因 $f(a)=f(b)=0$，
\[
\frac{\int_a^b f(x)\,dx-0}{-\frac1{12}(b-a)^3-0}
=\frac{F(b)-F(a)}{G(b)-G(a)}
\stackrel{\text{Cauchy 中值定理}}{=}\frac{F'(\xi_1)}{G'(\xi_1)}
\]
\[
=\frac{\frac12f(\xi_1)-\frac12(\xi_1-a)f'(\xi_1)-0}
{-\frac14(\xi_1-a)^2-0}
\stackrel{\text{Cauchy 中值定理}}{=}
\frac{-\frac12f''(\xi)(\xi-a)}{-\frac12(\xi-a)}
=f''(\xi),
\]
其中 $a\le\xi\le\xi_1\le b$. 此式同乘 $-\frac1{12}(b-a)^3=\frac1{12}(a-b)^3$，即得所求.

\noindent\textbf{证 II}\quad 先把习题 3.2.21 作为引理，这时函数 $F(x)=\int_a^x f(t)\,dt$ 满足条件，代入结果，立即可得 $\exists\xi\in(a,b)$，使得
\[
\int_a^b f(t)\,dt
=F(b)
=F(a)+\frac12(b-a)[F'(a)+F'(b)]-\frac{F'''(\xi)}{12}(b-a)^3
=\frac12(b-a)[f(a)+f(b)]-\frac{f''(\xi)}{12}(b-a)^3
=\frac{f''(\xi)}{12}(a-b)^3.
\]

\noindent\textbf{练习 2}\quad 已知 $f(x)$ 在 $[a,b]$ 内二阶连续可导，满足
\[
\left|\int_a^b f(x)\,dx\right|<\int_a^b|f(x)|\,dx.
\tag{1}
\]
设
\[
\max_{x\in[a,b]}|f'(x)|=M_1,\quad
\max_{x\in[a,b]}|f''(x)|=M_2,
\tag{2}
\]
求证：
\[
\left|\int_a^b f(x)\,dx\right|
\le\frac{M_1}{2}(b-a)^2+\frac{M_2}{6}(b-a)^3.
\tag{3}
\]
（南开大学）

\noindent\textbf{提示}\quad 由已知的不等式 (1)，可知 $f(x)$ 变号，故 $\exists x_0\in[a,b]$ 使得 $f(x_0)=0$. $f(x)$ 在 $x_0$ 处按 Taylor 公式展开，求积分便可得欲求的不等式.

\noindent\textbf{再提示}\quad
\[
f(x)=f(x_0)+f'(x_0)(x-x_0)+\frac12f''(\xi)(x-x_0)^2
\quad(f(x_0)=0).
\]

\noindent\textbf{证}\quad 由提示知，存在 $x_0\in(a,b)$，使得 $f(x_0)=0$. 于是
\[
\begin{aligned}
\left|\int_a^b f(x)\,dx\right|
&\le |f'(x_0)|\int_a^b|x-x_0|\,dx
+\frac12\int_a^b|f''(\xi)|\,|x-x_0|^2\,dx\\
&\le M_1\left[\int_a^{x_0}(x_0-x)\,dx+\int_{x_0}^b(x-x_0)\,dx\right]
+\frac{M_2}{2}\int_a^b(x-x_0)^2\,dx\\
&=\frac{M_1}{2}\left[(b-x_0)^2+(x_0-a)^2\right]
+\frac{M_2}{6}\left[(b-x_0)^3+(x_0-a)^3\right]\\
&\le\frac{M_1}{2}\left[(b-x_0)+(x_0-a)\right]^2
+\frac{M_2}{6}\left[(b-x_0)+(x_0-a)\right]^3\\
&=\frac{M_1}{2}(b-a)^2+\frac{M_2}{6}(b-a)^3.
\end{aligned}
\]
不等式 (3) 获证.

\noindent\textbf{例 4.3.7}\quad 设 $f(x)$ 在 $[a,b]$ 上连续可微，且 $f(a)=f(b)=0$，证明：
\[
\max_{a\le x\le b}|f'(x)|\ge
\frac{4}{(b-a)^2}\int_a^b|f(x)|\,dx.
\]

\noindent\textbf{提示}\quad
\[
\int_a^b|f(x)|\,dx
=\int_a^{\frac{a+b}{2}}|f(x)|\,dx
+\int_{\frac{a+b}{2}}^b|f(x)|\,dx.
\]

\noindent\textbf{再提示}\quad 右端两项分别在 $a,b$ 两点处应用微分中值公式.
\[
\begin{aligned}
\int_a^b|f(x)|\,dx
&=\int_a^{\frac{a+b}{2}}|f(a)+f'(\xi)(x-a)|\,dx
+\int_{\frac{a+b}{2}}^b|f(b)+f'(\xi)(x-b)|\,dx\\
&\le\max_{a\le x\le b}|f'(x)|
\left[\int_a^{\frac{a+b}{2}}(x-a)\,dx
+\int_{\frac{a+b}{2}}^b(b-x)\,dx\right]\\
&=\max_{a\le x\le b}|f'(x)|\cdot\frac14(b-a)^2.
\end{aligned}
\]
或者
\[
\begin{aligned}
\int_a^b|f(x)|\,dx
&=\int_a^{\frac{a+b}{2}}\left|\int_a^x f'(t)\,dt\right|\,dx
+\int_{\frac{a+b}{2}}^b\left|-\int_x^b f'(t)\,dt\right|\,dx\\
&\le\max_{a\le x\le b}|f'(x)|
\left(\int_a^{\frac{a+b}{2}}(x-a)\,dx
+\int_{\frac{a+b}{2}}^b(b-x)\,dx\right)\\
&=\max_{a\le x\le b}|f'(x)|\cdot\frac14(b-a)^2.
\end{aligned}
\]

\noindent 以上各例为已知函数的性质，要对积分进行估计. 下面看一个反问题：已知一个估计，看是否正确.

\noindent\textbf{☆例 4.3.8}\quad 在 $[0,2]$ 上是否存在这样的函数，连续可微，并且 $f(0)=f(2)=1$，$|f'(x)|\le1$，$\left|\int_0^2 f(x)\,dx\right|\le1$？（扬州大学）

\noindent\textbf{证}\quad （根据已知条件重新对积分 $\int_0^2 f(x)\,dx$ 进行估计.）
\[
\int_0^2 f(x)\,dx=\int_0^1 f(x)\,dx+\int_1^2 f(x)\,dx.
\]
右端第一项按 $x=0$ 展开，注意 $f(0)=1$ 及 $|f'(x)|\le1$，
\[
f(x)=f(0)+f'(\xi)x=1+f'(\xi)x\ge1-x\quad(\forall x\in[0,1]).
\tag{1}
\]
从而
\[
\int_0^1 f(x)\,dx\ge\int_0^1(1-x)\,dx=\frac12.
\]
类似，第二项有
\[
f(x)\ge x-1\quad(\text{当 }x\in[1,2]\text{ 时}).
\tag{2}
\]
所以
\[
\int_0^2 f(x)\,dx
\ge\int_0^1(1-x)\,dx+\int_1^2(x-1)\,dx
=\frac12+\frac12=1.
\]
（现证这种 $f$ 不存在.）假设这种 $f$ 存在，则由 $\left|\int_0^2 f(x)\,dx\right|\le1$ 及 $\int_0^2 f(x)\,dx\ge1$ 知
\[
\int_0^2 f(x)\,dx=1
=\int_0^1(1-x)\,dx+\int_1^2(x-1)\,dx,
\tag{3}
\]
记
\[
g(x)=
\begin{cases}
1-x, & \text{当 }0\le x\le1\text{ 时},\\
x-1, & \text{当 }1<x\le2\text{ 时}.
\end{cases}
\]
式 (1)、(2)、(3) 表明两连续函数 $f(x)\ge g(x)$，且积分值相等. 因此有 $f(x)\equiv g(x)$ 在 $[0,2]$ 上，与 $f$ 的可微性矛盾，所以 $f$ 不存在.

\noindent 以上，我们主要讨论积分估计的基本方法. 下面我们来看积分估计的某些应用.

\noindent\textbf{关于函数值的估计}

\noindent\textbf{☆例 4.3.9}\quad 设 $f(x)$ 在 $[0,1]$ 上连续，
\[
\int_0^1 f(x)\,dx=0,\quad
\int_0^1 xf(x)\,dx=0,\cdots,\quad
\int_0^1 x^{n-1}f(x)\,dx=0,\quad
\int_0^1 x^n f(x)\,dx=1\quad(n>1).
\]
求证：在 $[0,1]$ 的某一部分上 $|f(x)|\ge2^n(n+1)$.（吉林大学）

\noindent\textbf{证}\quad 由已知条件，对任意 $\alpha$，恒有
\[
\int_0^1(x-\alpha)^n f(x)\,dx=1.
\tag{1}
\]
（反证法）假设 $[0,1]$ 上处处都有 $|f(x)|<2^n(n+1)$. 若能选取恰当的 $\alpha$，由此得出估计 $\left|\int_0^1(x-\alpha)^n f(x)\,dx\right|<1$，便找到了矛盾. 事实上，这时有
\[
\begin{aligned}
\left|\int_0^1(x-\alpha)^n f(x)\,dx\right|
&<2^n(n+1)\int_0^1|x-\alpha|^n\,dx\quad\left(\text{取 }\alpha=\frac12\right)\\
&=2^n(n+1)\int_0^1\left|x-\frac12\right|^n\,dx\\
&=2^n(n+1)\left[\int_0^{\frac12}\left(\frac12-x\right)^n\,dx
+\int_{\frac12}^1\left(x-\frac12\right)^n\,dx\right]=1.
\end{aligned}
\]
证毕.

\noindent 下面两例通过积分值的估计解决零点存在性问题.

\noindent\textbf{例 4.3.10}\quad 设 $f(x)$ 在 $\left[0,\frac\pi2\right]$ 上连续，
\[
\int_0^{\frac\pi2} f(x)\sin x\,dx
=\int_0^{\frac\pi2} f(x)\cos x\,dx=0.
\]
试证：$f(x)$ 在 $\left(0,\frac\pi2\right)$ 内至少有两个零（值）点.

\noindent\textbf{证}\quad $1^\circ$ 若 $f(x)$ 在 $\left(0,\frac\pi2\right)$ 无零点，因 $f(x)$ 连续，$f(x)$ 在 $\left(0,\frac\pi2\right)$ 恒保持同号，例如 $f(x)>0$（或 $f(x)<0$），则得估计
\[
\int_0^{\frac\pi2} f(x)\sin x\,dx>0\quad(\text{或 }<0),
\]
与已知条件矛盾. 可见 $\left(0,\frac\pi2\right)$ 中至少有一个零点 $x_0$.

\noindent $2^\circ$ 若 $f(x)$ 除 $x_0$ 外在 $\left(0,\frac\pi2\right)$ 内再无零点，则 $f(x)$ 在 $(0,x_0)$ 与 $\left(x_0,\frac\pi2\right)$ 内分别保持不变号. 若 $f$ 在此二区间符号相异，则 $f(x)\sin(x-x_0)$ 恒正（或恒负），
\[
\int_0^{\frac\pi2} f(x)\sin(x-x_0)\,dx>0\quad(\text{或 }<0).
\]
但由已知条件
\[
\int_0^{\frac\pi2} f(x)\sin(x-x_0)\,dx
=\cos x_0\int_0^{\frac\pi2} f(x)\sin x\,dx
-\sin x_0\int_0^{\frac\pi2} f(x)\cos x\,dx=0,
\]
矛盾. 若 $f$ 在两区间上符号相同，则 $f(x)\cos(x-x_0)$ 恒正（或恒负），同样可推出矛盾.

\noindent\textbf{*☆例 4.3.11}\quad 证明方程
\[
\int_0^x e^{-t}\left(1+\frac{t}{1!}+\frac{t^2}{2!}+\cdots+\frac{t^{100}}{100!}\right)\,dt=50
\]
在开区间 $(50,100)$ 内有根 $\alpha$.（国外赛题）

\noindent\textbf{证}\quad 作为上限 $x$ 的函数，
\[
F(x)=\int_0^x e^{-t}\left(1+\frac{t}{1!}+\frac{t^2}{2!}+\cdots+\frac{t^{100}}{100!}\right)\,dt
\]
连续. 只要证明了 $F(50)<50,F(100)>50$，则由连续函数介值定理，知 $\exists\alpha\in(50,100)$ 使得 $F(\alpha)=50$. 因为
\[
e^{-t}\left(1+\frac{t}{1!}+\frac{t^2}{2!}+\cdots+\frac{t^{100}}{100!}\right)<1,
\]
故 $F(50)<50$ 明显. 剩下的只要证明估计式：$F(100)>50$.

反复使用分部积分法：
\[
\begin{aligned}
F(100)
&=\int_0^{100}e^{-t}\left(1+\frac{t}{1!}+\cdots+\frac{t^{100}}{100!}\right)\,dt\\
&=-\int_0^{100}\left(1+\frac{t}{1!}+\cdots+\frac{t^{100}}{100!}\right)\,de^{-t}\\
&=1-e^{-100}\left(1+\frac{100}{1!}+\cdots+\frac{100^{100}}{100!}\right)
+\int_0^{100}e^{-t}\left(1+\frac{t}{1!}+\cdots+\frac{t^{99}}{99!}\right)\,dt\\
&=\cdots\\
&=1-e^{-100}\left(1+\frac{100}{1!}+\cdots+\frac{100^{99}}{99!}+\frac{100^{100}}{100!}\right)\\
&\quad+1-e^{-100}\left(1+\frac{100}{1!}+\cdots+\frac{100^{99}}{99!}\right)+\cdots
+1-e^{-100}\left(1+\frac{100}{1!}\right)+1-e^{-100}\\
&>101-e^{-100}\cdot\frac{102}{2}
\left(1+\frac{100}{1!}+\cdots+\frac{100^{101}}{101!}\right)>101-51=50.
\end{aligned}
\]
证毕.

\noindent\textbf{☆二、积分不等式}

我们把联系两个以上的定积分的不等式，称为积分不等式. 关于积分不等式，有不少著名的结果，我们将在下节讨论. 这里只介绍证明积分不等式的若干基本方法.

\noindent\textbf{a. 用微分学的方法证明积分不等式}

\noindent\textbf{☆例 4.3.12}\quad 设 $f(x)$ 在 $[0,1]$ 上可微，且当 $x\in(0,1)$ 时，$0<f'(x)<1,f(0)=0$. 试证：
\[
\left(\int_0^1 f(x)\,dx\right)^2>\int_0^1 f^3(x)\,dx.
\]
（上海交通大学）

\noindent\textbf{证 I}\quad 问题在于证明 $\left(\int_0^1 f(x)\,dx\right)^2-\int_0^1 f^3(x)\,dx>0$. 令
\[
F(x)=\left(\int_0^x f(t)\,dt\right)^2-\int_0^x f^3(t)\,dt.
\]
因 $F(0)=0$，故只要证明在 $(0,1)$ 内有 $F'(x)>0$. 事实上，
\[
F'(x)=2f(x)\int_0^x f(t)\,dt-f^3(x)
=f(x)\left(2\int_0^x f(t)\,dt-f^2(x)\right).
\tag{1}
\]
已知 $f(0)=0,0<f'(x)<1\ (x\in(0,1))$，故当 $x\in(0,1)$ 时，$f(x)>0$.（以下证 (1) 中另一因子大于零.）记
\[
g(x)=2\int_0^x f(t)\,dt-f^2(x),
\]
则 $g(0)=0$，
\[
g'(x)=2f(x)-2f(x)f'(x)=2f(x)(1-f'(x))>0,
\]
于是
\[
g(x)=2\int_0^x f(t)\,dt-f^2(x)>0\quad(x\in(0,1)).
\]
$F'(x)>0$ 获证.

\noindent\textbf{证 II}\quad 问题在于证明
\[
\frac{\left(\int_0^1 f(x)\,dx\right)^2}{\int_0^1 f^3(x)\,dx}>1.
\tag{2}
\]
令
\[
F(x)=\left(\int_0^x f(t)\,dt\right)^2,\quad
G(x)=\int_0^x f^3(t)\,dt,
\]
则式 (2) 左端（利用 Cauchy 中值定理）
\[
\begin{aligned}
\frac{\left(\int_0^1 f(x)\,dx\right)^2}{\int_0^1 f^3(x)\,dx}
&=\frac{F(1)-F(0)}{G(1)-G(0)}
\overset{\exists0<\xi<1}{=}
\frac{F'(\xi)}{G'(\xi)}
=\frac{2f(\xi)\int_0^\xi f(t)\,dt}{f^3(\xi)}\\
&=\frac{2\int_0^\xi f(t)\,dt}{f^2(\xi)}
=\frac{2\int_0^\xi f(t)\,dt-2\int_0^\xi f(0)\,dt}{f^2(\xi)-f^2(0)}\\
&=\frac{2f(\eta)}{2f(\eta)f'(\eta)}
=\frac1{f'(\eta)}>1\quad(0<\eta<\xi<1).
\end{aligned}
\]

\noindent\textbf{b. 利用被积函数的不等式证明积分不等式}

\noindent\textbf{☆例 4.3.13}\quad 试证
\[
\int_0^1\frac{\cos x}{\sqrt{1-x^2}}\,dx
>\int_0^1\frac{\sin x}{\sqrt{1-x^2}}\,dx.
\tag{1}
\]
（国外赛题）

\noindent\textbf{分析}\quad 令 $t=\arcsin x$，
\[
\int_0^1\frac{\cos x}{\sqrt{1-x^2}}\,dx
=\int_0^{\frac\pi2}\cos(\sin t)\,dt.
\]
令 $t=\arccos x$，
\[
\int_0^1\frac{\sin x}{\sqrt{1-x^2}}\,dx
=\int_0^{\frac\pi2}\sin(\cos t)\,dt.
\]
欲证的不等式化为
\[
\int_0^{\frac\pi2}\cos(\sin t)\,dt
>\int_0^{\frac\pi2}\sin(\cos t)\,dt.
\tag{2}
\]
为此只要证明
\[
\cos(\sin t)>\sin(\cos t)\quad\left(\text{当 }t\in\left(0,\frac\pi2\right)\text{ 时}\right).
\tag{3}
\]
但已知 $\left(0,\frac\pi2\right]$ 上有 $\sin x<x,\cos x$ 严减，故 $\sin(\cos t)<\cos t<\cos(\sin t)$. 证毕.

\noindent\textbf{例 4.3.14}\quad 函数 $f(x)$ 在闭区间 $[0,1]$ 上有连续的一阶导数，证明：
\[
\int_0^1|f(x)|\,dx
\le\max\left\{\int_0^1|f'(x)|\,dx,\left|\int_0^1 f(x)\,dx\right|\right\}.
\]
（国外赛题）

\noindent\textbf{证}\quad $1^\circ$ 若
\[
\left|\int_0^1 f(x)\,dx\right|=\int_0^1|f(x)|\,dx,
\]
问题自明.

\noindent $2^\circ$ 若
\[
\left|\int_0^1 f(x)\,dx\right|<\int_0^1|f(x)|\,dx,
\]
则 $f(x)$ 在 $[0,1]$ 上变号，由 $f$ 的连续性，知 $\exists x_0\in(0,1)$ 使得 $f(x_0)=0$，于是
\[
|f(x)|=|f(x)-f(x_0)|
=\left|\int_{x_0}^x f'(x)\,dx\right|
\le\int_{x_0}^x|f'(x)|\,dx
\le\int_0^1|f'(x)|\,dx.
\]
取积分知 $\int_0^1|f(x)|\,dx\le\int_0^1|f'(x)|\,dx$. 原不等式获证.

\noindent\textbf{☆例 4.3.15}\quad 函数 $f(x)$ 在 $[0,1]$ 上单调不增，证明：对于任何 $\alpha\in(0,1)$，
\[
\int_0^\alpha f(x)\,dx\ge\alpha\int_0^1 f(x)\,dx.
\tag{1}
\]
（国外赛题，华中科技大学）

\noindent\textbf{证}\quad 式 (1) 即
\[
\int_0^\alpha f(x)\,dx\ge
\alpha\int_0^\alpha f(x)\,dx+\alpha\int_\alpha^1 f(x)\,dx,
\]
亦即
\[
(1-\alpha)\int_0^\alpha f(x)\,dx\ge
\alpha\int_\alpha^1 f(x)\,dx
\quad\text{或}\quad
\frac1\alpha\int_0^\alpha f(x)\,dx
\ge\frac1{1-\alpha}\int_\alpha^1 f(x)\,dx.
\]
但 $f(x)$ 单调不增，故
\[
\frac1\alpha\int_0^\alpha f(x)\,dx
\ge f(\alpha)\ge
\frac1{1-\alpha}\int_\alpha^1 f(x)\,dx.
\]

\noindent\textbf{**☆练习 1}\quad 设函数 $f(x)\ge0$，在 $[0,1]$ 上连续递减，试证若 $0<\alpha<\beta<1$，则
\[
\int_0^\alpha f(x)\,dx\ge
\frac{\alpha}{\beta-\alpha}\int_\alpha^\beta f(x)\,dx
>\frac{\alpha}{\beta}\int_0^\beta f(x)\,dx.
\]

\noindent\textbf{证}\quad $[0,\alpha]$ 上 $f(x)\ge f(\alpha)$，故
\[
\int_0^\alpha f(x)\,dx
\ge\int_0^\alpha f(\alpha)\,dx
=\alpha f(\alpha)
\ge\frac{\alpha}{\beta-\alpha}\int_\alpha^\beta f(x)\,dx
>\frac{\alpha}{\beta}\int_0^\beta f(x)\,dx.
\]

\noindent\textbf{**☆练习 2}\quad 设 $f(x)$ 是 $[a,b]$ 上连续递增函数，则如下不等式成立：
\[
\int_a^b xf(x)\,dx\ge\frac{a+b}{2}\int_a^b f(x)\,dx.
\tag{1}
\]
（上海交通大学）

\noindent\textbf{证}\quad 记 $c=\frac{a+b}{2}$，式 (1) 等价于
\[
I=\int_a^b(c-x)f(x)\,dx\le0.
\tag{2}
\]
\[
I=\int_a^b(c-x)f(x)\,dx
=\int_a^c(c-x)f(x)\,dx+\int_c^b(c-x)f(x)\,dx
=I_1+I_2.
\]
令 $h=\frac{b-a}{2}$，则
\[
I_1=\int_a^c(c-x)f(x)\,dx
\overset{x=c-t}{=}
\int_0^h t f(c-t)\,dt,
\]
\[
I_2=\int_c^b(c-x)f(x)\,dx
\overset{x=c+t}{=}
\int_0^h -t f(c+t)\,dt.
\]
故
\[
I=I_1+I_2=\int_0^h t[f(c-t)-f(c+t)]\,dt\le0\quad(\text{因 }f\text{ 递增}).
\]
式 (2) 获证，亦即式 (1) 成立.

\noindent\textbf{例 4.3.16}\quad 设 $a,b>0,f(x)\ge0$，且 $f(x)$ 在 $[a,b]$ 上可积，$\int_a^b xf(x)\,dx=0$. 试证：
\[
\int_a^b x^2 f(x)\,dx\le ab\int_a^b f(x)\,dx.
\]

\noindent\textbf{提示}\quad
\[
\int_a^b[(x+a)(b-x)]f(x)\,dx\ge0.
\]

\noindent\textbf{例 4.3.17}\quad 设 $f(x)$ 与 $g(x)$ 为拟序的，即 $\forall x_1,x_2$，有
\[
\left(f(x_1)-f(x_2)\right)\left(g(x_1)-g(x_2)\right)\ge0.
\tag{1}
\]
试证：
\[
\int_a^b f(x)\,dx\int_a^b g(x)\,dx
\le(b-a)\int_a^b f(x)g(x)\,dx.
\tag{2}
\]

\noindent\textbf{提示}\quad 在式 (1) 里先对 $x_1$ 于 $[a,b]$ 上取积分，然后对 $x_2$ 在 $[a,b]$ 上取积分.

\noindent\textbf{注}\quad 若将式 (1) 里的不等号反向，则 $f,g$ 称为反序的. 若 $f,g$ 反序，则式 (2) 不等号反向.

该题是著名的 Chebyshev 定理的积分形式. 其级数形式见下一节习题 4.4.10.

\noindent\textbf{例 4.3.18}\quad 设 $f'(x)$ 在 $[a,b]$ 上连续，试证
\[
\max_{a\le x\le b}|f(x)|
\le\left|\frac1{b-a}\int_a^b f(x)\,dx\right|
+\int_a^b|f'(x)|\,dx.
\]

\noindent\textbf{提示}\quad $\exists\xi,x_0\in[a,b]$ 使得
\[
|f(\xi)|=\left|\frac1{b-a}\int_a^b f(x)\,dx\right|,
\]
\[
\max_{a\le x\le b}|f(x)|
=|f(x_0)|
=\left|\int_\xi^{x_0} f'(t)\,dt+f(\xi)\right|.
\]

\noindent\textbf{c. 在不等式两端取变限积分证明新的不等式}

\noindent\textbf{☆例 4.3.19}\quad 证明：$x>0$ 时，
\[
x-\frac{x^3}{6}<\sin x<x-\frac{x^3}{6}+\frac{x^5}{120}.
\]
（吉林大学）

\noindent\textbf{证}\quad 已知 $\cos x\le1$（$x>0$，只有 $x=2n\pi$ 时等号才成立）. 在此式两端同时取 $[0,x]$ 上的积分，得 $\sin x<x\ (x>0)$.

再次取 $[0,x]$ 上的积分，得 $1-\cos x<\frac{x^2}{2}\ (x>0)$.

第三次取 $[0,x]$ 上的积分，得 $x-\sin x<\frac{x^3}{6}$，即 $x-\frac{x^3}{6}<\sin x\ (x>0)$.

继续在 $[0,x]$ 上积分两次，可得 $\sin x<x-\frac{x^3}{6}+\frac{x^5}{120}$. 证毕.

\noindent\textbf{注}\quad 上面是用积分法证明，对偶地还可用微分法证明，如用例 3.4.3 中的方法.

\noindent\textbf{三、另一些问题}

以上我们讨论了关于定积分的几类较典型的问题. 但是问题是多种多样、错综复杂的，不是几种类型所能概括的. 本段主要讨论一些灵活多变、带综合性的问题，以作上述内容的补充.

\noindent\textbf{例 4.3.20}\quad 设函数 $f(x)$ 二次可微，证明在 $(a,b)$ 内存在一点 $\xi$，使得
\[
f''(\xi)=\frac{24}{(b-a)^3}\int_a^b\left(f(x)-f\left(\frac{a+b}{2}\right)\right)\,dx.
\tag{1}
\]

\noindent\textbf{证}\quad 记 $x_0=\frac{a+b}{2}$，将被积函数在 $x=x_0$ 处按 Taylor 公式展开，
\[
f(x)-f(x_0)=(x-x_0)f'(x_0)+\frac{(x-x_0)^2}{2}f''(\eta),
\tag{2}
\]
其中 $\eta$ 在 $x_0$ 与 $x$ 之间. 在区间 $[a,b]$ 上取积分，注意式 (2) 右端第一项的积分为零. 因此
\[
\int_a^b(f(x)-f(x_0))\,dx
=\frac12\int_a^b(x-x_0)^2f''(\eta)\,dx.
\tag{3}
\]
此式右端中 $f''(\eta)$ 虽然不一定连续，但导数具有介值性（例 3.2.24），因而积分第一中值定理仍然成立. 故 $\exists\xi\in(a,b)$，使得
\[
\int_a^b(x-x_0)^2f''(\eta)\,dx
=f''(\xi)\int_a^b(x-x_0)^2\,dx
=\frac{(b-a)^3}{12}f''(\xi).
\]
代入式 (3) 即得欲证式 (1).

\noindent\textbf{☆例 4.3.21}\quad 设 $f(x)$ 在 $(0,+\infty)$ 中单调不减，且 $\forall A>0$，$f(x)$ 在 $[0,A]$ 上可积. 试证 $\lim_{x\to+\infty}\frac1x\int_0^x f(t)\,dt=C$ 的充要条件是 $\lim_{x\to+\infty}f(x)=C$（其中 $C$ 为有限数或 $+\infty$）.（武汉大学）

\noindent\textbf{证}\quad $1^\circ$（充分性）\textbf{方法 I}（拟合法）\quad 因 $\lim_{x\to+\infty}\frac1x\int_0^x C\,dt=C$，因此只要证明
\[
\lim_{x\to+\infty}\frac1x\int_0^x(f(t)-C)\,dt=0.
\tag{1}
\]
根据已知条件，$\lim_{x\to+\infty}f(x)=C,f\nearrow$，因此 $\forall\varepsilon>0$，$\exists A>0$，当 $x>A$ 时，有
\[
0\le C-f(x)<\frac{\varepsilon}{2}.
\]
于是
\[
0\le\frac1x\int_0^x(C-f(t))\,dt
=\frac1x\int_0^A(C-f(t))\,dt+\frac{\varepsilon}{2x}\int_A^x dt.
\]
将 $A$ 固定，令 $x\to+\infty$，当 $x$ 充分大时，就有 $\frac1x\int_0^A(C-f(t))\,dt<\frac{\varepsilon}{2}$. 于是
\[
0\le\frac1x\int_0^x(C-f(t))\,dt
<\frac{\varepsilon}{2}+\frac{\varepsilon}{2}=\varepsilon.
\]
式 (1) 获证.

\noindent\textbf{方法 II}\quad 已知 $\lim_{x\to+\infty}f(x)=C,f\nearrow$，故当 $x>A$ 时，有 $f(A)\le f(x)\le C$.
\[
\frac1x\int_0^A f(t)\,dt+\frac1x\int_A^x f(A)\,dt
\le\frac1x\int_0^x f(t)\,dt
\le\frac1x\int_0^x C\,dt=C.
\tag{2}
\]
固定 $A$，令 $x\to+\infty$ 得
\[
f(A)\le\underline{\lim_{x\to+\infty}}\frac1x\int_0^x f(t)\,dt
\le\overline{\lim_{x\to+\infty}}\frac1x\int_0^x f(t)\,dt
\le C.
\]
在此式中令 $A\to+\infty$（而 $\lim_{A\to+\infty}f(A)=C$），由两边夹法则得 $\lim_{x\to+\infty}\frac1x\int_0^x f(t)\,dt=C$.

\noindent $2^\circ$（必要性）现在已知 $\lim_{x\to+\infty}\frac1x\int_0^x f(t)\,dt=C$，要证 $\lim_{x\to+\infty}f(x)=C$. 由于 $f(x)\nearrow$，有
\[
\lim_{x\to+\infty}f(x)=\sup_{0\le x<+\infty} f(x)\stackrel{\text{记}}{=}M,
\tag{3}
\]
（故问题归结为证明 $M=C$.）由式 (3) 知 $\forall x>0$，有
\[
\frac1x\int_0^x f(t)\,dt\le\frac1x\int_0^x M\,dt=M.
\]
因为
\[
C=\lim_{x\to+\infty}\frac1x\int_0^x f(t)\,dt\le M,
\tag{4}
\]
可见当 $C=+\infty$ 时，$M=+\infty=C$.

（下面只需从 $0\le C<+\infty$ 导出 $C\ge M$，则有 $C=M$.）按确界定义，$\forall M_n:0<M_n<M$，$\exists x_n>0$，使得 $M_n\le f(x_n)\le M$. 因此，当 $x>x_n$ 时，有 $M_n\le f(x_n)\le f(x)\le M$. 于是
\[
\frac1x\int_0^x f(t)\,dt
=\frac1x\int_0^{x_n} f(t)\,dt+\frac1x\int_{x_n}^x f(t)\,dt
\ge\frac1x\int_0^{x_n} f(t)\,dt+\frac1x M_n(x-x_n).
\]
固定 $n$，令 $x\to+\infty$，得 $C\ge M_n$. 注意，不论 $C$ 是有限或无穷，都可使 $M_n\le M$，且 $M_n\to M\ (n\to+\infty)$. 故在不等式 $C\ge M_n$ 里，令 $n\to+\infty$ 取极限，即得 $C\ge M$. 但由式 (4)：$C\le M$. 故等式 $C=M$ 获证.

\noindent\textbf{练习}\quad 设函数 $f(x)$ 在任何有限区间上可积，且 $\lim_{x\to+\infty}f(x)=C$. 求证：
\[
\lim_{x\to+\infty}\frac1x\int_0^x f(t)\,dt=C.
\]
（武汉大学）

\noindent\textbf{提示}\quad 类似地：$\forall\varepsilon>0,\exists A>0$，当 $x>A$ 时，有 $|C-f(t)|<\frac{\varepsilon}{2}$. 于是
\[
0\le\left|\frac1x\int_0^x(C-f(t))\,dt\right|
\le\frac1x\int_0^A(|C|+|f(t)|)\,dt+\frac{\varepsilon}{2x}\int_A^x dt.
\]
由 $f(x)$ 在 $[0,A]$ 上可积，存在 $M_A>0$，使得 $[0,A]$ 上 $|f(x)|\le M_A$. 固定 $A$，令 $x\to+\infty$，当 $x$ 充分大时，可使
\[
\frac1x\int_0^A(|C|+|f(t)|)\,dt<\frac{\varepsilon}{2},
\]
从而
\[
\frac1x\int_0^A(C-f(t))\,dt
\le\frac1x\int_0^A(|C|+|f(t)|)\,dt<\frac{\varepsilon}{2}.
\]
于是
\[
0\le\left|\frac1x\int_0^x(C-f(t))\,dt\right|
<\frac{\varepsilon}{2}+\frac{\varepsilon}{2}=\varepsilon.
\]

\noindent\textbf{☆注}\quad 若 $f(x)$ 连续，则 $F(x)=\int_0^x f(t)\,dt$ 可导，且 $F'(x)=f(x)$.

若 $f(x)$ 在任一有限区间上可积，只能断言 $F(x)$ 连续，不能断言 $F(x)$ 一定可导. 因此这里（包括上例）不能应用 L'Hospital 法则.

\noindent\textbf{*例 4.3.22}\quad 设 $f(x)$ 为 $n$ 次多项式，且
\[
\int_0^1 x^k f(x)\,dx=0\quad(k=1,2,\cdots,n).
\tag{1}
\]
试证：1) $f(0)=(n+1)^2\int_0^1 f(x)\,dx$；
\[
\text{2) }\quad
\int_0^1 f^2(x)\,dx
=\left[(n+1)\int_0^1 f(x)\,dx\right]^2.
\]
（东北师范大学）

\noindent\textbf{证}\quad 既然 $f(x)$ 为 $n$ 次多项式，不妨设
\[
f(x)=a_0+a_1x+\cdots+a_nx^n.
\tag{2}
\]
于是
\[
f(0)=a_0.
\tag{3}
\]
利用已知条件 (1)，得
\[
\int_0^1 f^2(x)\,dx
=\int_0^1(a_0+a_1x+\cdots+a_nx^n)f(x)\,dx
=a_0\int_0^1 f(x)\,dx.
\tag{4}
\]
可见，只要证明了
\[
a_0=(n+1)^2\int_0^1 f(x)\,dx,
\tag{5}
\]
代入式 (3)、(4)，即得结论 1) 与 2). 下面我们来证明式 (5).

由式 (2)，
\[
\int_0^1 x^k f(x)\,dx
=\frac{a_0}{k+1}+\frac{a_1}{k+2}+\cdots+\frac{a_n}{n+k+1}\quad(\text{通分后})
\]
\[
=\frac{Q(k)}{(k+1)(k+2)\cdots(k+n+1)}.
\tag{6}
\]
其中 $Q(k)$ 是关于 $k$ 的 $n$ 次多项式. 根据式 (1)、(6) 知 $Q(k)=0\ (k=1,2,\cdots,n)$，因此
\[
Q(k)=c(k-1)(k-2)\cdots(k-n)\quad(\text{其中 }c\text{ 为某一常数}).
\tag{7}
\]
将 (7) 代入 (6) 的后一等式，同乘 $k+1$，并令 $k=-1$，则得
\[
a_0=(-1)^n(n+1)c;
\tag{8}
\]
将 (7) 代入 (6)，并令 $k=0$，可得
\[
\int_0^1 f(x)\,dx=\frac{(-1)^n}{n+1}c.
\tag{9}
\]
在 (8)、(9) 中消去 $c$，即得 $a_0=(n+1)^2\int_0^1 f(x)\,dx$. 证毕.

\noindent\textbf{用积分解函数方程}

\noindent\textbf{例 4.3.23}\quad 设 $f(x)$ 在任意有限区间上可积且满足方程
\[
f(x+y)=f(x)+f(y).
\tag{1}
\]
试证：$f(x)=ax$，其中 $a=f(1)$.

\noindent\textbf{证}\quad 要证 $f(x)=ax$，当 $x\ne0$ 时即要证 $\frac{f(x)}x=$ 常数，或 $\forall x,y\ne0$，$\frac{f(x)}x=\frac{f(y)}y$，亦即
\[
f(x)y=xf(y).
\tag{2}
\]
为此在已知方程 $f(t+y)=f(t)+f(y)$ 两边对 $t$ 取积分：
\[
\int_0^x f(t+y)\,dt
=\int_0^x f(t)\,dt+f(y)\cdot x.
\tag{3}
\]
但
\[
\int_0^x f(t+y)\,dt
=\int_y^{x+y}f(u)\,du
=\int_0^{x+y}f(t)\,dt-\int_0^y f(t)\,dt,
\]
故
\[
xf(y)=\int_0^{x+y}f(t)\,dt-\int_0^y f(t)\,dt-\int_0^x f(t)\,dt.
\]
此式右端 $x,y$ 以对称的形式出现，$x,y$ 互换知 $xf(y)=yf(x)$，从而
\[
f(x)=ax\quad(\text{当 }x\ne0\text{ 时}).
\tag{4}
\]
在 (1) 中令 $x=0,y=1$，得 $f(0)=0$. 可见式 (4) 对于 $x=0$ 也成立. 最后，在 (4) 中令 $x=1$，可得 $a=f(1)$. 证毕.

\noindent\textbf{new 练习}\quad 设 $f(x)$ 在 $[a,b]$ 上可积，且 $\forall\alpha,\beta:a\le\alpha<\beta\le b$，有
\[
\frac1{\beta-\alpha}\int_\alpha^\beta f(t)\,dt
=\frac{f(\beta)+f(\alpha)}2
\]
（任一内闭区间：积分均值等于端点值的平均）. 求证：$f(x)=kx+c$（其中 $k,c$ 为常数）.

\noindent\textbf{提示}\quad （$\forall x\in[a,b]$）. 根据已知条件可写出如下两式（两式相加即得）：
\[
\int_a^x f(t)\,dt=\frac{x-a}{2}[f(a)+f(x)],\quad
\int_x^b f(t)\,dt=\frac{b-x}{2}[f(x)+f(b)].
\]
\[
f(x)=kx+c,\quad\text{其中 }k=\frac{f(b)-f(a)}{b-a},\quad
c=\frac{af(a)-bf(b)}{b-a}+\frac{2}{b-a}\int_a^b f(t)\,dt.
\]

\noindent\textbf{函数线性相关的充要条件}

\noindent\textbf{定义}\quad 设 $f_1(x),f_2(x),\cdots,f_n(x)$ 在 $[a,b]$ 上有定义，那么，当且仅当存在不全为零的常数 $c_1,c_2,\cdots,c_n$，使得
\[
\sum_{i=1}^n c_i f_i(x)=0\quad(\forall x\in[a,b])
\]
成立时，函数 $f_1(x),f_2(x),\cdots,f_n(x)$ 称为在 $[a,b]$ 上线性相关，否则称为线性无关.

\noindent\textbf{**例 4.3.24}\quad 设 $f_1(x),f_2(x),\cdots,f_n(x)$ 在 $[a,b]$ 上连续，记
\[
a_{ij}=\int_a^b f_i(x)f_j(x)\,dx\quad(i,j=1,2,\cdots,n),
\]
试证：函数 $f_1(x),f_2(x),\cdots,f_n(x)$ 在 $[a,b]$ 上线性相关的充要条件是行列式
\[
\operatorname{Det}(a_{ij})=
\begin{vmatrix}
a_{11}&a_{12}&\cdots&a_{1n}\\
a_{21}&a_{22}&\cdots&a_{2n}\\
\vdots&\vdots&&\vdots\\
a_{n1}&a_{n2}&\cdots&a_{nn}
\end{vmatrix}=0.
\]

\noindent\textbf{证}\quad 若 $f_1(x),f_2(x),\cdots,f_n(x)$ 线性相关，则按定义，$\exists c_1,c_2,\cdots,c_n$（不全为零），使得
\[
g(x)=\sum_{i=1}^n c_i f_i(x)\equiv0\quad(\forall x\in[a,b]).
\tag{1}
\]
由此
\[
\int_a^b g(x)f_j(x)\,dx
=\sum_{i=1}^n c_i\int_a^b f_i(x)f_j(x)\,dx
=\sum_{i=1}^n a_{ij}c_i=0\quad(j=1,2,\cdots,n),
\tag{2}
\]
即关于 $c_i(i=1,2,\cdots,n)$ 的方程组 (2) 有非零解，故系数行列式
\[
\operatorname{Det}(a_{ij})=0.
\tag{3}
\]
反之，若 (3) 成立，则 $\exists c_1,c_2,\cdots,c_n$（不全为零）使得式 (2) 成立，从而
\[
\int_a^b g^2(x)\,dx
=\int_a^b g(x)\sum_{j=1}^n c_j f_j(x)\,dx
=\sum_{j=1}^n c_j\int_a^b g(x)f_j(x)\,dx=0.
\]
由练习 4.2.6 知 $g(x)=\sum_{i=1}^n c_i f_i(x)\equiv0\ (\forall x\in[a,b])$. 因此 $f_1(x),f_2(x),\cdots,f_n(x)$ 在 $[a,b]$ 上线性相关.

\noindent 利用特征函数的积分表示区间的长度.

\noindent\textbf{**例 4.3.25}\quad 设 $[\alpha_i,\beta_i]\ (i=1,2,\cdots,n)$ 为 $[0,1]$ 上 $n$ 个区间，且 $[0,1]$ 上每个点至少属于这些区间里的 $q$ 个. 证明这些区间里至少有一个，其长度 $\ge\frac{q}{n}$.

\noindent\textbf{证}\quad 用 $f_i(x)$ 表示第 $i$ 个区间 $[\alpha_i,\beta_i]$ 的特征函数. 即
\[
f_i(x)=
\begin{cases}
1, & \text{当 }x\in[\alpha_i,\beta_i]\text{ 时},\\
0, & \text{当 }x\notin[\alpha_i,\beta_i]\text{ 时}.
\end{cases}
\]
则 $\int_0^1 f_i(x)\,dx=\beta_i-\alpha_i$（第 $i$ 个区间的长度），且函数 $f(x)=\sum_{i=1}^n f_i(x)=k$（当 $x$ 属于 $[\alpha_i,\beta_i]\ (i=1,2,\cdots,n)$ 中的 $k$ 个时）. 于是已知条件可表达为
\[
f(x)=k\ge q.
\tag{1}
\]
此外，
\[
\int_0^1 f(x)\,dx
=\int_0^1\sum_{i=1}^n f_i(x)\,dx
=\sum_{i=1}^n\int_0^1 f_i(x)\,dx
=\sum_{i=1}^n\int_{\alpha_i}^{\beta_i}dx
=\sum_{i=1}^n(\beta_i-\alpha_i)
\tag{2}
\]
为区间 $[\alpha_i,\beta_i]_{i=1}^n$ 的总长度. 由式 (1) 知总长度
\[
\sum_{i=1}^n(\beta_i-\alpha_i)=\int_0^1 f(x)\,dx\ge q.
\tag{3}
\]
假若每个区间的长度 $\beta_i-\alpha_i<\frac{q}{n}$，则这些区间的总长度
\[
\sum_{i=1}^n(\beta_i-\alpha_i)<n\cdot\frac{q}{n}=q,
\]
与式 (3) 矛盾. 因此，$\exists[\alpha_i,\beta_i]$ 使得 $\beta_i-\alpha_i\ge\frac{q}{n}$.

\noindent\textbf{凸函数的积分性质}

\noindent\textbf{例 4.3.26}\quad （Hadamard 定理）设 $f(x)$ 是 $[a,b]$ 上连续的凸函数. 试证：$\forall x_1,x_2\in[a,b],x_1<x_2$，有
\[
f\left(\frac{x_1+x_2}{2}\right)
\le\frac1{x_2-x_1}\int_{x_1}^{x_2}f(t)\,dt
\le\frac{f(x_1)+f(x_2)}2.
\]
（中南大学）

\noindent\textbf{证}\quad 令 $t=x_1+\lambda(x_2-x_1),\lambda\in(0,1)$，则
\[
\frac1{x_2-x_1}\int_{x_1}^{x_2}f(t)\,dt
=\int_0^1 f[x_1+\lambda(x_2-x_1)]\,d\lambda.
\tag{1}
\]
同理，令 $t=x_2-\lambda(x_2-x_1)$，亦有
\[
\frac1{x_2-x_1}\int_{x_1}^{x_2}f(t)\,dt
=\int_0^1 f[x_2-\lambda(x_2-x_1)]\,d\lambda.
\]
从而
\[
\begin{aligned}
\frac1{x_2-x_1}\int_{x_1}^{x_2}f(t)\,dt
&=\frac12\int_0^1\{f[x_1+\lambda(x_2-x_1)]
+f[x_2-\lambda(x_2-x_1)]\}\,d\lambda.
\end{aligned}
\tag{2}
\]
注意 $x_1+\lambda(x_2-x_1)$ 与 $x_2-\lambda(x_2-x_1)$ 关于中点 $\frac{x_1+x_2}{2}$ 对称. 由于 $f(x)$ 是凸函数，
\[
\frac12\{f[x_1+\lambda(x_2-x_1)]
+f[x_2-\lambda(x_2-x_1)]\}
\ge f\left(\frac{x_1+x_2}{2}\right).
\]
故由 (2) 得
\[
\frac1{x_2-x_1}\int_{x_1}^{x_2}f(t)\,dt
\ge f\left(\frac{x_1+x_2}{2}\right).
\]
另外，由 (1)，应用 $f(x)$ 的凸性，
\[
\begin{aligned}
\frac1{x_2-x_1}\int_{x_1}^{x_2}f(t)\,dt
&=\int_0^1 f[\lambda x_2+(1-\lambda)x_1]\,d\lambda\\
&\le\int_0^1[\lambda f(x_2)+(1-\lambda)f(x_1)]\,d\lambda\\
&=f(x_2)\cdot\frac{\lambda^2}{2}\bigg|_0^1
+f(x_1)\cdot\left[-\frac{(1-\lambda)^2}{2}\right]\bigg|_0^1
=\frac{f(x_1)+f(x_2)}2.
\end{aligned}
\]

\noindent\textbf{例 4.3.27}\quad 设 $f(x)$ 是 $[0,+\infty)$ 上的凸函数，求证
\[
F(x)=\frac1x\int_0^x f(t)\,dt
\tag{1}
\]
为 $(0,+\infty)$ 上的凸函数.

\noindent\textbf{证}\quad $f(x)$ 为 $[0,+\infty)$ 上的凸函数，因此它在 $(0,+\infty)$ 内连续（$\S3.4$ 定理 3 的推论 5），$f(x)$ 在 $[0,x]$ 上有界（当 $x\ge0$ 时）（例 3.4.9）. 由此知积分 (1) 有意义. 注意到：$\forall x>0$，令 $u=\frac{t}{x}$ 时，
\[
F(x)=\frac1x\int_0^x f(t)\,dt
=\int_0^1 f\left(x\cdot\frac{t}{x}\right)\,d\frac{t}{x}
=\int_0^1 f(xu)\,du.
\tag{2}
\]
则 $\forall\lambda\in(0,1),\forall x_1,x_2>0$，恒有
\[
\begin{aligned}
F[\lambda x_1+(1-\lambda)x_2]
&=\int_0^1 f\{[\lambda x_1+(1-\lambda)x_2]u\}\,du\quad(\text{因为 (2)})\\
&=\int_0^1 f[\lambda x_1u+(1-\lambda)x_2u]\,du\\
&\le\int_0^1[\lambda f(x_1u)+(1-\lambda)f(x_2u)]\,du\quad(f\text{ 的凸性})\\
&=\lambda F(x_1)+(1-\lambda)F(x_2),
\end{aligned}
\]
所以 $F$ 是 $(0,+\infty)$ 上的凸函数.

\noindent\textbf{例 4.3.28}\quad 设函数 $g(x)$ 在 $[a,b]$ 上递增，试证：$\forall c\in(a,b)$，函数
\[
f(x)=\int_c^x g(x)\,dx
\]
为凸函数.

\noindent\textbf{证}\quad 因 $g(x)$ 递增，积分有意义，且 $\forall x_1<x_2<x_3$，
\[
\frac{f(x_2)-f(x_1)}{x_2-x_1}
=\frac1{x_2-x_1}\int_{x_1}^{x_2}g(x)\,dx
\le g(x_2)
\le\frac1{x_3-x_2}\int_{x_2}^{x_3}g(x)\,dx
=\frac{f(x_3)-f(x_2)}{x_3-x_2}.
\]
由 $\S3.4$ 定理 3，知 $f(x)$ 为凸函数.

\noindent\textbf{例 4.3.29}\quad 设 $f(x)$ 为 $[a,b]$ 上的凸函数，试证：$\forall c,x\in(a,b)$ 有
\[
f(x)-f(c)=\int_c^x f'_{-}(t)\,dt
=\int_c^x f'_{+}(t)\,dt.
\tag{1}
\]

\noindent\textbf{证}\quad 因 $f(x)$ 为凸函数，由 $\S3.4$ 定理 3 的推论 4，$f'_{-}(t),f'_{+}(t)$ 存在且递增（$t\in(a,b)$），故式 (1) 中的积分有意义. 对 $[c,x]$ 任作一分割 $c=x_0<x_1<x_2<\cdots<x_n=x$，有
\[
f(x)-f(c)=\sum_{i=1}^n[f(x_i)-f(x_{i-1})].
\tag{2}
\]
但参看 $\S3.4$ 定理 4（用该定理的证法），有
\[
f'_{-}(x_{i-1})(x_i-x_{i-1})
\le f(x_i)-f(x_{i-1})
\le f'_{-}(x_i)(x_i-x_{i-1}).
\]
于是由式 (2) 知
\[
\sum_{i=1}^n f'_{-}(x_{i-1})(x_i-x_{i-1})
\le f(x)-f(c)
\le\sum_{i=1}^n f'_{-}(x_i)(x_i-x_{i-1}).
\]
将分割无限分细，令 $\lambda=\max(x_i-x_{i-1})\to0$，取极限可知
\[
\int_c^x f'_{-}(x)\,dx=f(x)-f(c).
\]
同理，有
\[
\int_c^x f'_{+}(x)\,dx=f(x)-f(c).
\]
证毕.

\noindent\textbf{例 4.3.30}\quad 设 $f(x),p(x)$ 在区间 $[a,b]$ 上连续，$p(x)\ge0$，$\int_a^b p(x)\,dx>0$，且 $m\le f(x)\le M$；$\varphi(x)$ 在 $[m,M]$ 上有定义，并有二阶导数，$\varphi''(x)>0$. 试证：
\[
\varphi\left(\frac{\int_a^b p(x)f(x)\,dx}{\int_a^b p(x)\,dx}\right)
\le
\frac{\int_a^b p(x)\varphi(f(x))\,dx}{\int_a^b p(x)\,dx}.
\]
（北京理工大学）

\noindent\textbf{证 I}\quad （利用积分和.）将 $[a,b]$ $n$ 等分，记 $x_i=a+\frac{i}{n}(b-a),p_i=p(x_i)$，
\[
f_i=f(x_i)\quad(i=1,2,\cdots,n).
\]
因为 $\varphi''(x)>0,\varphi(x)$ 为凸函数，由 $\S3.4$ 定理 6，知
\[
\varphi\left(\frac{p_1f_1+p_2f_2+\cdots+p_nf_n}{p_1+p_2+\cdots+p_n}\right)
\le
\frac{p_1\varphi(f_1)+p_2\varphi(f_2)+\cdots+p_n\varphi(f_n)}
{p_1+p_2+\cdots+p_n},
\]
即
\[
\varphi\left(
\frac{\sum p(x_i)f(x_i)\frac{b-a}{n}}
{\sum p(x_i)\frac{b-a}{n}}
\right)
\le
\frac{\sum p(x_i)\varphi(f(x_i))\frac{b-a}{n}}
{\sum p(x_i)\frac{b-a}{n}}.
\]
令 $n\to\infty$，取极限，便得欲证的不等式.

\noindent\textbf{证 II}\quad （利用 Taylor 公式，参考例 3.3.3 之证法.）记
\[
x_0=\frac{\int_a^b p(x)f(x)\,dx}{\int_a^b p(x)\,dx},
\tag{1}
\]
则
\[
\varphi(y)-\varphi(x_0)
=\varphi'(x_0)(y-x_0)+\frac12\varphi''(\xi)(y-x_0)^2.
\]
注意 $\varphi''(\xi)>0$，所以
\[
\varphi(y)-\varphi(x_0)>\varphi'(x_0)(y-x_0).
\]
在此式中令 $y=f(x)$，然后两边同乘 $\frac{p(x)}{\int_a^b p(x)\,dx}$；再在 $[a,b]$ 上取积分，并注意式 (1)，得
\[
\frac{\int_a^b p(x)\varphi(f(x))\,dx}{\int_a^b p(x)\,dx}
-\varphi(x_0)\frac{\int_a^b p(x)\,dx}{\int_a^b p(x)\,dx}
>
\frac{\int_a^b p(x)(f(x)-x_0)\,dx}{\int_a^b p(x)\,dx}\cdot\varphi'(x_0)=0,
\]
所以
\[
\varphi(x_0)<\frac{\int_a^b p(x)\varphi(f(x))\,dx}{\int_a^b p(x)\,dx}.
\]
证毕.

\noindent\textbf{各种技巧的灵活应用}

\noindent\textbf{单调性的妙用}

\noindent\textbf{☆例 4.3.31}\quad 设函数 $f(x)\ge0$，在 $[0,1]$ 上连续减小，$0<\alpha<\beta<1$，试证：
\[
\int_0^\alpha f(x)\,dx
\ge\frac{\alpha}{\beta-\alpha}\int_\alpha^\beta f(x)\,dx
>\frac{\alpha}{\beta}\int_\alpha^\beta f(x)\,dx.
\]

\noindent\textbf{证}\quad 在 $[0,\alpha]$ 上 $f(x)\ge f(\alpha)$，故
\[
\frac1\alpha\int_0^\alpha f(x)\,dx
\ge\frac1\alpha\int_0^\alpha f(\alpha)\,dx
=f(\alpha)
\ge\frac1{\beta-\alpha}\int_\alpha^\beta f(x)\,dx
>\frac1\beta\int_\alpha^\beta f(x)\,dx,
\]
待证不等式得证.

\noindent\textbf{利用对称性}

\noindent\textbf{☆例 4.3.32}\quad 设 $f(x)$ 为 $[a,b]$ 上的连续递增函数，则成立不等式：
\[
\int_a^b xf(x)\,dx
\ge\frac{a+b}{2}\int_a^b f(x)\,dx.
\tag{1}
\]
（上海交通大学）

\noindent\textbf{证}\quad (1) 式等价于
\[
\int_a^b(x-c)f(x)\,dx\ge0,
\tag{2}
\]
其中 $c=\frac{a+b}{2}$ 为 $[a,b]$ 的中点，注意 $g(x)=x-c$ 关于 $x=c$ 有点对称性，记 $h=\frac{b-a}{2}$，令 $x-c=t$，则
\[
\int_a^b(x-c)f(x)\,dx
=\int_{-h}^h t f(c+t)\,dt
=\left(\int_0^h+\int_{-h}^0\right) t f(c+t)\,dt,
\tag{3}
\]
其中
\[
\int_{-h}^0 t f(c+t)\,dt
\overset{\text{令 }\tau=-t}{=}
\int_h^0(-\tau)f(c-\tau)\,d(-\tau)
=-\int_0^h\tau f(c-\tau)\,d\tau
\overset{\text{将 }\tau\text{ 仍记作 }t}{=}
-\int_0^h t f(c-t)\,dt.
\]
于是
\[
\left(\int_0^h+\int_{-h}^0\right)t f(c+t)\,dt
=\int_0^h t[f(c+t)-f(c-t)]\,dt\ge0.
\]
（因 $f\nearrow$，两对称点上 $f(c+t)-f(c-t)\ge0$.）

\noindent\textbf{求导变成微分方程}

\noindent\textbf{☆例 4.3.33}\quad 设 $f(x)$ 在 $[0,+\infty)$ 上可微，且满足
\[
\int_0^x f(t)\,dt=\frac{x}{3}\int_0^x f(t)\,dt,\quad x>0,
\]
求 $f(x)$.（北京大学）

\noindent\textbf{解}\quad 原式表明当 $x>0,x\ne3$ 时，$\int_0^x f(t)\,dt=0$. 因 $\int_0^x f(t)\,dt$ 关于 $x$ 连续，故在
\[
[0,+\infty)\text{ 上 }\int_0^x f(t)\,dt=0.
\]
若 $f$ 不变号，由此即知 $f(x)\equiv0$.

将原式两边同时求导，可得
\[
(3-x)f'(x)=2f(x).
\]
（分离变量）积分得
\[
f(x)=\frac{c}{(3-x)^2}\quad(c\text{ 为任意常数}).
\]
可见 $f(x)$ 不变号，故 $c\stackrel{\text{必}}{=}0$. 即 $f(x)\equiv0$.

\noindent\textbf{巧用极值原理}

\noindent\textbf{☆例 4.3.34}\quad 设 $f(x)$ 在 $\mathbb{R}$ 上连续，又 $\varphi(x)=f(x)\int_0^x f(t)\,dt$ 单调递减，证明：$f(x)\equiv0,x\in\mathbb{R}$.（上海交通大学）

\noindent\textbf{证}\quad 已知
\[
\varphi(x)=f(x)\int_0^x f(t)\,dt
=\left[\left(\frac12\int_0^x f(t)\,dt\right)^2\right]',
\]
且 $\varphi(0)=0$，故函数
\[
F(x)=\left(\frac12\int_0^x f(t)\,dt\right)^2
\]
的导数
\[
F'(x)
\begin{cases}
\ge0, & \text{当 }x<0\text{ 时},\\
=0, & \text{当 }x=0\text{ 时},\\
\le0, & \text{当 }x>0\text{ 时}.
\end{cases}
\]
由此 $F(x)$ 在 $x=0$ 处取极大也是最大，即 $0\le F(x)\le F(0)=0$. 因此，$\int_0^x f(t)\,dt\equiv0$.

因 $f(x)$ 连续，故
\[
f(x)=\left(\int_0^x f(t)\,dt\right)'=0\quad(\forall x\in\mathbb{R}).
\]

\noindent\textbf{被积函数零点问题}

\noindent\textbf{☆例 4.3.35}\quad 设 $f(x)$ 在 $[a,b]$ 上连续，且
\[
\int_a^b f(x)\,dx=0,\quad \int_a^b xf(x)\,dx=0.
\]
证明：至少存在两点 $x_1,x_2\in(a,b)$，使得 $f(x_1)=f(x_2)=0$.（湖北大学）

\noindent\textbf{证}\quad $1^\circ$（利用 Rolle 定理）记 $F(x)=\int_a^x f(t)\,dt$，则有 $F(a)=F(b)=0$，因此 $\exists x_1\in(a,b)$，使得 $f(x_1)=F'(x_1)=0$.

\noindent $2^\circ$ 假若 $x\ne x_1$ 时 $f(x)>0$，则 $F'(x)>0,F$ 严格递增，$0=F(a)<F(x_1)<F(b)=0$，矛盾. 类似可证 $f(x)<0$ 不可能（当 $x\ne x_1$ 时）.

\noindent $3^\circ$ 有 $2^\circ$ 的结论，就可断言 $f(x)$ 必有第二个零点，因为不然的话，$f(x)$ 在 $(a,x_1)$ 内（和 $(x_1,b)$ 内）不能变号，且 $x_1$ 的两侧只能异号，从而 $(x-x_1)f(x)$ 在 $x_1$ 两侧保持同号，于是
\[
0\ne\int_a^b(x-x_1)f(x)\,dx
=\int_a^b xf(x)\,dx-x_1\int_a^b f(x)\,dx=0,
\]
矛盾.

\noindent\textbf{对数导数的妙用}

\noindent 用积分解决微分学的问题.

\noindent\textbf{☆例 4.3.36}\quad 已知 $f(x)$ 在 $(2,+\infty)$ 内可导，$f(x)>0$，且 $\frac{d}{dx}(xf(x))\le-kf(x)$，$k$ 为常数，试证在此区间上 $f(x)\le Ax^{-(k+1)}$，其中 $A$ 为与 $x$ 无关的常数.（兰州大学）

\noindent\textbf{提示}\quad 利用对数导数公式：当 $f(x)>0$ 时，有 $(\ln|f(x)|)'=\frac{f'(x)}{f(x)}$.

\noindent\textbf{再提示}\quad
\[
f(x)+xf'(x)\le-kf(x)
\Rightarrow\frac{f'(x)}{f(x)}\le-\frac{k+1}{x}
\]
\[
\Rightarrow\int_2^x\frac{f'(x)}{f(x)}\,dx
\le-\int_2^x\frac{k+1}{x}\,dx
\]
\[
\Rightarrow \ln f(x)\bigg|_2^x
\le-[(k+1)\ln x]\bigg|_2^x
\]
\[
\Rightarrow\ln\frac{f(x)}{f(2)}
\le(k+1)\ln\frac2x
=\ln\left(\frac2x\right)^{k+1}
\]
\[
\Rightarrow f(x)\le Ax^{-(k+1)}
\quad(A=f(2)\cdot2^{k+1}).
\]

\noindent\textbf{☆例 4.3.37}\quad 设 $f(x)$ 在 $(0,+\infty)$ 内连续可微，$f(0)=1$，当 $x\ge0$ 时，$f(x)>|f'(x)|$. 证明：当 $x>0$ 时，$e^x>f(x)$.（中国科学院）

\noindent\textbf{证}\quad 已知当 $x>0$ 时，
\[
(\ln f(x))'=\frac{f'(x)}{f(x)}<1.
\]
又由于 $f(0)=1$，知
\[
\ln f(x)=\int_0^x(\ln f(x))'\,dx<x.
\]
因此，$f(x)<e^x$.

\noindent\textbf{定积分的实际应用}

\noindent\textbf{new ☆例 4.3.38}\quad 设 $f:[0,1]\to[0,1]$ 有连续二阶导数，$f(0)=f(1)=0,f''(x)<0$. 曲线 $L=\{(x,f(x))\mid x\in[0,1]\}$ 的弧长记为 $\ell$，试证：$\ell<3$.（北京大学）

\noindent\textbf{提示}\quad 弧长公式：
\[
\ell=\int_0^1\sqrt{1+f'^2(x)}\,dx.
\]
因为在 $[0,1]$ 上 $f''(x)<0$，故 $f'(x)$ 严格递减，$f(x)$ 严格凸向上；加之 $f:[0,1]\to[0,1],f(0)=f(1)=0$，说明 $f(x)$ 是单峰函数：有唯一最大值点 $c\in(0,1)$，使 $0<f(c)\le1$.

\noindent\textbf{再提示}\quad 因 $\sqrt{1+(f'(x))^2}\le1+|f'(x)|$（直角三角形两直角边（长度）之和大于斜边），又有 $f(c)>0$，故
\[
\begin{aligned}
\ell
&=\int_0^1\sqrt{1+f'^2(x)}\,dx
\le\int_0^1(1+|f'(x)|)\,dx\\
&=\int_0^c(1+f'(x))\,dx+\int_c^1(1-f'(x))\,dx
=1+2f(c)\le3.
\end{aligned}
\]

\noindent\textbf{new ☆四、积分作为上下限的函数}

\noindent\textbf{要点}\quad 1）若函数 $f$ 在 $[a,b]$ 上可积，则 $F(x)=\int_a^x f(t)\,dt$ 在 $[a,b]$ 上连续.

2）在区间 $I$ 上，若函数 $f$ 连续，$\psi(x)$ 和 $\varphi(x)$ 可导，则
\[
F(x)=\int_{\varphi(x)}^{\psi(x)} f(t)\,dt
\]
在 $I$ 上也可导，且
\[
F'(x)=f(\psi(x))\cdot\psi'(x)-f(\varphi(x))\cdot\varphi'(x)\quad(x\in I).
\]

\noindent\textbf{new 例 4.3.39}\quad 设
\[
f(x)=\int_x^{x^2}\frac1t\ln\frac{t-1}{32}\,dt\quad(x\in(1,+\infty)).
\]
问函数 $f(x)$ 在何处取最小值？（中国科学技术大学）

\noindent\textbf{解}\quad
\[
\begin{aligned}
f'(x)
&=2x\cdot\frac1{x^2}\ln\frac{x^2-1}{32}
-\frac1x\ln\frac{x-1}{32}\\
&=\frac1x\ln\frac{(x-1)(x+1)^2}{32}
=\frac1x\ln\left[1+\frac{(x-3)(x^2+4x+11)}{32}\right]\quad(x>1).
\end{aligned}
\]
令 $f'(x)=0$，得唯一稳定点 $x=3$；并且当从左到右穿过 $x=3$ 时，$f'(x)$ 由负变为正，故 $f(3)$ 是 $f$ 的最小值：
\[
f_{\min}=\int_3^9\frac1t\ln\frac{t-1}{32}\,dt.
\]

\noindent\textbf{new *例 4.3.40}\quad 设 $f(x)$ 在 $[0,+\infty)$ 上二阶可导，$f(x)\le f''(x)$，且 $f(0)\ge0,f'(0)\ge0$. 求证：在 $[0,+\infty)$ 上，
\[
f(x)\ge f(0)+f'(0)x.
\tag{1}
\]
（中国科学技术大学）

\noindent\textbf{证}\quad 已知 $f(x)\le f''(x)$，两端同时积分得
\[
\int_0^x f(u_1)\,du_1
\le\int_0^x f''(u_1)\,du_1
=f'(x)-f'(0).
\]
再积分一次，得
\[
\int_0^x\left(\int_0^u f(u_1)\,du_1\right)\,du
\le\int_0^x[f'(u)-f'(0)]\,du
=f(x)-f(0)-f'(0)x,
\]
亦即
\[
f(x)\ge f(0)+f'(0)x+\int_0^x du\int_0^u f(u_1)\,du_1.
\tag{2}
\]
又因 $f(0)\ge0,f'(0)\ge0$，故
\[
f(x)\ge\int_0^x du\int_0^u f(u_1)\,du_1\quad(\forall x\ge0).
\tag{3}
\]
利用式 (3)，代入式 (2) 右端的积分中，反复迭代可得
\[
f(x)\ge f(0)+f'(0)x
+\int_0^x du\int_0^u f(u_1)\,du_1\cdots
\int_0^{u_{n-2}}f(u_{n-1})\,du_{n-1}\int_0^{u_{n-1}}f(u_n)\,du_n.
\tag{4}
\]
函数序列
\[
G_n(x)=\int_0^x du\int_0^u f(u_1)\,du_1\cdots
\int_0^{u_{n-2}}f(u_{n-1})\,du_{n-1}\int_0^{u_{n-1}}f(u_n)\,du_n
\quad(n=0,1,2,\cdots).
\]
\[
G_n'(0)=G_n''(0)=\cdots=G_n^{(n)}(0)=0.
\]
应用 Taylor 公式，$\exists\xi\in(0,x)$，
\[
G_n(x)=\frac{f(\xi)}{(n+1)!}x^{n+1}.
\]
记 $M(x)=\max f(\xi)\ (0<\xi<x)$，$\forall x>0$，
\[
|G_n(x)|\le\frac{M(x)}{(n+1)!}x^{n+1}\to0\quad(n\to\infty).
\]
因此在式 (4) 里令 $n\to\infty$，取极限，可得欲证的不等式 (1).

\noindent\textbf{new **例 4.3.41}\quad 设 $f(x),g(x),\varphi(x)$ 均为 $[a,b]$ 上的连续函数，且 $g(x)$ 单调递增，$\varphi(x)\ge0$，若对于任意 $x\in[a,b]\ (a>0)$，有
\[
f(x)\le g(x)+\int_a^x\varphi(t)f(t)\,dt.
\tag{1}
\]
试证：对于任意 $x\in[a,b]$，有
\[
f(x)\le g(x)e^{\int_a^x\varphi(s)\,ds}.
\tag{2}
\]
（浙江大学）

\noindent\textbf{证}\quad $1^\circ$ 将式 (1) 两边同乘以 $\varphi(x)$（$\varphi(x)\ge0$），变形得
\[
g(x)\varphi(x)\ge \varphi(x)f(x)-\varphi(x)\int_a^x\varphi(t)f(t)\,dt.
\]
令
\[
F(x)=\int_a^x\varphi(t)f(t)\,dt,
\]
则 $F'(x)=\varphi(x)f(x)$，上式可写为
\[
g(x)\varphi(x)\ge F'(x)-\varphi(x)F(x).
\]
两边同乘 $e^{-\int_a^x\varphi(s)\,ds}$ 得
\[
g(x)\varphi(x)e^{-\int_a^x\varphi(s)\,ds}
\ge F'(x)e^{-\int_a^x\varphi(s)\,ds}
-\varphi(x)e^{-\int_a^x\varphi(s)\,ds}F(x)
=\left(e^{-\int_a^x\varphi(s)\,ds}F(x)\right)'_x.
\tag{3}
\]
两端 $x$ 改成 $t$，同时求积分得（注意到 $F(a)=0$）
\[
\int_a^x g(t)\varphi(t)e^{-\int_a^t\varphi(s)\,ds}\,dt
\ge e^{-\int_a^x\varphi(s)\,ds}F(x).
\tag{4}
\]
因 $g\nearrow$，故 $t<x$ 时 $g(t)\le g(x)$，而 $\varphi(t)e^{-\int_a^t\varphi(s)\,ds}\ge0$，所以
\[
g(x)\int_a^x\varphi(t)e^{-\int_a^t\varphi(s)\,ds}\,dt
\ge \int_a^x g(t)\varphi(t)e^{-\int_a^t\varphi(s)\,ds}\,dt
\ge e^{-\int_a^x\varphi(s)\,ds}F(x),
\]
即
\[
g(x)\int_a^x\varphi(t)e^{-\int_a^t\varphi(s)\,ds}\,dt
\ge e^{-\int_a^x\varphi(s)\,ds}F(x).
\tag{5}
\]
两边同乘 $e^{\int_a^x\varphi(s)\,ds}$，（并注意到：当 $a\le t\le x$ 时，$e^{\int_a^x\varphi(s)\,ds}\cdot e^{-\int_a^t\varphi(s)\,ds}=e^{\int_t^x\varphi(s)\,ds}$）知
\[
e^{\int_a^x\varphi(s)\,ds}\cdot g(x)\int_a^x\varphi(t)e^{-\int_a^t\varphi(s)\,ds}\,dt
=g(x)\int_a^x\varphi(t)e^{\int_t^x\varphi(s)\,ds}\,dt
=-g(x)\int_a^x\left(e^{\int_t^x\varphi(s)\,ds}\right)'_t\,dt
\]
\[
=-g(x)+g(x)e^{\int_a^x\varphi(s)\,ds}.
\]
因此，式 (5) 变为
\[
-g(x)+g(x)e^{\int_a^x\varphi(s)\,ds}\ge F(x).
\tag{6}
\]
于是
\[
f(x)\stackrel{\text{式(1)}}{\le}g(x)+F(x)
\stackrel{\text{式(6)}}{\le}g(x)e^{\int_a^x\varphi(s)\,ds},
\]
式 (2) 获证.

\noindent\textbf{点评}\quad 该题技巧性高，极富挑战性. 张祖锦的解法十分精巧，编写时基本上采用了他的思路.

\noindent\textbf{new ☆练习 1}\quad 求积分的极限：
\[
\lim_{x\to+\infty}x^m\int_0^{\frac1x}\sin t^2\,dt
\]
（其中 $m$ 为任意整数）.（南开大学）
\[
\left(0\ (m=1,2\text{ 或 }m\le0),\quad \frac13\ (m=3),\quad +\infty\ (m>3)\right)
\]

\noindent\textbf{提示}\quad
\[
\lim_{x\to+\infty}x^m\int_0^{\frac1x}\sin t^2\,dt
\overset{u=\frac1x}{=}
\lim_{u\to0^+}\frac{\int_0^u\sin t^2\,dt}{u^m}
\overset{\text{L'Hospital 法则}}{=}
\lim_{u\to0^+}\frac{\sin u^2}{m\cdot u^{m-1}}
\overset{\text{等价代换}}{=}
\lim_{u\to0^+}\frac1m u^{3-m}\quad(m=1,2,3).
\]

\noindent\textbf{练习 2}\quad 设
\[
f(x)=\int_x^{x^2}\left(1+\frac1{2t}\right)^t\sin\frac1{\sqrt{t}}\,dt\quad(x>0),
\]
求 $\lim_{n\to\infty}f(n)\sin\frac1n$.（福建师范大学）

\noindent\textbf{提示}\quad
\[
\begin{aligned}
\lim_{n\to\infty}f(n)\sin\frac1n
&=\lim_{x\to+\infty}f(x)\sin\frac1x
\overset{\sin\frac1x\sim\frac1x}{=}
\lim_{x\to+\infty}\frac1x\int_x^{x^2}
\left(1+\frac1{2t}\right)^t\sin\frac1{\sqrt{t}}\,dt\\
&\overset{\text{L'Hospital 法则}}{=}
\lim_{x\to+\infty}\left[
2x\left(1+\frac1{2x^2}\right)^{x^2}\sin\frac1{\sqrt{x^2}}
-\left(1+\frac1{2x}\right)^x\sin\frac1{\sqrt{x}}\right]\\
&=2e^{\frac12}.
\end{aligned}
\]

\noindent\textbf{new ☆五、（第一）积分中值定理}

注意：因为一般不讲第二积分中值定理，因此说积分中值定理，就是指第一积分中值定理.

\noindent\textbf{（第一积分中值定理）}\quad 设 $f(x)$ 是 $[a,b]$ 上的连续函数，则 $\exists\xi\in[a,b]$ 使得
\[
f(\xi)=\frac1{b-a}\int_a^b f(x)\,dx,
\tag{A}
\]
亦即
\[
\int_a^b f(x)\,dx=f(\xi)\int_a^b dx=f(\xi)(b-a)\quad(m\le\xi\le M).
\]

\noindent\textbf{（强化形式）}\quad $\exists\xi\in(a,b)$ 使得
\[
f(\xi)=\frac1{b-a}\int_a^b f(x)\,dx,
\tag{B}
\]
亦即
\[
\int_a^b f(x)\,dx=f(\xi)\int_a^b dx=f(\xi)(b-a)\quad(m<\xi<M).
\]

\noindent\textbf{（推广形式）}\quad 设 $f(x)$ 在 $[a,b]$ 上连续，$g(x)$ 在 $[a,b]$ 上可积，且不变号，则 $\exists\xi\in(a,b)$ 使得
\[
\int_a^b f(x)g(x)\,dx=f(\xi)\int_a^b g(x)\,dx.
\tag{C}
\]
注意：当 $g(x)\equiv1$ 时，回到式 (B) 的情况.

\noindent\textbf{new **例 4.3.42}\quad 证明（第一）积分中值定理：公式 (A)，(B)，(C).

\noindent\textbf{证}\quad $1^\circ$ 回顾公式 (A) 的证明：因 $f(x)$ 连续，故 $f(x)$ 有最大值和最小值（记为 $M$ 和 $m$）：
\[
m=\frac1{b-a}\int_a^b m\,dx
\le\frac1{b-a}\int_a^b f(x)\,dx
\le\frac1{b-a}\int_a^b M\,dx=M,
\tag{1}
\]

亦即
\[
\int_a^b f(x)\,dx=\mu\int_a^b dx=\mu(b-a)\quad(m\le\mu\le M).
\]
再利用连续函数的介值性即得公式 (A).

$2^\circ$ 证明公式 (B)：在 $[a,b]$ 上作
\[
F(x)=\int_a^x f(x)\,dx,
\]
并用 Lagrange 定理，$\exists\xi:a<\xi<b$，使得
\[
\int_a^b f(x)\,dx=F(b)-F(a)\stackrel{\text{Lagrange 定理}}{=}F'(\xi)(b-a)
=f(\xi)(b-a)\quad(a<\xi<b).
\]
（注意：Lagrange 找得到的 $\xi$ 在区间内部！）再同除以 $(b-a)$，即得公式 (B).

$3^\circ$ 证明公式 (C)：因为 $f(x)$ 在 $[a,b]$ 上连续，故 $f(x)$ 有最大值和最小值（记为 $M$ 和 $m$），即 $\forall x\in[a,b]$，有 $m\le f(x)\le M$. 因此
\[
m\int_a^b g(x)\,dx\le\int_a^b f(x)g(x)\,dx\le M\int_a^b g(x)\,dx.
\tag{2}
\]
① 若 $m=M$，则 $f(x)$ 为常数：$m=f(x)=M$，式 (2) 变成等式，$(a,b)$ 的每一个 $x$ 都可作为 $\xi$ 满足公式 (C)（属于平凡情况无须讨论）.

② 若 $\displaystyle\int_a^b g(x)\,dx=0$，也是平凡的. 因为此时式 (2) 左、右两端都为零，所以
\[
\int_a^b f(x)g(x)\,dx=0.
\]
公式 (C) 变成 $0=0,\ \forall\xi\in[a,b]$，公式都成立.

③ 下设
\[
G=\int_a^b g(x)\,dx>0
\]
（$G<0$ 类似可证），此时式 (2) 等价于
\[
m\le \mu\stackrel{\text{记}}{=}
\frac{\displaystyle\int_a^b f(x)g(x)\,dx}{\displaystyle\int_a^b g(x)\,dx}\le M.
\tag{3}
\]
下面分 4 种情况进行讨论：

i）当 $m=\mu=M$ 时，$f(x)$ 为常数，即 $\forall x\in(a,b)$ 都可作为 $\xi$.

ii）当 $m<\mu<M$ 时，利用 $f$ 连续的介值性，知 $\exists\xi\in(a,b)$ 使得
\[
f(\xi)=\mu=
\frac{\displaystyle\int_a^b f(x)g(x)\,dx}{\displaystyle\int_a^b g(x)\,dx},
\]
满足式 (3)，故公式 (C) 成立.

iii）当 $m<\mu=M$，即
\[
\mu=\frac{\displaystyle\int_a^b f(x)g(x)\,dx}{\displaystyle\int_a^b g(x)\,dx}=M
\]
时，两边同乘 $\displaystyle\int_a^b g(x)\,dx$，得
\[
\int_a^b f(x)g(x)\,dx=\int_a^b Mg(x)\,dx,
\]
亦即
\[
\int_a^b(M-f(x))g(x)\,dx=0.
\tag{4}
\]
因 $(M-f(x))g(x)\ge0,\ \forall[x',x'']\subset[a,b]$，恒有
\[
\int_a^b(M-f(x))g(x)\,dx\ge
\int_{x'}^{x''}(M-f(x))g(x)\,dx.
\tag{5}
\]
由于 $g(x)$ 可积，$g(x)$ 在 $[a,b]$ 上的 Darboux 下和有极限：
\[
\lim_{\lambda\to0}\sum_{i=1}^n m_i\Delta x_i
=\int_a^b g(x)\,dx\stackrel{\text{记}}{=}G>0
\]
（$\lambda$ 是最大“小区间”的长度）. 因此，当分划足够细密时，有
\[
\sum_{i=1}^n m_i\Delta x_i\ge\frac G2.
\tag{6}
\]
从而至少有一个小区间（设为第 $k$ 个小区间：$[x_{k-1},x_k]$）使得对应的
\[
m_k\ge\frac{G}{2(b-a)}.
\tag{7}
\]
（一定存在，否则，对每个 $m_i<\dfrac{G}{2(b-a)}$，有
\[
\sum_{i=1}^n m_i\Delta x_i<
\frac{G}{2(b-a)}\sum_{i=1}^n\Delta x_i=\frac G2,
\]
与式 (6) 矛盾.）用小区间 $[x_{k-1},x_k]$ 取代式 (5) 里的 $[x',x'']$，则
\[
0\stackrel{\text{式(4)}}{=}\int_a^b(M-f(x))g(x)\,dx
\stackrel{\text{式(5)}}{\ge}
\int_{x_{k-1}}^{x_k}(M-f(x))g(x)\,dx
\ge\int_{x_{k-1}}^{x_k}(M-f(x))m_k\,dx\ge0.
\]
此式说明：
\[
\int_{x_{k-1}}^{x_k}[M-f(x)]m_k\,dx=0.
\]
故 $\forall x\in[x_{k-1},x_k]$，有 $f(x)\equiv M\stackrel{\text{前设}}{=}\mu$. 即：
\[
\forall x\in[x_{k-1},x_k]
\]
都可选作 $\xi$，使得 $f(\xi)=\mu$，满足式 (3)，公式 (2) 获证.

iv）当 $m=\mu<M$ 时，类似 iii）可证.

\noindent\textbf{注 1}\quad $f(x)$ 在 $[a,b]$ 上可积，则 $f(x)$ 有界，从而有上、下确界 $(M,m)$. 若 $f(x)$ 还有介值性，那么虽然不能断言 $f(x)$ 有最大、最小值，但能断言 $f(x)$ 取尽 $m$ 和 $M$ 之间的一切中间值. 因此，从公式 (A) 和 (C) 的证明看出，替代连续性，第一积分中值定理对具有介值性的可积函数仍然成立. 因为导函数具有介值性，可见导函数若在 $[a,b]$ 上可积，就能在 $[a,b]$ 上应用第一积分中值定理（如例 4.3.20）. 如果 $f(x)$ 没有介值性，只能写作
\[
\int_a^b f(t)\,dt=\mu(b-a),\quad
\inf_{a\le x\le b}f(x)<\mu<\sup_{a\le x\le b}f(x).
\]

\noindent\textbf{注 2}\quad 跟公式 (A) 不同，（强化后的）公式 (B) 强调 $\xi$ 一定落在区间内部（是内点，不是端点），这一改进，有时成为问题的关键，十分重要，如下例的练习 2.

\noindent\textbf{new ☆例 4.3.43}\quad 设 $f(x)$ 在 $(0,1)$ 上可微，且满足
\[
f(1)-2\int_0^{\frac12}xf(x)\,dx=0.
\tag{1}
\]
求证：在 $(0,1)$ 内至少有一点 $\xi$，使得
\[
f'(\xi)=-\frac{f(\xi)}{\xi}.
\tag{2}
\]
（中国人民大学，重庆大学）

\noindent\textbf{证 I}\quad
\[
f(1)\stackrel{\text{式(1)}}{=}2\int_0^{\frac12}xf(x)\,dx
\stackrel{\text{积分中值定理}}{=}\eta f(\eta)\quad
\left(\eta\in\left[0,\frac12\right]\right).
\]
函数 $xf(x)$ 在 $[\eta,1]$ 上可应用 Rolle 定理，$\exists\xi\in(\eta,1)$，使得
\[
(xf(x))'_{x=\xi}=0.
\]
知式 (2) 成立.

\noindent\textbf{证 II}\quad
\[
G(x)\stackrel{\text{记}}{=}\int_0^x(f(1)-tf(t))\,dt,
\]
有 $G(0)=G\left(\dfrac12\right)=0$. 对 $G(x)$ 应用 Rolle 定理，$\exists\eta\in\left(0,\dfrac12\right)$，使得
\[
G'(\eta)=f(1)-\eta f(\eta)=0.
\]
同上，再在 $[\eta,1]$ 上对 $xf(x)$ 应用 Rolle 定理，得 $\exists\xi\in(\eta,1)$，使得
\[
f'(\xi)+\xi f(\xi)=0.
\]
即式 (2) 成立.

\noindent\textbf{new 练习 1}\quad 设 $f(x)$ 在区间 $[0,1]$ 上连续，在 $(0,1)$ 上可导，满足
\[
f(1)=2\int_0^{\frac12}e^{1-x^2}f(x)\,dx.
\]
求证：存在一点 $\xi\in(0,1)$ 使得
\[
f'(\xi)=2\xi f(\xi).
\]
（北京大学）

\noindent\textbf{提示}\quad 一方面，利用积分中值定理，有
\[
\int_0^{\frac12}e^{1-x^2}f(x)\,dx=\frac12 e^{1-\eta^2}f(\eta),\quad
\eta\in\left(0,\frac12\right).
\]
另一方面，
\[
2\int_0^{\frac12}e^{1-x^2}f(x)\,dx\stackrel{\text{题设}}{=}f(1)
=\left(e^{1-x^2}f(x)\right)_{x=1}.
\]
因此，在 $[\eta,1]$ 上可对 $e^{1-x^2}f(x)$ 应用 Rolle 定理，得
$f'(\xi)=2\xi f(\xi)\ (\xi\in(0,1))$.

\noindent\textbf{new 练习 2}\quad 设 $f(x)$ 在 $[0,1]$ 上连续，在 $(0,1)$ 内三次可微，并且满足条件：
\[
1)\ \int_0^{\frac14}f(t)\,dt=\frac14f(0);\qquad
2)\ \int_{\frac14}^1f(t)\,dt=\frac34f(1).
\]
试证：若 $f(0)=f(1)$，则存在 $\xi\in(0,1)$，使得 $f'''(\xi)=0$.

\noindent\textbf{提示}\quad 利用（强化后的）积分中值定理的公式 (B)：由条件 1），2）可知
\[
\exists\xi_1\in\left(0,\frac14\right),\quad
\frac14f(0)=\int_0^{\frac14}f(t)\,dt=\frac14f(\xi_1),
\]
即 $f(0)=f(\xi_1)$；
\[
\exists\xi_2\in\left(\frac14,1\right),\quad
\frac34f(1)=\int_{\frac14}^1f(t)\,dt=\frac34f(\xi_2),
\]
即 $f(\xi_2)=f(1)$.
于是，$f(x)$ 有四个等值点，使得
\[
f(0)=f(\xi_1)=f(\xi_2)=f(1).
\]
应用 Rolle 定理，在相邻两等值点之间，导函数有一零点. 因此，$f'(x)$ 有至少三个（不同的）零点，使得
$f'(\eta_1)=f'(\eta_2)=f'(\eta_3)=0$. 如此反复应用 Rolle 定理，可得到 $f''(x)$ 至少有两个（不同的）零点 $\zeta_1,\zeta_2\in(0,1)$，最后得到至少一个 $\xi\in(0,1)$，使得 $f'''(\xi)=0$.

\noindent\textbf{注}\quad 如果不用强化形式，改用公式 (A)，则上面第一步里的 $\xi_1$ 和 $\xi_2$ 可能相等：$\xi_1=\dfrac14=\xi_2$，如此只能得到 $f''(x)$ 的一个零点，得不出 $f'''(x)$ 的零点.

\noindent\textbf{new 练习 3}\quad 设 $f(x)$ 在 $[0,1]$ 上连续，在 $(0,1)$ 内三次可微，并且满足条件：
\[
1)\ f(1)=f\left(\frac14\right)=0;\qquad
2)\ \int_{\frac14}^1f(t)\,dt=\frac34f(1).
\]
试证：存在 $\xi\in(0,1)$，使得 $f'''(\xi)=0$.（中国科学院）

\noindent\textbf{解 I}\quad 在 $\left[\dfrac14,1\right]$ 上，对
\[
g(x)=\int_{\frac14}^x f(t)\,dt-f(1)\left(x-\frac14\right)
\]
应用 Rolle 定理，$\exists\eta\in\left(\dfrac14,1\right)$，
\[
f(\eta)=f(1)=f\left(\frac14\right)=0.
\]
$f(x)$ 有三个零点. 反复应用 Rolle 定理，可得 $f'(\xi_1)=f'(\xi_2)=0$，从而得
$f''(\xi)=0$.

\noindent\textbf{解 II}\quad 如果在条件 2）里直接应用第一积分中值定理（强化形式），能直接得到 $\exists\eta\in\left(\dfrac14,1\right)$，使得
\[
f(\eta)=f(1).
\]
无须作辅助函数. 然后同上提示：$f(\eta)=f(1)=0$，可得 $f'(\xi_1)=f'(\xi_2)=0$，从而可得 $f''(\xi)=0$.

\subsection*{六、第二积分中值定理}

\noindent\textbf{new **例 4.3.44}\quad 证明：（第二积分中值定理）

1）在 $[a,b]$ 上，若 $f(x)$ 非负单调递减，$g(x)$ 可积，则 $\exists\xi\in[a,b]$ 使得
\[
I=\int_a^b f(x)g(x)\,dx=f(a)\int_a^\xi g(x)\,dx;
\tag{A}
\]
2）在 $[a,b]$ 上，若 $f(x)$ 非负单调递增，$g(x)$ 可积，则 $\exists\xi\in[a,b]$ 使得
\[
I=\int_a^b f(x)g(x)\,dx=f(b)\int_\xi^b g(x)\,dx;
\tag{B}
\]
3）在 $[a,b]$ 上，若 $f(x)$ 单调，$g(x)$ 可积，则 $\exists\xi\in[a,b]$ 使得
\[
I=\int_a^b f(x)g(x)\,dx=f(a)\int_a^\xi g(x)\,dx+f(b)\int_\xi^b g(x)\,dx.
\tag{C}
\]
公式 (A), (B), (C) 被统称为 Bonnet 公式（其中的单调性不要求是严格的）.

\noindent\textbf{证}\quad 1）（证明公式 (A).）作分划 $a=x_0<x_1<x_2<\cdots<x_n=b$，于是
\[
I=\int_a^b f(x)g(x)\,dx=\sum_{i=0}^{n-1}\int_{x_i}^{x_{i+1}}f(x)g(x)\,dx.
\]
将 $f(x)$ 写为 $(f(x)-f(x_i))+f(x_i)$，则
\[
I=\int_a^b f(x)g(x)\,dx
=\sum_{i=0}^{n-1}\int_{x_i}^{x_{i+1}}(f(x)-f(x_i))g(x)\,dx
+\sum_{i=0}^{n-1}\int_{x_i}^{x_{i+1}}f(x_i)g(x)\,dx
\stackrel{\text{记}}{=}I_1+I_2.
\tag{1}
\]
在 $[a,b]$ 上 $g(x)$ 可积，可知 $g(x)$ 有界：$\exists L>0,\ |g(x)|\le L\ (\forall x\in[a,b])$. 又因 $f\searrow$，故
\[
0\le\sum\omega_i'\le f(a)-f(b)\stackrel{\text{记}}{=}R
\]
（这里 $\omega_i'$ 表示 $f(x)$ 在第 $i$ 个小区间上的振幅），$\lambda\stackrel{\text{记}}{=}\max_i\Delta x_i,\ \Delta x_i=x_{i+1}-x_i$，得
\[
|I_1|\le\sum_{i=0}^{n-1}\int_{x_i}^{x_{i+1}}|f(x)-f(x_i)||g(x)|\,dx
\le L\sum_{i=0}^{n-1}\omega_i'\Delta x_i\le LR\lambda\to0\quad(\lambda\to0).
\]
令
\[
G(x)=\int_a^x g(x)\,dx,
\]
则
\[
\int_{x_i}^{x_{i+1}}g(x)\,dx=G(x_{i+1})-G(x_i),\quad G(x_0)=G(a)=0.
\]
于是
\[
I_2=\sum_{i=0}^{n-1}f(x_i)(G(x_{i+1})-G(x_i))
\]
\[
\stackrel{\text{Abel 变换}}{=}
\sum_{i=1}^{n-1}G(x_i)(f(x_{i-1})-f(x_i))+G(x_n)f(x_{n-1}).
\tag{2}
\]
（注 Abel 变换是一种和式的变形方法. 本书多次要用到它，这里是第一次，可能要多次才能熟练掌握. 下面我们来具体推导一次. 为了简洁，令 $f_i=f(x_i),\ G_i=G(x_i)$，那么 $G_0=G(x_0)=G(a)=0$. 于是
\[
I_2=\sum_{i=0}^{n-1}f_i(G_{i+1}-G_i)
\]
\[
=f_0G_1-f_0G_0+f_1G_2-f_1G_1+f_2G_3-f_2G_2+\cdots
+f_{n-3}G_{n-2}-f_{n-3}G_{n-3}+f_{n-2}G_{n-1}-f_{n-2}G_{n-2}+f_{n-1}G_n-f_{n-1}G_{n-1}
\]
\[
=-f_0G_0+(f_0G_1-f_1G_1)+(f_1G_2-f_2G_2)+\cdots
+(f_{n-2}G_{n-1}-f_{n-1}G_{n-1})+f_{n-1}G_n
\]
\[
\stackrel{G_0=0}{=}
\sum_{i=1}^{n-1}G_i(f_{i-1}-f_i)+G_nf_{n-1}
=\sum_{i=1}^{n-1}G(x_i)(f(x_{i-1})-f(x_i))+G(x_n)f(x_{n-1}),
\]
此即式 (2).）

因为 $f$ 单减非负，若“$f(a)=0$”，则 $f(x)\equiv0$，公式 (A) 自然成立. 因此，下面默认 $f(a)\ne0$. 因 $g(x)$ 可积，可知
\[
G(x)=\int_a^x g(x)\,dx
\]
连续，故在 $[a,b]$ 上 $G(x)$ 有最大值 $(M)$ 和最小值 $(m)$：
\[
m\le G_i\le M\quad(i=0,1,2,\cdots,n).
\tag{3}
\]
已知 $f(x)\ge0$ 单减，故 $f_{i-1}-f_i\ge0$，故由 (3) 得
\[
m(f_{i-1}-f_i)\le G_i(f_{i-1}-f_i)\le M(f_{i-1}-f_i)\quad(i=1,2,\cdots,n),
\]
\[
mf(a)=m\left[\sum_{i=1}^{n-1}(f_{i-1}-f_i)+f_{n-1}\right]\le I_2
\le M\left[\sum_{i=1}^{n-1}(f_{i-1}-f_i)+f_{n-1}\right]=Mf(a),
\]
即 $mf(a)\le I_2\le Mf(a)$，亦即
\[
m\le\frac{I_2}{f(a)}\le M.
\tag{4}
\]
而 $G(x)=\displaystyle\int_a^x g(x)\,dx$ 连续，有介值性：对 (4)，$\exists\xi\in[a,b]$ 使得
\[
G(\xi)=\int_a^\xi g(x)\,dx=\frac{I_2}{f(a)}.
\]
两边同乘 $f(a)$ 即得公式 (A)，证毕.

2）（证明公式 (B).）现知 $f(x)\ge0$ 单调递增，令 $x=b-t$，则 $f(b-t)\ge0$ 单减，符合公式 (A) 的条件. 故

\[
\int_a^b f(x)g(x)\,dx\stackrel{\text{令 }x=b-t}{=}
\int_0^{b-a}f(b-t)g(b-t)\,dt
\stackrel{\text{式(A)}}{=}f(b-0)\int_0^{\xi_1}g(b-t)\,dt
\]
\[
=-f(b-0)\int_b^{b-\xi_1}g(x)\,dx
\stackrel{\xi=b-\xi_1\in[a,b]}{=}f(b-0)\int_\xi^b g(x)\,dx
\]
（公式 (B) 获证）.

3）（证明公式 (C).）当 $f(x)$ 单调时，现设为递减（递增的情况类似可证）. 令
\[
F(x)=f(x)-f(b),
\]
则 $F(x)\ge0$，满足公式 (A) 的条件. 应用公式 (A)，得
\[
\int_a^b f(x)g(x)\,dx-f(b)\int_a^b g(x)\,dx
=\int_a^b F(x)g(x)\,dx
\stackrel{\text{式(A)}}{=}F(a)\int_a^\xi g(x)\,dx
\]
\[
=f(a)\int_a^\xi g(x)\,dx-f(b)\int_a^\xi g(x)\,dx.
\]
移项即得公式 (C)：
\[
\int_a^b f(x)g(x)\,dx
=f(a)\int_a^\xi g(x)\,dx
+f(b)\left(\int_a^b g(x)\,dx-\int_a^\xi g(x)\,dx\right)
\]
\[
=f(a)\int_a^\xi g(x)\,dx+f(b)\int_\xi^b g(x)\,dx.
\]

\noindent\textbf{注 1}\quad 公式 (A) 中的 $f(a)$ 可改写为 $f(a+0)$ 或比 $f(a+0)$ 大的任意常数 $A$，即 $\exists\xi\in[a,b]$，使得
\[
I=\int_a^b f(x)g(x)\,dx=A\int_a^\xi g(x)\,dx
\quad(\xi\text{ 的大小依赖于 }A\text{ 的选取}).
\tag{A'}
\]
（因为：$f(x)$ 单减，在端点 $a$ 处存在右极限 $f(a+0)$. 当 $f(a)=f(a+0)$ 或改为比 $f(a+0)$ 大的某数 $A$ 时，显然不会改变 $f(x)$ 的单减性，因此定理仍成立.）

\noindent\textbf{注 2}\quad 同理，公式 (B) 中的“$f(b)$”可以改换为“$f(b-0)$”或比 $f(b-0)$ 大的任一常数 $B$.

\noindent\textbf{注 3}\quad 公式 (C) 不要求 $f(x)$ 非负，这给应用带来很大方便.

\noindent\textbf{new 例 4.3.45}\quad 证明：若 $f(x)$ 在 $[0,1]$ 上严格单调下降，则

1）$\exists\theta\in(0,1)$，使得
\[
\int_0^1 f(x)\,dx=\theta f(0)+(1-\theta)f(1);
\]
2）$\forall c>f(0)$，$\exists\theta\in(0,1)$，使得
\[
\int_0^1 f(x)\,dx=\theta c+(1-\theta)f(1).
\]

\noindent\textbf{提示}\quad $g(x)\equiv1$，应用 Bonnet 公式 (C).

\noindent\textbf{new 例 4.3.46}\quad 设 $f:[a,b]\to[0,1]$ 可导，$f'(x)$ 在 $[a,b]$ 上单调下降，$f'(b)>0$. 求证：
\[
\int_a^b \cos f(x)\,dx\le\frac2{f'(b)}.
\]
（浙江大学）

\noindent\textbf{证 I}\quad
\[
\left|\int_a^b\cos f(x)\,dx\right|
\stackrel{\substack{\text{令 }y=f(x),\,g=f^{-1}\\ dx=g'(y)\,dy}}{=}
\left|\int_A^B\cos y\cdot g'(y)\,dy\right|
\quad(A=f(a),\ B=f(b)).
\tag{1}
\]
因 $f'(x)\downarrow,\ f'(x)\ge f'(b)>0$，知 $g'(y)>0$ 且 $g'(y)\downarrow$. 应用第二积分中值定理：
\[
\left|\int_A^B\cos y\cdot g'(y)\,dy\right|
=\left|g'(A)\int_A^\xi\cos y\,dy\right|
=\left|\frac1{f'(a)}(\sin\xi-\sin A)\right|
\le\frac2{f'(b)}.
\]

\noindent\textbf{证 II}\quad
\[
\int_a^b\cos f(x)\,dx
\stackrel{\text{令 }x=f^{-1}(y)}{=}
\int_{f(a)}^{f(b)}\cos y\cdot\frac1{f'(x)}\,dy\quad(x=f^{-1}(y)).
\]
而 $y=f(x)\in[0,1]$，$\cos y\ge0$ 不变号，可用第一积分中值定理：$\exists y_1\in[f(a),f(b)]$，$a\le f^{-1}(y_1)\le b$，使得（$f'(x)\downarrow$）
\[
\int_{f(a)}^{f(b)}\cos y\frac1{f'(x)}\,dy
=\frac1{f'(f^{-1}(y_1))}\int_{f(a)}^{f(b)}\cos y\,dy
\le\frac2{f'(b)}.
\]
（这是因为：$0\le\displaystyle\int_{f(a)}^{f(b)}\cos y\,dy\le|\sin f(b)|+|\sin f(a)|\le2$.）证毕.

\noindent\textbf{new 练习 1}\quad 设 $x>0,c>0$，试证：
\[
\left|\int_x^{x+c}\sin t^2\,dt\right|<\frac1x.
\]
（华中师范大学）

\noindent\textbf{证 I}\quad 令 $u=t^2$，在 $[x^2,(x+c)^2]$（$x>0,c>0$）上，$\dfrac1{\sqrt u}>0$ 且 $\downarrow$. 于是
\[
\left|\int_x^{x+c}\sin t^2\,dt\right|
=\left|\frac12\int_{x^2}^{(x+c)^2}\frac{\sin u}{\sqrt u}\,du\right|
\stackrel{\text{第二积分中值定理}}{=}
\left|\frac1{2\sqrt{x^2}}\int_{x^2}^{\xi}\sin u\,du\right|
\le\frac1{2x}|\cos x^2-\cos\xi|\le\frac1x.
\]

\noindent\textbf{证 II}\quad
\[
\int_x^{x+c}\sin t^2\,dt
\stackrel{t=\sqrt u}{=}\frac12\int_{x^2}^{(x+c)^2}\frac{\sin u}{\sqrt u}\,du
=-\frac12\int_{x^2}^{(x+c)^2}\frac1{\sqrt u}\,d\cos u
\]
\[
=\frac{\cos x^2}{2|x|}-\frac{\cos(x+c)^2}{2|x+c|}
-\frac14\int_{x^2}^{(x+c)^2}\frac{\cos u}{\sqrt{u^3}}\,du,
\]
\[
\left|\int_x^{x+c}\sin t^2\,dt\right|
<\frac1{2x}+\frac1{2(x+c)}+\frac14\int_{x^2}^{(x+c)^2}\frac1{\sqrt{u^3}}\,du
=\frac1x.
\]
（其中
\[
\frac14\int_{x^2}^{(x+c)^2}\frac1{\sqrt{u^3}}\,du
=\frac12\int_{x^2}^{(x+c)^2}\frac1u\,d\sqrt u
=\frac1{2x}-\frac1{2(x+c)}.
\]
）

\noindent\textbf{new 练习 2}\quad 设 $e^2<a<b$，求证：
\[
\int_a^b\frac{dx}{\ln x}<\frac{2b}{\ln b}.
\]
（北京大学）

\noindent\textbf{提示}\quad 方法 I：利用第二积分中值定理；方法 II：常数变量化（张卫建议）.

\noindent\textbf{证 I}\quad
\[
\int_a^b\frac{dx}{\ln x}\stackrel{u=\sqrt x}{=}
\int_{\sqrt a}^{\sqrt b}\frac{u\,du}{\ln u}.
\]
\[
\left(\frac u{\ln u}\right)'=\frac{\ln u-1}{(\ln u)^2}>0
\quad(u=\sqrt x>\sqrt a>e),
\]
$\dfrac u{\ln u}\nearrow$，应用第二积分中值定理得
\[
\int_{\sqrt a}^{\sqrt b}\frac{u\,du}{\ln u}
\stackrel{\text{Bonnet 公式(B)}}{=}
\frac{\sqrt b}{\ln\sqrt b}\int_\xi^{\sqrt b}du
=\frac{2\sqrt b}{\ln b}(\sqrt b-\xi)<\frac{2b}{\ln b}.
\]

\noindent\textbf{证 II}\quad （将常数 $b$ 视为变量 $(b>a)$）求证关于 $b$ 的不等式：
\[
F(b)\stackrel{\text{记}}{=}\frac{2b}{\ln b}-\int_a^b\frac1{\ln x}\,dx>0.
\tag{1}
\]
因 $e^2<a<b$，故 $F(a)\stackrel{\text{记}}{=}\dfrac{2a}{\ln a}\ge2>0$. 而
\[
F'(b)=\left(\frac{2b}{\ln b}\right)'_b-\frac1{\ln b}
=\frac1{\ln b}-\frac2{(\ln b)^2}
=\frac{\ln b-2}{(\ln b)^2}>0.
\]
故 $b>a$ 时，不等式 (1) 成立，原式获证.

\noindent\textbf{new 练习 3}\quad 设 $x\ne0$，
\[
f(x)=\int_0^x\sin\frac1t\,dt,\quad f(0)=0,
\]
求证：$f'(0)=0$.（北京大学习题）

\noindent\textbf{提示}\quad 因为
\[
f'(0)=\lim_{x\to0}\frac{f(x)-f(0)}x=\lim_{x\to0}\frac{f(x)}x,
\]
参看例 1.3.23 中 2）.

\noindent\textbf{再提示}\quad 根据例 1.3.23 中 2），只要证明：
\[
\lim_{x\to0}\frac{f(x)-f(x/2)}x=0,
\]
则有 $f'(0)=\displaystyle\lim_{x\to0}\dfrac{f(x)}x=0$.

事实上，
\[
\frac{f(x)-f(x/2)}x
=\frac1x\int_{x/2}^x\sin\frac1t\,dt
\stackrel{u=\frac1t}{=}
-\frac1x\int_{2/x}^{1/x}\sin u\frac1{u^2}\,du
\]
\[
=\frac1x\int_{1/x}^{2/x}\frac1{u^2}\sin u\,du
\stackrel{\text{第二积分中值定理}}{=}
x\int_{1/x}^{\xi}\sin u\,du\to0\quad(x\to0).
\]

\noindent\textbf{单元练习 4.3}

\noindent\textbf{4.3.1}\quad 证明：
\[
1)\ \sqrt2e^{-\frac12}<\int_{-\frac1{\sqrt2}}^{\frac1{\sqrt2}}e^{-x^2}\,dx<\sqrt2;
\quad
2)\ 0<\frac\pi2-\int_0^{\frac\pi2}\frac{\sin x}{x}\,dx<\frac{\pi^3}{144};
\]
\[
3)\ \frac29\pi^2\le\int_{\frac\pi6}^{\frac\pi2}\frac{2x}{\sin x}\,dx\le\frac49\pi^2;
\quad
4)\ 当\ 0\le x\le\frac\pi2\ 时，\sin x\le x-\frac1{3\pi}x^3.
\]

\noindent\textbf{提示}\quad （利用已知不等式求积分.）

1）$-\dfrac1{\sqrt2}\le x\le\dfrac1{\sqrt2}\Rightarrow0\le x^2\le\dfrac12
\Rightarrow e^{-x^2}\le e^{-\frac12}\Rightarrow e^{-\frac12}<e^{-x^2}<1$.

对最后的不等式求积分得
\[
\sqrt2e^{-\frac12}
=\int_{-\frac1{\sqrt2}}^{\frac1{\sqrt2}}e^{-\frac12}\,dx
<\int_{-\frac1{\sqrt2}}^{\frac1{\sqrt2}}e^{-x^2}\,dx
<\int_{-\frac1{\sqrt2}}^{\frac1{\sqrt2}}dx=\sqrt2.
\]
2）$x\in\left(0,\dfrac\pi2\right)$：
\[
\sin x>x-\frac{x^3}{3!}\Rightarrow
\frac{\sin x}{x}>1-\frac{x^2}{6}
\Rightarrow 1-\frac{\sin x}{x}<\frac{x^2}{6}
\]
\[
\Rightarrow\frac\pi2-\int_0^{\frac\pi2}\frac{\sin x}{x}\,dx<\frac{\pi^3}{144}.
\]
类似地，$\sin x<x\Rightarrow\displaystyle\int_0^{\frac\pi2}\frac{\sin x}{x}\,dx<\frac\pi2
\Rightarrow0<\frac\pi2-\int_0^{\frac\pi2}\frac{\sin x}{x}\,dx$.

3）当 $\dfrac\pi6\le x\le\dfrac\pi2$ 时，有
\[
2x\le\frac{2x}{\sin x}\le4x,
\]
对该不等式求积分可得.

4）已有
\[
\sin x<x-\frac{x^3}{6}+\frac{x^5}{120},
\]
因此只需证明：当 $0\le x\le\dfrac\pi2$ 时，
\[
x-\frac{x^3}{6}+\frac{x^5}{120}\le x-\frac{x^3}{3\pi}.
\tag{1}
\]
但此式等价于
\[
\pi\left(\frac12-\frac{x^2}{40}\right)>1.
\tag{2}
\]
而由 $0\le x\le\dfrac\pi2<2$ 可得 $x^2\le4$，即 $\dfrac{x^2}{40}<\dfrac1{10}$，故
\[
\pi\left(\frac12-\frac{x^2}{40}\right)>\pi\frac{5-1}{10}>1.
\]
式 (2) 获证，原不等式成立.

\noindent\textbf{new ☆4.3.2}\quad 设 $m,n\ge0$ 为整数，求积分
\[
I(m,n)=\int_0^1x^m(1-x)^n\,dx.
\]
（中国科学院）
\[
\left(\frac{m!\,n!}{(m+n+1)!}\right)
\]

\noindent\textbf{提示}\quad 建立递推公式，
\[
I(m,n)=\int_0^1(1-x)^n\,d\frac{x^{m+1}}{m+1}
\stackrel{\text{分部积分}}{=}
\frac n{m+1}I(m+1,n-1)=\cdots
=\frac{n!}{(m+n)\cdots(m+1)}I(m+n,0).
\]
因 $I(m+n,0)=\dfrac1{m+n+1}$，故
\[
I(m,n)=\frac{m!\,n!}{(m+n)!}\cdot\frac1{m+n+1}
=\frac{m!\,n!}{(m+n+1)!}.
\]

\noindent\textbf{☆4.3.3}\quad 求证：
\[
f(x)=\int_0^x(t-t^2)\sin^{2n}t\,dt
\]
（$n$ 为正整数）在 $x\ge0$ 上的最大值不超过
\[
\frac1{(2n+2)(2n+3)}.
\]
（西北大学）

\noindent\textbf{提示}\quad
\[
f'(x)\begin{cases}
>0,&0<x<1,\\
=0,&x=1,\\
<0,&x>1,
\end{cases}
\Rightarrow f(x)\le\max_{x>0}f(x)=f(1).
\]

\noindent\textbf{☆4.3.4}\quad 把满足下述条件 1）和 2）的实函数 $f$ 的全体记作 $F$：

1）$f(x)$ 在闭区间 $[0,1]$ 上连续，并且非负；2）$f(0)=0,f(1)=1$.

试证：
\[
\inf_{f\in F}\int_0^1f(x)\,dx=0,
\]
但不存在 $\varphi\in F$，使得
\[
\int_0^1\varphi(x)\,dx=0.
\]
（厦门大学）

\noindent\textbf{提示}\quad $\forall f\in F$，有 $\displaystyle\int_0^1f(x)\,dx\ge0$，又 $\exists f_n\in F$（如 $f_n(x)=x^n$），使得
\[
\lim_{n\to\infty}\int_0^1x^n\,dx=0,
\]
故 $\displaystyle\inf_{f\in F}\int_0^1f(x)\,dx=0$. 但 $\forall\varphi\in F$，由连续非负、$\varphi(1)=1$，易证
\[
\int_0^1\varphi(x)\,dx>0.
\]

\noindent\textbf{☆4.3.5}\quad 若 $f'(x)$ 在 $[0,2\pi]$ 上连续，且 $f'(x)\ge0$，则对任意正整数 $n$，有
\[
\left|\int_0^{2\pi}f(x)\sin nx\,dx\right|
\le\frac{2[f(2\pi)-f(0)]}{n}.
\]
（东北师范大学）

\noindent\textbf{提示}\quad
\[
\left|\int_0^{2\pi}f(x)\sin nx\,dx\right|
=\left|\frac1n\int_0^{2\pi}f(x)\,d\cos nx\right|
\]
\[
\le\frac1n[f(2\pi)-f(0)]
+\frac1n\left|\int_0^{2\pi}f'(x)\cos nx\,dx\right|
\le\frac{2[f(2\pi)-f(0)]}{n}.
\]

\noindent\textbf{4.3.6}\quad $f(x)$ 在 $[a,b]$ 上可导，$f'(x)\downarrow$，$|f'(x)|\ge m>0$，试证：
\[
\left|\int_a^b\cos f(x)\,dx\right|\le\frac2m.
\]

\noindent\textbf{提示}\quad 因 $y=f(x)$ 严格单调，有反函数 $g(y)=f^{-1}(y)$，且
\[
g'(y)=\frac1{f'(x)}\bigg|_{x=g(y)}.
\]
导数有介值性，既然 $|f'(x)|\ge m>0$，说明只有两种可能：$f'(x)\ge m>0$ 或 $f'(x)\le-m<0$.

再提示
\[
\left|\int_a^b\cos f(x)\,dx\right|
\stackrel{\substack{\text{令 }y=f(x),\,x=f^{-1}\\ dx=g'(y)\,dy}}{=}
\left|\int_A^B\cos y\cdot g'(y)\,dy\right|
\quad(A=f(a),\ B=f(b)).
\tag{1}
\]
$1^\circ$ 当 $f'(x)\ge m>0$ 时，$g'(y)>0$. 因
\[
g'(y)=\frac1{f'(x)}\bigg|_{x=g(y)},
\]
故由 $f'(x)\downarrow$ 知 $g'(y)\nearrow$. 应用第二积分中值定理：
\[
\left|\int_A^B\cos y\cdot g'(y)\,dy\right|
=\left|g'(B)\int_\xi^B\cos y\,dy\right|
=\left|\frac1{f'(b)}(\sin B-\sin\xi)\right|\le\frac2m.
\]
$2^\circ$ 当 $-f'(x)\ge m>0$ 时，因 $\cos x$ 是偶函数，
\[
\left|\int_a^b\cos f(x)\,dx\right|
=\left|\int_a^b\cos(-f(x))\,dx\right|.
\]
此时令 $F(x)=-f(x)$，则 $F'(x)=-f'(x)\ge m>0$. 反函数
\[
g'(y)=\frac1{-f'(x)}\bigg|_{x=g(y)}
\]
同样可应用第二积分中值定理，得证.

\noindent\textbf{说明}\quad 下面的 4.3.7, 4.3.9, 4.3.10 三题，曾在“博士数学论坛”（现为“博士数学家园”）上引起热烈讨论，出现不少精美的解法，令人赞赏，下面希望尽量弘扬他们的思想方法写出解答.

\noindent\textbf{*4.3.7}\quad $f(x)\ne0$，在 $[a,b]$ 内可导，$f(a)=f(b)=0$. 证明：至少存在一点 $\xi\in[a,b]$，使得
\[
|f'(\xi)|>\frac4{(b-a)^2}\int_a^b|f(x)|\,dx.
\tag{1}
\]

\noindent\textbf{证}\quad 假设结论不成立，即 $\forall\xi\in[a,b]$，有
\[
|f'(\xi)|\le\frac4{(b-a)^2}\int_a^b|f(x)|\,dx\stackrel{\text{记}}{=}k
\quad(k\ge0\text{ 常数}).
\tag{2}
\]
（由此推出矛盾）将区间 $[a,b]$ 的中点记为 $c$，即 $c=\dfrac{a+b}2$，则
\[
\int_a^b|f(t)|\,dt\le\int_a^c|f(t)|\,dt+\int_c^b|f(t)|\,dt.
\tag{3}
\]
应用 Lagrange 定理，当 $x\in[a,c]$ 时，
\[
|f(x)|=|f(x)-f(a)|
\stackrel{\exists\xi_1\in(a,x)}{=}|f'(\xi_1)(x-a)|
\stackrel{\text{式(2)}}{\le}k(x-a).
\tag{4}
\]
同理，对 $x\in[c,b]$，有
\[
|f(x)|\le k(b-x).
\tag{5}
\]
将 (4), (5) 代入 (3)，得
\[
\int_a^b|f(t)|\,dt
\le\int_a^c k(t-a)\,dt+\int_c^b k(b-t)\,dt
=\frac k2\left[(c-a)^2+(b-c)^2\right]
=\frac{k(b-a)^2}{4}
\stackrel{\text{式(2)}}{=}\int_a^b|f(t)|\,dt.
\tag{6}
\]

此式：“左端$\le$右端”，且两端是同一数 $\displaystyle\int_a^b|f(t)|\,dt$，故其中的“$\le$”应改为“$=$”. 由于 $f(x)$ 连续，利用积分性质（例 4.2.12）知：式 (4) 和 (5) 也应是等号. 故
\[
k(x-a)\stackrel{x\in[a,c]}{=}|f(x)|\stackrel{x\in[c,b]}{=}k(b-x)
\quad(k\ge0\text{ 为常数}).
\tag{7}
\]
由此可见，$k\ne0$（因为若 $k=0$，则 $f(x)\equiv0$，与题设 $f(x)\ne0$ 矛盾），故只能 $k>0$. 说明 $f(x)$ 在 $(a,b)$ 内无零点，因而不会变号.

假设 $f(x)>0$（$f(x)<0$ 类似可证），则
\[
\text{当 }x\in(a,c)\text{ 时， }f(x)=k(x-a),\quad f'(x)=k;
\]
\[
\text{当 }x\in(c,b)\text{ 时， }f(x)=k(b-x),\quad f'(x)=-k.
\]
这表明：导函数 $f'$ 在中点 $c$ 处有第一类间断，与导数特性矛盾！

或者利用 Lagrange 定理，知
\[
f'_-(c)=\lim_{0<h\to0}\frac{f(c-h)-f(c)}{-h}
=\lim_{0<\theta<h\to0}\frac{f'(c-\theta)(-h)}{-h}=k,
\]
\[
f'_+(c)=\lim_{0<h\to0}\frac{f(c+h)-f(c)}{h}
=\lim_{0<\theta<h\to0}\frac{f'(c+\theta)h}{h}=-k.
\]
因此，$f'_-(c)\ne f'_+(c)$，即 $f'(c)$ 不存在，与题设矛盾. 证毕.

网上有很多类似的解法，很难说此种解法应归于谁.

\noindent\textbf{类题}\quad 设 $f(x)$ 在 $[a,b]$ 连续，在 $(a,b)$ 内可导，$f\left(\dfrac{a+b}2\right)=0$，且 $f(x)$ 不是常数. 试证：
\[
\exists\xi\in(a,b),\quad
|f'(\xi)|\ge\frac4{(b-a)^2}|f(b)-f(a)|.
\]
（南开大学）

\noindent\textbf{4.3.8}\quad 将条件 $f(x)\ne0$ 换为 $f''(x)<0$，重新证明例 4.3.5.

\noindent\textbf{*4.3.9}\quad 证明
\[
\int_0^{\frac\pi2}t\left(\frac{\sin nt}{\sin t}\right)^4dt<\frac{\pi^2n^2}{4}.
\]

\noindent\textbf{证 I}\quad （Gingkuan 和 tian27456 等人的证法.）只需证明 $n\ge2$ 的情况. 令
\[
\int_0^{\frac\pi2}t\left(\frac{\sin nt}{\sin t}\right)^4dt
=\int_0^{\frac\pi{2n}}t\left(\frac{\sin nt}{\sin t}\right)^4dt
+\int_{\frac\pi{2n}}^{\frac\pi2}t\left(\frac{\sin nt}{\sin t}\right)^4dt
\stackrel{\text{记}}{=}I+J.
\]
对 $I$，用数学归纳法易知
\[
\left|\frac{\sin nx}{\sin x}\right|\le n,
\]
从而
\[
I\le\frac{n^2\pi^2}{8}.
\]
对 $J$，利用 $|\sin nx|\le1$ 及 $\dfrac2\pi t\le\sin t\ (0\le t\le\dfrac\pi2)$，可得
\[
J\le\frac{n^2\pi^2}{8}.
\]
于是
\[
\int_0^{\frac\pi2}t\left(\frac{\sin nt}{\sin t}\right)^4dt
=I+J\le\frac{n^2\pi^2}{4}.
\]

\noindent\textbf{证 II}\quad （此证由张卫提供，作者 Hansschwarzkopf 等.）取 $c\in\left(0,\dfrac\pi2\right)$，
\[
\int_0^{\frac\pi2}t\left(\frac{\sin nt}{\sin t}\right)^4dt
=\int_0^c t\left(\frac{\sin nt}{\sin t}\right)^4dt
+\int_c^{\frac\pi2}t\left(\frac{\sin nt}{\sin t}\right)^4dt
\stackrel{\text{记}}{=}I_1+I_2.
\]
其中（利用 $\left|\dfrac{\sin nx}{\sin x}\right|\le n$）
\[
I_1=\int_0^c t\left(\frac{\sin nt}{\sin t}\right)^4dt
\le n^4\int_0^c t\,dt=\frac{n^4c^2}{2},
\]
\[
I_2=\int_c^{\frac\pi2}t\left(\frac{\sin nt}{\sin t}\right)^4dt
\le\left(\frac\pi2\right)^4\int_c^{\frac\pi2}t\cdot\frac1{t^4}\,dt
=\frac{\pi^4}{32c^2}-\frac{\pi^2}{8}.
\]
因此
\[
0<\int_0^{\frac\pi2}t\left(\frac{\sin nt}{\sin t}\right)^4dt
=I_1+I_2<\frac{n^4c^2}{2}+\frac{\pi^4}{32c^2}-\frac{\pi^2}{8}.
\tag{1}
\]
“不等式 $\dfrac12(u^2+v^2)\ge uv$ 当且仅当 $u=v$ 时才变为等号”，应用到上式右端：
\[
\frac{n^4c^2}{2}+\frac{\pi^4}{32c^2}
=\frac12\left[(n^2c)^2+\left(\frac{\pi^2}{4c}\right)^2\right]\ge n^2\frac{\pi^2}{4},
\tag{2}
\]
当且仅当 $n^2c=\dfrac{\pi^2}{4c}$ 亦即 $c=\dfrac\pi{2n}$ 时，式 (2) 的“$\ge$”才变成“$=$”，式 (1) 右端达最小：
\[
\min_{c\in(0,\frac\pi2)}
\left(\frac{n^4c^2}{2}+\frac{\pi^4}{32c^2}-\frac{\pi^2}{8}\right)
=\left(\frac{n^4c^2}{2}+\frac{\pi^4}{32c^2}-\frac{\pi^2}{8}\right)_{c=\frac\pi{2n}}
=\frac{n^2\pi^2}{4}-\frac{\pi^2}{8}.
\]
因此
\[
0<\int_0^{\frac\pi2}t\left(\frac{\sin nt}{\sin t}\right)^4dt
=I_1+I_2\le\frac{n^2\pi^2}{4}-\frac{\pi^2}{8}
<\frac{n^2\pi^2}{4}.
\]

\noindent\textbf{评}\quad 证 II 实际给出了“更精准一点”的估计. 因此题目可改为求证下式成立：
\[
\int_0^{\frac\pi2}t\left(\frac{\sin nt}{\sin t}\right)^4dt
\le\frac{n^2\pi^2}{4}-\frac{\pi^2}{8}.
\]
但原题也有优点：结论简洁，易于证明.

\noindent\textbf{*4.3.10}\quad 对自然数 $n\ge2$，证明：
\[
\frac1\pi\int_0^{\frac\pi2}\left|\frac{\sin(2n+1)t}{\sin t}\right|dt
<\frac{2+\ln n}{2}.
\]

\noindent\textbf{证}\quad （此证由张卫提供，作者 Hansschwar Ekpof 等）
\[
\text{原式左端}=\frac1\pi\int_0^{\frac\pi2}
\frac{|\sin(2n+1)x|}{\sin x}\,dx
\]
\[
\le\frac1\pi\int_0^{\frac\pi2}\left(\frac1{\sin x}-\frac1x\right)
|\sin(2n+1)x|\,dx
+\frac1\pi\int_0^{\frac\pi2}\frac{|\sin(2n+1)x|}{x}\,dx
\stackrel{\text{记}}{=}I_1+I_2.
\tag{1}
\]
（下面放大 $I_1+I_2$，并证明放大后还小于等于原式右端.）

$2^\circ$（放大 $I_1$）在 $\left[0,\dfrac\pi2\right]$ 上，
\[
f(x)\stackrel{\text{记}}{=}\left(\frac1{\sin x}-\frac1x\right)\nearrow
\]
（因 $f'\ge0$）. 故
\[
I_1\le f\left(\frac\pi2\right)\cdot\frac1\pi
\int_0^{\frac\pi2}|\sin(2n+1)x|\,dx
\le f\left(\frac\pi2\right)\frac12
=\frac12-\frac1\pi.
\tag{2}
\]

$3^\circ$（将 $I_2$ 放大.）
\[
I_2=\frac1\pi\int_0^{\frac\pi2}\frac{|\sin(2n+1)x|}{x}\,dx
\stackrel{t=(2n+1)x}{=}
\frac1\pi\int_0^{(2n+1)\pi/2}\frac{|\sin t|}{t}\,dt
\]
\[
=\frac1\pi\int_0^{\frac\pi2}\frac{\sin t}{t}\,dt
+\frac1\pi\sum_{k=1}^{2n}\int_{k\pi/2}^{(k+1)\pi/2}\frac{|\sin t|}{t}\,dt
\]
\[
\le\frac12+\frac1\pi\sum_{k=1}^{2n}\int_{k\pi/2}^{(k+1)\pi/2}\frac{|\sin t|}{k\pi/2}\,dt
\le\frac12+\frac2{\pi^2}\sum_{k=1}^{2n}\frac1k
\]
\[
\le\frac12+\frac2{\pi^2}+\frac2{\pi^2}\sum_{k=2}^{2n}\int_{k-1}^k\frac{dt}{t}
=\frac12+\frac2{\pi^2}+\frac2{\pi^2}\ln(2n).
\tag{3}
\]
$4^\circ$ 式 (2) 和 (3) 代入式 (1)，得（$n\ge2$ 时）
\[
\text{原式左端}\le
1-\frac1\pi+\frac2{\pi^2}+\frac2{\pi^2}\ln2+\frac2{\pi^2}\ln n
\le 1+\frac{\ln n}{2}
=\text{原式右端}.
\]
（其中最后的不等式只需验证当 $n=2$ 时成立即可，因为当 $n$ 增大时，$\dfrac{\ln n}{2}$ 比 $\dfrac2{\pi^2}\ln n$ 增加得更快. 事实上，当 $n=2$ 时，最后的不等式等价于：$(\pi^2-8)\ln2+(2\pi-4)\ge0$，显然成立.）

\noindent\textbf{new 类题}\quad 设 $f,g$ 在 $[a,b]$ 连续，证明：存在 $\xi\in(a,b)$，使得
\[
g(\xi)\int_a^\xi f(x)\,dx=f(\xi)\int_\xi^b g(x)\,dx.
\]
（中国科学技术大学）

\noindent\textbf{提示}\quad 在 $[a,b]$ 上，
\[
F(x)\stackrel{\text{记}}{=}\int_a^x f(t)\,dt\cdot\int_x^b g(t)\,dt,
\]
$F(a)=F(b)=0$，应用 Rolle 定理.

\noindent\textbf{4.3.11}\quad 函数 $f(x)$ 在 $[a,b]$ 上连续，并且对于任何区间 $[\alpha,\beta]\ (a\le\alpha<\beta\le b)$，不等式
\[
\left|\int_\alpha^\beta f(x)\,dx\right|
\le M|\beta-\alpha|^{1+\delta}\quad(M,\delta\text{ 是正常数})
\]
成立. 证明：在 $[a,b]$ 上，$f(x)\equiv0$.（国外赛题）

\noindent\textbf{提示}\quad $\forall x_0\in(a,b)$，当 $n$ 充分大时，
\[
(\alpha_n,\beta_n)\stackrel{\text{记}}{=}\left(x_0-\frac1{2n},x_0+\frac1{2n}\right)\subset[a,b].
\]
因 $f(x)$ 连续，$\exists x_n:f(x_n)\stackrel{\text{记}}{=}\max_{\alpha_n\le x\le\beta_n}f(x)$，$(\alpha_n,\beta_n)$ 可作为 $(\alpha,\beta)$，满足已知条件.

再提示
\[
\frac1n|f(x_n)|
=\left|\int_{\alpha_n}^{\beta_n}f(x_n)\,dx\right|
\le\left|\int_{\alpha_n}^{\beta_n}f(x)\,dx\right|
\le M(\beta_n-\alpha_n)^{1+\delta}=\frac{M}{n^{1+\delta}}.
\]
因此，任何内点上：
\[
f(x_0)=\lim_{n\to\infty}f(x_n)=0.
\]
类似可证 $x_0$ 为端点的情况.

\noindent\textbf{☆4.3.12}\quad 证明：若 $f(x)$ 为 $[0,1]$ 上的连续函数，且对一切 $x\in[0,1]$，有
\[
\int_0^x f(u)\,du\ge f(x)\ge0,
\]
则 $f(x)\equiv0$.（上海师范大学）

\noindent\textbf{提示}\quad 方法 I：记
\[
F(x)=\int_0^x f(t)\,dt,
\]
则 $0\le F'(x)=f(x)\le F(x)$. 用例 3.2.18，知 $F(x)\equiv0$. 又因 $f$ 连续、非负，故 $f(x)\equiv0$（于 $[0,1]$ 上）.

\noindent\textbf{方法 II}\quad $\exists M>0$，使 $|f(x)|\le M$（于 $[0,1]$ 上）. $\forall x\in[0,1)$，有
\[
0\le f(x)\le\int_0^x f(u)\,du=f(\xi_1)x\quad(0\le\xi_1\le x)
\]
（反复利用此结果）
\[
\le f(\xi_2)\xi_1x\le\cdots\le f(\xi_n)\xi_{n-1}\cdots\xi_1x\le Mx^n\to0.
\]
故在 $[0,1)$ 上，$f(x)\equiv0$. 由连续性知 $f(1)=0$.

\noindent\textbf{4.3.13}\quad 证明：如果在 $(-\infty,+\infty)$ 上的连续函数 $f(x)$ 满足
\[
\int_x^{x+1}f(t)\,dt=0,
\]
那么 $f(x)$ 是周期函数.

\noindent\textbf{☆4.3.14}\quad 设 $f(x)$ 处处连续，
\[
F(x)=\frac1{2\delta}\int_{-\delta}^{\delta}f(x+t)\,dt,
\]
其中 $\delta$ 为任何正数，证明：

1）$F(x)$ 对任何 $x$ 有连续导数；

2）在任意闭区间 $[a,b]$ 上，当 $\delta$ 足够小时，可使 $F(x)$ 与 $f(x)$ 一致逼近（即任给 $\varepsilon>0$，对一切 $x\in[a,b]$，均有 $|F(x)-f(x)|<\varepsilon$）.（华东师范大学）

\noindent\textbf{提示}\quad 1）令 $u=x+t$，知
\[
F(x)=\frac1{2\delta}\int_{x-\delta}^{x+\delta}f(u)\,du.
\]
又因 $f$ 连续，故
\[
F'(x)=\frac1{2\delta}[f(x+\delta)-f(x-\delta)]
\]
也连续.

2）由 $f$ 连续知 $f(x)$ 在 $[a,b]$ 上一致连续，故 $\forall\varepsilon>0$，$\exists\delta>0$，当 $|\xi|<\delta$ 时，$\forall x\in[a,b]$，有
\[
|f(x+\xi)-f(x)|<\varepsilon.
\]
因此
\[
|F(x)-f(x)|
=\left|\frac1{2\delta}\int_{-\delta}^{\delta}f(x+t)\,dt-f(x)\right|
=|f(x+\xi)-f(x)|<\varepsilon.
\]

\noindent\textbf{☆4.3.15}\quad $[a,b]$ 上的连续函数序列 $\varphi_1,\varphi_2,\cdots,\varphi_n,\cdots$ 满足
\[
\int_a^b\varphi_n^2(x)\,dx=1.
\]
证明：存在自然数 $N$ 及定数 $c_1,c_2,\cdots,c_N$，使
\[
\sum_{k=1}^N c_k^2=1,\quad
\max_{x\in[a,b]}\left|\sum_{k=1}^N c_k\varphi_k(x)\right|>100.
\]
（扬州师范学院）

\noindent\textbf{提示}\quad 可对
\[
\int_a^b[\varphi_1^2(x)+\cdots+\varphi_N^2(x)]\,dx=N
\]
应用积分中值定理.

\noindent\textbf{再提示}\quad 因 $\exists\xi\in[a,b]$，使得
\[
\sum_{i=1}^N\varphi_i^2(\xi)=\frac N{b-a}.
\]
取 $N>100^2(b-a)$，令
\[
c_i=\sqrt{\frac{b-a}{N}}\varphi_i(\xi)\quad(i=1,2,\cdots,N),
\]
则 $\sum_{i=1}^N c_i^2=1$ 且
\[
\sum_{i=1}^N c_i\varphi_i(\xi)>100.
\]

\noindent\textbf{4.3.16}\quad 按牛顿二项式展开及代换 $x=\sin t$ 两种方法计算积分
\[
\int_0^1(1-x^2)^n\,dx
\]
（$n$ 为正整数），并由此说明：
\[
1-\frac{C_n^1}{3}+\frac{C_n^2}{5}-\frac{C_n^3}{7}+\cdots
+\frac{(-1)^{n-1}}{2n-1}C_n^{n-1}+\frac{(-1)^n}{2n+1}
=\frac{(2n)!!}{(2n+1)!!}.
\]

\noindent\textbf{4.3.17}\quad 设在 $\left(0,\dfrac\pi2\right)$ 内连续函数 $f(x)>0$，且满足
\[
f^2(x)=\int_0^x f(t)\frac{\tan t}{\sqrt{1+2\tan^2t}}\,dt.
\]
求 $f(x)$ 的初等函数表达式.（复旦大学）

\noindent\textbf{提示}\quad 两端同时对 $x$ 求导.

\noindent\textbf{☆4.3.18}\quad 设
\[
\lim_{x\to0}\frac1{bx-\sin x}\int_0^x\frac{t^2}{\sqrt{a+t^2}}\,dt=1,
\]
试求正常数 $a$ 与 $b$.（华中师范大学）
\[
\langle a=4,b=1\rangle
\]

\noindent\textbf{提示}\quad 可用 $\dfrac00$ 型 L'Hospital 法则.

\noindent\textbf{再提示}\quad
\[
1=\text{原极限}\stackrel{\text{L'Hospital 法则}}{=}
\lim_{x\to0}\frac{\dfrac{x^2}{\sqrt{a+x^2}}}{b-\cos x}
\]
\[
=\begin{cases}
0,&b\ne1\ (\text{矛盾}),\\[0.4em]
\displaystyle\lim_{x\to0}\frac{\dfrac{x^2}{\sqrt{a+x^2}}}{\frac12x^2}
=\frac2{\sqrt a},&b=1.
\end{cases}
\]
故 $\sqrt a=2,a=4,b=1$.

\noindent\textbf{☆4.3.19}\quad 求
\[
\lim_{x\to+\infty}\int_x^{x+2}t\left(\sin\frac3t\right)f(t)\,dt,
\]
其中 $f(x)$ 可微，且已知 $\displaystyle\lim_{t\to+\infty}f(t)=1$.（中国科学技术大学）

\noindent\textbf{提示}\quad 可用积分中值定理.
\[
\text{原式}=\lim_{\xi\to+\infty}\xi\sin\frac3\xi\cdot f(\xi)\cdot2=6
\quad(x<\xi<x+2,\ x\to+\infty\text{ 时，}\xi\to+\infty).
\]

\noindent\textbf{4.3.20}\quad 设 $a>0$，函数 $f(x)$ 在 $[0,a]$ 上连续可微，证明：
\[
|f(0)|\le\frac1a\int_0^a|f(x)|\,dx+\int_0^a|f'(x)|\,dx.
\]
（华中师范大学）

\noindent\textbf{提示}\quad
\[
\frac1a\int_0^a f(x)\,dx=f(\xi)=f(0)+\int_0^\xi f'(x)\,dx,
\]
\[
|f(0)|\le\left|\frac1a\int_0^a f(x)\,dx\right|
+\left|\int_0^\xi f'(x)\,dx\right|\le\cdots.
\]

\noindent\textbf{☆4.3.21}\quad 设 $f(x)$ 的一阶导数在 $[0,1]$ 上连续，且 $f(0)=f(1)=0$，求证：
\[
\left|\int_0^1 f(x)\,dx\right|
\le\frac14\max_{0\le x\le1}|f'(x)|.
\]
（清华大学）

\noindent\textbf{提示}\quad 可先通过平移 $x-\dfrac12=t$，再分部积分并放大.

\noindent\textbf{再提示}\quad 令 $x-\dfrac12=t$，
\[
\left|\int_0^1f(x)\,dx\right|
=\left|\int_{-\frac12}^{\frac12}f\left(t+\frac12\right)\,dt\right|
\]
\[
=\left|\left.tf\left(t+\frac12\right)\right|_{-\frac12}^{\frac12}
-\int_{-\frac12}^{\frac12}tf'\left(t+\frac12\right)\,dt\right|
\le M\int_{-\frac12}^{\frac12}|t|\,dt=\frac M4
\]
（$M=\displaystyle\max_{0\le x\le1}|f'(x)|$）.

\noindent\textbf{☆4.3.22}\quad 设 $f\in C([0,1])$（即 $f$ 在 $[0,1]$ 上连续），且在 $(0,1)$ 上可微，若有
\[
8\int_{\frac78}^1f(x)\,dx=f(0),
\]
证明：存在 $\xi\in(0,1)$，使得 $f'(\xi)=0$.（北京大学）

\noindent\textbf{提示}\quad 可先用积分中值定理，再用 Rolle 定理：
\[
8\int_{\frac78}^1f(x)\,dx=8f(\eta)\left(1-\frac78\right)=f(0),
\]
即 $f(\eta)=f(0)$. 故 $\exists\xi:f'(\xi)=0$，其中 $0<\xi<\eta\le1$.

\noindent\textbf{4.3.23}\quad 设函数 $f(x)$ 在 $[a,b]$ 上连续，$f(x)>0$，又
\[
F(x)=\int_a^x f(t)\,dt+\int_b^x\frac1{f(t)}\,dt.
\]
试证：

1）$F'(x)\ge2$；\qquad 2）$F(x)=0$ 在 $[a,b]$ 中有且仅有一个实根.（华中师范大学）

\noindent\textbf{提示}\quad 1）（平均不等式）
\[
f(x)+\frac1{f(x)}\ge2\sqrt{f(x)\frac1{f(x)}}=2.
\]
2）$F$ 在端点异号，且 $F'>0$，$F$ 严$\nearrow$.

\noindent\textbf{*4.3.24}\quad 设
\[
f(x)=\int_x^{x+1}\sin t^2\,dt,
\]
求证：$x>0$ 时，
\[
|f(x)|<\frac1x.
\]
（北京工业大学）

\noindent\textbf{提示}\quad 令 $t=\sqrt{\tau}$，
\[
f(x)=\int_{x^2}^{(x+1)^2}\sin\tau\cdot\frac1{2\sqrt\tau}\,d\tau
=\frac12\int_{x^2}^{(x+1)^2}\frac1{\sqrt\tau}\,d(-\cos\tau).
\]
分部积分，（注意 $|\cos\tau|\le1$）可得 $x>0$ 时，
\[
|f(x)|<\frac1{2x}+\frac1{2(x+1)}
+\frac14\int_{x^2}^{(x+1)^2}\tau^{-\frac32}\,d\tau=\frac1x.
\]
此题还可利用第二积分中值定理进行证明.

\noindent\textbf{※4.3.25}\quad 设 $f(x)$ 在 $(a,b)$ 内连续，
\[
\lim_{h\to0}\frac1{h^3}\int_0^h[f(x+u)+f(x-u)-2f(x)]\,du=0
\quad(x\in[a,b]),
\]
试证 $f(x)$ 为线性函数.

\noindent\textbf{提示}\quad 利用练习 3.2.34.

\noindent\textbf{4.3.26}\quad 设 $f(x)$ 是 $[-\pi,\pi]$ 上的凸函数，$f'(x)$ 有界. 求证：
\[
a_{2n}=\frac1\pi\int_{-\pi}^{\pi}f(x)\cos2nx\,dx\ge0,
\]
\[
a_{2n+1}=\frac1\pi\int_{-\pi}^{\pi}f(x)\cos(2n+1)x\,dx\le0.
\]

\noindent\textbf{证 I}\quad
\[
a_{2n}=\frac1\pi\int_{-\pi}^{\pi}f(x)\cos2nx\,dx
=-\frac1{2n\pi}\int_{-\pi}^{\pi}f'(x)\sin2nx\,dx.
\]
令 $t=2nx$，得
\[
a_{2n}=-\frac1{4n^2\pi}\int_{-2n\pi}^{2n\pi}f'\left(\frac t{2n}\right)\sin t\,dt
\]
\[
=-\frac1{4n^2\pi}\sum_{k=-n}^{n-1}
\int_{2k\pi}^{2(k+1)\pi}f'\left(\frac t{2n}\right)\sin t\,dt.
\]
于是
\[
a_{2n}=-\frac1{4n^2\pi}\sum_{k=-n}^{n-1}
\int_0^{2\pi}f'\left(\frac{t+2k\pi}{2n}\right)\sin t\,dt
\]
\[
=-\frac1{4n^2\pi}\sum_{k=-n}^{n-1}
\int_0^\pi\left[
f'\left(\frac{t+2k\pi}{2n}\right)
-f'\left(\frac{t+(2k+1)\pi}{2n}\right)\right]\sin t\,dt.
\]
因为 $f$ 为凸函数，故 $f'(x)\nearrow$，从而
\[
f'\left(\frac{t+2k\pi}{2n}\right)
\le f'\left(\frac{t+(2k+1)\pi}{2n}\right).
\]
因此
\[
a_{2n}=\frac1{4n^2\pi}\sum_{k=-n}^{n-1}
\int_0^\pi\left[
f'\left(\frac{t+(2k+1)\pi}{2n}\right)
-f'\left(\frac{t+2k\pi}{2n}\right)\right]\sin t\,dt\ge0.
\]
同理可证
\[
a_{2n+1}\le0.
\]

\noindent\textbf{证 II}\quad（使用第二积分中值定理）
\[
a_{2n}=-\frac1{2n\pi}\int_{-\pi}^{\pi}f'(x)\sin2nx\,dx
\]
\[
=-\frac1{2n\pi}\sum_{k=-n}^{n-1}
\int_{\frac{k\pi}{n}}^{\frac{(k+1)\pi}{n}}f'(x)\sin2nx\,dx.
\]
由第二积分中值定理知，存在
$\xi_k\in\left[\dfrac{k\pi}{n},\dfrac{(k+1)\pi}{n}\right]$，使得
\[
\int_{\frac{k\pi}{n}}^{\frac{(k+1)\pi}{n}}f'(x)\sin2nx\,dx
=f'\left(\frac{k\pi}{n}\right)
\int_{\frac{k\pi}{n}}^{\xi_k}\sin2nx\,dx
+f'\left(\frac{(k+1)\pi}{n}\right)
\int_{\xi_k}^{\frac{(k+1)\pi}{n}}\sin2nx\,dx
\]
\[
=\left[f'\left(\frac{k\pi}{n}\right)
-f'\left(\frac{(k+1)\pi}{n}\right)\right]
\int_{\frac{k\pi}{n}}^{\xi_k}\sin2nx\,dx.
\]
若 $k$ 为偶数，则
\[
\int_{\frac{k\pi}{n}}^{\xi_k}\sin2nx\,dx\ge0;
\]
若 $k$ 为奇数，则
\[
\int_{\frac{k\pi}{n}}^{\xi_k}\sin2nx\,dx\le0.
\]
综合起来，得 $a_{2n}\ge0$.

同理可证 $a_{2n+1}\le0$.

\noindent\textbf{注}\quad 以上两题的已知条件可放松为：“$f'(x)$ 单增有界”（无须凸函数知识）.

\noindent\textbf{new 类题 1}\quad 设 $f(x)\ge0$ 在 $[-\pi,\pi]$ 上单调下降，求证：
\[
\int_{-\pi}^{\pi}f(x)\cos x\,dx\ge0.
\]

\noindent\textbf{new 类题 2}\quad 设 $f(x)$ 在 $[-\pi,\pi]$ 上单调下降，求证：
\[
\int_{-\pi}^{\pi}f(x)\sin x\,dx\ge0.
\]

\noindent\textbf{4.3.27}\quad 已知 $f(x)$ 在 $[a,b]$ 上二阶可导，且 $f''(x)>0$，证明：对任意
$\lambda\in(0,1)$，有
\[
f(\lambda a+(1-\lambda)b)<\lambda f(a)+(1-\lambda)f(b).
\]

\noindent\textbf{证}\quad 令
\[
g(\lambda)=\lambda f(a)+(1-\lambda)f(b)-f(\lambda a+(1-\lambda)b).
\]
则
\[
g''(\lambda)=-(a-b)^2f''(\lambda a+(1-\lambda)b)<0.
\]
因为 $g(0)=g(1)=0$，所以 $g(\lambda)>0\ (0<\lambda<1)$，即结论成立.

\noindent\textbf{4.3.28}\quad 设 $f(x)$ 在 $[a,b]$ 上连续，试证：$f(x)$ 为凸的充分必要条件是
\[
f(x)\le\frac1{2h}\int_{-h}^h f(x+t)\,dt,\quad
\forall [x-h,x+h]\subset[a,b].
\]

\noindent\textbf{证}\quad（必要性）
\[
\frac1{2h}\int_{-h}^hf(x+t)\,dt
=\frac1{2h}\left(\int_0^hf(x+t)\,dt+\int_{-h}^0f(x+t)\,dt\right)
\]
\[
=\frac1h\int_0^h\frac{f(x+t)+f(x-t)}2\,dt
\ge\frac{f(x)}h\int_0^hdt=f(x).
\]

（充分性）已知：$h>0$，$\forall [x-h,x+h]\subset[a,b]$，有
\[
f(x)h\le\frac12\int_{-h}^hf(x+t)\,dt.
\]
两端同时对 $h$ 求导，得
\[
f(x)\le\frac{f(x+h)+f(x-h)}2,
\]
故 $f(x)$ 为 $[a,b]$ 上的凸函数.

\noindent\textbf{new 4.3.29}\quad 设 $f(x)$ 在 $[0,2]$ 上可积，且
\[
\int_0^2f(x)\,dx=0.
\]
求证：存在 $\alpha\in[0,1]$，使得
\[
\int_\alpha^{\alpha+1}f(x)\,dx=0.
\]
（南开大学）

\noindent\textbf{提示}\quad 一方面，在 $[0,1]$ 上 $f(x)$ 可积，可得
\[
F(x)=\int_x^{x+1}f(t)\,dt
\]
连续，有介值性.

另一方面，
\[
F(0)+F(1)=\int_0^2f(x)\,dx=0,
\]
表明 $F(0)$ 与 $F(1)$ 异号，或同时为零，故 $F(x)$ 在 $[0,1]$ 里有零点，即
$\exists\alpha\in[0,1]$，使得
\[
F(\alpha)=\int_\alpha^{\alpha+1}f(x)\,dx=0.
\]
结论获证.

\noindent\textbf{new *4.3.30}\quad 设 $(0,+\infty)$ 内 $f(x)$ 可导，$\sqrt{x}f'(x)$ 有界.

1）证明：$f(x)$ 在 $(0,+\infty)$ 内一致连续；

2）证明：极限 $f(0^+)=\lim\limits_{x\to0^+}f(x)$ 存在，$\lim\limits_{x\to+\infty}f(x)$ 未必存在；

3）将 $\sqrt{x}f'(x)$ 有界的条件改为：$\lim\limits_{x\to+\infty}\sqrt{x}f(x)$ 和
$\lim\limits_{x\to0^+}\sqrt{x}f(x)$ 都存在，那么 $f(x)$ 是否仍在 $(0,+\infty)$ 内一致连续？说明理由.（华中师范大学）

\noindent\textbf{证}\quad 1）证法 I\quad 已知 $\exists M>0$，使得
\[
|\sqrt{x}f'(x)|\le M.
\]
$\forall x_1,x_2\in(0,+\infty)\ (x_1<x_2)$，则有
\[
|f(x_2)-f(x_1)|=\left|\int_{x_1}^{x_2}f'(x)\,dx\right|
\le\int_{x_1}^{x_2}\frac{|\sqrt{x}f'(x)|}{\sqrt{x}}\,dx
\le M\int_{x_1}^{x_2}\frac{dx}{\sqrt{x}}
\]
\[
=2M(\sqrt{x_2}-\sqrt{x_1})\le2M\sqrt{x_2-x_1}. \tag{1}
\]
（最后的一个不等式，通过平方、移项，立即可证.）
$\forall\varepsilon>0$，根据式 (1)，取
\[
\delta=\left(\frac{\varepsilon}{2M}\right)^2,
\]
则 $\forall x_2>x_1>0$，当 $|x_2-x_1|<\delta$ 时，有
\[
|f(x_2)-f(x_1)|\le2M\sqrt{x_2-x_1}<\varepsilon.
\]
证毕.

\noindent\textbf{证法 II}\quad 利用 Cauchy 微分中值定理：
\[
\left|\frac{f(x_2)-f(x_1)}{\sqrt{x_2}-\sqrt{x_1}}\right|
=\left|\frac{f'(\xi)}{\frac1{2\sqrt{\xi}}}\right|\le2M.
\]
同样可得式 (1).

\noindent\textbf{注}\quad 下面的方法只能证明 $f(x)$ 在 $(0,+\infty)$ 的任意内闭区间
$[\alpha,+\infty)$ 上一致连续，未能证明 $f(x)$ 在 $(0,+\infty)$ 上一致连续：

（利用 Lagrange 定理.）$\forall\alpha,x_1,x_2,\ 0<\alpha<x_1<x_2$，有
\[
|f(x_2)-f(x_1)|
=|f'(\xi)(x_2-x_1)|
\le\left|\frac1{\sqrt{\alpha}}\sqrt{\xi}f'(\xi)\right|\,|x_2-x_1|
\le\frac M{\sqrt{\alpha}}|x_2-x_1|.
\]
（$M$ 定义同证法 I.）所以，$\forall\varepsilon>0$，取
$\delta=\dfrac{\varepsilon\sqrt{\alpha}}M$，可使得 $\forall x_1,x_2\in[\alpha,+\infty)$，当
$|x_2-x_1|<\delta$ 时，有
\[
|f(x_2)-f(x_1)|\le\frac M{\sqrt{\alpha}}|x_2-x_1|<\varepsilon.
\]

2）已证 $f(x)$ 在 $(0,+\infty)$ 上一致连续. 如例 2.2.6 所证：极限
\[
\lim_{x\to0^+}f(x)
\]
存在.

当 $x\to+\infty$ 时，极限可以存在也可以不存在. 例如：
\[
f_1(x)\equiv A\quad(\text{常数}),\qquad f_2(x)=\sin\sqrt{x}.
\]
这里 $f_1,f_2$ 都满足题中的条件，极限 $\lim\limits_{x\to+\infty}f_1(x)=A$（存在），而
$\lim\limits_{x\to+\infty}f_2(x)$ 不存在.

3）不能说：“若 $\lim\limits_{x\to+\infty}\sqrt{x}f(x)$ 和
$\lim\limits_{x\to0^+}\sqrt{x}f(x)$ 都存在，则 $f(x)$ 在 $(0,+\infty)$ 上一致连续”. 例如，
\[
f(x)=\sin\frac1x,
\]
这时极限
\[
\lim_{x\to+\infty}\sqrt{x}f(x)=\lim_{x\to0^+}\sqrt{x}f(x)=0
\]
都存在，但 $f(x)=\sin\dfrac1x$ 在 $(0,+\infty)$ 内不一致连续. 因为如令
\[
x_n=(2n\pi)^{-1},\quad x'_n=\left(2n\pi+\frac\pi2\right)^{-1},
\]
当 $n\to+\infty$ 时，虽然 $|x_n-x'_n|\to0$，但
\[
|f(x_n)-f(x'_n)|=1
\]
（不趋向 0）. 所以非一致连续（见例 2.2.3）.

\noindent\textbf{new **4.3.31}\quad 若函数 $f(x)$ 在区间 $[0,1]$ 上 Riemann 可积，并且对
$[0,1]$ 中任意有限个两两不相交的闭区间序列 $[a_i,b_i]$，都有
\[
\left|\sum\int_{a_i}^{b_i}f(x)\,dx\right|\le1.
\]
求证：
\[
\int_0^1|f(x)|\,dx\le2.
\]
（北京大学）

\noindent\textbf{证}\quad 作
\[
f_+(x)=
\begin{cases}
f(x),& f(x)>0,\\
0,& f(x)\le0,
\end{cases}
\qquad
f_-(x)=
\begin{cases}
0,& f(x)>0,\\
-f(x),& f(x)\le0.
\end{cases}
\]
则
\[
f(x)=f_+(x)-f_-(x),\quad |f(x)|=f_+(x)+f_-(x).
\]
因 $f(x)$ 可积，知 $|f(x)|,f_+(x),f_-(x)$ 也都可积（因
$\omega^{|\ |},\omega^+,\omega^-$ 都 $\le\omega^f$）.
对 $[0,1]$ 作分划：
\[
0=x_0<x_1<x_2<\cdots<x_n=1,
\]
记
\[
\lambda=\max_{k=0,\cdots,n-1}\Delta x_k
=\max_{k=0,\cdots,n-1}(x_{k+1}-x_k).
\]
由可积性：$\forall\varepsilon>0$，$\exists\delta_1>0$，当 $\lambda<\delta_1$ 时，有
\[
\left|\int_0^1 f_+(x)\,dx-\sum_{k=0}^{n-1}f_+(\xi_k)\Delta x_k\right|<\frac\varepsilon4. \tag{1}
\]
同理，$\forall\varepsilon>0$，$\exists\delta_2>0$，当 $\lambda<\delta_2$ 时，有
\[
\left|\int_0^1 f_-(x)\,dx-\sum_{k=0}^{n-1}f_-(\eta_k)\Delta x_k\right|<\frac\varepsilon4. \tag{2}
\]
另外，对 $f(x)$，$\exists\delta_3>0$，当 $\lambda<\delta_3$ 时，有
\[
\sum_{k=0}^{n-1}\omega_k^f\Delta x_k<\frac\varepsilon4, \tag{3}
\]
其中 $\omega_k$ 表示 $f(x)$ 在 $[x_k,x_{k+1}]$ 上的振幅. 现令
$\delta=\min\{\delta_1,\delta_2,\delta_3\}$，（对三者）作统一的分划，使得当
$\lambda<\delta$ 时式 (1)，(2) 和 (3) 同时成立.

定义两个指标集：
\[
I=\{k\mid f_+(\xi_k)\ge0\}\quad\text{和}\quad
J=\{k\mid f_-(\eta_k)>0\},
\]
那么
\[
\sum_{k=0}^{n-1}f_+(\xi_k)\Delta x_k
=\sum_{k\in I}f_+(\xi_k)\Delta x_k
=\sum_{k\in I}f(\xi_k)\Delta x_k
\]
\[
=\sum_{k\in I}\left(f(\xi_k)\Delta x_k-\int_{x_k}^{x_{k+1}}f(x)\,dx\right)
+\sum_{k\in I}\int_{x_k}^{x_{k+1}}f(x)\,dx
\]
\[
=\sum_{k\in I}(f(\xi_k)\Delta x_k-\mu_k\Delta x_k)
+\sum_{k\in I}\int_{x_k}^{x_{k+1}}f(x)\,dx
\]
\[
\le\sum_{k=0}^{n-1}\omega_k^f\Delta x_k+1
<\frac\varepsilon4+1.
\]
亦即
\[
\sum_{k=0}^{n-1}f_+(\xi_k)\Delta x_k<\frac\varepsilon4+1. \tag{4}
\]
同理，
\[
\sum_{k=0}^{n-1}f_-(\eta_k)\Delta x_k<\frac\varepsilon4+1. \tag{5}
\]
最后
\[
\int_0^1|f(x)|\,dx=\int_0^1(f_+(x)+f_-(x))\,dx=I_1+I_2,
\]
其中
\[
I_1=\left(\int_0^1f_+(x)\,dx-\sum_{k=0}^{n-1}f_+(\xi_k)\Delta x_k\right)
+\sum_{k=0}^{n-1}f_+(\xi_k)\Delta x_k
<\frac\varepsilon4+\frac\varepsilon4+1=\frac\varepsilon2+1.
\]
同理，$I_2<\dfrac\varepsilon2+1$. 因此
\[
\int_0^1|f(x)|\,dx<\left(\frac\varepsilon2+1\right)
+\left(\frac\varepsilon2+1\right)=\varepsilon+2.
\]
因为 $\varepsilon>0$ 的任意性. 令 $\varepsilon\to0$，得
\[
\int_0^1|f(x)|\,dx\le2.
\]
（证毕.）

\noindent\textbf{new 4.3.32}\quad（积分等式）设 $a>0$，$f(x)$ 是定义在 $[-a,a]$ 上的连续偶函数. 试证：
\[
\int_{-a}^a\frac{f(x)}{1+e^x}\,dx=\int_0^a f(x)\,dx.
\]
（中国科学院）

\noindent\textbf{证}\quad
\[
\int_{-a}^a\frac{f(x)}{1+e^x}\,dx
\stackrel{x=-y}{\underset{f\text{ 为偶}}=}
\int_a^{-a}\frac{f(y)}{1+e^{-y}}(-dy)
=\int_{-a}^a\frac{f(y)}{1+e^{-y}}\,dy
\]
\[
=\int_{-a}^a\frac{f(x)}{1+e^{-x}}\,dx
\stackrel{\text{分子、分母同乘 }e^x}{=}
\int_{-a}^a\frac{f(x)e^x}{e^x+1}\,dx.
\]
故
\[
\int_{-a}^a\frac{f(x)}{1+e^x}\,dx
=\frac12\left(
\int_{-a}^a\frac{f(x)}{1+e^x}\,dx
+\int_{-a}^a\frac{f(x)e^x}{e^x+1}\,dx\right)
\]
\[
=\frac12\int_{-a}^a f(x)\,dx
=\int_0^a f(x)\,dx
\quad(\text{因为 }f(x)\text{ 为偶函数}).
\]

\noindent\textbf{new ※4.3.33}\quad（为了叙述方便引入一个名词）称 $f$ 是区间 $[a,b]$ 上的 $F$ 函数，意指它满足条件：

① $f$ 在 $[a,b]$ 上有连续导数，② $(a,b)$ 内二阶可导，③ 区间中点
$c=\dfrac{a+b}2$ 上，$f'(c)=0$.

试问下面关于 $F$ 函数的论断是否成立？说明理由：

1）若 $f$ 是区间 $[a,b]$ 上的 $F$ 函数，则必存在 $\xi\in(a,b)$，使得
\[
|f''(\xi)|\ge\frac4{(b-a)^2}|f(b)-f(a)|; \tag{1}
\]

2）试证：式 (1) 中的 4 是最好的，即：若将式 (1) 里的 4 换成任一较大的 $M$，则结论 1）不再成立.

3）$\forall f\in F$，若 $f(x)\ne$ 常数，则 $\exists\xi\in(a,b)$，使得
\[
|f''(\xi)|>\frac4{(b-a)^2}|f(b)-f(a)|. \tag{2}
\]
（南开大学）

\noindent\textbf{注}\quad 肯定形式变成否定形式的法则是：将逻辑量词“$\forall$”改成“$\exists$”，将“$\exists$”改成“$\forall$”，然后不等式反向. 例如：本题

结论 1）意指：$\forall f\in F$，$\exists\xi\in(a,b)$ 使得不等式 (1) 成立；

\noindent 结论 2）意指：若将式 (1) 中的 4 换成 $\forall M=4+\varepsilon(\varepsilon>0)$，结论 1）不再成立. 即 $\forall \xi\in(a,b)$，
$\exists f\in F$ 使得（换后）不等式 (1) 不再成立，亦即下式成立：
\[
|f''(\xi)|<\frac{M}{(b-a)^2}|f(b)-f(a)|\quad(\forall \xi\in(a,b)). \tag{3}
\]

\noindent\textbf{提示}\quad 1）参考例 3.3.6 的证明方法；2）构造反例；3）利用 Lagrange 定理.

\noindent\textbf{解}\quad 1）应用 Taylor 公式（已知区间中点处 $f'(c)=0$），
\[
f(a)=f(c)+\frac{f''(\eta_1)}{2!}\left(\frac{b-a}{2}\right)^2,\quad
f(b)=f(c)+\frac{f''(\eta_2)}{2!}\left(\frac{b-a}{2}\right)^2. \tag{4}
\]
由式 (4) 得
\[
\begin{aligned}
|f(b)-f(a)|
&=\left|\frac{f''(\eta_1)-f''(\eta_2)}2\right|
\left(\frac{b-a}{2}\right)^2\\
&\le \frac{|f''(\eta_1)|+|f''(\eta_2)|}{2}
\left(\frac{b-a}{2}\right)^2\\
&\le |f''(\xi)|\frac{(b-a)^2}{4},
\end{aligned}
\]
其中 $|f''(\xi)|\stackrel{\text{记}}{=}\max_{k=1,2}\{|f''(\eta_k)|\}$.
再同乘 $\dfrac4{(b-a)^2}$ 即得式 (1)，结论 1）获证.

2）为了证明结论 1）不再成立，只需构造一个函数 $f\in F$ 使得 $\forall \xi\in(a,b)$，式 (3) 成立.
最简便是利用幂函数 $f(x)=x^\alpha$，则 $f'(x)=\alpha x^{\alpha-1}$，
$f''(x)=\alpha(\alpha-1)x^{\alpha-2}$. 取区间 $[a,b]=[-1,1]$，
则中点 $0$ 上：
\[
f'(0)=\alpha x^{\alpha-1}\big|_{x=0}=0.
\]
$F$ 的三个条件都成立，剩下只需选择 $\alpha$ 使式 (3) 成立. 此时，式 (3) 是
\[
|\alpha(\alpha-1)x^{\alpha-2}|
<\frac{4+\varepsilon}{[1-(-1)]^2}|1-(-1)^\alpha|
\quad(\forall x\in(-1,1)). \tag{3'}
\]
若取 $\alpha=2+\dfrac1m$（$m=2k+1$ 为奇数），则式 (3') 里：
$|x^{\alpha-2}|<1$，$(-1)^\alpha=(-1)^{2+\frac1m}=-1$. 于是
\[
\text{式 (3') 左端：}\quad
|\alpha(\alpha-1)x^{\alpha-2}|<|\alpha(\alpha-1)|
=2+\frac3m+\frac1{m^2}<2+\frac4m;
\]
\[
\text{式 (3') 右端：}\quad
\frac{4+\varepsilon}{[1-(-1)]^2}|1-(-1)^\alpha|
=2+\frac\varepsilon2.
\]
因式 (3) 与 (3') 等价，只要令 $f(x)=x^{2+\frac1m}$，并让 $m$ 为充分大的奇数，就可以使得
\[
\text{式 (3') 左端}<2+\frac4m<2+\frac\varepsilon2=\text{式 (3') 右端},
\]
亦即 $m>\dfrac8\varepsilon$，就能保证式 (3') 成立，结论 2）获证.

3）（反证法）假设结论 3）不成立，即 $\forall \xi\in[a,b]$，有
\[
|f''(\xi)|\le \frac4{(b-a)^2}|f(b)-f(a)|
=\frac4{(b-a)^2}\left|\int_a^b f'(x)\,dx\right|
\le \frac4{(b-a)^2}\int_a^b |f'(x)|\,dx
\stackrel{\text{记}}{=} k\ge0. \tag{5}
\]
区间 $[a,b]$ 的中点 $c=\dfrac{a+b}{2}$ 处 $f'(c)=0$. 故 $\forall x\in[a,b]$，应用 Lagrange 定理得
\[
|f'(x)|=|f'(x)-f'(c)|=|f''(\xi)(x-c)|\le k|x-c|. \tag{6}
\]
于是
\[
\int_a^b |f'(t)|\,dt
\le \int_a^c k(c-t)\,dt+\int_c^b k(t-c)\,dt
=\frac k2[(c-a)^2+(b-c)^2]
=\frac{k(b-a)^2}{4}
\stackrel{\text{式(5)}}{=}\int_a^b |f'(t)|\,dt.
\]
此不等式两端是同一积分，故上式只能是等号. 根据积分性质（例 4.2.12），若非负函数积分相等，则被积函数也相等.
因此式 (6) 里的“$\le$”也只能是“$=$”. 于是式 (6) 变为
\[
|f'(x)|=
\begin{cases}
k(c-x),&x\le c,\\
k(x-c),&x>c,
\end{cases}
\quad k\ge0. \tag{7}
\]
亦即
\[
f'(x)=
\begin{cases}
k(c-x),&x\le c,\\
k(x-c),&x>c,
\end{cases}
\quad\text{或}\quad
f'(x)=
\begin{cases}
-k(c-x),&x\le c,\\
-k(x-c),&x>c.
\end{cases}
\]
可见
\[
f''_+(c)=\lim_{x\to c^+}\frac{k(x-c)-0}{x-c}=k,\quad
f''_-(c)=\lim_{x\to c^-}\frac{k(c-x)-0}{x-c}=-k,
\]
或
\[
f''_+(c)=\lim_{x\to c^+}\frac{-k(x-c)-0}{x-c}=-k,\quad
f''_-(c)=\lim_{x\to c^-}\frac{-k(c-x)-0}{x-c}=k.
\]
如果 $k\ne0$，则 $f''_+(c)\ne f''_-(c)$，说明 $f''(c)$ 不存在，与题设矛盾.
如果 $k=0$，式 (7) 表明：$\forall x\in[a,b]$，$f'(x)\equiv0$，$f(x)$ 为常值函数，也与题设矛盾.

\section*{§4.4\quad 几个著名的不等式}

本节讨论几个著名的不等式，这些不等式不仅本身是重要的，而且证明这些不等式的方法也十分典型. 因此，本节较系统地介绍这些不等式，并着重讨论它们的证明、变形与应用.

\subsection*{一、Cauchy 不等式及 Schwarz 不等式}

\noindent\textbf{a. Cauchy 不等式}

\noindent\textbf{定理 1}\quad 设 $a_i,b_i$ 为任意实数（$i=1,2,\cdots,n$），则
\[
\left(\sum_{i=1}^n a_i b_i\right)^2
\le \sum_{i=1}^n a_i^2\cdot \sum_{i=1}^n b_i^2, \tag{1}
\]
其中等号当且仅当 $a_i$ 与 $b_i$ 成比例时成立.

式 (1) 称为 Cauchy 不等式.（注意此定理以后应用很广泛.）

\noindent\textbf{证 I}\quad（判别式法）因为
\[
0\le \sum_{i=1}^n(a_i x+b_i)^2
=\left(\sum_{i=1}^n a_i^2\right)x^2
+2\left(\sum_{i=1}^n a_i b_i\right)x
+\left(\sum_{i=1}^n b_i^2\right),
\]
关于 $x$ 的二次三项式保持非负，故判别式
\[
\left(\sum_{i=1}^n a_i b_i\right)^2
-\sum_{i=1}^n a_i^2\cdot\sum_{i=1}^n b_i^2\le0.
\]

\noindent\textbf{证 II}\quad（配方法）因为
\[
\begin{aligned}
&\sum_{i=1}^n a_i^2\cdot\sum_{j=1}^n b_j^2
-\left(\sum_{i=1}^n a_i b_i\right)^2\\
={}&\sum_{i=1}^n a_i^2\cdot\sum_{j=1}^n b_j^2
-\sum_{i=1}^n a_i b_i\cdot\sum_{j=1}^n a_j b_j\\
={}&\sum_{i=1}^n\sum_{j=1}^n a_i^2b_j^2
-\sum_{i=1}^n\sum_{j=1}^n a_i b_i a_j b_j\\
={}&\frac12\sum_{i,j=1}^n(a_i b_j-a_j b_i)^2\ge0.
\end{aligned}
\]
故式 (1) 获证，等号当且仅当 $a_i b_j=a_j b_i$（$i,j=1,2,\cdots,n$）时成立.

\noindent\textbf{证 III}\quad（利用二次型）因为
\[
0\le\sum_{i=1}^n(a_i x+b_i y)^2
=\left(\sum_{i=1}^n a_i^2\right)x^2
+2\left(\sum_{i=1}^n a_i b_i\right)xy
+\left(\sum_{i=1}^n b_i^2\right)y^2,
\]
即关于 $x,y$ 的二次型非负定，故
\[
\begin{vmatrix}
\sum_{i=1}^n a_i^2&\sum_{i=1}^n a_i b_i\\
\sum_{i=1}^n a_i b_i&\sum_{i=1}^n b_i^2
\end{vmatrix}\ge0,
\]
此即式 (1).

\noindent\textbf{注}\quad 用证 III，易将结果进行推广. 因
\[
0\le\sum_{i=1}^n(a_{i1}x_1+a_{i2}x_2+\cdots+a_{im}x_m)^2
=\sum_{k,j=1}^m\left(\sum_{i=1}^n a_{ik}a_{ij}\right)x_kx_j,
\]
此式右边为 $x_1,x_2,\cdots,x_m$ 的二次型，此式表明该二次型非负定，因此系数行列式
\[
\operatorname{Det}\left(\sum_{i=1}^n a_{ik}a_{ij}\right)
=
\begin{vmatrix}
\sum_{i=1}^n a_{i1}^2&\sum_{i=1}^n a_{i1}a_{i2}&\cdots&\sum_{i=1}^n a_{i1}a_{im}\\
\sum_{i=1}^n a_{i2}a_{i1}&\sum_{i=1}^n a_{i2}^2&\cdots&\sum_{i=1}^n a_{i2}a_{im}\\
\vdots&\vdots&\ddots&\vdots\\
\sum_{i=1}^n a_{im}a_{i1}&\sum_{i=1}^n a_{im}a_{i2}&\cdots&\sum_{i=1}^n a_{im}^2
\end{vmatrix}\ge0, \tag{2}
\]
等号当且仅当 $(a_{11},a_{21},\cdots,a_{n1})$，$(a_{12},a_{22},\cdots,a_{n2})$，$\cdots$，
$(a_{1m},a_{2m},\cdots,a_{nm})$ 线性相关（即 $\exists$ 不全为零的常数 $x_1,\cdots,x_m$ 使得
\[
a_{i1}x_1+a_{i2}x_2+\cdots+a_{im}x_m=0\quad(i=1,2,\cdots,n)
\]
成立）时取得. 式 (2) 是 Cauchy 不等式的推广形式.

\noindent\textbf{b. Schwarz 不等式}

Cauchy 不等式的积分形式称为 Schwarz 不等式，它可以通过积分定义，直接由 Cauchy 不等式推得；也可仿照 Cauchy 不等式的证法类似证明.

\noindent\textbf{定理 2}\quad 若 $f(x),g(x)$ 在 $[a,b]$ 上可积，则
\[
\left(\int_a^b f(x)g(x)\,dx\right)^2
\le \int_a^b f^2(x)\,dx\int_a^b g^2(x)\,dx. \tag{1}
\]
若 $f(x),g(x)$ 在 $[a,b]$ 上连续，其中等号当且仅当存在常数 $\alpha,\beta$，使得
$\alpha f(x)\equiv\beta g(x)$ 时成立（$\alpha,\beta$ 不同时为零）.（南京理工大学等）

（注意此定理以后经常用到.）

\noindent\textbf{证 I}\quad 将 $[a,b]$ $n$ 等分，令
\[
x_i=a+\frac in(b-a),
\]
应用 Cauchy 不等式，
\[
\left(\frac1n\sum_{i=1}^n f(x_i)g(x_i)\right)^2
\le \frac1n\sum_{i=1}^n f^2(x_i)\cdot \frac1n\sum_{i=1}^n g^2(x_i).
\]
令 $n\to\infty$，取极限即得式 (1).

\noindent\textbf{证 II}\quad
\[
\begin{aligned}
&\int_a^b f^2(x)\,dx\int_a^b g^2(x)\,dx
-\left(\int_a^b f(x)g(x)\,dx\right)^2\\
={}&\frac12\int_a^b f^2(x)\,dx\int_a^b g^2(y)\,dy
+\frac12\int_a^b f^2(y)\,dy\int_a^b g^2(x)\,dx\\
&-\int_a^b f(x)g(x)\,dx\int_a^b f(y)g(y)\,dy\\
={}&\frac12\int_a^b dy\int_a^b
[f^2(x)g^2(y)+f^2(y)g^2(x)-2f(x)g(x)f(y)g(y)]\,dx\\
={}&\frac12\int_a^b dy\int_a^b[f(x)g(y)-g(x)f(y)]^2\,dx\ge0,
\end{aligned}
\]
这就证明了式 (1). 由此看出，若 $f(x),g(x)$ 连续，等号当且仅当存在常数 $\alpha,\beta$（不全为零）使得
$\alpha f(x)\equiv\beta g(x)$ 时成立.

还可用本节定理 1 中证 I 和证 III 类似的方法证明.

类似可以推广到一般的情况. 若函数 $f_i(x),g_i(x)$（$i=1,2,\cdots,m$）在 $[a,b]$ 上可积，则
\[
\operatorname{Det}\left(\int_a^b f_i(x)f_j(x)\,dx\right)\ge0.
\]
若 $f_i(x)$ 在 $[a,b]$ 上连续，其中等号当且仅当 $f_i(x)$（$i=1,2,\cdots,m$）线性相关（即 $\exists$ 不全为零的常数
$\alpha_1,\alpha_2,\cdots,\alpha_m$ 使得
$\alpha_1f_1(x)+\alpha_2f_2(x)+\cdots+\alpha_mf_m(x)\equiv0$）时成立.

\noindent\textbf{c. Schwarz 不等式的应用}

应用 Schwarz 不等式，可证明另外一些不等式. 使用时，要注意恰当地选取函数 $f(x)$ 与 $g(x)$.

\noindent\textbf{例 4.4.1}\quad 已知 $f(x)\ge0$，在 $[a,b]$ 上连续，$\displaystyle\int_a^b f(x)\,dx=1$，$k$ 为任意实数，求证：
\[
\left(\int_a^b f(x)\cos kx\,dx\right)^2+
\left(\int_a^b f(x)\sin kx\,dx\right)^2\le1. \tag{1}
\]
（中国科学技术大学）

\noindent\textbf{证}\quad 式 (1) 左端第一项应用 Schwarz 不等式，得
\begin{equation}
\begin{aligned}
\left(\int_a^b f(x)\cos kx\,dx\right)^2
&=\left[\int_a^b \sqrt{f(x)}(\sqrt{f(x)}\cos kx)\,dx\right]^2\\
&\le \int_a^b f(x)\,dx\cdot\int_a^b f(x)\cos^2kx\,dx
=\int_a^b f(x)\cos^2kx\,dx.
\end{aligned}
\tag{2}
\end{equation}
同理
\[
\left(\int_a^b f(x)\sin kx\,dx\right)^2
\le\int_a^b f(x)\sin^2kx\,dx. \tag{3}
\]
(2)+(3) 即得式 (1).

\noindent\textbf{练习}\quad 1）设 $f(x)$ 在 $[a,b]$ 上连续，证明不等式
\[
\left(\int_a^b f(x)\,dx\right)^2\le(b-a)\int_a^b f^2(x)\,dx;
\]
（北京大学）

2）设 $h(x)$ 是 $[a,b]$ 上的正值连续函数，求证
\[
\int_a^b h(t)\,dt\cdot\int_a^b\frac1{h(t)}\,dt\ge(b-a)^2;
\]
（中国科学院，哈尔滨工业大学）

\noindent\textbf{new 3）}\quad 设实数函数 $f(x)$ 及一阶导函数在 $[a,b]$ 上连续，$f(a)=0$，试证：
\[
\text{i）}\quad
\max_{x\in[a,b]}f(x)\le\sqrt{b-a}\int_a^b |f'(t)|^2\,dt;
\quad
\text{ii）}\quad
\int_a^b f^2(t)\,dt\le\frac12(b-a)^2\int_a^b |f'(t)|^2\,dt.
\]
（中国科学院）

\noindent\textbf{new 4）}\quad 设 $f(x)$ 在 $[0,1]$ 上连续，且 $\displaystyle\int_0^1 f(x)\,dx=1$. 试证
\[
\int_0^1(1+x^2)f^2(x)\,dx\ge\frac4\pi.
\]

\noindent\textbf{提示}\quad 1）对 $\left(\displaystyle\int_a^b 1\cdot f(x)\,dx\right)^2$ 用 Schwarz 不等式；

2）$(b-a)^2=\left(\displaystyle\int_a^b \sqrt{h(t)}\cdot\dfrac1{\sqrt{h(t)}}\,dt\right)^2$. 有时需要对积分作适当变形，才能用 Schwarz 不等式.

3）i）因 $f$ 在 $[a,b]$ 上达最大，故 $\exists x_0\in[a,b]$ 使得
\[
\max_{x\in[a,b]}f(x)=f(x_0)\stackrel{f(a)=0}{=}\int_a^{x_0}f'(t)\,dt=\int_a^{x_0}1\cdot f'(t)\,dt.
\]
再用 Schwarz 不等式可得.
\[
\text{ii）}\quad f^2(x)=\left(\int_a^x f'(t)\,dt\right)^2\quad(\text{应用 Schwarz 不等式})
\]
\[
\le \int_a^x(f'(t))^2\,dt\cdot\int_a^x(1)^2\,dt
\le (x-a)\int_a^b(f'(t))^2\,dt.
\]
两端同时积分可得.

4）
\[
1=\left(\int_0^1 f(x)\,dx\right)^2
=\left[\int_0^1\frac1{\sqrt{1+x^2}}(\sqrt{1+x^2}f(x))\,dx\right]^2\quad(\text{应用 Schwarz 不等式})
\]
\[
\le \int_0^1\left(\frac1{\sqrt{1+x^2}}\right)^2\,dx\cdot
\int_0^1(\sqrt{1+x^2}f(x))^2\,dx
=\frac\pi4\int_0^1(1+x^2)(f(x))^2\,dx.
\]

\noindent\textbf{例 4.4.2}\quad 设函数 $g(x)$ 在 $[0,a]$ 上连续可微，$g(0)=0$，试证：
\[
\int_0^a |g(x)g'(x)|\,dx\le\frac a2\int_0^a |g'(x)|^2\,dx, \tag{1}
\]
其中等号当且仅当 $g(x)=cx$（$c$ 为常数）时成立.（北京师范大学）

\noindent\textbf{证}\quad $1^\circ$ 记
\[
f(x)=\int_0^x |g'(t)|\,dt\quad(0\le x\le a),
\]
则 $f'(x)=|g'(x)|$，由 $g(0)=0$ 知
\[
|g(x)|=|g(x)-g(0)|=\left|\int_0^x g'(t)\,dt\right|
\le \int_0^x |g'(t)|\,dt=f(x),
\]
因此
\[
\int_0^a |g(x)g'(x)|\,dx
\le \int_0^a f(x)f'(x)\,dx
=\int_0^a f(x)\,df(x)
\]
\[
=\frac12 f^2(x)\bigg|_0^a
=\frac12\left(\int_0^a |g'(t)|\,dt\right)^2
=\frac12\left(\int_0^a 1\cdot |g'(t)|\,dt\right)^2
\quad(\text{应用 Schwarz 不等式})
\]
\[
\le \frac12\int_0^a 1^2\,dx\cdot\int_0^a |g'(t)|^2\,dt
=\frac a2\int_0^a g'^2(t)\,dt.
\]
式 (1) 获证.

$2^\circ$ 当 $g=cx$ 时，式 (1) 明显成立，只需证明必要性. 如上已证
\[
\int_0^a |g(x)g'(x)|\,dx
\le \frac12\left(\int_0^a |g'(t)|\,dt\right)^2
\le \frac a2\int_0^a g'^2(t)\,dt.
\]
若式 (1) 中等号成立，则
\[
\left(\int_0^a |g'(t)|\,dt\right)^2=a\int_0^a g'^2(t)\,dt. \tag{2}
\]
记
\[
A=\int_0^a g'^2(t)\,dt,\quad B=\int_0^a |g'(t)|\,dt,
\]
于是式 (2) 相当于方程
\[
\int_0^a (1+\lambda |g'(x)|)^2\,dx=A\lambda^2+2B\lambda+a=0 \tag{3}
\]
的判别式 $\Delta=0$. 因而二次方程 (3) 有唯一根：
\[
\lambda_0=-\frac BA\quad(A\ne0). \tag{4}
\]
但 $g'(x)$ 在 $[0,a]$ 上连续，（(4) 代入 (3)）由
\[
\int_0^a\left(1-\frac BA|g'(x)|\right)^2\,dx=0
\]
可得 $B|g'(x)|=A$.
当 $A\ne0$ 时，$B$ 也不为零，故
\[
g'(x)=\pm\frac AB,\quad g(x)=\pm\frac ABx+c_1.
\]
又由于 $g(0)=0$，$c_1=0$，所以 $g(x)=cx$（$c=\pm A/B$ 为常数）.

最后，假若 $A=0$，即 $\displaystyle\int_0^a g'^2(x)\,dx=0$，由 $g'(x)$ 连续知 $g'(x)\equiv0$（在 $[0,a]$ 上）.
从而 $g(x)=c_2$，但 $g(0)=0$，所以 $g(x)\equiv0$，属于 $g(x)=cx$ 中 $c=0$ 的特殊情况.
总之，不论 $A$ 是否为 0，当式 (1) 等号成立时，$g(x)=cx$（$c$ 为常数）. 必要性获证.

\noindent\textbf{注}\quad 对任意区间 $[a,b]$，若 $g'$ 连续，$g(a)=0$，则有
\[
\int_a^b |g(x)g'(x)|\,dx\le\frac{b-a}{2}\int_a^b (g'(x))^2\,dx,
\]
其中等号当且仅当 $g(x)=c(x-a)$ 时成立（$c$ 为常数）.

\noindent\textbf{例 4.4.3}\quad 假设函数 $f(x)$ 在闭区间 $[a,b]$（$b>a$）上有连续 $n$ 阶导数 $f^{(n)}(x)$，
并且 $f^{(k)}(a)=0$，$k=0,1,\cdots,n-1$. 求证：
\[
\left[\int_a^b (f^{(k)}(x))^2\,dx\right]^{\frac12}
\le \left(\frac12\right)^{\frac{m-k}{2}}(b-a)^{m-k}
\left[\int_a^b (f^{(m)}(x))^2\,dx\right]^{\frac12}, \tag{1}
\]
这里 $0\le k<m\le n$.（中山大学）

\noindent\textbf{分析}\quad $1^\circ$ 先设法证明最简单而又必须证明的情况. 令 $n=1$（此时 $k=0,m=1$），
我们所要证明的结论是
\[
\text{假若 }\varphi(x)\text{ 在 }[a,b]\text{ 上有连续的导数，}\varphi(a)=0,\text{ 则必有}
\]
\[
\left[\int_a^b(\varphi(x))^2\,dx\right]^{\frac12}
\le \left(\frac12\right)^{\frac12}(b-a)
\left[\int_a^b(\varphi'(x))^2\,dx\right]^{\frac12}. \tag{2}
\]
为把 $\varphi$ 与 $\varphi'$ 联系起来，用公式
\[
\varphi(x)=\int_a^x\varphi'(x)\,dx.
\]
应用 Schwarz 公式，
\[
(\varphi(x))^2=\left(\int_a^x\varphi'(t)\,dt\right)^2
\le \int_a^x1^2\,dt\cdot\int_a^x(\varphi'(t))^2\,dt
=(x-a)\int_a^x\varphi'^2(t)\,dt. \tag{3}
\]
两边同时积分，
\[
\begin{aligned}
\int_a^b(\varphi(x))^2\,dx
&\le\int_a^b\left[(x-a)\int_a^x\varphi'^2(t)\,dt\right]dx\\
&=\frac12\int_a^b\left(\int_a^x\varphi'^2(t)\,dt\right)d(x-a)^2\quad(\text{应用分部积分法})\\
&=\frac12(x-a)^2\left(\int_a^x\varphi'^2(t)\,dt\right)\bigg|_{x=a}^{x=b}
-\frac12\int_a^b(x-a)^2\varphi'^2(x)\,dx\\
&\le\frac{(b-a)^2}{2}\int_a^b\varphi'^2(t)\,dt.
\end{aligned}
\]
两边同时开方，便得欲证的式 (2).

$2^\circ$ 回到一般情况，令 $\varphi(x)=f^{(k)}(x)$，反复应用我们刚刚证得的不等式 (2) $m-k$ 次，便可得欲证的不等式 (1).

下例用 Schwarz 不等式求极限.

\noindent\textbf{例 4.4.4}\quad 设 $f(x),g(x)$ 在 $[a,b]$ 上连续，$f(x)\ne0$，$g(x)$ 有正下界. 记
\[
d_n=\int_a^b |f(x)|^n g(x)\,dx,\quad n=1,2,\cdots.
\]
试证
\[
\lim_{n\to\infty}\frac{d_{n+1}}{d_n}=\max_{a\le x\le b}|f(x)|.
\]
（南开大学，四川大学）

\noindent\textbf{证}\quad $1^\circ$（为了分析 $\left\{\dfrac{d_{n+1}}{d_n}\right\}$ 的变化状态，我们先研究 $d_n$ 相邻项之间的关系.）
\[
\begin{aligned}
d_n&=\int_a^b |f(x)|^n g(x)\,dx\\
&=\int_a^b \sqrt{g(x)}\,|f(x)|^{\frac{n-1}{2}}\cdot
\sqrt{g(x)}\,|f(x)|^{\frac{n+1}{2}}\,dx\quad(\text{应用 Schwarz 不等式})\\
&\le\left(\int_a^b g(x)|f(x)|^{n-1}\,dx\right)^{\frac12}
\left(\int_a^b g(x)|f(x)|^{n+1}\,dx\right)^{\frac12}
=d_{n-1}^{\frac12}\cdot d_{n+1}^{\frac12}.
\end{aligned}
\]
因 $d_n>0$，平方得 $d_n^2\le d_{n-1}d_{n+1}$，即
\[
\frac{d_{n+1}}{d_n}\ge\frac{d_n}{d_{n-1}},
\]
故 $\dfrac{d_{n+1}}{d_n}$ 递增.

$2^\circ$ 因 $f(x)$ 在 $[a,b]$ 上连续，$\exists M>0$，使得 $|f(x)|\le M$ 于 $[a,b]$ 上，故
\[
0\le\frac{d_{n+1}}{d_n}
=\frac{\int_a^b g(x)|f(x)|^{n+1}\,dx}{\int_a^b g(x)|f(x)|^n\,dx}
\le \frac{M\int_a^b g(x)|f(x)|^n\,dx}{\int_a^b g(x)|f(x)|^n\,dx}=M.
\]

$3^\circ$ 既然 $\left\{\dfrac{d_{n+1}}{d_n}\right\}$ 单调有界，所以有极限. 根据例 1.3.4 中 2）的结果，
\[
\lim_{n\to\infty}\frac{d_{n+1}}{d_n}
=\lim_{n\to\infty}\sqrt[n]{d_n}
=\lim_{n\to\infty}\left(\int_a^b g(x)|f(x)|^n\,dx\right)^{\frac1n}
=\max_{a\le x\le b}|f(x)|.
\]

\noindent\textbf{二、平均值不等式}

\noindent\textbf{a. 基本形式}

\noindent\textbf{定理 3}\quad 对任意 $n$ 个实数 $a_i\ge0$（$i=1,2,\cdots,n$），恒有
\[
\sqrt[n]{a_1a_2\cdots a_n}\le \frac{a_1+a_2+\cdots+a_n}{n} \tag{1}
\]
（即几何平均值 $\le$ 算术平均值），其中等号当且仅当 $a_1=a_2=\cdots=a_n$ 时成立（证明见第一章 §1.1）.

\noindent\textbf{例 4.4.5}\quad 设正值函数 $f(x)$ 在 $[0,1]$ 上连续，试证：
\[
\exp\left(\int_0^1\ln f(x)\,dx\right)\le\int_0^1 f(x)\,dx.
\]
（中国科学院）

\noindent\textbf{证 I}\quad 由条件知 $f(x),\ln f(x)$ 在 $[0,1]$ 上可积. 将 $[0,1]$ $n$ 等分，作积分和，
\[
\int_0^1 f(x)\,dx=\lim_{n\to\infty}\frac1n\sum_{i=1}^n f\left(\frac in\right),
\]
\[
\int_0^1\ln f(x)\,dx
=\lim_{n\to\infty}\frac1n\sum_{i=1}^n\ln f\left(\frac in\right)
=\lim_{n\to\infty}\ln\left[\prod_{i=1}^n f\left(\frac in\right)\right]^{\frac1n}.
\]
所以
\[
\exp\left(\int_0^1\ln f(x)\,dx\right)
=\exp\left[\lim_{n\to\infty}\ln\left(\prod_{i=1}^n f\left(\frac in\right)\right)^{\frac1n}\right]
=\lim_{n\to\infty}\left[\prod_{i=1}^n f\left(\frac in\right)\right]^{\frac1n}.
\]
应用定理 3，
\[
\left[\prod_{i=1}^n f\left(\frac in\right)\right]^{\frac1n}
\le\frac1n\sum_{i=1}^n f\left(\frac in\right),
\]
故
\[
\exp\left(\int_0^1\ln f(x)\,dx\right)\le\int_0^1 f(x)\,dx.
\]

\noindent\textbf{证 II}\quad 取对数，原式等价于
\[
\int_0^1\ln f(x)\,dx\le\ln\int_0^1 f(x)\,dx\stackrel{\text{记}}{=}\ln S,
\]
即要证
\[
\int_0^1(\ln f(x)-\ln S)\,dx\le0.
\]
事实上，
\[
\begin{aligned}
\int_0^1(\ln f(x)-\ln S)\,dx
&=\int_0^1\ln\frac{f(x)}S\,dx\\
&=\int_0^1\ln\left[1+\left(\frac{f(x)}S-1\right)\right]dx
\quad(\text{利用 }\ln(1+x)<x\ (x>-1))\\
&\le\int_0^1\left(\frac{f(x)}S-1\right)dx
=\frac1S\int_0^1 f(x)\,dx-1=0.
\end{aligned}
\]
证毕.

\noindent\textbf{b. 平均值不等式的推广形式}

\noindent\textbf{定义 1}\quad 设 $a_i\ge0$（$i=1,2,\cdots,n$），记
\[
M_r(a)=\left(\frac1n\sum_{i=1}^n a_i^r\right)^{\frac1r}\quad(r>0);
\]
称 $M_r(a)$ 为 $a_1,\cdots,a_n$ 的 $r$ 次幂平均. 它与算术平均值的关系是
\[
M_1(a)=\frac{a_1+a_2+\cdots+a_n}{n}=A(a),\quad
M_r(a)=(A(a^r))^{\frac1r}.
\]

\noindent\textbf{定义 2（加权平均）}\quad $p_i>0$（$i=1,2,\cdots,n$），记
\[
M_r(a,p)=\left(\frac{\sum_{i=1}^n p_i a_i^r}{\sum_{i=1}^n p_i}\right)^{\frac1r}.
\]
\[
G(a,p)=\left(\prod_{i=1}^n a_i^{p_i}\right)^{1/\sum_{i=1}^n p_i}
=(a_1^{p_1}a_2^{p_2}\cdots a_n^{p_n})^{\frac1{p_1+p_2+\cdots+p_n}}.
\]
$M_r(a,p)$ 和 $G(a,p)$ 分别称为 $a_1,a_2,\cdots,a_n$ 的加权（$r$ 次幂）算术平均和加权几何平均，
$p_i$（$i=1,2,\cdots,n$）称为权数.

若令
\[
q_i=\frac{p_i}{p_1+\cdots+p_n},
\]
则 $\sum_{i=1}^n q_i=1$，这时
\[
M_r(a,q)=\left(\sum_{i=1}^n q_i a_i^r\right)^{\frac1r},\quad
G(a,q)=\prod_{i=1}^n a_i^{q_i}=a_1^{q_1}a_2^{q_2}\cdots a_n^{q_n}.
\]
将 $p_i$ 改成 $q_i$（$i=1,2,\cdots,n$）称为权数的标准化.（为了简洁，以下我们将连加和连乘符号中的标号略去.）

\noindent\textbf{引理 1}\quad 设 $r>0$，$a_1,a_2,\cdots,a_n$ 不全相等，则
\[
M_r(a,q)>M_{\frac r2}(a,q).
\]

\noindent\textbf{证}\quad
\[
\begin{aligned}
M_{\frac r2}(a,q)
&=\left(\sum q_i a_i^{\frac r2}\right)^{\frac2r}
=\left(\sum \sqrt{q_i}\sqrt{q_i a_i^r}\right)^{\frac2r}\quad(\text{应用 Cauchy 不等式})\\
&\le\left(\sum q_i\sum q_i a_i^r\right)^{\frac1r}\quad(\text{根据 }\sum q_i=1)\\
&=\left(\sum q_i a_i^r\right)^{\frac1r}=M_r(a,q).
\end{aligned}
\]

\noindent\textbf{引理 2}\quad $G(a,q)=\displaystyle\lim_{r\to0^+}M_r(a,q)$.

\noindent\textbf{证}\quad 若 $a_i>0$（$i=1,2,\cdots,n$），则
\[
a_i^r=e^{r\ln a_i}=1+r\ln a_i+o(r^2).
\]
所以
\[
\begin{aligned}
M_r(a,q)
&=\left(\sum q_i a_i^r\right)^{\frac1r}
=e^{\frac1r\ln\sum q_i a_i^r}\\
&=\exp\left[\frac1r\ln\sum q_i(1+r\ln a_i+o(r^2))\right]\\
&=\exp\left[\frac1r\ln\left(\sum q_i+r\sum q_i\ln a_i+o(r^2)\sum q_i\right)\right]\\
&=\exp\left[\frac1r\ln\left(1+r\sum q_i\ln a_i+o(r^2)\right)\right].
\end{aligned}
\]
如此
\[
\lim_{r\to0^+}M_r(a,q)
=\exp\left(\sum q_i\ln a_i\right)=\prod a_i^{q_i}=G(a,q).
\]

利用这两个引理，立即可得平均值不等式的推广形式.

\noindent\textbf{定理 4}\quad 设 $a_1,a_2,\cdots,a_n$ 不全相等，则有 $G(a,q)<M_1(a,q)$，即
\[
a_1^{q_1}a_2^{q_2}\cdots a_n^{q_n}
<q_1a_1+\cdots+q_na_n\quad(q_i>0,\ \sum q_i=1),
\]
亦即
\[
(a_1^{p_1}a_2^{p_2}\cdots a_n^{p_n})^{\frac1{p_1+p_2+\cdots+p_n}}
<\frac{p_1a_1+p_2a_2+\cdots+p_na_n}{p_1+p_2+\cdots+p_n},
\]
只有 $a_1,a_2,\cdots,a_n$ 全相等时“$<$”才成为“$=$”.

\noindent\textbf{证}\quad 由引理 1 知
\[
M_1(a,q)>M_{\frac12}(a,q)>M_{\frac1{2^2}}(a,q)>\cdots>
M_{\frac1{2^k}}(a,q)>\lim_{r\to0^+}M_r(a,q).
\]
又由引理 2 知 $\displaystyle\lim_{r\to0^+}M_r(a,q)=G(a,q)$，故 $G(a,q)<M_1(a,q)$.

显然，若 $q_1=q_2=\cdots=q_n=\dfrac1n$，则回到定理 3. 可见定理 4 是定理 3 的推广.

\noindent\textbf{c. 平均值不等式的积分形式}

\noindent\textbf{定义 3}\quad 设函数 $f(x)$ 及 $p(x)>0$ 在 $[a,b]$ 上有定义，且下面所出现的积分有意义，记
\[
A(f)=\frac{\int_a^b p(x)f(x)\,dx}{\int_a^b p(x)\,dx},\quad
M_r(f)=\left(\frac{\int_a^b p(x)f^r(x)\,dx}{\int_a^b p(x)\,dx}\right)^{\frac1r}\quad(r>0).
\]
若 $f(x)>0$，记
\[
G(f)=\exp\left(\frac{\int_a^b p(x)\ln f(x)\,dx}{\int_a^b p(x)\,dx}\right).
\]
它们分别称为 $f(x)$ 的加权算术平均，加权（$r$ 次幂）算术平均和加权几何平均，其中 $p(x)>0$ 称为权函数.

若用
\[
q(x)=\frac{p(x)}{\int_a^b p(x)\,dx}
\]
取代 $p(x)$，则 $\displaystyle\int_a^b q(x)\,dx=1$，$q(x)$ 称为 $p(x)$ 的标准化.

\noindent\textbf{注}\quad 由上述定义，明显可看出

1）$\ln G(f)=A(\ln f)$（当 $f(x)>0$ 时）.

2）若 $\alpha,\beta$ 为常数，则 $A(\alpha f\pm\beta g)=\alpha A(f)\pm\beta A(g)$.

3）$M_r(f)=[A(f^r)]^{\frac1r}$，$A(f)=M_1(f)$.

4）$G(f)$ 是加权几何平均 $G(a,p)$：
\[
\left(\prod_{i=1}^n a_i^{p_i}\right)^{\frac1{p_1+\cdots+p_n}}
=\exp\left[\ln\left(\prod_{i=1}^n a_i^{p_i}\right)^{\frac1{\sum p_i}}\right]
=\exp\left(\frac{\sum p_i\ln a_i}{\sum p_i}\right)
\]
的积分形式.

5）若 $p(x)\equiv1$，则
\[
A(f)=\frac1{b-a}\int_a^b f(x)\,dx,\quad
M_r(f)=\left(\frac1{b-a}\int_a^b f^r(x)\,dx\right)^{\frac1r},\quad
G(f)=\exp\left(\frac1{b-a}\int_a^b\ln f(x)\,dx\right).
\]

\noindent\textbf{定理 5（平均值不等式的积分形式）}\quad 设 $r>0$，$f(x)>0$，所证的积分有意义，则 $G(f)\le A(f)$，
$G(f)\le M_r(f)$，即
\[
\exp\left(\int_a^b q(x)\ln f(x)\,dx\right)
\le\left(\int_a^b q(x)f^r(x)\,dx\right)^{\frac1r}.
\]
（$r>0$，包括 $r=1$ 的情况.）

\noindent\textbf{证}\quad $\forall r>0$，
\[
M_{\frac r2}(f)=\left(\int_a^b q(x)f^{\frac r2}(x)\,dx\right)^{\frac2r}
\]
\[
=\left(\int_a^b\sqrt{q(x)}\cdot\sqrt{q(x)}f^{\frac r2}(x)\,dx\right)^{\frac2r}\quad(\text{应用 Schwarz 不等式})
\]
\[
\le\left(\int_a^b q(x)\,dx\cdot\int_a^b q(x)f^r(x)\,dx\right)^{\frac1r}
=\left(\int_a^b q(x)f^r(x)\,dx\right)^{\frac1r}=M_r(f).
\]
但
\[
\lim_{r\to0^+}M_r(f)
=\lim_{r\to0^+}\left(\int_a^b q(x)f^r(x)\,dx\right)^{\frac1r}
=\max_{a\le x\le b}f(x)\stackrel{\text{记}}{=}\mu\quad(\text{见例 4.1.7}),
\]
且
\[
G(f)=\exp\left(\int_a^b q(x)\ln f(x)\,dx\right)
\le \exp\left(\ln\mu\int_a^b q(x)\,dx\right)=\mu,
\]
故
\[
M_r(f)\ge M_{\frac r2}(f)\ge M_{\frac r4}(f)\ge\cdots\ge
M_{\frac r{2^i}}(f)\ge\cdots\ge\lim_{r\to0^+}M_r(f)=\mu\ge G(f).
\]
$r=1$ 时，$A(f)\ge G(f)$. 证明过程里“$\ge$”中的等号当且仅当 $f(x)\equiv$ 常数时成立.

\subsection*{三、Hölder 不等式}

\noindent\textbf{a. 基本形式}

\noindent\textbf{定理 6（Hölder 不等式）}\quad 设 $a_i,b_i\ge0$（$i=1,2,\cdots,n$），$k,k'$ 为实数：
$\dfrac1k+\dfrac1{k'}=1$. 则当 $k>1$（从而 $k'>1$）时，
\[
\sum_{i=1}^n a_i b_i
\le \left(\sum_{i=1}^n a_i^k\right)^{\frac1k}
\left(\sum_{i=1}^n b_i^{k'}\right)^{\frac1{k'}}; \tag{1}
\]
当 $k<1$，$k\ne0$（从而 $k'<1$）时，
\[
\sum_{i=1}^n a_i b_i
\ge \left(\sum_{i=1}^n a_i^k\right)^{\frac1k}
\left(\sum_{i=1}^n b_i^{k'}\right)^{\frac1{k'}}; \tag{2}
\]
其中等号当且仅当 $a_i$ 与 $b_i$ 成比例（即 $\exists\alpha,\beta$ 不全为零使
$\alpha a_i^k=\beta b_i^{k'}$（$i=1,2,\cdots,n$））时成立.

\noindent\textbf{证}\quad $1^\circ$ 当 $k>1$ 时
\[
k'=\frac1{1-\frac1k}>1.
\]
下文中的 $\sum$ 的上、下标跟式 (1) 相同，省略.
\[
\frac{\sum a_i b_i}{(\sum a_i^k)^{\frac1k}(\sum b_i^{k'})^{\frac1{k'}}}
=\sum_i\left(\frac{a_i^k}{\sum a_i^k}\right)^{\frac1k}
\left(\frac{b_i^{k'}}{\sum b_i^{k'}}\right)^{\frac1{k'}}
\]
（应用定理 4）
\[
\le\sum_i\left[\frac1k\left(\frac{a_i^k}{\sum a_i^k}\right)
+\frac1{k'}\left(\frac{b_i^{k'}}{\sum b_i^{k'}}\right)\right]
=\frac1k\frac{\sum a_i^k}{\sum a_i^k}
+\frac1{k'}\frac{\sum b_i^{k'}}{\sum b_i^{k'}}=1.
\]
“$\le$”中的等号当且仅当
\[
\frac{a_i^k}{\sum a_i^k}=\frac{b_i^{k'}}{\sum b_i^{k'}}
\quad(i=1,2,\cdots,n)
\]
时成立.

$2^\circ$ 当 $k<1$ 时，注意到 $\dfrac1k+\dfrac1{k'}=1$，可得 $k'(1-k)+k=0$，故
\[
\sum a_i^k=\sum a_i^k b_i^{k+k'(1-k)}
=\sum (a_i b_i)^k(b_i^{k'})^{1-k}.
\]
利用刚证得的式 (1)，把 $(a_i b_i)^k$，$(b_i^{k'})^{1-k}$，$\dfrac1k$，$\dfrac1{1-k}$ 分别看作式 (1) 中的 $a_i,b_i,k,k'$，则
\[
\sum a_i^k\le(\sum a_i b_i)^k(\sum b_i^{k'})^{1-k},
\]
故
\[
\left(\sum a_i^k\right)^{\frac1k}
\le \sum a_i b_i\cdot\left(\sum b_i^{k'}\right)^{-\frac1{k'}},
\]
即
\[
\sum a_i b_i\ge
\left(\sum a_i^k\right)^{\frac1k}\left(\sum b_i^{k'}\right)^{\frac1{k'}}.
\]
等号当且仅当 $a_i^k$ 与 $b_i^{k'}$ 成比例时成立.

\noindent\textbf{new 例 4.4.6}\quad 设 $|a_i|$（$i=1,2,\cdots,2n$）是 $2n$ 个不等于 1 的整数，满足
$\displaystyle\prod_{i=1}^{2n}a_i=1$；又设 $\alpha>0$，$\beta=\alpha+1$. 问
$\displaystyle\sum_{i=1}^{2n}a_i^\alpha$ 与
$\displaystyle\sum_{i=1}^{2n}a_i^\beta$ 哪个较大？如何证明？（中国科学技术大学）

\noindent\textbf{提示}\quad 以 $a_1=2,a_2=\dfrac12$ 及 $\alpha=1,\beta=2$ 为例，猜想应为：
\[
\sum_{i=1}^{2n}a_i^\alpha\le\sum_{i=1}^{2n}a_i^\beta.
\]
可利用本节定理 6（Hölder 不等式），因 $\dfrac\alpha\beta+\dfrac1\beta=1$，设
\[
\frac1k=\frac\alpha\beta,\quad \frac1{k'}=\frac1\beta.
\]

\noindent\textbf{证}\quad 根据 Hölder 不等式，
\[
\begin{aligned}
\sum_{i=1}^{2n}a_i^\alpha
&=\sum_{i=1}^{2n}(a_i^\alpha\cdot1)
\le\left(\sum_{i=1}^{2n}(a_i^\alpha)^{\frac\beta\alpha}\right)^{\frac\alpha\beta}
\left(\sum_{i=1}^{2n}(1)^\beta\right)^{\frac1\beta}\\
&=\left(\sum_{i=1}^{2n}a_i^\beta\right)^{\frac\alpha\beta}(2n\cdot1)^{\frac1\beta}\\
&=\left(\sum_{i=1}^{2n}a_i^\beta\right)^{\frac\alpha\beta}
\left(2n\cdot\sqrt[2n]{\prod_{i=1}^{2n}a_i^\beta}\right)^{\frac1\beta}.
\end{aligned}
\]
利用平均值不等式（见例 1.1.7），有
\[
\sqrt[2n]{\prod_{i=1}^{2n}a_i^\beta}\le
\frac{\sum_{i=1}^{2n}a_i^\beta}{2n},
\]
故
\[
\sum_{i=1}^{2n}a_i^\alpha
\le\left(\sum_{i=1}^{2n}a_i^\beta\right)^{\frac\alpha\beta}
\left(\sum_{i=1}^{2n}a_i^\beta\right)^{\frac1\beta}
=\sum_{i=1}^{2n}a_i^\beta.
\]

\noindent\textbf{new 练习}\quad 设 $a,b,c,d$ 是 4 个不等于 1 的整数，满足 $abcd=1$，问如下两数哪个较大：
\[
S=a^{2010}+b^{2010}+c^{2010}+d^{2010},\quad
R=a^{2011}+b^{2011}+c^{2011}+d^{2011}.
\]
说明理由，（中国科学技术大学）\hfill《$S\le R$》

\noindent\textbf{b. Hölder 不等式的积分形式}

\noindent\textbf{定理 7}\quad 设 $f(x),g(x)\ge0$，并使得所讨论的积分有意义，$k,k'\ne0,1$，为共轭实数
（即 $\dfrac1k+\dfrac1{k'}=1$），则
\[
k>1\text{ 时，}\quad
\int_a^b f(x)g(x)\,dx
\le\left(\int_a^b f^k(x)\,dx\right)^{\frac1k}
\left(\int_a^b g^{k'}(x)\,dx\right)^{\frac1{k'}}; \tag{1}
\]
\[
k<1\text{ 时，}\quad
\int_a^b f(x)g(x)\,dx
\ge\left(\int_a^b f^k(x)\,dx\right)^{\frac1k}
\left(\int_a^b g^{k'}(x)\,dx\right)^{\frac1{k'}}. \tag{2}
\]
若 $f(x),g(x)$ 连续，则其中的等号当且仅当 $f^k(x)$ 与 $g^{k'}(x)$ 成比例（即 $\exists\alpha,\beta$ 不全为零，使得
$\alpha f^k(x)\equiv\beta g^{k'}(x)$，$\forall x\in[a,b]$）时成立.

证明与定理 6 完全类似，只要把“$\sum$”改为积分符号“$\displaystyle\int_a^b$”即可. 那里应用定理 4，这里应用定理 5.
不等式 (1)、(2) 亦可应用积分和的极限，从定理 6 推得.

\noindent\textbf{例 4.4.7}\quad 试证明：
\[
\int_0^\pi x a^{\sin x}\,dx\cdot\int_0^{\frac\pi2}a^{-\cos x}\,dx
\ge\frac{\pi^3}4\quad(a>0). \tag{1}
\]
（广西大学）

\noindent\textbf{证}\quad 令 $x=t+\dfrac\pi2$，
\[
\int_0^\pi x a^{\sin x}\,dx
=\pi\int_0^{\frac\pi2}a^{\cos t}\,dt.
\]
于是原式左端
\[
\int_0^\pi x a^{\sin x}\,dx\cdot\int_0^{\frac\pi2}a^{-\cos x}\,dx
=\pi\int_0^{\frac\pi2}a^{\cos x}\,dx\cdot\int_0^{\frac\pi2}a^{-\cos x}\,dx
\]
（应用定理 7）
\[
\ge \pi\left(\int_0^{\frac\pi2}a^{\frac{\cos x}{2}}a^{-\frac{\cos x}{2}}\,dx\right)^2
=\pi\cdot\frac{\pi^2}4=\frac{\pi^3}4.
\]

\noindent\textbf{new 练习}\quad 设 $f(x)$ 在 $[a,b]$ 连续，且 $\displaystyle\min_{x\in[a,b]}f(x)=1$. 试证：
\[
\lim_{n\to\infty}\left[\int_a^b\frac{dx}{(f(x))^n}\right]^{\frac1n}=1.
\]
（浙江大学）

\noindent\textbf{证}\quad 因连续性，$\exists x_0\in[a,b]$，使得 $f(x_0)=\min f(x)=1$.
暂设 $x_0$ 为内点，则 $\forall\varepsilon>0$，$\exists\delta>0$，当
$x\in U=U(x_0,\delta)=(x_0-\delta,x_0+\delta)$ 时，有 $1\le f(x)\le1+\varepsilon$，从而
\[
(b-a)^{\frac1n}\ge
\left[\int_a^b\frac{dx}{(f(x))^n}\right]^{\frac1n}
\ge \left[\int_U\frac{dx}{(f(x))^n}\right]^{\frac1n}
\ge \frac1{1+\varepsilon}(2\delta)^{\frac1n}.
\]
令 $n\to\infty$，得
\[
1\ge\lim_{n\to\infty}\left[\int_a^b\frac{dx}{(f(x))^n}\right]^{\frac1n}
\ge\frac1{1+\varepsilon}.
\]
由 $\varepsilon>0$ 的任意性，欲证的等式成立.（若 $x_0=a$ 或 $x_0=b$，上述推理保持有效，只需将双侧邻域改为单侧邻域，$(2\delta)^{1/n}$ 改为 $\delta^{1/n}$ 即可.）

\subsection*{四、Minkowski 不等式}

\noindent\textbf{a. 基本形式}

\noindent\textbf{定理 8（Minkowski 不等式）}\quad 对于任意实数 $r\ne0,1$ 及 $a_i,b_i\ge0$（$i=1,2,\cdots,n$），有
\[
r>1\text{ 时，}\quad
\left(\sum_{i=1}^n(a_i+b_i)^r\right)^{\frac1r}
\le\left(\sum_{i=1}^n a_i^r\right)^{\frac1r}
+\left(\sum_{i=1}^n b_i^r\right)^{\frac1r}; \tag{1}
\]
\[
r<1\text{ 时，}\quad
\left(\sum_{i=1}^n(a_i+b_i)^r\right)^{\frac1r}
\ge\left(\sum_{i=1}^n a_i^r\right)^{\frac1r}
+\left(\sum_{i=1}^n b_i^r\right)^{\frac1r}. \tag{2}
\]
其中等号当且仅当 $a_i$ 与 $b_i$ 成比例（即 $\exists\alpha,\beta$ 不全为零，使得 $\alpha a_i=\beta b_i$）时成立.

式 (1) 又称为距离不等式. $r=2,n=3$ 时，式 (1) 表示 $\mathbf{R}^3$ 中三角形任一边小于另外两边之和. 因此式 (1) 又称三角不等式.

\noindent\textbf{证}\quad $r>1$ 时，记 $s_i=a_i+b_i$，则
\[
\sum_{i=1}^n s_i^r
=\sum_{i=1}^n(a_i+b_i)^r
=\sum_{i=1}^n(a_i+b_i)(a_i+b_i)^{r-1}
=\sum a_i s_i^{r-1}+\sum b_i s_i^{r-1}.
\]
令 $k=r,k'=\dfrac r{r-1}$，则 $\dfrac1k+\dfrac1{k'}=1$. 对上式右端应用 Hölder 不等式，
\[
\sum s_i^r
\le \left(\sum a_i^r\right)^{\frac1r}\left(\sum s_i^r\right)^{\frac{r-1}{r}}
+\left(\sum b_i^r\right)^{\frac1r}\left(\sum s_i^r\right)^{\frac{r-1}{r}}.
\]
两边同乘 $\left(\sum s_i^r\right)^{\frac1r-1}$ 得
\[
\left(\sum s_i^r\right)^{\frac1r}
\le\left(\sum a_i^r\right)^{\frac1r}
+\left(\sum b_i^r\right)^{\frac1r},
\]
其中 $s_i=a_i+b_i$. 式 (1) 得证. 由定理 6 知，等号当且仅当 $a_i$ 与 $b_i$ 成比例时成立.
$r<1$ 的情况类似可证.

\noindent\textbf{b. Minkowski 不等式的积分形式}

\noindent\textbf{定理 9}\quad 设 $f(x),g(x)\ge0$，在 $[a,b]$ 上有定义，使下面积分有意义，则当 $r>1$ 时，
\[
\left(\int_a^b(f(x)+g(x))^r\,dx\right)^{\frac1r}
\le \left(\int_a^b f^r(x)\,dx\right)^{\frac1r}
+\left(\int_a^b g^r(x)\,dx\right)^{\frac1r}; \tag{1}
\]
当 $0<r<1$ 时，
\[
\left(\int_a^b(f(x)+g(x))^r\,dx\right)^{\frac1r}
\ge \left(\int_a^b f^r(x)\,dx\right)^{\frac1r}
+\left(\int_a^b g^r(x)\,dx\right)^{\frac1r}. \tag{2}
\]
\noindent\textbf{证}\quad 可仿照有限形式（定理 8）的证法，从 Hölder 不等式的积分形式推出. 不等式 (1)、(2) 亦可用积分和的极限，从定理 8 的不等式 (1)、(2) 推得.

\noindent\textbf{c. $n$ 元 Minkowski 不等式}

上面的不等式立即可写出它们的一般形式：

\noindent\textbf{定理 10}\quad 对于任意实数 $r\ne0$ 及 $a_{ij}\ge0$（$i=1,2,\cdots,n$，$j=1,2,\cdots,m$），有当 $r>1$ 时，
\[
\left[\sum_{i=1}^n(a_{i1}+a_{i2}+\cdots+a_{im})^r\right]^{\frac1r}
\le \left(\sum_{i=1}^n a_{i1}^r\right)^{\frac1r}
+\left(\sum_{i=1}^n a_{i2}^r\right)^{\frac1r}
+\cdots+\left(\sum_{i=1}^n a_{im}^r\right)^{\frac1r};
\]
当 $0<r<1$ 时，
\[
\left[\sum_{i=1}^n(a_{i1}+a_{i2}+\cdots+a_{im})^r\right]^{\frac1r}
\ge \left(\sum_{i=1}^n a_{i1}^r\right)^{\frac1r}
+\left(\sum_{i=1}^n a_{i2}^r\right)^{\frac1r}
+\cdots+\left(\sum_{i=1}^n a_{im}^r\right)^{\frac1r}.
\]
等号当且仅当 $\forall j,k$，$(a_{ik})_{i=1}^n$ 与 $(a_{ij})_{i=1}^n$ 成比例时成立.

\noindent\textbf{定理 11}\quad 设 $f_i(x)$（$i=1,2,\cdots,m$）在 $[a,b]$ 上有定义，下界为正，在 $[a,b]$ 上可积，则当 $r>1$ 时，
\[
\left[\int_a^b(f_1(x)+\cdots+f_m(x))^r\,dx\right]^{\frac1r}
\le\left(\int_a^b f_1^r(x)\,dx\right)^{\frac1r}
+\cdots+\left(\int_a^b f_m^r(x)\,dx\right)^{\frac1r}; \tag{1}
\]
当 $0<r<1$ 时，
\[
\left[\int_a^b(f_1(x)+\cdots+f_m(x))^r\,dx\right]^{\frac1r}
\ge\left(\int_a^b f_1^r(x)\,dx\right)^{\frac1r}
+\cdots+\left(\int_a^b f_m^r(x)\,dx\right)^{\frac1r}. \tag{2}
\]

\subsection*{五、Young 不等式}

著名的不等式还有很多，我们不准备一一介绍. 最后，只介绍一个在证法上有特点的 Young 不等式.

\noindent\textbf{定理 12}\quad 设 $f(x)$ 严格递增，于 $[0,+\infty)$ 上连续，$f(0)=0$，$a,b>0$，$f^{-1}(x)$ 表示 $f(x)$ 的反函数，则
\[
ab\le\int_0^a f(x)\,dx+\int_0^b f^{-1}(y)\,dy, \tag{1}
\]
其中等号当且仅当 $f(a)=b$ 时成立.

该式从几何上看，是十分清楚的：因积分等于曲边梯形的面积，可能发生的三种情况，如图 4.4.1 所示，这时
\[
\int_0^a f(x)\,dx=S_{OABO},\quad
\int_0^b f^{-1}(y)\,dy=S_{OCEO},\quad
ab=S_{OADEO}.
\]

\begin{center}
\includegraphics[width=0.78\textwidth]{4-4-1.png}
\end{center}

（其中 $S_{OABO}$ 表示图形 $OABO$ 的面积，如此，等等.）
从图形看，三种情况都有
\[
S_{OABO}+S_{OCEO}\ge S_{OADEO},
\]
并且等号只在第三种情况（$b=f(a)$）发生. 故从几何上看命题十分明显. 问题是分析上如何证明.

\noindent\textbf{证}\quad $1^\circ$ 我们先证明：
\[
\int_0^a f(x)\,dx+\int_0^{f(a)}f^{-1}(y)\,dy=af(a). \tag{2}
\]

因 $f(x)\nearrow$，于 $[0,a]$ 上连续，故 $f^{-1}(y)\nearrow$，于 $[0,f(a)]$ 上连续，因此式 (2) 中积分有意义. 将 $[0,a]$ $n$ 等分，记分点为
\[
0=x_0<x_1<\cdots<x_n=a,
\]
相应的点 $y_i=f(x_i)$（$i=0,1,\cdots,n$）构成区间 $[0,f(a)]$ 的一个分划：
\[
0=y_0<y_1<\cdots<y_n=f(a). \tag{3}
\]
因为 $f(x)$ 在 $[0,a]$ 上连续，故在 $[0,a]$ 上一致连续. 因此当 $n\to\infty$ 时，对于分划 (3)，有
\[
\max \Delta y_i=\max (y_i-y_{i-1})
=\max[f(x_i)-f(x_{i-1})]\to0.
\]
从而
\[
\begin{aligned}
&\int_0^a f(x)\,dx+\int_0^{f(a)}f^{-1}(y)\,dy\\
&=\lim_{n\to\infty}\left\{\sum_{i=1}^n f(x_i)\Delta x_i
+\sum_{i=1}^n f^{-1}(y_{i-1})\Delta y_i\right\}\\
&=\lim_{n\to\infty}\sum_{i=1}^n
\{f(x_i)(x_i-x_{i-1})+f^{-1}(f(x_{i-1}))(f(x_i)-f(x_{i-1}))\}\\
&=\lim_{n\to\infty}\sum_{i=1}^n
\{f(x_i)(x_i-x_{i-1})+x_{i-1}(f(x_i)-f(x_{i-1}))\}\\
&=\lim_{n\to\infty}\sum_{i=1}^n[x_i f(x_i)-x_{i-1}f(x_{i-1})]\\
&=\lim_{n\to\infty}[x_n f(x_n)-x_0 f(x_0)]
=\lim_{n\to\infty}[af(a)-0\cdot f(0)]=af(a).
\end{aligned}
\]
式 (2) 获证.

$2^\circ$ 由式 (2) 可知，若 $f(a)=b$，则式 (1) 中等号成立.

$3^\circ$ 若 $0<b<f(a)$，则由 $f(x)$ 的连续性可知，$\exists x_0\in(0,a)$ 使得 $f(x_0)=b$. 故
\[
\begin{aligned}
&\int_0^a f(x)\,dx+\int_0^b f^{-1}(y)\,dy\\
&=\int_0^{x_0}f(x)\,dx+\int_{x_0}^a f(x)\,dx+\int_0^b f^{-1}(y)\,dy\\
&>\int_0^{x_0}f(x)\,dx+\int_0^{f(x_0)}f^{-1}(y)\,dy+f(x_0)(a-x_0)\\
&>x_0f(x_0)+f(x_0)(a-x_0)\qquad\text{（应用式 (2)）}\\
&=af(x_0)=ab.
\end{aligned}
\]

$4^\circ$ $b>f(a)$ 时，只要把 $f(x)$ 看作 $f^{-1}(y)$ 的反函数，就可由 $3^\circ$ 的结论得到.

$5^\circ$ 联系 $2^\circ$，$3^\circ$，$4^\circ$ 时可知式 (1) 成立，当且仅当 $f(a)=b$ 时式 (1) 中的等号成立.

\noindent\textbf{例 4.4.8}\quad 证明：当 $a,b\ge1$ 时，不等式
\[
ab\le e^{a-1}+b\ln b
\]
成立.（安徽大学）

\noindent\textbf{提示}\quad $f(x)=e^x-1\nearrow$，连续，$f^{-1}(y)=\ln(1+y)$，因为 $a,b>1$，应用 Young 不等式，得
\[
(a-1)(b-1)\le\int_0^{a-1}f(x)\,dx+\int_0^{b-1}f^{-1}(y)\,dy,
\]
即得欲证不等式.

\noindent\textbf{例 4.4.9}\quad 设 $a,b>0$，$p>1$，$\frac1p+\frac1q=1$，试证：
\[
ab\le\frac{a^p}{p}+\frac{b^q}{q}.
\]

\noindent\textbf{证 I}\quad 因为 $p>1$，故 $f(x)=x^{p-1}\nearrow$，连续. 当 $x\ge0$ 时，
\[
f^{-1}(y)=y^{\frac1{p-1}}=y^{q-1}\quad\left(\frac1{p-1}=q-1\right).
\]
应用 Young 不等式，有
\[
ab\le\int_0^a f(x)\,dx+\int_0^b f^{-1}(y)\,dy
=\frac{a^p}{p}+\frac{b^q}{q}.
\]

\noindent\textbf{证 II}\quad 令
\[
f(x)=\frac{x^p}{p}+\frac{b^q}{q}-bx\quad(x>0),
\]
则
\[
f'(x)=x^{p-1}-b
\begin{cases}
>0,& x>b^{\frac1{p-1}},\\
<0,& x<b^{\frac1{p-1}}.
\end{cases}
\]
所以 $f(x)$ 在 $x=b^{\frac1{p-1}}$ 处最小，但 $f(b^{\frac1{p-1}})=0$，故 $f(a)\ge0$.

\noindent\textbf{注}\quad 例 4.4.9 的不等式有时也被称为 Young 不等式，常用于缩放，如例 5.1.22 的练习 2.

\subsection*{单元练习 4.4}

\noindent\textbf{4.4.1}\quad 证明：
\[
0.83<\int_0^1\frac{dx}{\sqrt{1+x^4}}<0.95.
\]

\noindent\textbf{提示}\quad 将区间 $[0,1]$ 二等分，作为分划，则 0.83 和 0.95 分别是积分的 Darboux 下和与 Darboux 上和.

\noindent\textbf{4.4.2}\quad 设 $f(x)$ 在 $[a,b]$ 上有连续导数，$f(a)=0$. 试证：
\[
M^2\le(b-a)\int_a^b f'^2(x)\,dx,
\]
其中
\[
M=\max_{a\le x\le b}|f(x)|.
\]

\noindent\textbf{提示}\quad 利用 Schwarz 不等式：
\[
f^2(x)=\left(\int_a^x |f'(t)|\,dt\right)^2
\le(x-a)\int_a^x(f'(t))^2\,dt,
\]
然后两端同时取最大值.

\noindent\textbf{4.4.3}\quad 设 $f(x)$ 在 $[0,1]$ 上有连续导数，且 $f(1)-f(0)=1$. 证明：
\[
\int_0^1 f'^2(x)\,dx\ge1.
\]
（国外赛题）

\noindent\textbf{提示}\quad 令 $F(x)=f(x)-f(0)$，则 $F'(x)=f'(x)$，$F(x)$ 满足上题的条件，可利用上题的结果（或用类似方法推证）.

\noindent\textbf{再提示}\quad
\[
\int_0^1(f'(x))^2\,dx=(1-0)\int_0^1(F'(x))^2\,dx
\ge\max_{0\le x\le1}F^2(x)\ge F^2(1)=1.
\]

\noindent\textbf{4.4.4}\quad 设 $f(x)$ 在 $[a,b]$ 上有连续导数（$0<a<b$），$f(a)=f(b)=0$，$\int_a^b f^2(x)\,dx=1$. 试证：
\[
\int_a^b x^2f'^2(x)\,dx>\frac14.
\]

\noindent\textbf{提示}\quad
\[
1=\left(\int_a^b f^2(x)\,dx\right)^2
=4\left(\int_a^b f(x)\cdot f'(x)x\,dx\right)^2,
\]
再用 Schwarz 不等式.

\noindent\textbf{4.4.5}\quad 证明：
\[
\ln\frac Pq\le\frac{P-q}{\sqrt{Pq}}\quad(0<q\le P).
\]

\noindent\textbf{提示}\quad 方法 I：用积分（Schwarz 不等式）；方法 II：用导数.

\noindent\textbf{再提示}\quad 用方法 I 时：若 $P=q$ 是平凡的，原式变为 $0=0$.

当 $0<q<P$ 时，
\[
\begin{aligned}
\ln\frac Pq&=\ln P-\ln q=\int_q^P\frac1x\cdot1\,dx\\
&\le\left(\int_q^P\left(\frac1x\right)^2\,dx\cdot\int_q^P1^2\,dx\right)^{\frac12}\\
&=\left[\left(\frac1q-\frac1P\right)(P-q)\right]^{\frac12}
=\frac{P-q}{\sqrt{Pq}}.
\end{aligned}
\]
用方法 II 时：原式可写为
\[
\ln\frac Pq\le\sqrt{\frac Pq}-\sqrt{\frac qP}.
\]
设 $x=\sqrt{\frac Pq}\ge1$（$0<q\le P$），则问题化为：当 $x\ge1$ 时，求证：
\[
2\ln x\le x-\frac1x.
\]
令 $F(x)=2\ln x-x+\frac1x$，因 $F(1)=0$，
\[
F'(x)=-\frac{(x-1)^2}{x^2}<0,
\]
故 $F(x)<0$（$x>1$），即 $2\ln x\le x-\frac1x$，证毕.

\noindent\textbf{4.4.6}\quad 设函数 $f(x)$ 在 $[a,b]$ 上有连续导数，$f(a)=f(b)=0$. 试证：
\[
\int_a^b |f(x)f'(x)|\,dx\le\frac{b-a}{4}\int_a^b f'^2(x)\,dx,
\]
并且 $\frac{b-a}{4}$ 不能再小.

\noindent\textbf{提示}\quad 令
\[
F(x)=\int_a^x |f'(t)|\,dt.
\]

\noindent\textbf{再提示}\quad $1^\circ$ 此时，$F'(x)=|f'(x)|$，
\[
F(x)\ge\left|\int_a^x f'(t)\,dt\right|=|f(x)|.
\]
将区间 $[a,b]$ 的中点记为 $c=\frac{a+b}{2}$，那么
\[
\int_a^b |f(x)f'(x)|\,dx
=\int_a^c |f(x)f'(x)|\,dx+\int_c^b |f(x)f'(x)|\,dx=I_1+I_2.
\]
而
\begin{equation}
\begin{aligned}
I_1&=\int_a^c |f(x)f'(x)|\,dx
\le\int_a^c F(x)F'(x)\,dx\\
&=\frac12F^2(c)-\frac12F^2(a)
=\frac12F^2(c)
=\frac12\left(\int_a^c |f'(x)|\cdot1\,dx\right)^2\\
&\le\frac{b-a}{4}\int_a^c f'^2(x)\,dx.
\end{aligned}
\tag{1}
\end{equation}
同理，因 $f(b)=0$，有
\[
I_2=\int_c^b |f(x)f'(x)|\,dx
\le\frac{b-a}{4}\int_c^b f'^2(x)\,dx. \tag{2}
\]
因此
\[
\int_a^b |f(x)f'(x)|\,dx=I_1+I_2
\le\frac{b-a}{4}\int_a^b f'^2(x)\,dx.
\]

$2^\circ$ （为什么 $\frac{b-a}{4}$ 不能再小？）作折线函数
\[
f_1(x)=
\begin{cases}
x-a,& x\in(a,c),\\
b-x,& x\in(c,b),
\end{cases}
\]
在 $[a,b]$ 上连续，满足 $f_1(a)=f_1(b)$，且使原不等式左、右两端都等于 $\frac{(b-a)^2}{4}$（变成了等式）. $f_1(x)$ 仅在 $x=c$ 处无导数，若在 $c$ 点的任意小邻域里，换成光滑函数就可保证可微性，邻域充分小，就可使积分值改变任意小，与 $\frac{b-a}{4}$ 之差能任意小. 可见 $\frac{b-a}{4}$ 不能再小.

\noindent\textbf{4.4.7}\quad 若 $u_1,u_2,\cdots,u_n\ge0$，$u_1\cdot u_2\cdots u_n=1$，则有
\[
u_1+u_2+\cdots+u_n\ge n.
\]
试证明这一结论，并由它导出定理 3（平均值不等式）.

\noindent\textbf{证}\quad $1^\circ$ 本节定理 4 指出：若 $p_i\ge1$（$i=1,2,\cdots,n$），$a_i\ge0$（$i=1,2,\cdots,n$），则有
\[
\left(a_1^{p_1}\cdot a_2^{p_2}\cdots a_n^{p_n}\right)^{\frac1{p_1+p_2+\cdots+p_n}}
\le\frac{p_1a_1+p_2a_2+\cdots+p_na_n}{p_1+p_2+\cdots+p_n}.
\]
令 $a_i=u_i>0$，并取 $p_i=1$（$i=1,2,\cdots,n$），应用上式得
\[
1=\sqrt[n]{u_1\cdot u_2\cdots u_n}
\le\frac1n(u_1+u_2+\cdots+u_n),
\]
故 $u_1+u_2+\cdots+u_n\ge n$. 本命题成立.

$2^\circ$ 反之，若本命题成立，即：“若正数之积 $u_1\cdot u_2\cdots u_n=1$，则 $u_1+u_2+\cdots+u_n\ge n$.” 因此，对任意 $a_i>0$（$i=1,2,\cdots,n$），令
\[
u_i=\frac{a_i}{\sqrt[n]{a_1a_2\cdots a_n}}\quad(i=1,2,\cdots,n),
\]
则 $u_i$ 满足条件：$u_1\cdot u_2\cdots u_n=1$. 应用本命题，得
\[
u_1+u_2+\cdots+u_n\ge n,
\]
亦即
\[
\sum_{i=1}^n\frac{a_i}{\sqrt[n]{a_1a_2\cdots a_n}}\ge n.
\]
故
\[
\frac1n\sum_{i=1}^n a_i\ge\sqrt[n]{a_1a_2\cdots a_n}.
\]
（即平均值不等式获证.）

\noindent\textbf{☆4.4.8}\quad 设 $x_1,x_2,\cdots,x_n$ 是正数，且 $n\ge1$. 证明：
\[
\sqrt[n]{x_1x_2\cdots x_n}
\ge
\frac1{\frac1n\left(\frac1{x_1}+\frac1{x_2}+\cdots+\frac1{x_n}\right)}.
\]
（中山大学）

\noindent\textbf{提示}\quad 参看例 3.4.8.

\noindent\textbf{4.4.9}\quad 设 $f(x)\nearrow$，连续（当 $x\ge0$ 时），$f(0)=0$，$a,b\ge0$，试证：
\[
ab\le af(a)+bf^{-1}(b).
\]

\noindent\textbf{证}\quad 因 $f(x)\nearrow$，记
\[
f(a)=\max_{0\le x\le a}f(x),\qquad
f^{-1}(b)=\max_{0\le x\le b}f^{-1}(x).
\]
利用 Young 不等式，得
\[
ab\le\int_0^a f(x)\,dx+\int_0^b f^{-1}(x)\,dx
\le af(a)+bf^{-1}(b).
\]

\noindent\textbf{4.4.10}\quad （Чебышев 定理）若 $\forall i,j$，有 $(a_i-a_j)(b_i-b_j)\ge0$，则 $a_i,b_i$ 称为拟序的；若恒有相反的不等式，则称之为反序的. 试证：$a_i,b_i$ 拟序时，
\[
\left(\sum_{i=1}^n a_i\right)\left(\sum_{i=1}^n b_i\right)
\le n\sum_{i=1}^n a_i b_i;
\]
$a_i,b_i$ 反序时此不等式反向，等号当且仅当 $a_1=a_2=\cdots=a_n$，或 $b_1=b_2=\cdots=b_n$ 时成立.

\noindent\textbf{注}\quad 拟序意即：当 $i\ne j$ 时 $(a_i-a_j)$ 与 $(b_i-b_j)$ 同号（反序则指异号）.

\noindent\textbf{证 I}\quad 证（原式）右端 $-$ 左端
\[
=\frac12\sum_{i=1}^n\sum_{j=1}^n(a_i-a_j)(b_i-b_j)
\begin{cases}
\ge0&\text{（拟序时）},\\
\le0&\text{（反序时）}.
\end{cases}
\]
事实上，
\[
\text{（原式）右端}=n\sum_{i=1}^n a_i b_i
=\sum_{i=1}^n na_i b_i
=\sum_{i=1}^n\sum_{j=1}^n a_i b_i
\quad\left(\text{或 }=\sum_{i=1}^n\sum_{j=1}^n a_jb_j\right),
\]
\[
\text{（原式）左端}=\sum_{i=1}^n a_i\sum_{j=1}^n b_j
=\sum_{i=1}^n\sum_{j=1}^n a_i b_j.
\]
因此
\[
\text{右端}-\text{左端}
=\sum_{i=1}^n\sum_{j=1}^n a_i b_i
-\sum_{i=1}^n\sum_{j=1}^n a_i b_j
=\sum_{i=1}^n\sum_{j=1}^n(a_i b_i-a_i b_j)
\]
或
\[
\text{右端}-\text{左端}
=\sum_{i=1}^n\sum_{j=1}^n a_j b_j
-\sum_{i=1}^n\sum_{j=1}^n a_i b_j
=\sum_{i=1}^n\sum_{j=1}^n(a_j b_j-a_j b_i).
\]
两式相加除以 2，得
\[
\text{右端}-\text{左端}
=\frac12\sum_{i=1}^n\sum_{j=1}^n(a_i b_i-a_i b_j+a_j b_j-a_j b_i)
=\frac12\sum_{i=1}^n\sum_{j=1}^n(a_i-a_j)(b_i-b_j)
\begin{cases}
\ge0&\text{（拟序时）},\\
\le0&\text{（反序时）}.
\end{cases}
\]

\noindent\textbf{证 II}\quad （用数学归纳法，只证“拟序”情况，“反序”类似.）

$1^\circ$ $n=1$ 时平凡；当 $n=2$ 时也明显，因为
\[
\text{右端}-\text{左端}
=2(a_1b_1+a_2b_2)-(a_1+a_2)(b_1+b_2)
=(a_2-a_1)(b_2-b_1)\ge0.
\]

$2^\circ$ 设 $n=k$ 时命题成立：
\[
k\sum_{j=1}^k a_jb_j-\sum_{j=1}^k a_j\sum_{j=1}^k b_j\ge0. \tag{1}
\]
要证 $n=k+1$ 也成立：
\[
(k+1)\sum_{j=1}^{k+1}a_jb_j-\sum_{j=1}^{k+1}a_j\sum_{j=1}^{k+1}b_j\ge0. \tag{2}
\]
实际上，
\[
\begin{aligned}
\text{式 (2) 左端}
&=(k+1)\left(\sum_{j=1}^k a_jb_j+a_{k+1}b_{k+1}\right)
-\left(\sum_{j=1}^k a_j+a_{k+1}\right)\left(\sum_{j=1}^k b_j+b_{k+1}\right)\\
&=\left(k\sum_{j=1}^k a_jb_j+ka_{k+1}b_{k+1}
+\sum_{j=1}^k a_jb_j+a_{k+1}b_{k+1}\right)\\
&\quad-\left(\sum_{j=1}^k a_j\sum_{j=1}^k b_j
+b_{k+1}\sum_{j=1}^k a_j+a_{k+1}\sum_{j=1}^k b_j+a_{k+1}b_{k+1}\right).
\end{aligned}
\]

\[
\begin{aligned}
&\ge\left(\sum_{j=1}^k a_{k+1}b_{k+1}
+\sum_{j=1}^k a_jb_j\right)
-\left(\sum_{j=1}^k a_jb_{k+1}
+\sum_{j=1}^k a_{k+1}b_j\right)\\
&=\sum_{j=1}^k(a_{k+1}-a_j)(b_{k+1}-b_j)\ge0.
\end{aligned}
\]
（因为拟序，每项非负和也非负.）

反序的情况只要将不等号反向，推理仍有效.

$3^\circ$ （证明等号成立的充要条件是：$a_1=a_2=\cdots=a_n$ 或 $b_1=b_2=\cdots=b_n$.）

已证：拟序时，原不等式可转化为
\[
\sum_{i=1}^n\sum_{j=1}^n(a_i-a_j)(b_i-b_j)\ge0.
\]
可见上述条件直接使不等式成为等式。充分性明显，只需证明必要性.

若不等式成为等式（即：每项非负，其和也为零），则每项必为零：
\[
(a_i-a_j)(b_i-b_j)=0\quad(\forall i,j).
\]
如果 $a_i-a_j,b_i-b_j$ 统统为零，问题已明. 若有一个 $a_i-a_j,b_i-b_j$ 不为零，例如：$a_i-a_j\ne0$，则可推出 $b_1=b_2=\cdots=b_n$. 事实上，由 $(a_i-a_j)(b_i-b_j)=0,\ a_i-a_j\ne0$，知 $b_i=b_j$. 对其他每个下标 $k\ne i,j$ 也应有 $b_k=b_i=b_j$，否则由 $b_k-b_j\ne0$ 和 $(a_k-a_j)(b_k-b_j)=0$ 可得 $a_k=a_j$；由 $b_k-b_i\ne0$ 和 $(a_k-a_i)(b_k-b_i)=0$ 可得 $a_k=a_i$，如此 $a_i=a_j$，与 $a_i-a_j\ne0$ 矛盾.

同理，若有某个 $b_i-b_j\ne0$，则得 $a_1=a_2=\cdots=a_n$.（必要性获证）（反序情况类似可证）.

\noindent\textbf{☆4.4.11}\quad （上题的积分形式）假设函数 $f(x)$ 和 $g(x)$ 都在 $[a,b]$ 上可积，且 $f$ 和 $g$ 为拟序的（即：同为递增（或同为递减）），那么有如下不等式成立：
\[
\int_a^b f(x)\,dx\cdot\int_a^b g(x)\,dx
\le(b-a)\int_a^b f(x)g(x)\,dx.
\]
若 $f$ 和 $g$ 为反序的（即其一递增，另一递减），则该不等式反向.

\noindent\textbf{提示}\quad 例如：作积分和时，将 $[a,b]$ $n$ 等分，每个小区间的 $\xi_i$ 都统一取右端点（或左端点），不等式两端都如此操作. 写出积分和后，（根据上题）有不等式成立，然后令 $n\to\infty$，取极限即得欲求的积分不等式.（也可先证明积分形式，再利用阶梯函数得出离散形式.）

\noindent\textbf{注}\quad 不等式
\[
\sum_{i=1}^n a_i\sum_{i=1}^n b_i\le n\sum_{i=1}^n a_i b_i,\quad
\int_a^b f(x)\,dx\cdot\int_a^b g(x)\,dx
\le(b-a)\int_a^b f(x)g(x)\,dx
\]
可分别改写为均值形式：
\[
\frac1n\sum_{i=1}^n a_i\cdot\frac1n\sum_{i=1}^n b_i
\le\frac1n\sum_{i=1}^n a_i b_i,
\]
和
\[
\frac1{b-a}\int_a^b f(x)\,dx\cdot
\frac1{b-a}\int_a^b g(x)\,dx
\le\frac1{b-a}\int_a^b f(x)g(x)\,dx.
\]
（可称为：“均值的乘积”小于“乘积的均值”.）

\noindent\textbf{☆4.4.12}\quad （Чебышев 定理（级数形式的推广））设 $q_i>0\ (i=1,2,\cdots,n)$ 满足 $\sum_{i=1}^n q_i=1$（$q_i$ 称为权数），则（当 $a_i,b_i$ 为拟序时）有不等式：
\[
\sum_{i=1}^n q_i a_i\cdot\sum_{i=1}^n q_i b_i
\le\sum_{i=1}^n q_i a_i b_i, \tag{1}
\]
或等价地：对 $p_i>0\ (i=1,2,\cdots,n)$，有
\[
\sum_{i=1}^n p_i a_i\cdot\sum_{i=1}^n p_i b_i
\le\sum_{i=1}^n p_i\sum_{i=1}^n p_i a_i b_i. \tag{2}
\]

\noindent\textbf{注}\quad 在式 (2) 里，若 $\sum_{i=1}^n p_i=1$，就回到式 (1)（因此式 (1) 可看成式 (2) 的特例）. 式 (2) 两端同除以 $\sum_{i=1}^n p_i$，并令
\[
q_i=\frac{p_i}{\sum_{i=1}^n p_i},
\]
则式 (2) 可从式 (1) 获证.

另外，当 $a_i$ 与 $b_i$ 反序时，以上各式中的“$\le$”应改为“$\ge$”.

\noindent\textbf{☆4.4.13}\quad （Чебышев 定理（积分形式的推广））设 $p(x)>0$，且 $p,f,g$ 都在 $[a,b]$ 可积，$f,g$ 同增（或同减），则有如下积分不等式成立：
\[
\int_a^b p(x)f(x)\,dx\cdot\int_a^b p(x)g(x)\,dx
\le\int_a^b p(x)\,dx\int_a^b p(x)f(x)g(x)\,dx. \tag{1}
\]

\noindent\textbf{证}\quad （参考 4.4.10 题的证 I，用类似方法证明式 (1).）
\[
\text{（式 (1)）右端}
=\int_a^b p(x)\,dx\int_a^b p(y)f(y)g(y)\,dy, \tag{2}
\]
或
\[
\text{（式 (1)）右端}
=\int_a^b p(y)\,dy\int_a^b p(x)f(x)g(x)\,dx, \tag{$2'$}
\]
\[
\text{（式 (1)）左端}
=\int_a^b p(x)f(x)\,dx\cdot\int_a^b p(y)g(y)\,dy, \tag{3}
\]
或
\[
\text{（式 (1)）左端}
=\int_a^b p(y)f(y)\,dy\cdot\int_a^b p(x)g(x)\,dx. \tag{$3'$}
\]
$(2)-(3)$ 得
\[
\begin{aligned}
\text{右端}-\text{左端}
&=\int_a^b p(x)\,dx\int_a^b p(y)f(y)g(y)\,dy
-\int_a^b p(x)f(x)\,dx\cdot\int_a^b p(y)g(y)\,dy\\
&=\int_a^b\int_a^b
[p(x)p(y)f(y)g(y)-p(x)f(x)p(y)g(y)]\,dxdy\\
&=\int_a^b\int_a^b[p(x)p(y)g(y)(f(y)-f(x))]\,dxdy. \qquad(4)
\end{aligned}
\]
$(2')-(3')$ 得
\[
\text{右端}-\text{左端}
=\int_a^b\int_a^b[p(x)p(y)g(x)(f(x)-f(y))]\,dxdy. \tag{5}
\]
$[(4)+(5)]\div2$ 得
\[
\text{右端}-\text{左端}
=\frac12\int_a^b\int_a^b
[p(x)p(y)(f(x)-f(y))(g(x)-g(y))]\,dxdy\ge0.
\]
式 (1) 证毕.

\section*{§4.5\quad 反常积分}

\noindent\textbf{导读}\quad （一元）反常积分是各类考试的热点之一，非数学院系学生可只侧重于计算.

本节内容包括：反常积分的计算，收敛性的判定，反常积分的极限，无穷限反常积分敛散性与无穷远处的状态，反常积分作为“积分和”的极限.

\subsection*{一、反常积分的计算}

\noindent\textbf{a. 三大基本方法}

Newton-Leibniz 公式，变量替换法，分部积分法，是计算反常积分的三大基本方法.

\noindent\textbf{要点}\quad 设 $\displaystyle\int_a^b f(x)\,dx$ 是反常积分，$b$ 为唯一的奇点（$b$ 为有限数或 $+\infty$），计算
\[
\int_a^b f(x)\,dx.
\]

1） （用 Newton-Leibniz 公式）若 $f(x)$ 在 $[a,b)$ 内连续，且 $F(x)$ 为 $f(x)$ 的原函数，则
\[
\int_a^b f(x)\,dx
=F(x)\Big|_a^{b-0}=F(b-0)-F(a).
\]

2） （变量替换法）若 $\varphi(t)$ 在 $[\alpha,\beta)$ 内单调，有连续的导数 $\varphi'(t)$，$\varphi(\alpha)=a$，$\varphi(\beta-0)=b$（$\beta$ 为有限数或无穷大），则
\[
\int_a^b f(x)\,dx
=\int_\alpha^\beta f(\varphi(t))\varphi'(t)\,dt.
\]

3） （分部积分法）设 $u=u(x),v=v(x)$ 在 $[a,b)$ 内有连续的导数，则
\[
\int_a^b u(x)v'(x)\,dx
=\int_a^b u\,dv
=u(x)v(x)\Big|_a^{b-0}-\int_a^b v(x)u'(x)\,dx.
\]

一般来说，变量替换与分部积分只把一个积分换为另一个积分，最后还是靠 Newton-Leibniz 公式算出积分值。但不善于用变量替换与分部积分，常常无法应用 Newton-Leibniz 公式.

\noindent\textbf{例 4.5.1}\quad 计算反常积分
\[
I=\int_{-\infty}^{+\infty}|t-x|^{\frac12}
\frac{y}{(t-x)^2+y^2}\,dt.
\]

\noindent\textbf{解}\quad （这里 $x,y$ 为参变量，$t$ 为积分变量）令 $t-x=u$，则
\[
I=\int_{-\infty}^{+\infty}|u|^{\frac12}\frac{y\,du}{u^2+y^2}
=2\int_0^{+\infty}\frac{u^{\frac12}y\,du}{u^2+y^2}.
\]
再令 $\dfrac{u^{\frac12}}{\sqrt y}=v$，则
\[
I=4\sqrt y\int_0^{+\infty}\frac{v^2\,dv}{v^4+1}. \tag{1}
\]
令 $v=\dfrac1w$，作变量替换，然后，仍把积分变量写成 $v$，则
\[
\int_0^{+\infty}\frac{v^2\,dv}{v^4+1}
=\int_0^{+\infty}\frac1{1+w^4}\,dw
=\int_0^{+\infty}\frac{dv}{1+v^4}.
\]
此式左端和右端相加，除以 2，知
\[
\begin{aligned}
\int_0^{+\infty}\frac{v^2\,dv}{v^4+1}
&=\frac12\int_0^{+\infty}\frac{1+v^2}{1+v^4}\,dv \quad\text{（拆项）}\\
&=\frac14\int_0^{+\infty}
\left(\frac1{1+v^2+\sqrt2v}+\frac1{1+v^2-\sqrt2v}\right)\,dv\\
&=\frac{\sqrt2}{4}\left[
\arctan\frac{v+\frac1{\sqrt2}}{\frac1{\sqrt2}}
+\arctan\frac{v-\frac1{\sqrt2}}{\frac1{\sqrt2}}
\right]_0^{+\infty}\\
&=\frac{\sqrt2}{4}\pi.
\end{aligned}
\]
代入式 (1) 得
\[
I=\sqrt2\pi\sqrt y.
\]

\noindent\textbf{附注}\quad 式 (2) 右端的积分可另解如下：
\[
\begin{aligned}
\int_0^{+\infty}\frac{1+v^2}{1+v^4}\,dv
&=\int_0^{+\infty}\frac{1+\frac1{v^2}}{v^2+\frac1{v^2}}\,dv
=\int_0^{+\infty}\frac{d\left(v-\frac1v\right)}
{\left(v-\frac1v\right)^2+2}\\
&=\frac1{\sqrt2}\arctan\left(v-\frac1v\right)\Big|_0^{+\infty}
=\frac{\pi}{\sqrt2}.
\end{aligned}
\]

\noindent\textbf{练习}\quad 假设 $f:\mathbb{R}\to\mathbb{R}$ 是有界连续函数，求
\[
\lim_{t\to+0}\int_{\mathbb{R}}\frac{t}{t^2+x^2}f(x)\,dx.
\]
（北京大学）\hfill《$\pi f(0)$》

\noindent\textbf{提示}\quad
\[
\lim_{t\to+0}\left|\int_1^{+\infty}\frac{t}{t^2+x^2}f(x)\,dx\right|
\le M\lim_{t\to+0}\int_1^{+\infty}\frac{t}{t^2+x^2}\,dx=0
\]
（$M$ 是 $|f(x)|$ 的上界），而
\[
\lim_{t\to+0}\int_0^1\frac{t}{t^2+x^2}f(x)\,dx
=\frac\pi2 f(0)
\]
（例 4.1.5）. 负半轴类似.

\noindent\textbf{例 4.5.2}\quad 证明等式
\[
\int_0^{+\infty}f\left(ax+\frac bx\right)\,dx
=\frac1a\int_0^{+\infty}f\left(\sqrt{t^2+4ab}\right)\,dt, \tag{1}
\]
其中 $a,b>0$（假定两积分有意义）.

\noindent\textbf{分析}\quad 比较该等式的两边，必须使得
\[
ax+\frac bx=\sqrt{t^2+4ab}. \tag{2}
\]
因 $a,b,x>0$，此即要求
\[
\left(ax+\frac bx\right)^2=t^2+4ab,
\]
亦即
\[
\left(ax-\frac bx\right)^2=t^2.
\]
故我们选取变换：
\[
ax-\frac bx=t.
\]

\noindent\textbf{证}\quad 令
\[
ax-\frac bx=t, \tag{3}
\]
此时式 (2) 成立. $(2)+(3)$ 可得
\[
x=\frac1{2a}\left(t+\sqrt{t^2+4ab}\right),\qquad
dx=\frac{t+\sqrt{t^2+4ab}}{2a\sqrt{t^2+4ab}}\,dt.
\]
于是式 (1) 左端的积分（设为 $I$）
\[
I=\int_0^{+\infty}f\left(ax+\frac bx\right)\,dx
=\frac1{2a}\left(\int_{-\infty}^0+\int_0^{+\infty}\right)
f\left(\sqrt{t^2+4ab}\right)
\frac{t+\sqrt{t^2+4ab}}{\sqrt{t^2+4ab}}\,dt,
\]
在右边第一个积分里令 $t=-u$，
\[
I=\frac1{2a}\left[
\int_0^{+\infty}f\left(\sqrt{u^2+4ab}\right)
\frac{\sqrt{u^2+4ab}-u}{\sqrt{u^2+4ab}}\,du
+\int_0^{+\infty}f\left(\sqrt{t^2+4ab}\right)
\frac{t+\sqrt{t^2+4ab}}{\sqrt{t^2+4ab}}\,dt
\right].
\]
再将 $u$ 改写成 $t$，两积分合并，
\[
I=\frac1a\int_0^{+\infty}f\left(\sqrt{t^2+4ab}\right)\,dt.
\]
式 (1) 获证.

\noindent\textbf{例 4.5.3}\quad 已知
\[
\int_0^{+\infty}e^{-x^2}\,dx=\frac{\sqrt\pi}{2},
\]
试证
\[
\int_0^{+\infty}e^{-a^2x^2-\frac{b^2}{x^2}}\,dx
=\frac{\sqrt\pi}{2a}e^{-2ab}.
\]

\noindent\textbf{提示}\quad 可利用上例的结果或方法. 注意
\[
\int_0^{+\infty}e^{-a^2x^2-\frac{b^2}{x^2}}\,dx
=e^{-2ab}\int_0^{+\infty}e^{-\left(ax+\frac bx\right)^2}\,dx
=\frac{\sqrt\pi}{2a}e^{-2ab}.
\]
利用分部积分法，常常可获得递推公式，或把困难的积分变成较易的积分.

\noindent\textbf{例 4.5.4}\quad 设 $m,n$ 为自然数，求
\[
\int_0^1 t^n(\ln t)^m\,dt.
\]
（北京师范大学）

\noindent\textbf{解}\quad
\[
\begin{aligned}
I_m&=\int_0^1 t^n(\ln t)^m\,dt
=\frac1{n+1}\int_0^1(\ln t)^m\,dt^{n+1}\\
&=\frac1{n+1}t^{n+1}(\ln t)^m\Big|_0^1
-\frac{m}{n+1}\int_0^1 t^n(\ln t)^{m-1}\,dt
=-\frac{m}{n+1}I_{m-1}.
\end{aligned}
\]
至此已得一递推公式. 反复使用此式，
\[
I_m=-\frac{m}{n+1}\left(-\frac{m-1}{n+1}\right)\cdots
\left(-\frac1{n+1}\right)I_0
=(-1)^m\frac{m!}{(n+1)^m}\int_0^1t^n\,dt
=(-1)^m\frac{m!}{(n+1)^{m+1}}.
\]

\noindent\textbf{☆例 4.5.5}\quad 计算积分
\[
\int_0^{\frac\pi2}\cos(2nx)\ln(\cos x)\,dx.
\]

\noindent\textbf{解}\quad （困难在于被积函数中有对数符号“$\ln$”，用分部积分法消去“$\ln$”.）

原式
\[
\begin{aligned}
&=\frac1{2n}\int_0^{\frac\pi2}\ln(\cos x)\,d\sin(2nx)\\
&=\frac1{2n}\sin(2nx)\ln(\cos x)\Big|_0^{\frac\pi2}
-\frac1{2n}\int_0^{\frac\pi2}\frac{\sin(2nx)(-\sin x)}{\cos x}\,dx\\
&=\frac1{2n}\int_0^{\frac\pi2}
\frac{\sin(2nx)\sin x}{\cos x}\,dx. \qquad(1)
\end{aligned}
\]
（我们看到，这里如果被积函数没有分母 $\cos x$，用积化和差公式立即可算出积分值，因此，我们希望设法应用公式
\[
\frac{\sin(2n+1)t}{\sin t}=1+2\sum_{k=1}^n\cos(2kt) \tag{2}
\]
将被积函数拆开.）因为
\[
\sin(2nx)\sin x=\cos(2nx)\cos x-\cos(2n+1)x,
\]
\[
\frac1{2n}\int_0^{\frac\pi2}\frac{\sin(2nx)\sin x}{\cos x}\,dx
=\frac1{2n}\int_0^{\frac\pi2}\cos(2nx)\,dx
-\frac1{2n}\int_0^{\frac\pi2}\frac{\cos(2n+1)x}{\cos x}\,dx,
\]
第一个积分为 0，第二个积分中令 $x=t-\dfrac\pi2$，故
\[
\begin{aligned}
\frac1{2n}\int_0^{\frac\pi2}\frac{\sin(2nx)\sin x}{\cos x}\,dx
&=\frac{(-1)^{n-1}}{2n}
\int_{\frac\pi2}^{\pi}\frac{\sin(2n+1)t}{\sin t}\,dt\quad\text{（利用公式 (2)）}\\
&=\frac{(-1)^{n-1}}{2n}
\int_{\frac\pi2}^{\pi}\left[1+2\sum_{k=1}^n\cos(2kt)\right]\,dt
=(-1)^{n-1}\frac{\pi}{4n}.
\end{aligned}
\]

\noindent\textbf{b. 其他方法}

\noindent\textbf{要点}\quad 计算反常积分，除上述三大基本方法之外，根据具体情况，需要灵活运用各种其他方法，其中比较常用的有：待定系数法，把有理分式化为部分分式，方程法，分段积分自我消去法，级数法，等等.

\noindent\textbf{※例 4.5.6}\quad 计算积分
\[
I_n=\int_1^{+\infty}\frac{dx}{x(x+1)\cdots(x+n)}.
\]

\noindent\textbf{解}\quad （拆为部分分式.）设
\[
\frac1{x(x+1)\cdots(x+n)}
=\frac{A_0}{x}+\frac{A_1}{x+1}+\cdots+\frac{A_k}{x+k}
+\cdots+\frac{A_n}{x+n}
\]
（$A_0,A_1,\cdots,A_n$ 为待定系数）. 将 $x(x+1)\cdots(x+n)$ 同乘等式两边，然后令 $x\to-k$，得
\[
\begin{aligned}
A_k&=\frac1{(-k)(-k+1)\cdots(-1)\cdot1\cdot2\cdots(-k+n)}\\
&=(-1)^k\frac1{k!(n-k)!}=(-1)^k\frac{C_n^k}{n!}\quad(k=0,1,2,\cdots,n),
\end{aligned}
\]
其中 $C_n^k=\dfrac{n!}{k!(n-k)!}$. 于是
\[
\begin{aligned}
I_n&=\int_1^{+\infty}\left(\sum_{k=0}^n(-1)^k\frac{C_n^k}{n!}\frac1{x+k}\right)dx\\
&=\sum_{k=0}^n(-1)^k\frac{C_n^k}{n!}\int_1^{+\infty}\frac1{x+k}\,dx\\
&=\frac1{n!}\sum_{k=0}^n(-1)^k C_n^k\ln(x+k)\Big|_1^{+\infty}.
\end{aligned}
\]
注意到
\[
\begin{aligned}
\sum_{k=0}^n(-1)^k C_n^k\ln(x+k)
&=\sum_{k=0}^n(-1)^k C_n^k\ln\left[x\left(1+\frac{k}{x}\right)\right]\\
&=\ln x\cdot\sum_{k=0}^n(-1)^k C_n^k
+\sum_{k=0}^n(-1)^k C_n^k\ln\left(1+\frac{k}{x}\right)\\
&=\ln x\cdot(1-1)^n+\sum_{k=0}^n(-1)^k C_n^k\ln\left(1+\frac{k}{x}\right)
\longrightarrow0\quad(x\to+\infty),
\end{aligned}
\]
因此
\[
I_n=\frac1{n!}\sum_{k=1}^n(-1)^{k+1}C_n^k\ln(1+k).
\]

\noindent\textbf{☆例 4.5.7}\quad 计算积分
\[
I=\int_0^{\frac\pi2}\ln\sin x\,dx.
\]
（武汉大学）

\noindent\textbf{提示}\quad $x=0$ 是奇点，当 $0<\lambda<1$ 时，$\int_0^{\frac\pi2}x^{-\lambda}\,dx$ 收敛. 利用 L'Hospital 法则，易知当 $x\to0^+$ 时，$x^{-\lambda}$ 比 $\ln\sin x$ 高阶，所以 $\int_0^{\frac\pi2}\ln\sin x\,dx$ 收敛.

\noindent\textbf{再提示}\quad
\[
\begin{aligned}
I&=\int_0^{\frac\pi2}\ln\sin x\,dx
=\int_0^{\frac\pi4}2\ln\sin2t\,dt\\
&=2\ln2\cdot\frac\pi4
+2\int_0^{\frac\pi4}\ln\sin t\,dt
+2\int_0^{\frac\pi4}\ln\cos t\,dt\\
&=\frac\pi2\ln2+2\int_0^{\frac\pi4}\ln\sin t\,dt
+2\int_{\frac\pi4}^{\frac\pi2}\ln\sin u\,du
\quad\left(\text{这里 }u=\frac\pi2-t\right)\\
&=\frac\pi2\ln2+2I.
\end{aligned}
\]
解方程得
\[
I=-\frac\pi2\ln2.
\]

\noindent\textbf{☆例 4.5.8}\quad 计算积分
\[
\int_0^{+\infty}\frac{\ln x}{1+x^2}\,dx.
\]
（北京航空航天大学）

\noindent\textbf{解 I}\quad
\[
\begin{aligned}
\int_0^{+\infty}\frac{\ln x}{1+x^2}\,dx
&=\int_0^1\frac{\ln x}{1+x^2}\,dx
+\int_1^{+\infty}\frac{\ln x}{1+x^2}\,dx
\quad\left(\text{第二个积分中令 }x=\frac1t\right)\\
&=\int_0^1\frac{\ln x}{1+x^2}\,dx
+\int_1^0\frac{\ln t}{1+t^2}\,dt=0.
\end{aligned}
\]

\noindent\textbf{解 II}\quad 令 $x=\dfrac1t$，
\[
I=\int_0^{+\infty}\frac{\ln x}{1+x^2}\,dx
=\int_{+\infty}^0
\frac{\ln\frac1t}{1+\frac1{t^2}}\left(-\frac1{t^2}\right)\,dt
=-\int_0^{+\infty}\frac{\ln t}{1+t^2}\,dt=-I,
\]
故
\[
I=0.
\]

\noindent\textbf{例 4.5.9}\quad 证明：
\[
\int_0^{+\infty}\frac{dx}{(1+x^2)(1+x^\alpha)}
\]
与 $\alpha$ 无关.（国外赛题）

\noindent\textbf{提示}\quad
\[
\int_0^{+\infty}\frac{dx}{(1+x^2)(1+x^\alpha)}
=\int_0^1\frac{dx}{(1+x^2)(1+x^\alpha)}
+\int_1^{+\infty}\frac{dx}{(1+x^2)(1+x^\alpha)}
=\int_0^{+\infty}\frac{dx}{1+x^2}.
\]

\noindent\textbf{例 4.5.10}\quad 求
\[
\max_{0\le s\le1}\int_0^1|\ln|s-t||\,dt.
\]

\noindent\textbf{解}\quad
\[
\begin{aligned}
I&=\int_0^1|\ln|s-t||\,dt
=-\int_0^s\ln(s-t)\,dt-\int_s^1\ln(t-s)\,dt\\
&=1-s\ln s-(1-s)\ln(1-s),
\end{aligned}
\]
\[
I_s'=\ln\left(\frac1s-1\right).
\]
令 $I_s'=0$，得 $s=\dfrac12$.
当 $s\nearrow$ 时，$I_s'$ 由正变负，所以
\[
\max_{0\le s\le1}\int_0^1|\ln|s-t||\,dt
=I\left(\frac12\right)=1+\ln2.
\]

\noindent\textbf{例 4.5.11}\quad 证明：
\[
\int_1^{+\infty}\left(\frac1{[x]}-\frac1x\right)\,dx
=\lim_{n\to\infty}\left(1+\frac12+\cdots+\frac1n-\ln n\right).
\]

\noindent\textbf{提示}\quad
$1^\circ$ 当 $x>2$ 时，
\[
\left|\frac1{[x]}-\frac1x\right|\le\frac1{(x-1)x},
\]
积分收敛；

$2^\circ$
\[
\int_1^{+\infty}\left(\frac1{[x]}-\frac1x\right)\,dx
=\lim_{n\to\infty}\int_1^n\left(\frac1{[x]}-\frac1x\right)\,dx.
\]

\subsection*{二、反常积分敛散性的判定（十二法）}

\noindent\textbf{要点}\quad （这里只就无穷限的反常积分进行叙述，对于无界函数反常积分，有类似的结果.）判定反常积分 $\displaystyle\int_a^{+\infty}f(x)\,dx$ 的敛散性的要点如下：

1） 若 $f(x)\ge0$，且 $\lim\limits_{x\to+\infty}f(x)=0$，可考察 $x\to+\infty$ 时无穷小量 $f(x)$ 的阶；若阶数 $\lambda>1$，则反常积分 $\displaystyle\int_a^{+\infty}f(x)\,dx$ 收敛；当 $\lambda\le1$ 时，发散.

2） 若 $f(x)\ge0$，可用比较判别法或比较判别法的极限形式进行判断.

3） 若 $f(x)\ge0$，可考察 $\displaystyle\int_a^A f(x)\,dx$ 是否有界.

4） 以上 $f(x)\ge0$ 的条件，只要对于充分大的 $x\ (x\ge a)$ 保持成立即可.

5） 因为 $\displaystyle\int_a^{+\infty} f(x)\,dx$ 与 $\displaystyle\int_a^{+\infty}-f(x)\,dx$ 同时敛散，故对 $f(x)\le0$ 有类似的方法.

6） 若 $x\to+\infty$ 时 $f(x)$ 无穷次变号，则以上判别法失效；可考虑用 Abel 判别法或 Dirichlet 判别法.

\noindent\textbf{Abel 判别法}\quad
若 $\displaystyle\int_a^{+\infty}f(x)\,dx$ 收敛，且 $x\nearrow+\infty$ 时，$g(x)$ 单调有界，则 $\displaystyle\int_a^{+\infty}f(x)g(x)\,dx$ 收敛.

\noindent\textbf{Dirichlet 判别法}\quad
若 $\exists M>0$，使
\[
\left|\int_a^A f(x)\,dx\right|\le M\quad(\forall A>a),
\]
且 $g(x)\nearrow0$（或 $g(x)\searrow0$），则
\[
\int_a^{+\infty}f(x)g(x)\,dx
\]
收敛.

7） 用 Abel 判别法与 Dirichlet 判别法判定为收敛，只是 $\displaystyle\int_a^{+\infty}f(x)\,dx$ 本身收敛，至于是绝对收敛还是条件收敛，还需进一步考虑 $\displaystyle\int_a^{+\infty}|f(x)|\,dx$ 收敛还是发散.

8） 若以上方法无效，可考虑用 Cauchy 准则来判断.

9） 用定义看极限 $\displaystyle\lim_{A\to+\infty}\int_a^A f(x)\,dx$ 是否存在；

10） 用分部积分法或变量替换法变成别的形式，看是否能判定它的敛散性.

11） 用级数方法判定积分的敛散性（见第五章）.

12） 用运算性质判断敛散性，例如
\[
\text{若 }\int_a^{+\infty}f(x)\,dx,\ \int_a^{+\infty}g(x)\,dx\text{ 收敛，则 }
\int_a^{+\infty}(f(x)\pm g(x))\,dx\text{ 也收敛.}
\]
\[
\text{若 }\int_a^{+\infty}f(x)\,dx\text{ 收敛，}\int_a^{+\infty}g(x)\,dx\text{ 发散，则 }
\int_a^{+\infty}(f(x)\pm g(x))\,dx\text{ 发散.}
\]

13） 对于无界函数的反常积分，以上各条都有类似结论，只是 1）要特别注意：对于无界函数的反常积分而言，此条应是 $x$ 趋向奇点时，$f(x)$ 为无穷大量. 非负函数的情况：若无穷大量的阶数 $\lambda<1$，则积分收敛；若阶数 $\lambda\ge1$，则积分发散.

\noindent\textbf{例 4.5.12}\quad 讨论
\[
\int_0^\pi\frac{\sin^{\alpha-1}x\,dx}{|1+k\cos x|^\alpha}
\]
的敛散性.

\noindent\textbf{解}\quad
$1^\circ$ 若 $|k|<1$，则积分以 $0$ 与 $\pi$ 为奇点，当 $x\to0^+$ 时，$\dfrac{\sin^{\alpha-1}x}{|1+k\cos x|^\alpha}$ 与 $\dfrac1{x^{1-\alpha}}$ 同阶；当 $x\to\pi$ 时，$\dfrac{\sin^{\alpha-1}x}{|1+k\cos x|^\alpha}$ 与 $\dfrac1{(\pi-x)^{1-\alpha}}$ 同阶. 故当且仅当 $\alpha>0$ 时，积分收敛.

$2^\circ$ 若 $k=1$，则积分仍以 $0,\pi$ 为奇点. 当 $x\to0^+$ 时，与 $1^\circ$ 中情况一样，收敛要求 $\alpha>0$. 对于奇点 $\pi$，将 $\cos x$ 在点 $\pi$ 展开，可知 $|1+\cos x|=|-1-\cos x|$ 与 $(\pi-x)^2$ 同阶；而 $\sin x=\sin(\pi-x)$ 与 $(\pi-x)$ 同阶，因此当 $x\to\pi$ 时，$\dfrac{\sin^{\alpha-1}x}{|1+\cos x|^\alpha}$ 与
\[
\frac1{(\pi-x)^{2\alpha+1-\alpha}}
=\frac1{(\pi-x)^{1+\alpha}}
\]
同阶. 故对于奇点 $\pi$，要求 $\alpha<0$. 可见 $k=1$ 时，$0,\pi$ 两奇点不能同时收敛. 故积分发散.

$3^\circ$ 若 $k>1$，记 $\theta=\arccos\left(-\dfrac1k\right)$，则积分以 $0,\theta,\pi$ 为奇点. 对 $0,\pi$，与 $1^\circ$ 中情况一样，收敛要求 $\alpha>0$. 对于奇点 $\theta$，将 $\cos x$ 在 $x=\theta$ 处展开，可知 $1+k\cos x=-k\left(-\dfrac1k-\cos x\right)=-k(\cos\theta-\cos x)$ 与 $|x-\theta|$ 同阶，因此 $\dfrac{\sin^{\alpha-1}x}{|1+k\cos x|^\alpha}$ 与 $\dfrac1{|x-\theta|^\alpha}$ 同阶，收敛要求 $\alpha<1$. 故 $k>1$ 时，积分当且仅当 $0<\alpha<1$ 时收敛.

$4^\circ$ 当积分作变换 $y=\pi-x$ 时，知
\[
\int_0^\pi\frac{\sin^{\alpha-1}x}{|1+k\cos x|^\alpha}\,dx
=\int_0^\pi\frac{\sin^{\alpha-1}x}{|1-k\cos x|^\alpha}\,dx,
\]
即表明该积分关于 $k$ 对称（是 $k$ 的偶函数）.

总结上述结果知：积分当且仅当 $|k|<1$ 且 $\alpha>0$，以及 $|k|>1$ 且 $0<\alpha<1$ 时收敛.

\noindent\textbf{例 4.5.13}\quad 设 $f(x)$ 在 $[1,+\infty)$ 上连续，对任意 $x\in[1,+\infty)$，有 $f(x)>0$. 另外，
\[
\lim_{x\to+\infty}\frac{\ln f(x)}{\ln x}=-\lambda.
\]
试证：若 $\lambda>1$，则 $\displaystyle\int_1^{+\infty}f(x)\,dx$ 收敛.（华东师范大学）

\noindent\textbf{证}\quad （用比较判别法.）因
\[
\lim_{x\to+\infty}\frac{\ln f(x)}{\ln x}=-\lambda,
\]
所以 $\forall\varepsilon>0$，$\exists A>1$，当 $x>A$ 时有
\[
\frac{\ln f(x)}{\ln x}<-\lambda+\varepsilon,
\]
即
\[
\ln f(x)<(-\lambda+\varepsilon)\ln x=\ln x^{-\lambda+\varepsilon}.
\]
所以
\[
0<f(x)<\frac1{x^{\lambda-\varepsilon}}\quad(\text{当 }x>A\text{ 时}).
\]
因 $\lambda>1$，故可取 $0<\varepsilon<\lambda-1$，于是 $\lambda-\varepsilon>1$. 根据比较判别法，积分 $\displaystyle\int_1^{+\infty}f(x)\,dx$ 收敛.

\noindent\textbf{例 4.5.14}\quad 设 $f(x)$ 在 $(-\infty,+\infty)$ 上有定义，$f(x)>0$，且在任意有限区间 $[-A,B]\ (A,B>0)$ 上可积；又有定数 $M$，使得
\[
\int_{-\infty}^{+\infty}f(x)e^{-\frac{|x|}{k}}\,dx<M
\]
对任意 $k>0$ 成立. 试证明 $\displaystyle\int_{-\infty}^{+\infty}f(x)\,dx$ 收敛.（新疆大学）

\noindent\textbf{证}\quad 要证明 $\displaystyle\int_{-\infty}^{+\infty}f(x)\,dx$ 收敛，因 $f(x)>0$，只要证明积分 $\displaystyle\int_{-A}^B f(x)\,dx$ 对 $A,B\ (A,B>0)$ 保持有界. 已知 $\exists M>0$，使得
\[
\int_{-\infty}^{+\infty}f(x)e^{-\frac{|x|}{k}}\,dx\le M\quad(\forall k>0).
\]
$\forall A,B>0$，记 $C=\max\{A,B\}$，取 $k>C$，则
\[
0\le\int_{-A}^B f(x)\,dx
\le\int_{-A}^B f(x)e^{\frac{C-|x|}{k}}\,dx
=e^{\frac Ck}\int_{-A}^B f(x)e^{-\frac{|x|}{k}}\,dx
\le e^{\frac Ck}\cdot M\le3M.
\]
故 $\displaystyle\int_{-\infty}^{+\infty}f(x)\,dx$ 收敛.

\noindent\textbf{例 4.5.15}\quad 设函数 $f(x)$ 在半开半闭区间 $(0,1]$ 里连续，且 $\lim\limits_{x\to0}f(x)=+\infty$. 对任何正整数 $N$，定义 $f_N(x)=\min\{f(x),N\}$，证明：反常积分 $\displaystyle\int_0^1 f(x)\,dx$ 收敛的充要条件是
\[
\lim_{N\to+\infty}\int_0^1 f_N(x)\,dx
\]
存在.（厦门大学）

\noindent\textbf{证}\quad $1^\circ$ 充分性. 因 $\lim\limits_{x\to0^+}f(x)=+\infty$，故 $\exists\delta>0$，当 $x\in(0,\delta)$ 时，$f(x)>0$. 故对 $\displaystyle\int_0^1 f(x)\,dx$ 的敛散性，可用非负函数的判别法进行判定. 下面我们来证明：当 $0<\alpha<\delta$ 时，$\displaystyle\int_\alpha^1f(x)\,dx$ 保持有上界.

事实上，因为 $f(x)$ 在 $[\alpha,1]$ 上连续，所以 $\exists M>0$，使得 $f(x)\le M\ (x\in[\alpha,1])$. 因而 $N>M$ 时，$[\alpha,1]$ 上恒有 $f_N(x)=\min\{f(x),N\}=f(x)$，从而
\[
\int_\alpha^1 f(x)\,dx=\int_\alpha^1 f_N(x)\,dx
\le\int_0^1 f_N(x)\,dx.
\]
令 $N\to+\infty$，取极限得
\[
\int_\alpha^1 f(x)\,dx
\le\lim_{N\to+\infty}\int_0^1 f_N(x)\,dx<+\infty.
\]
故 $\displaystyle\int_0^1 f(x)\,dx$ 收敛.

$2^\circ$ 必要性. 只要注意到 $f_N(x)$ 对 $N$ 递增，且 $f_N(x)\le f(x)$，立刻可用单调有界原理证得
\[
\lim_{N\to+\infty}\int_0^1 f_N(x)\,dx
\]
存在.

下面讨论在奇点附近无穷次变号的例子.

\noindent\textbf{☆例 4.5.16}\quad 证明积分
\[
\int_0^1\left(x\sin\frac1{x^2}-\frac1x\cos\frac1{x^2}\right)dx
\]
有意义.

\noindent\textbf{证 I}\quad
$1^\circ$ 对 $\displaystyle\int_0^1x\sin\dfrac1{x^2}\,dx$，因当 $x\to0$ 时 $x\sin\dfrac1{x^2}\to0$，故该积分为正常积分，只要补充在 $x=0$ 处为 0，则在 $[0,1]$ 上连续，所以该积分有意义.

$2^\circ$ 考虑第二项的积分. 首先
\[
\int_0^1\frac1{x^2}\cos\frac1{x^2}\,dx
=-\int_0^1\cos\frac1{x^2}\,d\frac1x
\stackrel{u=\frac1x}{=}\int_1^{+\infty}\cos u^2\,du
\]
\[
\stackrel{u=\sqrt t}{=}\int_1^{+\infty}\cos t\,d\sqrt t
=\frac12\int_1^{+\infty}\frac{\cos t}{\sqrt t}\,dt,
\]
\[
\left|\int_0^A\cos t\,dt\right|\le2,\qquad \frac1{\sqrt t}\searrow0.
\]
根据 Dirichlet 判别法，此积分收敛.

其次，原积分第二项
\[
\int_0^1\frac1x\cos\frac1{x^2}\,dx
=\int_0^1 x\cdot\frac1{x^2}\cos\frac1{x^2}\,dx.
\]
$\displaystyle\int_0^1\frac1{x^2}\cos\dfrac1{x^2}\,dx$ 收敛，因子 $x$ 单调有界. 故由 Abel 判别法，此积分收敛. 总之原积分有意义.

\noindent\textbf{证 II}\quad
\[
\int_0^1\left(x\sin\frac1{x^2}-\frac1x\cos\frac1{x^2}\right)dx
=\frac12x^2\sin\frac1{x^2}\Big|_{0+}^1=\frac12\sin1,
\]
故该积分有意义.

\noindent\textbf{☆例 4.5.17}\quad 积分
\[
\int_0^{+\infty}\left[\left(1-\frac{\sin x}{x}\right)^{-\frac13}-1\right]dx
\]
是否收敛？是否绝对收敛？证明所述结论.（北京大学）

\noindent\textbf{解}\quad
\[
\int_0^{+\infty}\left[\left(1-\frac{\sin x}{x}\right)^{-\frac13}-1\right]dx
=\int_0^1\left(1-\frac{\sin x}{x}\right)^{-\frac13}dx-\int_0^1dx
+\int_1^{+\infty}\left[\left(1-\frac{\sin x}{x}\right)^{-\frac13}-1\right]dx,
\]
其中 $\displaystyle\int_0^1\left(1-\frac{\sin x}{x}\right)^{-\frac13}dx$ 以 $x=0$ 为奇点，
\[
\left(1-\frac{\sin x}{x}\right)^{-\frac13}
=\left[\frac1{3!}x^2+o(x^2)\right]^{-\frac13}
\]
与 $\dfrac1{x^{2/3}}$ 同阶，故 $\displaystyle\int_0^1\left(1-\frac{\sin x}{x}\right)^{-\frac13}dx$ 收敛. 因为 $\left(1-\dfrac{\sin x}{x}\right)>0$，收敛即为绝对收敛.

其次对积分
\[
\int_1^{+\infty}\left[\left(1-\frac{\sin x}{x}\right)^{-\frac13}-1\right]dx,
\]
因为当 $x>1$ 时，$\left|\dfrac{\sin x}{x}\right|<1$，可利用 $(1+x)^\alpha$ 的 Taylor 公式，有
\[
\left(1-\frac{\sin x}{x}\right)^{-\frac13}-1
=\frac13\frac{\sin x}{x}+o\left(\frac1{x^2}\right),
\]
于是
\[
\int_1^{+\infty}\left[\left(1-\frac{\sin x}{x}\right)^{-\frac13}-1\right]dx
=\frac13\int_1^{+\infty}\frac{\sin x}{x}\,dx
+\int_1^{+\infty}o\left(\frac1{x^2}\right)\,dx.
\]
而 $\displaystyle\int_1^{+\infty}\frac{\sin x}{x}\,dx$ 条件收敛，$\displaystyle\int_1^{+\infty}o\left(\frac1{x^2}\right)\,dx$ 绝对收敛. 故原积分条件（不绝对）收敛.

\noindent\textbf{例 4.5.18}\quad 设 $\alpha,\beta$ 为实数，试讨论积分
\[
I=\int_0^{+\infty}x^\alpha\sin(x^\beta)\,dx
\]
的敛散性.（中国科学院）

\noindent\textbf{解}\quad 若 $\beta=0$，则
\[
I=\sin1\int_0^1x^\alpha\,dx+\sin1\int_1^{+\infty}x^\alpha\,dx.
\]
不论 $\alpha<-1$ 还是 $\alpha\ge-1$，积分都发散.

若 $\beta\ne0$，则
\[
I\stackrel{t=x^\beta}{=}
\begin{cases}
\displaystyle\int_0^{+\infty}t^{\frac\alpha\beta}\sin t\cdot\frac1\beta t^{\frac1\beta-1}\,dt,&\beta>0,\\[1.1em]
\displaystyle\int_{+\infty}^0t^{\frac\alpha\beta}\sin t\cdot\frac1\beta t^{\frac1\beta-1}\,dt,&\beta<0
\end{cases}
\]
\[
=\frac1{|\beta|}\int_0^{+\infty}t^{\frac{\alpha+1-\beta}{\beta}}\sin t\,dt
\quad\left(\text{记 }\mu=\frac{\alpha+1-\beta}{\beta}\right)
\]
\[
=\frac1{|\beta|}\int_0^1t^\mu\sin t\,dt
+\frac1{|\beta|}\int_1^{+\infty}t^\mu\sin t\,dt
\equiv I_1+I_2.
\]
对于 $I_1=\dfrac1{|\beta|}\int_0^1t^\mu\sin t\,dt$，因为
\[
\lim_{t\to0^+}\frac{t^\mu\sin t}{t^{\mu+1}}=1,
\]
故 $I_1$ 与 $\displaystyle\int_0^1t^{\mu+1}\,dt$ 同时敛散. 因而当且仅当 $-\mu-1<1$（即 $\mu>-2$），亦即 $\dfrac{\alpha+1}{\beta}>-1$ 时 $I_1$ 收敛. 因为被积函数为正，收敛亦为绝对收敛.

对于 $I_2=\dfrac1{|\beta|}\int_1^{+\infty}t^\mu\sin t\,dt$，我们只需讨论 $-\mu-1<1$（即 $\mu>-2$）时的情况.

i） 当 $-2<\mu<-1$（即 $-1<\dfrac{\alpha+1}{\beta}<0$）时，因 $|t^\mu\sin t|\le t^\mu$，且 $\displaystyle\int_1^{+\infty}t^\mu\,dt$ 收敛，所以此时 $I_2$ 绝对收敛.

ii） 当 $-1\le\mu<0$（即 $0\le\dfrac{\alpha+1}{\beta}<1$）时，随 $x\nearrow+\infty$，$x^\mu\searrow0$，且
\[
\left|\int_1^A\sin t\,dt\right|
=|-\cos A+\cos1|\le2
\]
（有界）. 由 Dirichlet 判别法，$I_2$ 收敛，且由
\[
|t^\mu\sin t|\ge t^\mu\sin^2t=\frac{t^\mu}{2}-\frac{t^\mu\cos2t}{2},
\]
知 $I_2$ 非绝对收敛.

iii） 当 $\mu\ge0$（即 $\dfrac{\alpha+1}{\beta}\ge1$）时，因 $\forall k\in\mathbb{N}$，有
\[
\left|\int_{2k\pi}^{(2k+1)\pi}t^\mu\sin t\,dt\right|
\ge\int_0^\pi\sin t\,dt=2,
\]
所以 $I_2$ 发散.

总之，原积分当 $-1<\dfrac{\alpha+1}{\beta}<0$ 时，绝对收敛；当 $0\le\dfrac{\alpha+1}{\beta}<1$ 时，条件收敛；其他情况发散.

\noindent\textbf{例 4.5.19}\quad 讨论
\[
I=\int_0^{+\infty}\frac{\sin\left(x+\frac1x\right)}{x^\alpha}\,dx
\]
的绝对收敛性与条件收敛性.

\noindent\textbf{证}\quad 积分 $I$ 的反常点为 $0$ 和 $+\infty$. 将 $I$ 拆成两项：
\[
I=\int_0^{+\infty}\frac{\sin\left(x+\frac1x\right)}{x^\alpha}\,dx
=\int_0^1\frac{\sin\left(x+\frac1x\right)}{x^\alpha}\,dx
+\int_1^{+\infty}\frac{\sin\left(x+\frac1x\right)}{x^\alpha}\,dx
\equiv I_1+I_2.
\]

i） 显然当 $\alpha\le0$ 时 $I_2$ 发散. 作变换 $x=\dfrac1t$，易知 $\alpha\ge2$ 时 $I_1$ 发散.

ii） 考虑 $0<\alpha<2$ 的情况，将积分 $I$ 写成
\[
\begin{aligned}
I&=\int_0^{+\infty}
\frac{\left(1-\frac1{x^2}\right)\sin\left(x+\frac1x\right)}
{x^\alpha\left(1-\frac1{x^2}\right)}\,dx\\
&=\left(\int_0^1+\int_1^{+\infty}\right)
\frac{\left(1-\frac1{x^2}\right)\sin\left(x+\frac1x\right)}
{x^\alpha\left(1-\frac1{x^2}\right)}\,dx
\equiv I_3+I_4,
\end{aligned}
\]
其中
\[
\begin{aligned}
\left|\int_1^A\left(1-\frac1{x^2}\right)
\sin\left(x+\frac1x\right)\,dx\right|
&=\left|\int_1^A\sin\left(x+\frac1x\right)d\left(x+\frac1x\right)\right|\\
&=\left|\cos2-\cos\left(A+\frac1A\right)\right|\le2.
\end{aligned}
\]
关于 $A>1$ 有界，且 $x\nearrow+\infty$ 时，
\[
\frac1{x^\alpha\left(1-\frac1{x^2}\right)}
=\frac1{x^\alpha-\frac1{x^{2-\alpha}}}\searrow0.
\]
由 Dirichlet 判别法知，$I_4$ 当 $0<\alpha<2$ 时收敛. 作变换 $x=\dfrac1t$，类似可知 $I_3$ 也收敛. 故 $I$ 在 $(0,2)$ 内收敛.

iii） 证明 $I$ 对 $\alpha\in(0,2)$ 非绝对收敛.
\[
\left|\frac{\sin\left(x+\frac1x\right)}{x^\alpha}\right|
\ge\frac{\sin^2\left(x+\frac1x\right)}{x^\alpha}
=\frac1{2x^\alpha}-\frac1{2x^\alpha}\cos2\left(x+\frac1x\right).
\]
当 $0<\alpha\le1$ 时，此不等式右端第 1 项 $\dfrac1{2x^\alpha}$ 在 $[1,+\infty)$ 上的积分发散；第 2 项 $\dfrac1{2x^\alpha}\cos2\left(x+\dfrac1x\right)$ 在 $[1,+\infty)$ 上的积分收敛（证法与 ii）类似）. 因而正值函数
\[
\frac{\sin^2\left(x+\frac1x\right)}{x^\alpha}
\]
在 $[1,+\infty)$ 上的积分发散，从而
\[
\left|\frac{\sin\left(x+\frac1x\right)}{x^\alpha}\right|
\]
在 $[0,+\infty)$ 上的积分发散，故 $I$ 当 $0<\alpha\le1$ 时非绝对收敛.

类似可证当 $1<\alpha<2$ 时也非绝对收敛.

总结 i）、ii）、iii），积分 $I$ 当且仅当 $\alpha\in(0,2)$ 时才条件收敛.

\noindent\textbf{☆例 4.5.20}\quad 设 $f(x)$ 在 $(a,+\infty)$ 内可微，$f'(x)$ 可积，且当 $x\to+\infty$ 时，$f(x)\searrow0$；又积分 $\displaystyle\int_a^{+\infty}f(x)\,dx$ 收敛. 试证：
\[
\int_a^{+\infty}xf'(x)\,dx
\]
收敛.（辽宁师范大学，北京大学，哈尔滨工业大学）

\noindent\textbf{证}\quad
\[
\int_a^{+\infty}xf'(x)\,dx
=\int_a^{+\infty}x\,df(x)
=xf(x)\Big|_a^{+\infty}-\int_a^{+\infty}f(x)\,dx.
\]
已知 $\displaystyle\int_a^{+\infty}f(x)\,dx$ 收敛，故 $\displaystyle\int_a^{+\infty}xf'(x)\,dx$ 收敛与否取决于极限 $\lim\limits_{x\to+\infty}xf(x)$ 是否存在.

因为 $\displaystyle\int_a^{+\infty}f(x)\,dx$ 收敛，利用 Cauchy 准则，知：$\forall\varepsilon>0$，$\exists A>0$，当 $x>A$ 时，有
\[
\int_x^{2x}f(t)\,dt<\frac\varepsilon2.
\]
因 $f(x)\searrow0$，$[\frac x2,x]$ 上 $f$ 的最小值为 $f(x)$，所以当 $x>2A$ 时，
\[
0\le xf(x)=2f(x)\int_{\frac x2}^xdt
\le2\int_{\frac x2}^xf(t)\,dt<\varepsilon.
\]
表明存在极限
\[
\lim_{x\to+\infty}xf(x)=0.
\]
证毕.

\noindent\textbf{例 4.5.21}\quad 设 $f(x)>0\searrow$，试证：
\[
\int_a^{+\infty}f(x)\,dx
\]
与
\[
\int_a^{+\infty}f(x)\sin^2x\,dx
\]
同时敛散.

\noindent\textbf{证}\quad 因 $f(x)>0\searrow$，故 $f(x)\searrow0$ 或 $f(x)\searrow A(>0)$.

$1^\circ$ 若 $f(x)\searrow0$，则由 Dirichlet 判别法知，$\displaystyle\int_a^{+\infty}f(x)\cos2x\,dx$ 收敛，从而由
\[
\int_a^{+\infty}f(x)\sin^2x\,dx
=\int_a^{+\infty}f(x)\frac{1-\cos2x}{2}\,dx
=\frac12\int_a^{+\infty}f(x)\,dx
-\frac12\int_a^{+\infty}f(x)\cos2x\,dx
\]
知，$\displaystyle\int_a^{+\infty}f(x)\sin^2x\,dx$ 与 $\displaystyle\int_a^{+\infty}f(x)\,dx$ 同时敛散.

$2^\circ$ 若 $f(x)\searrow A(>0)$，则易证两积分发散. 总之两积分同时敛散.

\noindent\textbf{例 4.5.22}\quad 讨论如下积分的敛散性：
\[
1)\ \int_2^{+\infty}\frac{\sin^2x}{x^p(x^p+\sin x)}\,dx\quad(p>0);
\]
\[
2)\ \int_2^{+\infty}\frac{\sin x}{x^p+\sin x}\,dx\quad(p>0);
\]
\[
3)\ \int_0^{+\infty}\frac{\sin x}{x^p+\sin x}\,dx\quad(p>0).
\]

\noindent\textbf{解}\quad
1） 为非负函数的积分，可用比较判别法，由不等式
\[
\frac{\sin^2x}{x^p(x^p+1)}
<\frac{\sin^2x}{x^p(x^p+\sin x)}
<\frac1{x^p(x^p-1)}
\]
知
\[
\text{若 }p>\frac12,\text{ 则积分 }\int_2^{+\infty}\frac1{x^p(x^p-1)}\,dx\text{ 收敛，从而 }
\int_2^{+\infty}\frac{\sin^2x}{x^p(x^p+\sin x)}\,dx\text{ 收敛；}
\]
\[
\text{若 }p\le\frac12,\text{ 由积分 }\int_2^{+\infty}\frac{dx}{x^p(x^p+1)}\text{ 发散，由上例可知 }
\int_2^{+\infty}\frac{\sin^2x}{x^p(x^p+1)}\,dx\text{ 亦发散，}
\]
从而
\[
\int_2^{+\infty}\frac{\sin^2x}{x^p(x^p+\sin x)}\,dx
\]
发散.

2） 利用 1）之结果及等式
\[
\frac{\sin x}{x^p+\sin x}
=\frac{\sin x}{x^p}
-\frac{\sin^2x}{x^p(x^p+\sin x)}
\]
可知，积分
\[
\int_2^{+\infty}\frac{\sin x}{x^p+\sin x}\,dx
\]
当且仅当 $p>\dfrac12$ 时收敛.

3） 因 $x\to0^+$ 时，$\dfrac{\sin x}{x^p+\sin x}\to C(p)$（与 $p$ 有关的常数），故 $0$ 不是奇点，敛散性与 2）相同.

\noindent\textbf{例 4.5.23}\quad 证明如下积分收敛：
\[
\int_0^{+\infty}x\sin x^4\sin x\,dx.
\]

\noindent\textbf{证}\quad 设 $A''>A'>A$，利用分部积分法，
\[
\begin{aligned}
\int_{A'}^{A''}x\sin x^4\sin x\,dx
&=-\frac{\sin x\cos x^4}{4x^2}\Big|_{A'}^{A''}
+\frac14\int_{A'}^{A''}\frac{\cos x^4\cos x}{x^2}\,dx\\
&\quad-\frac12\int_{A'}^{A''}\frac{\cos x^4\sin x}{x^3}\,dx
\longrightarrow0\quad(A\to+\infty).
\end{aligned}
\]
故积分收敛.

利用级数判断反常积分的敛散性问题，请见下章例 5.1.50 等.

\subsection*{三、无穷限反常积分的敛散性与无穷远处的极限}

本段我们来讨论 $\displaystyle\int_a^{+\infty}f(x)\,dx$ 收敛与 $\lim\limits_{x\to+\infty}f(x)=0$ 的关系.

1） 我们知道，$\displaystyle\int_a^{+\infty}f(x)\,dx$ 收敛一般不意味着 $f(x)\to0\ (x\to+\infty)$. 例如
\[
\int_0^{+\infty}\sin x^2\,dx
=\int_0^{+\infty}\frac{\sin t}{2\sqrt t}\,dt
\quad(x=\sqrt t)
\]
收敛，但 $\sin x^2\nrightarrow0$（当 $x\to+\infty$ 时）.

2） $\displaystyle\int_a^{+\infty}f(x)\,dx$ 收敛，并且 $f(x)\ge0$，仍不能断言 $f(x)\to0\ (x\to+\infty)$. 例如
\[
f(x)=
\begin{cases}
\dfrac1{1+x^2},&x\ne\text{整数},\\
1,&x=\text{整数}.
\end{cases}
\]

3） $\displaystyle\int_a^{+\infty}f(x)\,dx$ 收敛，$f(x)\ge0$，$f(x)$ 连续，还可能 $f(x)\nrightarrow0\ (x\to+\infty)$. 例如：$n=1,2,\cdots$，
\[
f(x)=
\begin{cases}
1,&x=n,\\
0,&x=n+\dfrac1{2^n},\\
\text{直线段},&x\in\left[n-\dfrac1{2^n},n\right]\text{ 或 }x\in\left[n,n+\dfrac1{2^n}\right],\\
0,&\text{其余},
\end{cases}
\]
如图 4.5.1，此函数可以简单表示为
\[
f(x)=
\begin{cases}
1-2^n|x-n|,&x\in\left[n-\dfrac1{2^n},n+\dfrac1{2^n}\right]\ (n=1,2,\cdots),\\
0,&\text{其余}.
\end{cases}
\]
此时
\[
\int_0^{+\infty}f(x)\,dx
=\sum_{n=1}^\infty\frac12\cdot\frac2{2^n}\cdot1
=\sum_{n=1}^\infty\frac1{2^n}=1,
\]
收敛，$f(x)\ge0$，连续，但 $f(x)\nrightarrow0$（当 $x\to+\infty$ 时）.

\begin{center}
\includegraphics[width=0.75\textwidth]{4-5-1.png}
\end{center}

4） 上述条件将 $f(x)\ge0$ 改为 $f(x)>0$，依然不能肯定 $f(x)\to0\ (x\to+\infty)$. 这时只要考虑函数
\[
f(x)=\max\left\{\frac1{x^2},\varphi(x)\right\},
\]
其中 $\varphi(x)$ 按 3）中的 $f(x)$ 同样的方式定义.

5） 若 $f(x)$ 单调，$\displaystyle\int_a^{+\infty}f(x)\,dx$ 收敛，则
\[
\lim_{x\to+\infty}f(x)=0.
\]
（自证）

6） 若 $f(x)$ 在 $[a,+\infty)$ 上一致连续（或更强些，设 $f(x)$ 有有界导数），则可由
\[
\int_a^{+\infty}f(x)\,dx
\]
收敛推出
\[
\lim_{x\to+\infty}f(x)=0.
\]

\noindent\textbf{☆例 4.5.24}\quad 试证明：若 $f(x)$ 在 $[a,+\infty)$ 上一致连续，且反常积分 $\displaystyle\int_a^{+\infty}f(x)\,dx$ 收敛，则
\[
\lim_{x\to+\infty}f(x)=0.
\]
（武汉大学）

\noindent\textbf{证}\quad （反证法）若当 $x\to+\infty$ 时，$f(x)\nrightarrow0$，则 $\exists\varepsilon_0>0$，使得 $\forall A>0$，$\exists x_1>A$：
\[
|f(x_1)|\ge\varepsilon_0.
\]
又因为 $f(x)$ 在 $[a,+\infty)$ 上一致连续，对 $\dfrac{\varepsilon_0}2>0$，$\exists\delta>0$，当 $|x'-x''|<\delta$ 时，有
\[
|f(x')-f(x'')|<\frac{\varepsilon_0}2.
\]
故当 $x\in[x_1,x_1+\delta]$ 时，有
\[
|f(x)|\ge\big||f(x_1)|-|f(x_1)-f(x)|\big|>\frac{\varepsilon_0}2, \tag{1}
\]
并且 $f(x)$ 与 $f(x_1)$ 同号（因为不然的话，$|f(x)-f(x_1)|>\varepsilon_0$，产生矛盾）. 若 $f(x_1)>0$，则 $f(x)>0$，从而由式 (1) 知 $f(x)>\dfrac{\varepsilon_0}2$，故
\[
\left|\int_{x_1}^{x_1+\delta}f(x)\,dx\right|
\ge\frac{\varepsilon_0}2\int_{x_1}^{x_1+\delta}dx
=\frac{\varepsilon_0}2\delta.
\]
同理，若 $f(x_1)<0$，亦有
\[
\left|\int_{x_1}^{x_1+\delta}f(x)\,dx\right|
\ge\frac{\varepsilon_0}2\delta.
\]
这就证明了：对 $\dfrac{\varepsilon_0}2\delta>0$，$\forall A$，$\exists x_1+\delta>x_1>A$，使得
\[
\left|\int_{x_1}^{x_1+\delta}f(x)\,dx\right|
\ge\frac{\varepsilon_0}2\delta.
\]
根据 Cauchy 准则，此即表明 $\displaystyle\int_a^{+\infty}f(x)\,dx$ 发散，矛盾. 证毕.

以上关于“$\displaystyle\int_a^{+\infty}f(x)\,dx$ 收敛得出 $f(x)\to0\ (x\to+\infty)$”的讨论并没有完，如下例.

\noindent\textbf{☆例 4.5.25}\quad 证明：若 $f(x)$ 连续可微，积分 $\displaystyle\int_a^{+\infty}f(x)\,dx$ 和 $\displaystyle\int_a^{+\infty}f'(x)\,dx$ 都收敛，则当 $x\to+\infty$ 时，有 $f(x)\to0$.（新疆大学）

\noindent\textbf{证 I}\quad 要证明 $x\to+\infty$ 时 $f(x)$ 有极限，根据 Heine 定理，我们只要证明 $\forall\{x_n\}\to+\infty$，恒有 $\{f(x_n)\}$ 收敛. 事实上，已知积分 $\displaystyle\int_a^{+\infty}f'(x)\,dx$ 收敛，根据 Cauchy 准则，$\forall\varepsilon>0$，$\exists A>a$，以致 $\forall x_1,x_2>A$，恒有
\[
\left|\int_{x_1}^{x_2}f'(x)\,dx\right|
=|f(x_2)-f(x_1)|<\varepsilon.
\]
如此 $\forall\{x_n\}\to+\infty$，$\exists N>0$，当 $n,m>N$ 时，有 $x_n,x_m>A$，从而
\[
\left|\int_{x_n}^{x_m}f'(x)\,dx\right|
=|f(x_n)-f(x_m)|<\varepsilon.
\]
这即表明 $\{f(x_n)\}$ 收敛. 故由 Heine 定理，极限 $\lim\limits_{x\to+\infty}f(x)=\alpha$ 存在.

现在来证 $\alpha=0$.

若 $\alpha>0$，则由保号性，$\exists A>0$，当 $x>A$ 时，有 $f(x)>\dfrac\alpha2>0$，从而 $A>A$ 时，
\[
\int_A^{2A}f(x)\,dx\ge\frac\alpha2A\longrightarrow+\infty\quad(A\to+\infty\text{ 时}),
\]
这与 $\displaystyle\int_a^{+\infty}f(x)\,dx$ 收敛矛盾.

同理可论 $\alpha<0$ 也不可能. 故
\[
\lim_{x\to+\infty}f(x)=\alpha=0.
\]

\noindent\textbf{证 II}\quad （反证法）
$1^\circ$ 若 $\lim\limits_{x\to+\infty}f(x)\ne0$，则 $\exists\varepsilon_0>0$ 及 $x_n\to+\infty$，使得 $|f(x_n)|>\varepsilon_0$. 设 $\{f(x_n)\}$ 中有无穷多项为正（无穷多项为负，类似可证），则可将负项去掉. 若把 $\{x_n\}$ 看成是剩下的点列，于是
\[
f(x_n)>\varepsilon_0\quad(n=1,2,\cdots).
\]

$2^\circ$ 因 $\displaystyle\int_a^{+\infty}f(x)\,dx$ 收敛，可知 $\exists\{x'_m\}\to+\infty$ 使得
\[
f(x'_m)<\frac{\varepsilon_0}2\quad(m=1,2,\cdots)
\]
（因为不然的话，$\exists A>0$，当 $x>A$ 时，恒有 $f(x)\ge\dfrac{\varepsilon_0}2$，于是当 $A>A$ 时，
\[
\int_A^{2A}f(x)\,dx\ge\frac{\varepsilon_0}2A\longrightarrow+\infty\quad(A\to+\infty\text{ 时}),
\]
与 $\displaystyle\int_a^{+\infty}f(x)\,dx$ 收敛矛盾）.

$3^\circ$ 于是，$\forall n,m$，有
\[
\left|\int_{x'_m}^{x_n}f'(x)\,dx\right|
=|f(x_n)-f(x'_m)|\ge\frac{\varepsilon_0}2>0,
\]
与 $\displaystyle\int_a^{+\infty}f'(x)\,dx$ 收敛矛盾.

当函数单调时，问题常常变得很简单，结果一般更深入，如下例，不仅得出 $f(x)\to0$（当 $x\to+\infty$ 时），而且对“阶”作了估计.

\noindent\textbf{☆例 4.5.26}\quad 设 $[a,+\infty)$ 上 $f(x)\searrow$，且 $\displaystyle\int_a^{+\infty}f(x)\,dx$ 收敛，试证：
\[
\lim_{x\to+\infty}xf(x)=0.
\]
（内蒙古大学）

\noindent\textbf{证}\quad 首先，我们有 $f(x)\ge0$. 因为，若存在某个 $x_1$，使得 $f(x_1)<0$，则当 $x>x_1$ 时，恒有 $f(x)<f(x_1)<0$，从而 $\displaystyle\int_a^{+\infty}f(x)\,dx$ 发散，与已知条件矛盾.

其次，由 $\displaystyle\int_a^{+\infty}f(x)\,dx$ 收敛，知 $\forall\varepsilon>0$，$\exists A>a$，当 $A''>A'>A$ 时，有
\[
\int_{A'}^{A''}f(x)\,dx<\frac\varepsilon2.
\]
故 $\forall x>2A$，有
\[
0\le xf(x)\le2\int_{\frac x2}^xf(t)\,dt<\varepsilon.
\]
此即
\[
\lim_{x\to+\infty}xf(x)=0.
\]

\noindent\textbf{例 4.5.27}\quad 设 $\displaystyle\int_a^{+\infty}f(x)\,dx$ 收敛，$xf(x)$ 在 $[a,+\infty)$ 上单调下降，求证
\[
\lim_{x\to+\infty}xf(x)\ln x=0.
\]

\noindent\textbf{证}\quad 当 $x>1$ 时，
\[
\begin{aligned}
\int_{\sqrt x}^{x}f(t)\,dt
&=\int_{\sqrt x}^{x}tf(t)\frac1t\,dt
=\xi f(\xi)\int_{\sqrt x}^{x}\frac1t\,dt\quad(\sqrt x\le\xi\le x)\\
&=\xi f(\xi)\ln\sqrt x
\ge\frac12xf(x)\ln x\ge0.
\end{aligned}
\]
根据 Cauchy 准则，当 $x\to+\infty$ 时，$\displaystyle\int_{\sqrt x}^{x}f(t)\,dt\to0$. 所以，
\[
\lim_{x\to+\infty}xf(x)\ln x=0.
\]

\noindent\textbf{练习}\quad 若 $\displaystyle\int_a^{+\infty}f(x)\,dx$ 收敛，$\dfrac{f(x)}x$ 在 $[a,+\infty)$ 上单调下降（$a>0$）. 求证：
\[
\lim_{x\to+\infty}xf(x)=0.
\]
（南开大学）

\noindent\textbf{提示}\quad 先证 $f(x)\ge0$，再利用
\[
\int_{\frac x2}^{x}f(t)\,dt
\ge\frac{f(x)}x\int_{\frac x2}^{x}t\,dt
=\frac38xf(x)\ge0.
\]

\noindent\textbf{再提示}\quad
$1^\circ$ （用反证法证 $f(x)\ge0$.）若 $\exists x_0>a>0$ 使得 $f(x_0)<0$，则
\[
\int_{x_0}^A f(x)\,dx
\stackrel{x=tx_0}{=}x_0\int_1^{A/x_0}\frac{f(tx_0)}{tx_0}\cdot tx_0\,dt
\quad\left(\text{因为在 }[a,+\infty)\text{ 上 }\frac{f(x)}x\searrow\right)
\]
\[
\le x_0\int_1^{A/x_0}\frac{f(x_0)}{x_0}\cdot tx_0\,dt
=x_0f(x_0)\int_1^{A/x_0}t\,dt\to-\infty\quad(A\to+\infty),
\]
与 $\displaystyle\int_a^{+\infty}f(x)\,dx$ 收敛矛盾. 所以只能 $f(x)\ge0$.

$2^\circ$ 由 $\dfrac{f(x)}x\searrow$ 知：$\forall\varepsilon>0$，有
\[
\int_{\frac x2}^{x}f(t)\,dt
=\int_{\frac x2}^{x}\frac{f(t)}t\cdot t\,dt
\ge\int_{\frac x2}^{x}\frac{f(x)}x\cdot t\,dt
=\frac38xf(x)\ge0.
\]
而 $\displaystyle\int_a^{+\infty}f(x)\,dx$ 收敛，由 Cauchy 准则，知 $\displaystyle\int_{\frac x2}^{x}f(t)\,dt\to0\ (x\to+\infty)$. 故 $\lim\limits_{x\to+\infty}xf(x)=0$.

\subsection*{四、反常积分的极限}

\noindent\textbf{要点}\quad 反常积分作为某参数的函数，可以提出求极限的问题，这种问题，原则上还是应用第一章介绍的求极限的各种方法，不同的是现在要充分利用反常积分的定义以及各种性质. 至于在积分号下取极限，要用到含参变量反常积分的理论，这方面的内容，我们移到含参变量反常积分里叙述，见第七章例 7.1.25 等.

\noindent\textbf{☆例 4.5.28}\quad 设 $f(x)$ 对一切 $b\ (0<b<+\infty)$ 在 $[0,b]$ 上可积，且
\[
\lim_{x\to+\infty}f(x)=\alpha,
\]
证明：
\[
\lim_{t\to0^+}t\int_0^{+\infty}e^{-tx}f(x)\,dx=\alpha. \tag{1}
\]
（中山大学）

\noindent\textbf{证 I}\quad （拟合法）注意到 $t\ne0$ 时，
\[
\int_0^{+\infty}te^{-tx}\,dx=1.
\]
故我们可把 $\alpha$ 改写成
\[
\alpha=\int_0^{+\infty}\alpha te^{-tx}\,dx.
\]
于是
\[
\left|t\int_0^{+\infty}e^{-tx}f(x)\,dx-\alpha\right|
=\left|\int_0^{+\infty}te^{-tx}(f(x)-\alpha)\,dx\right|
\le\int_0^{+\infty}te^{-tx}|f(x)-\alpha|\,dx.
\]
（我们来证明右端积分，当 $t$ 充分接近 $0^+$ 时，能任意小.）因 $\lim\limits_{x\to+\infty}f(x)=\alpha$，所以 $\forall\varepsilon>0$，$\exists A>0$，当 $x>A$ 时，有 $|f(x)-\alpha|<\dfrac\varepsilon2$. 取 $A=A$，于是
\[
\begin{aligned}
\int_0^{+\infty}te^{-tx}|f(x)-\alpha|\,dx
&=\int_0^Ate^{-tx}|f(x)-\alpha|\,dx
+\int_A^{+\infty}te^{-tx}|f(x)-\alpha|\,dx\\
&\le\int_0^Ate^{-tx}(|f(x)|+|\alpha|)\,dx
+\frac\varepsilon2\int_A^{+\infty}te^{-tx}\,dx\\
&\le M\int_0^Ate^{-tx}\,dx+\frac\varepsilon2
=M(1-e^{-tA})+\frac\varepsilon2.
\end{aligned}
\]
（因 $f(x)$ 在 $[0,A]$ 上有界，所以 $\exists M>0$，使得 $|f(x)|+|\alpha|\le M$.）因当 $t\to0^+$ 时，$1-e^{-tA}\to0$，所以对 $\varepsilon>0$，$\exists\delta>0$，使得当 $0<t<\delta$ 时，有 $M(1-e^{-tA})<\dfrac\varepsilon2$，故
\[
\int_0^{+\infty}te^{-tx}|f(x)-\alpha|\,dx
<\frac\varepsilon2+\frac\varepsilon2=\varepsilon.
\]

\noindent\textbf{证 II}\quad （利用上、下极限.）要证明式 (1)，只要证明：$\forall\varepsilon>0$，有
\[
\varlimsup_{t\to0^+}t\int_0^{+\infty}e^{-tx}f(x)\,dx\le\alpha+\varepsilon \tag{2}
\]
及
\[
\varliminf_{t\to0^+}t\int_0^{+\infty}e^{-tx}f(x)\,dx\ge\alpha-\varepsilon. \tag{3}
\]
因为 $\lim\limits_{x\to+\infty}f(x)=\alpha$，所以 $\forall\varepsilon>0$，$\exists A>0$，使当 $x>A$ 时，$\alpha-\varepsilon<f(x)<\alpha+\varepsilon$. 故
\[
\begin{aligned}
t\int_0^{+\infty}e^{-tx}f(x)\,dx
&=t\int_0^Ae^{-tx}f(x)\,dx+\int_A^{+\infty}te^{-tx}f(x)\,dx\\
&<t\int_0^Ae^{-tx}f(x)\,dx
+(\alpha+\varepsilon)\int_A^{+\infty}te^{-tx}\,dx\\
&\le M(1-e^{-tA})+(\alpha+\varepsilon)e^{-tA}.
\end{aligned}
\]
（因 $f(x)$ 在 $[0,A]$ 上有界，即 $\exists M>0$，使得 $x\in[0,A]$ 时，有 $f(x)\le M$.）令 $t\to0^+$，在不等式两边同时取极限，得
\[
\varlimsup_{t\to0^+}t\int_0^{+\infty}e^{-tx}f(x)\,dx
\le\lim_{t\to0^+}[M(1-e^{-tA})+(\alpha+\varepsilon)e^{-tA}]
=\alpha+\varepsilon.
\]
这便证明了式 (2). 类似可证式 (3).

\noindent\textbf{例 4.5.29}\quad 求极限
\[
\lim_{x\to0^+}x^\alpha\int_x^1\frac{f(t)}{t^{\alpha+1}}\,dt,
\]
其中 $\alpha>0$，$f(x)$ 为闭区间 $[0,1]$ 上的连续函数.

\noindent\textbf{证 I}\quad （应用 L'Hospital 法则.）
\[
\lim_{x\to0^+}x^\alpha\int_x^1\frac{f(t)}{t^{\alpha+1}}\,dt
=\lim_{x\to0^+}
\frac{\displaystyle\int_x^1\frac{f(t)}{t^{\alpha+1}}\,dt}
{\displaystyle\frac1{x^\alpha}}.
\]
因 $x\to0$ 时，$\dfrac1{x^\alpha}\to+\infty$，使用 L'Hospital 法则，
\[
\text{上式}=\lim_{x\to0^+}
\frac{-\dfrac{f(x)}{x^{\alpha+1}}}
{-\alpha\dfrac1{x^{\alpha+1}}}
=\lim_{x\to0^+}\frac{f(x)}\alpha=\frac{f(0)}\alpha.
\]

\noindent\textbf{证 II}\quad
\[
x^\alpha\int_x^1\frac{f(t)}{t^{\alpha+1}}\,dt
=x^\alpha\int_x^1\frac{f(t)-f(0)}{t^{\alpha+1}}\,dt
+x^\alpha\int_x^1\frac{f(0)}{t^{\alpha+1}}\,dt. \tag{1}
\]
我们来证明右端第一项的极限为 $0$，第二项的极限为 $\dfrac{f(0)}\alpha$.

因 $f(x)$ 在 $x=0$ 处连续，所以 $\forall\varepsilon>0$，$\exists\delta>0\ (\delta<1)$，当 $0<x<\delta$ 时，有
$|f(x)-f(0)|<\dfrac{\varepsilon\alpha}2$. 故
\[
\begin{aligned}
\left|x^\alpha\int_x^1\frac{f(t)-f(0)}{t^{\alpha+1}}\,dt\right|
&\le\left|x^\alpha\int_x^\delta\frac{f(t)-f(0)}{t^{\alpha+1}}\,dt\right|
+\left|x^\alpha\int_\delta^1\frac{f(t)-f(0)}{t^{\alpha+1}}\,dt\right|\\
&\le x^\alpha\int_x^\delta\frac{|f(t)-f(0)|}{t^{\alpha+1}}\,dt+x^\alpha M
\quad\left(M\stackrel{\text{记}}{=}\left|\int_\delta^1\frac{f(x)-f(0)}{t^{\alpha+1}}\,dt\right|\right)\\
&\le\frac{\varepsilon\alpha}2\cdot\frac1\alpha\left(1-\frac{x^\alpha}{\delta^\alpha}\right)+x^\alpha M\\
&\le\frac\varepsilon2+x^\alpha M<\varepsilon\quad(\text{当 }x>0\text{ 充分小时}).
\end{aligned}
\]
（式 (1) 右端第二项）
\[
x^\alpha\int_x^1\frac{f(0)}{t^{\alpha+1}}\,dt
=x^\alpha f(0)\left(\frac1{\alpha x^\alpha}-\frac1\alpha\right)
\longrightarrow\frac{f(0)}\alpha\quad(x\to0^+).
\]

\noindent\textbf{☆例 4.5.30}\quad 设当 $x>a$ 时，$\varphi(x)>0$，$f(x),\varphi(x)$ 在任何有限区间 $[a,b]$ 上可积，
$\displaystyle\int_a^{+\infty}\varphi(t)\,dt$ 发散，当 $x\to+\infty$ 时，$f(x)=o(\varphi(x))$，证明：
\begin{equation}
\int_a^{+\infty}f(t)\,dt=o\left(\int_a^{+\infty}\varphi(t)\,dt\right).
\tag{1}
\end{equation}
（清华大学）

\noindent\textbf{证}\quad 要证明式 (1)，即要证明：$\forall\varepsilon>0$，$\exists\Delta>0$，当 $x>\Delta$ 时，
\[
\left|\frac{\displaystyle\int_a^x f(t)\,dt}{\displaystyle\int_a^x\varphi(t)\,dt}\right|<\varepsilon.
\]
因 $\varphi(x)>0$，此即
\[
\left|\int_a^x f(t)\,dt\right|<\varepsilon\int_a^x\varphi(t)\,dt.
\]
已知当 $x\to+\infty$ 时，$f(x)=o(\varphi(x))$. 所以 $\exists A>a$，使得当 $x>A$ 时，有
$|f(x)|<\dfrac\varepsilon2\varphi(x)$. 于是
\begin{equation}
\begin{aligned}
\left|\int_a^x f(t)\,dt\right|
&\le\left|\int_a^A f(t)\,dt\right|+\int_A^x|f(t)|\,dt\\
&\le\left|\int_a^A f(t)\,dt\right|+\frac\varepsilon2\int_A^x\varphi(t)\,dt\\
&\le\left|\int_a^A f(t)\,dt\right|+\frac\varepsilon2\int_a^x\varphi(t)\,dt.
\end{aligned}
\tag{2}
\end{equation}
又因 $\displaystyle\int_a^{+\infty}\varphi(t)\,dt$ 发散，故 $x$ 充分大时，能使
\[
\int_a^x\varphi(t)\,dt>\frac2\varepsilon\left|\int_a^A f(t)\,dt\right|,
\]
亦即
\[
\left|\int_a^A f(t)\,dt\right|<\frac\varepsilon2\int_a^x\varphi(t)\,dt.
\]
所以
\[
\left|\int_a^x f(t)\,dt\right|\le\varepsilon\int_a^x\varphi(t)\,dt.
\]
证毕.

\noindent\textbf{例 4.5.31}\quad 设 $k>0,\ a<\xi<b$，证明当 $n\to+\infty$ 时，
\[
\int_a^b e^{-kn(x-\xi)^2}\,dx\sim\sqrt{\frac\pi{kn}}.
\]
（已知 $\displaystyle\int_{-\infty}^{+\infty}e^{-t^2}\,dt=\sqrt\pi$.）

\noindent\textbf{证}\quad 问题等价于要证明
\[
\lim_{n\to+\infty}\sqrt{\frac{kn}\pi}\int_a^b e^{-kn(x-\xi)^2}\,dx=1.
\]
令 $\sqrt{kn}(x-\xi)=t$，作变换，则 $x=\xi+\dfrac{t}{\sqrt{kn}}$，$dx=\dfrac1{\sqrt{kn}}\,dt$，于是上述极限成为
\[
\lim_{n\to+\infty}\sqrt{\frac{kn}\pi}\int_a^b e^{-kn(x-\xi)^2}\,dx
=\lim_{n\to+\infty}\frac1{\sqrt\pi}\int_{\sqrt{kn}(a-\xi)}^{\sqrt{kn}(b-\xi)}e^{-t^2}\,dt
=\frac1{\sqrt\pi}\int_{-\infty}^{+\infty}e^{-t^2}\,dt=1.
\]
（因为 $a-\xi<0,\ b-\xi>0$，当 $n\to+\infty$ 时，上限 $\sqrt{kn}(b-\xi)\to+\infty$，下限
$\sqrt{kn}(a-\xi)\to-\infty$.）

（利用正常积分现有结果）

\noindent\textbf{例 4.5.32（Riemann 定理）}\quad 设 $f(x)$ 在 $[a,+\infty)$ 上绝对可积，$g(x)$ 是周期为 $T$ 的函数，在 $[0,T]$ 上正常可积，则
\begin{equation}
\lim_{n\to\infty}\int_a^{+\infty}f(x)g(nx)\,dx
=\frac1T\int_0^T g(x)\,dx\int_a^{+\infty}f(x)\,dx. \tag{1}
\end{equation}

\noindent\textbf{分析}\quad 要证明式 (1)，即要证明 $n$ 充分大时，
\begin{equation}
\left|\int_a^{+\infty}f(x)g(nx)\,dx-\frac1T\int_0^Tg(x)\,dx\int_a^{+\infty}f(x)\,dx\right| \tag{2}
\end{equation}
能任意小. 用关系 $\displaystyle\int_a^{+\infty}\cdots=\int_a^A\cdots+\int_A^{+\infty}\cdots$ 将积分拆开，再放大：
\[
\begin{aligned}
&\left|\int_a^{+\infty}f(x)g(nx)\,dx-\frac1T\int_0^Tg(x)\,dx\int_a^{+\infty}f(x)\,dx\right|\\
\le{}&
\left|\int_a^A f(x)g(nx)\,dx-\frac1T\int_0^Tg(x)\,dx\int_a^A f(x)\,dx\right|\\
&+\left|\int_A^{+\infty}f(x)g(nx)\,dx\right|
+\frac1T\left|\int_0^Tg(x)\,dx\right|\left|\int_A^{+\infty}f(x)\,dx\right|\\
\le{}&
\left|\int_a^A f(x)g(nx)\,dx-\frac1T\int_0^Tg(x)\,dx\int_a^A f(x)\,dx\right|\\
&+M\int_A^{+\infty}|f(x)|\,dx
+\frac1T\left|\int_0^Tg(x)\,dx\right|\int_A^{+\infty}|f(x)|\,dx.
\end{aligned}
\]
（因 $g(x)$ 有界，$\exists M>0$ 使得 $|g(x)|\le M$.）由于 $f(x)$ 绝对可积，故当 $A$ 充分大时，可
使上式第二、三项任意小. 然后，将 $A$ 固定，令 $n$ 充分大，因该命题对于正常积分已成立（见例 4.1.10），故当 $n$ 充分大时，第一项亦可任意小. 命题获证.

请读者写出简明严格的证明.

\noindent\textbf{例 4.5.33}\quad 设 $\varphi(x)$ 为有界周期函数，周期为 $T$，且 $\dfrac1T\int_0^T\varphi(x)\,dx=C$，试证：
\begin{equation}
\lim_{n\to\infty}n\int_n^{+\infty}\frac{\varphi(t)}{t^2}\,dt=C. \tag{1}
\end{equation}

\noindent\textbf{证 I}\quad （利用上例结果.）
\[
n\int_n^{+\infty}\frac{\varphi(t)}{t^2}\,dt
=\int_n^{+\infty}\frac{\varphi\left(n\cdot\dfrac tn\right)d\left(\dfrac tn\right)}
{\left(\dfrac tn\right)^2}
\stackrel{u=t/n}{=}\int_1^{+\infty}\frac1{u^2}\varphi(nu)\,du
\]
\[
\longrightarrow\frac1T\int_0^T\varphi(x)\,dx\int_1^{+\infty}\frac{du}{u^2}
=C\quad(\text{当 }n\to\infty\text{ 时}).
\]

\noindent\textbf{证 II}\quad （利用积分第二中值定理：若 $f(x)$ 在 $[a,b]$ 上可积，$g(x)$ 单调减小，且
$g(b)\ge0$，则 $\exists\xi\in[a,b]$ 使得 $\displaystyle\int_a^b f(x)g(x)\,dx
=g(a+0)\int_a^\xi f(x)\,dx$.）

注意到 $\displaystyle n\int_n^{+\infty}\frac{dt}{t^2}=1$，知
$C=n\int_n^{+\infty}\dfrac{C\,dt}{t^2}$. 故只要证明
\[
\lim_{n\to\infty}n\int_n^{+\infty}\frac{\varphi(t)-C}{t^2}\,dt=0.
\]
根据已知条件，$\displaystyle\int_n^A(\varphi(t)-C)\,dt$ 有界，即 $\exists M>0$ 使得
\[
\left|\int_n^A(\varphi(t)-C)\,dt\right|\le M.
\]
于是利用第二中值定理，
\[
\left|n\int_n^A\frac{\varphi(t)-C}{t^2}\,dt\right|
=n\cdot\frac1{n^2}\left|\int_n^\xi(\varphi(t)-C)\,dt\right|
\le\frac Mn\quad(\forall A>n),
\]
从而
\[
\left|n\int_n^{+\infty}\frac{\varphi(t)-C}{t^2}\,dt\right|\le\frac Mn.
\]
此式对每个固定的 $n$ 都成立，令 $n\to+\infty$，取极限，知式 (1) 成立.

（利用两边夹法则）

\noindent\textbf{例 4.5.34}\quad 设 $\varphi(x)$ 连续，$\lim\limits_{x\to+\infty}\varphi(x)=1$，记
\[
\psi(x)=\int_x^{+\infty}\frac{\varphi(t)\,dt}{t^{1+\alpha}}\quad(\alpha>0).
\]
求证：当 $x\to+\infty$ 时，$\psi(x)\sim\dfrac{c}{x^\alpha}$，并求 $c$ 之值.

\noindent\textbf{提示}\quad $x$ 充分大时，
\[
x^\alpha\int_x^{+\infty}\frac{1-\varepsilon}{t^{1+\alpha}}\,dt
\le x^\alpha\int_x^{+\infty}\frac{\varphi(t)\,dt}{t^{1+\alpha}}
\le x^\alpha\int_x^{+\infty}\frac{1+\varepsilon}{t^{1+\alpha}}\,dt.
\]
或直接对 $\displaystyle\lim_{x\to+\infty}\frac{\psi(x)}{x^{-\alpha}}$ 用 L'Hospital 法则.

（利用分部积分法）

\noindent\textbf{例 4.5.35}\quad 设 $\varphi(x)=\displaystyle\int_0^x\cos\frac1t\,dt$，求 $\varphi'(0)$.

\noindent\textbf{解}\quad
\begin{equation}
\varphi'(0)=\lim_{x\to0}\frac{\varphi(x)-\varphi(0)}{x-0}
=\lim_{x\to0}\frac{\displaystyle\int_0^x\cos\frac1t\,dt}{x}. \tag{1}
\end{equation}
但
\[
\begin{aligned}
\int_0^x\cos\frac1t\,dt
&\stackrel{1/t=u}{=}\int_{1/x}^{+\infty}\frac{\cos u}{u^2}\,du\\
&=\left.\frac1{u^2}\sin u\right|_{1/x}^{+\infty}
+\int_{1/x}^{+\infty}\frac{2\sin u}{u^3}\,du\\
&=-x^2\sin\frac1x+\int_{1/x}^{+\infty}\frac{2\sin u}{u^3}\,du.
\end{aligned}
\]
故
\[
\left|\frac{\displaystyle\int_0^x\cos\frac1t\,dt}{x}\right|
\le |x|\left|\sin\frac1x\right|+\frac1{|x|}\int_{1/x}^{+\infty}\frac2{u^3}\,du
=|x|\left|\sin\frac1x\right|+|x|\to0\quad(\text{当 }x\to0\text{ 时}).
\]
所以 $\varphi'(0)=0$.

\noindent\textbf{注}\quad 本例式 (1) 中的极限可以用 L'Hospital 法则，然而此极限使用 L'Hospital 法则后，所得极限
$\lim\limits_{x\to0}\cos\dfrac1x$ 反而不存在. 注意这并不矛盾，因为 L'Hospital 法则是在分子、分母求导之后所得极限存在时得到的.
当求导后的极限不存在时，原式仍可能有极限，本例就是如此. 所以求导后极限不存在，只能说明此时 L'Hospital 法则失效，不能说明原式无极限.

\noindent\textbf{例 4.5.36}\quad 设 $\rho(t)$ 为全数轴的连续函数：
\[
1)\ \rho(t)=0,\ |t|\ge1;\qquad
2)\ \int_{-\infty}^{+\infty}\rho(t)\,dt=0;\qquad
3)\ \int_{-\infty}^{+\infty}t\rho(t)\,dt=1.
\]
又设 $f(x)$ 是全数轴可微函数，证明：
\[
\lim_{\lambda\to0}\int_{-\infty}^{+\infty}\frac1{\lambda^2}\rho\left(\frac{t-x}{\lambda}\right)f(t)\,dt=f'(x).
\]

\noindent\textbf{提示}\quad 用拟合法归结为
\[
\lim_{\lambda\to0}\int_{-\infty}^{+\infty}
u\rho(u)\left[\frac{f(x+\lambda u)-f(x)}{\lambda u}-f'(x)\right]\,du=0.
\]

\subsection*{五、反常积分作为``积分和''的极限}

我们知道，反常积分不能像正常积分那样定义为积分和的极限. 本段我们将看到，对于单调无界反常积分，我们可以用类似于积分和的和式来逼近.

\noindent\textbf{例 4.5.37}\quad 设 $f(x)$ 在 $(0,1)$ 上单调，$x=0,1$ 为奇点，$\displaystyle\int_0^1f(x)\,dx$ 存在，证明：
\[
\lim_{n\to\infty}
\frac{f\left(\frac1n\right)+f\left(\frac2n\right)+\cdots+f\left(\frac{n-1}n\right)}n
=\int_0^1f(x)\,dx.
\]

\noindent\textbf{证}\quad 若 $f(x)\nearrow$，则有
\begin{equation}
\int_0^{1-\frac1n}f(x)\,dx
\le\frac{f\left(\frac1n\right)+f\left(\frac2n\right)+\cdots+f\left(\frac{n-1}n\right)}n
\le\int_{\frac1n}^1f(x)\,dx. \tag{1}
\end{equation}
（图 4.5.2 中阴影部分的面积 $=\dfrac1n\left[f\left(\dfrac1n\right)+f\left(\dfrac2n\right)+\cdots+f\left(\dfrac{n-1}n\right)\right]$.）在式 (1) 中令
$n\to\infty$，取极限，即得所需的等式.

若 $f(x)\searrow$，替换 $f(x)$，可以考虑函数 $-f(x)$.

\begin{center}
\includegraphics[width=0.78\textwidth]{4-5-2.png}
\end{center}

\noindent\textbf{注}\quad $1^\circ$ 单调性条件可减弱为只在奇点某邻域内单调.

$2^\circ$ 本例的结论，可推广到更一般的情况.

若 $f(x)$ 在 $(0,1)$ 内单调，$0,1$ 为奇点，反常积分 $\displaystyle\int_0^1f(x)\,dx$ 收敛，$\varphi(x)$ 在 $[0,1]$ 上正常可积，则
\[
\lim_{n\to\infty}\frac1n\sum_{i=1}^{n-1}\varphi\left(\frac in\right)f\left(\frac in\right)
=\int_0^1\varphi(x)f(x)\,dx.
\]

\noindent\textbf{例 4.5.38}\quad 设 $f(x)$ 单调，$x=0$ 为奇点，$\displaystyle\int_0^1f(x)\,dx$ 收敛，证明：
\[
\lim_{n\to\infty}\frac1n\sum_{i=0}^{n-1}f\left(\varepsilon_n+\frac in\right)
=\int_0^1f(x)\,dx\quad\left(\frac\theta n<\varepsilon_n<\frac1n\right),
\]
其中 $\theta\in(0,1)$ 为常数.

\noindent\textbf{证}\quad 设 $f(x)\nearrow$（否则考虑 $-f(x)$），利用两边夹法则，在不等式中令 $n\to+\infty$，取极限即得. 这里
$\dfrac1n\displaystyle\sum_{i=0}^{n-1}f\left(\varepsilon_n+\dfrac in\right)$ 是图 4.5.3 中阴影部分的面积.

其中右边的不等式可由
\[
\varepsilon_nf(\varepsilon_n)\le\int_0^{\varepsilon_n}f(x)\,dx
\quad\text{及}\quad
\frac1n f(\varepsilon_n)\le\frac1{n\varepsilon_n}\int_0^{\varepsilon_n}f(x)\,dx
\le\frac1\theta\int_0^{\varepsilon_n}f(x)\,dx
\]
得到.

\begin{center}
\includegraphics[width=0.78\textwidth]{4-5-3.png}
\end{center}

我们来看上述结果的一个应用.

\noindent\textbf{例 4.5.39}\quad 设 $0<a<d$，
\[
A_n=\frac{a+(a+d)+\cdots+[a+(n-1)d]}n,\quad
G_n=\sqrt[n]{a(a+d)\cdots[a+(n-1)d]},
\]
试证：$\displaystyle\lim_{n\to\infty}\frac{G_n}{A_n}=\frac2e$.

\noindent\textbf{证}\quad 首先将 $\dfrac{G_n}{A_n}$ 变形，
\[
\frac{G_n}{A_n}
=\frac{\sqrt[n]{a(a+d)\cdots[a+(n-1)d]}}
{\dfrac{a+(a+d)+\cdots+[a+(n-1)d]}n}
=\frac{\sqrt[n]{a(a+d)\cdots[a+(n-1)d]}}
{a+\dfrac12(n-1)d}.
\]
分子分母同除以 $nd$，并记 $c=\dfrac ad$，则
\begin{equation}
\frac{G_n}{A_n}
=\frac{\sqrt[n]{\dfrac cn\cdot\dfrac{c+1}n\cdots\dfrac{c+(n-1)}n}}
{\dfrac cn+\dfrac{n-1}{2n}}. \tag{1}
\end{equation}
这里分母趋于 $\dfrac12$，而分子的对数
\[
\ln\sqrt[n]{\frac cn\cdot\frac{c+1}n\cdots\frac{c+(n-1)}n}
=\frac1n\left[\ln\frac cn+\ln\left(\frac cn+\frac1n\right)+\cdots+
\ln\left(\frac cn+\frac{n-1}n\right)\right]
\]
\[
=\frac1n\sum_{i=0}^{n-1}\ln\left(\varepsilon_n+\frac in\right)
\quad\left(\text{其中 }\varepsilon_n=\frac cn=\frac1n\frac ad\right).
\]
利用上例结论，
\[
\frac1n\sum_{i=0}^{n-1}\ln\left(\varepsilon_n+\frac in\right)
\to\int_0^1\ln x\,dx=-1\quad(n\to\infty).
\]
故式 (1) 分子
\[
\sqrt[n]{\frac cn\cdot\frac{c+1}n\cdots\frac{c+(n-1)}n}
\to e^{-1}\quad(n\to\infty).
\]
从而 $\displaystyle\lim_{n\to\infty}\frac{G_n}{A_n}=\frac2e$. 证毕.

\section*{六、综合性问题}

\noindent\textbf{☆例 4.5.40}\quad 假定函数 $f(x)$ 当 $x>0$ 时连续并且非负；对任何正数 $M$，定义
\begin{equation}
f_M(x)=\min\{f(x),M\}. \tag{1}
\end{equation}
证明：如果
\[
\lim_{0<\eta\to0}\lim_{M\to+\infty}\int_0^\eta f_M(x)\,dx
\]
存在，那么
\begin{equation}
\lim_{0<\eta\to0}\lim_{M\to+\infty}\int_0^\eta f_M(x)\,dx=0. \tag{2}
\end{equation}
（厦门大学）

\noindent\textbf{证}\quad
$1^\circ$ 我们看到，要证明式 (2)，只要对某个 $\eta_1>0$，证明 $\displaystyle\int_0^{\eta_1}f(x)\,dx$ 收敛.
这是因为，若 $\displaystyle\int_0^{\eta_1}f(x)\,dx$ 收敛，则 $\forall\varepsilon>0$，$\exists\delta>0$，当 $0<\eta<\delta$ 时，有
$0<\displaystyle\int_0^\eta f(x)\,dx<\dfrac\varepsilon2$，从而
\[
0\le\int_0^\eta f_M(x)\,dx\le\int_0^\eta f(x)\,dx<\frac\varepsilon2,
\]
所以
\[
\lim_{M\to+\infty}\int_0^\eta f_M(x)\,dx\le\frac\varepsilon2<\varepsilon.
\]
这就证明了式 (2).

$2^\circ$ 我们来证明：$\exists\eta_1>0$，使得 $\displaystyle\int_0^{\eta_1}f(x)\,dx$ 收敛.
因 $f(x)\ge0$，要证明 $\displaystyle\int_0^{\eta_1}f(x)\,dx$ 收敛，只要证明 $0<\eta\to0$ 时，
$\displaystyle\int_\eta^{\eta_1}f(x)\,dx$ 保持有界.

已知 $\displaystyle\lim_{0<\eta\to0}\lim_{M\to+\infty}\int_0^\eta f_M(x)\,dx$ 存在，故对 $\eta_1>0$（充分小），极限
\[
\lim_{M\to+\infty}\int_0^{\eta_1}f_M(x)\,dx=\varphi(\eta_1)
\]
存在（有限）. 但当 $M\nearrow$ 时 $f_M(x)\nearrow$，$\displaystyle\int_0^{\eta_1}f_M(x)\,dx\nearrow$. 故
\begin{equation}
\varphi(\eta_1)=\lim_{M\to+\infty}\int_0^{\eta_1}f_M(x)\,dx
=\sup_{M>0}\int_0^{\eta_1}f_M(x)\,dx. \tag{3}
\end{equation}
$\forall 0<\eta<\eta_1$，因为 $f(x)$ 连续，$f(x)$ 在区间 $[\eta,\eta_1]$ 上有上确界 $M_1\ge0$，从而在 $[\eta,\eta_1]$ 上
$f_{M_1}(x)=f(x)$，
\[
0\le\int_\eta^{\eta_1}f(x)\,dx
=\int_\eta^{\eta_1}f_{M_1}(x)\,dx
\le\int_0^{\eta_1}f_{M_1}(x)\,dx
\le\varphi(\eta_1).
\]
即 $\displaystyle\int_\eta^{\eta_1}f(x)\,dx$ 关于 $\eta\in(0,\eta_1)$ 有界. 故
$\displaystyle\int_0^{\eta_1}f(x)\,dx$ 收敛，证毕.

\noindent\textbf{☆例 4.5.41}\quad 设 $f(x)$ 在每个有限区间 $[a,b]$ 上可积，并且
$\lim\limits_{x\to+\infty}f(x)=A,\ \lim\limits_{x\to-\infty}f(x)=B$ 存在. 求证：对任何一个实数 $a>0$，
\[
\int_{-\infty}^{+\infty}[f(x+a)-f(x)]\,dx
\]
存在，并求出它的值.（同济大学）

\noindent\textbf{解}\quad 问题在于：证明极限
\[
\lim_{\substack{\alpha\to-\infty\\ \beta\to+\infty}}
\int_\alpha^\beta[f(x+a)-f(x)]\,dx
\]
存在，并求其值. 事实上，
\[
\begin{aligned}
\int_\alpha^\beta[f(x+a)-f(x)]\,dx
&=\int_\alpha^\beta f(x+a)\,dx-\int_\alpha^\beta f(x)\,dx\quad(\text{令 }x+a=t)\\
&=\int_{\alpha+a}^{\beta+a}f(t)\,dt-\int_\alpha^\beta f(x)\,dx\\
&=\int_\beta^{\beta+a}f(x)\,dx-\int_\alpha^{\alpha+a}f(x)\,dx\\
&=\int_\beta^{\beta+a}[A+(f(x)-A)]\,dx-\int_\alpha^{\alpha+a}[B+(f(x)-B)]\,dx\\
&=Aa-Ba+\int_\beta^{\beta+a}(f(x)-A)\,dx-\int_\alpha^{\alpha+a}(f(x)-B)\,dx.
\end{aligned}
\]
利用已知条件 $\lim\limits_{x\to+\infty}f(x)=A,\ \lim\limits_{x\to-\infty}f(x)=B$，易知上式右端后两项当
$\alpha\to-\infty,\ \beta\to+\infty$ 时极限为零. 故
\[
\lim_{\substack{\alpha\to-\infty\\ \beta\to+\infty}}
\int_\alpha^\beta[f(x+a)-f(x)]\,dx=a(A-B).
\]

\noindent\textbf{例 4.5.42}\quad 设 $f(x)$ 在 $[0,1]$ 上有连续导数，且
\[
|f'(x)|\le\frac{M}{(1-x)^{1-\alpha}}\quad(0<\alpha\le1).
\]
求证：$f(x)$ 在 $[0,1]$ 上满足条件
\[
|f(x_2)-f(x_1)|\le\frac M\alpha |x_2-x_1|^\alpha.
\]

\noindent\textbf{证}\quad
\[
\begin{aligned}
|f(x_2)-f(x_1)|
&=\left|\int_{x_1}^{x_2}f'(x)\,dx\right|
\le\left|\int_{x_1}^{x_2}|f'(x)|\,dx\right|\\
&\le\left|\int_{x_1}^{x_2}\frac M{(1-x)^{1-\alpha}}\,dx\right|
=\frac M\alpha\left|(1-x_2)^\alpha-(1-x_1)^\alpha\right|.
\end{aligned}
\]
因 $\forall x,y$，有 $|x+y|^\alpha\le |x|^\alpha+|y|^\alpha$，可知 $\forall x,y$，有
\[
\bigl||x|^\alpha-|y|^\alpha\bigr|\le |x-y|^\alpha,
\]
故
\[
|f(x_2)-f(x_1)|\le\frac M\alpha |x_2-x_1|^\alpha.
\]

\noindent\textbf{例 4.5.43}\quad 设
\[
F(x)=\int_0^x\left(\frac1t-\left[\frac1t\right]\right)\,dt,
\]
其中 $\left[\dfrac1t\right]$ 表示取不大于 $\dfrac1t$ 的最大整数. 试证：$F'(0)=\dfrac12$.

\noindent\textbf{提示}\quad 参看例 4.5.33.

\noindent\textbf{new ☆例 4.5.44}\quad
$1)$ 设 $f\in C[a,+\infty)$（意思：$f(x)$ 在 $[a,+\infty)$ 上连续），且
$\displaystyle\int_a^{+\infty}|f(x)|\,dx$ 收敛. 试证：必存在 $|x_n|\subset[a,+\infty)$ 满足：
$\lim\limits_{n\to+\infty}x_n=+\infty$，且
\[
\lim_{n\to+\infty}x_nf(x_n)=0;
\]

$2)$ 若 $f(x)$ 放弃连续条件，但在 $[a,+\infty)$ 任意内闭区间上可积，$\displaystyle\int_a^{+\infty}f(x)\,dx$ 条件收敛（不绝对收敛），则 1）的结论可以不成立，举例说明.（华东师范大学）

\noindent\textbf{提示}\quad
$1)$ 由于 $\displaystyle\int_a^{+\infty}|f(x)|\,dx$ 收敛，由 Cauchy 准则知，$\forall\varepsilon_n=\dfrac1n>0$，$\exists A_n>n>0$，
$\forall A'_n>A_n>n$，有
\begin{equation}
2\int_{A'_n}^{2A'_n}|f(x)|\,dx<\varepsilon_n=\frac1n. \tag{1}
\end{equation}
又因 $f(x)$ 连续，可知 $|f(x)|$ 连续. 根据积分中值定理，$\exists\xi_n:A'_n<\xi_n<2A'_n$，使得
\[
2\int_{A'_n}^{2A'_n}|f(x)|\,dx
\stackrel{\text{积分中值定理}}{=}2|f(\xi_n)|\int_{A'_n}^{2A'_n}dx
=2A'_n|f(\xi_n)|>\xi_n|f(\xi_n)|.
\]
故
\[
\xi_n|f(\xi_n)|<2\int_{A'_n}^{2A'_n}|f(x)|\,dx<\frac1n.
\]
可见 $\lim\limits_{n\to\infty}\xi_n|f(\xi_n)|=0$. 但 $\xi_n>n$，只要令 $x_n=\xi_n$，则 $|x_n|\to+\infty$，满足一切要求.

$2)$ 令
\[
a_0=a,\quad a_n=a+\sum_{k=1}^n\frac1k\quad(n=1,2,\cdots).
\]
则 $a_n$ 严格增至 $+\infty$. 定义
\begin{equation}
f(x)=(-1)^n\quad(a_n\le x<a_{n+1})\quad(n=0,1,2,\cdots). \tag{2}
\end{equation}
此 $f(x)$ 不连续，可作反例，说明如果 $f(x)$ 不连续，则 1）的结论可以不成立. 事实上：
\[
i)\quad \int_a^{+\infty}f(x)\,dx
=\sum_{k=1}^{+\infty}(-1)^{k-1}\frac1k
\]
收敛（显然非绝对收敛，因 $|f(x)|\equiv1$）.
\[
ii)\quad \forall |x_n|,\ \text{若 }|x_n|\to+\infty,\ \text{则根据 }f(x)\text{ 的定义式 (2)：}
\]
\[
\lim_{n\to+\infty}|x_nf(x_n)|=\lim_{n\to+\infty}x_n=+\infty,
\]
说明原结论不再成立.

\noindent\textbf{new *例 4.5.45}\quad 函数 $f$ 和 $g$ 在 $[0,+\infty)$ 上非负，连续，单调递减，反常积分
\[
\int_0^{+\infty}f(x)\,dx\quad\text{和}\quad \int_0^{+\infty}g(x)\,dx
\]
发散. 如果定义：$h(x)=\min\{g(x),f(x)\},\ x\in[0,+\infty)$. 问
\[
\int_0^{+\infty}h(x)\,dx
\]
是否一定发散？为什么？（北京大学）

\noindent\textbf{提示}\quad （否！）因为 $\displaystyle\int_0^{+\infty}\frac1{(1+x)^2}\,dx$ 收敛，只要构造两个非负、连续、递减函数 $g$ 和 $f$，使得
$\displaystyle\int_0^{+\infty}f(x)\,dx$ 和 $\displaystyle\int_0^{+\infty}g(x)\,dx$ 发散，且
\[
h(x)=\min\{g(x),f(x)\}=\frac1{(1+x)^2},
\]
就可作为反例，说明命题不成立.

\noindent\textbf{再提示}\quad 如图 4.5.4 所示，作阶梯函数，让 $f(x),g(x)$ 同时从点 $A$ 出发，一个走折线另一个走曲线，每下一步台阶，两者交换一次（即：走折线改走曲线，走曲线改走折线）. 这样下去，直到 $+\infty$. 让每步台阶的下方图形面积大于某一固定的正数 $\varepsilon_0>0$.
例如取 $\varepsilon_0=\dfrac12$，只需台面足够宽（即 $BC$ 足够长），总能使（线段 $BC$）下方图形面积大于 $\varepsilon_0=\dfrac12$.

\begin{center}
\includegraphics[width=0.45\textwidth]{4-5-4.png}
\end{center}

以 $BCD$ 作为一步台阶的代表，如点
\[
B\left(x_n,\frac1{(1+x_n)^2}\right),\quad
C\left(x'_{n+1},\frac1{(1+x_n)^2}\right),\quad
D\left(x_{n+1},\frac1{(1+x_{n+1})^2}\right)\quad(n=1,2,\cdots).
\]
当 $x_n$ 确定后，令 $x'_{n+1}>\dfrac12(1+x_n)^2+x_n$，则 $BC$ 下方图形的面积：
\[
\int_{x_n}^{x'_{n+1}}\frac1{(1+x_n)^2}\,dx
=(x'_{n+1}-x_n)\frac1{(1+x_n)^2}\ge\frac12.
\]
然后取 $x_{n+1}=x'_{n+1}+m\ (m>0)\ (n=1,2,\cdots)$，在 $(x'_{n+1},x_{n+1})$ 段，以直线连接（保证折线连续）. 于是，$\forall A>0$，
\[
\int_A^{+\infty}f(x)\,dx>\frac12;
\]
同样，$\displaystyle\int_A^{+\infty}g(x)\,dx>\frac12$. 故 $\displaystyle\int_0^{+\infty}f(x)\,dx$ 和
$\displaystyle\int_0^{+\infty}g(x)\,dx$ 发散. 但是，
\[
\int_0^{+\infty}\min\{f(x),g(x)\}\,dx
=\int_0^{+\infty}\frac{dx}{(1+x)^2}
\]
收敛.

\noindent\textbf{点评}\quad 此题曾引起网友们的广泛关注，大家的表达方式虽然不同，但本质上惊人地一致. 如果您学过级数，不妨考虑将本题离散化，写出级数的相应结果.

\noindent\textbf{new *例 4.5.46}\quad 假设 $f(x)>0$，积分 $\displaystyle\int_1^{+\infty}f(x)\,dx$ 收敛. 试证：
$\displaystyle\int_1^{+\infty}(f(x))^{1-\frac1x}\,dx$ 收敛.

\noindent\textbf{提示}\quad 利用例 4.4.9 的不等式进行放大，使放大后的积分收敛.

\noindent\textbf{再提示}\quad $\forall b>0$，
\[
(f(x))^{1-\frac1x}
=\frac1b\left[(f(x))^{1-\frac1x}\cdot b\right]
\le\frac1b\left[\frac1p(f(x))^{\frac{x-1}x p}+\frac1q b^q\right]
\quad\left(\forall p,q>1:\frac1p+\frac1q=1\right)
\]
\[
\le\frac1b\left[(f(x))^{\frac{x-1}x p}+b^q\right].
\]
（对 $x>1$，取 $p=\dfrac x{x-1},\ q=x>1$，则满足 $\dfrac1p+\dfrac1q=1$.）由上式得
\[
(f(x))^{1-\frac1x}\le\frac1b f(x)+b^{x-1}.
\]
取 $b=e^{-1}>0$，则得
\[
0\le\int_1^A(f(x))^{1-\frac1x}\,dx
\le e\int_1^A[f(x)+e^{-x}]\,dx.
\]
故由 $\displaystyle\int_1^{+\infty}f(x)\,dx$ 和 $\displaystyle\int_1^{+\infty}e^{-x}\,dx$ 收敛，知
$\displaystyle\int_1^{+\infty}(f(x))^{1-\frac1x}\,dx$ 收敛.（证毕.）

\noindent\textbf{new *例 4.5.47}\quad 设 $f(x)$ 是 $[0,+\infty)$ 上的（正值）连续函数，且
$\displaystyle\int_0^{+\infty}\frac1{f(x)}\,dx<+\infty$，求证：
\[
\lim_{A\to+\infty}\frac1{A^2}\int_0^A f(x)\,dx=+\infty.
\]
（北京大学，华中师范大学）

\noindent\textbf{提示}\quad
\[
\frac14=\frac1{A^2}\left(\int_{A/2}^A\frac1{\sqrt{f(x)}}\cdot\sqrt{f(x)}\,dx\right)^2
\stackrel{\text{Schwarz不等式}}{\le}
\frac1{A^2}\int_0^A f(x)\,dx\cdot\int_{A/2}^A\frac1{f(x)}\,dx.
\]
因此
\[
\frac1{A^2}\int_0^A f(x)\,dx
\ge\frac14\frac1{\displaystyle\int_{A/2}^A\frac1{f(x)}\,dx}
\to+\infty\quad(A\to+\infty).
\]
问题得证.

\noindent\textbf{注}\quad
$i)$ 若题目添加条件：$\dfrac x{f(x)}=o(1)\ (x\to+\infty)$，则证法可大为简化：（用强化的 L'Hospital，例 3.2.30）
\[
\lim_{x\to+\infty}\frac1{x^2}\int_1^x f(t)\,dt
=\lim_{x\to+\infty}\frac{f(x)}{2x}=+\infty.
\]

$ii)$ 或有 $f(x)\nearrow$，则可导出条件：$\dfrac x{f(x)}=o(1)\ (x\to+\infty)$.（利用 Cauchy 准则）
$\forall n>0$，$\exists A>1$，使得 $\forall x>A$，有
\[
\frac1n>\int_x^{2x}\frac1{f(t)}\,dt
\ge\frac1{f(x)}\int_x^{2x}dt=\frac x{f(x)},
\]
即 $\dfrac{f(x)}x>n$. 因此，$\lim\limits_{x\to+\infty}\dfrac{f(x)}x=+\infty$，故
$\dfrac x{f(x)}=o(1)\ (x\to+\infty)$.

\noindent\textbf{new *例 4.5.48}\quad 已知双曲函数：
\[
\sinh x=\frac{e^x-e^{-x}}2,\quad \cosh x=\frac{e^x+e^{-x}}2.
\]
若 $\sinh x\cdot\sinh y=1$，试计算积分
\[
\int_0^{+\infty}y(x)\,dx.
\]
（浙江大学）

\noindent\textbf{解}\quad 由 $\sinh x\cdot\sinh y=1$ 可得
\[
\sinh y\cdot\cosh x\,dx+\sinh x\cdot\cosh y\,dy=0.
\]
故
\[
dx=-\frac{\sinh x\cdot\cosh y\,dy}{\sinh y\cdot\sqrt{1+\sinh^2x}}
=-\frac{\sinh x\cdot\cosh y\,dy}{\sqrt{\sinh^2y+1}}
=-\frac{dy}{\sinh y}.
\]
因此
\[
\begin{aligned}
\int_0^{+\infty}y(x)\,dx
&=-\int_{+\infty}^0 y\cdot\frac{dy}{\sinh y}
=\int_0^{+\infty}y\cdot\frac{dy}{\sinh y}
=\int_0^{+\infty}\frac{2y\,dy}{e^y-e^{-y}}\\
&=\int_0^{+\infty}\frac{2ye^{-y}\,dy}{1-e^{-2y}}
\stackrel{s=e^{-y}}{=}\int_0^1\frac{-2\ln s\,ds}{1-s^2}\\
&=-2\int_0^1\ln s\cdot\sum_{k=0}^{\infty}s^{2k}\,ds
=-2\sum_{k=0}^{\infty}\int_0^1s^{2k}\ln s\,ds\\
&\stackrel{\text{例 }4.5.4}{=}2\sum_{k=0}^{\infty}\frac1{(1+2k)^2}
\stackrel{\text{例 }5.4.8}{=}\frac{\pi^2}4.
\end{aligned}
\]

\section*{单元练习 4.5}

\noindent\textbf{☆4.5.1}\quad 计算：
\[
1)\ \int_a^b\frac{dx}{\sqrt{(x-a)(b-x)}}\quad(b>a);\qquad
2)\ \int_{-1}^1\frac1{(a-x)\sqrt{1-x^2}}\,dx\quad(a>1).
\]
\[
\left\langle \pi \right\rangle\qquad
\left\langle \frac{\pi}{\sqrt{a^2-1}} \right\rangle
\]

\noindent\textbf{提示}\quad
$1)$ 原式
\[
=\int_a^b\frac{2\,d\sqrt{x-a}}{\sqrt{b-x}}
=2\int_a^b\frac{d\sqrt{x-a}}{\sqrt{(\sqrt{b-a})^2-(\sqrt{x-a})^2}}
\]
\[
=2\int_a^b\frac{d\sqrt{\dfrac{x-a}{b-a}}}
{\sqrt{1-\left(\sqrt{\dfrac{x-a}{b-a}}\right)^2}}
=2\arcsin\left.\frac{\sqrt{x-a}}{\sqrt{b-a}}\right|_a^b=\pi.
\]
$2)$ 可令 $x=\sin t$.

\noindent\textbf{☆4.5.2}\quad 计算
\[
\int_{-\infty}^{+\infty}\frac{dx}{(x^2+2x+2)^n}.
\]
（中国科学院）
\[
\left\langle \frac{(2n-3)!!}{(2n-2)!!}\pi \right\rangle
\]

\noindent\textbf{提示}\quad 原式 $=2\displaystyle\int_0^{+\infty}\dfrac{dt}{(t^2+1)^n}=2I_n$. 分部积分可建立 $I_n$ 的递推公式：
\[
I_{n+1}=\frac{2n-1}{2n}I_n.
\]
或利用变换 $t=\tan\theta$，其中
\[
I_n\equiv\int_0^{+\infty}\frac{dt}{(t^2+1)^n}
=\int_0^{\pi/2}(\cos\theta)^{2n-2}\,d\theta,
\]
再直接引用 Wallis 公式，
\[
I_n=\frac{(2n-3)!!}{(2n-2)!!}\cdot\frac\pi2.
\]

\noindent\textbf{4.5.3}\quad 求
\[
\int_0^{+\infty}f(x^p+x^{-p})\frac{\ln x}{1+x^2}\,dx
\]
（函数 $f(x)$ 连续）.
\[
\left\langle 0 \right\rangle
\]

\noindent\textbf{提示}\quad 可参看例 4.5.8.

\noindent\textbf{4.5.4}\quad 计算
\[
\int_0^1\frac{\arcsin x}{x}\,dx.
\]
\[
\left\langle\frac\pi2\ln2\right\rangle
\]

\noindent\textbf{提示}\quad 可令 $x=\sin t$，参看例 4.5.7.

\noindent\textbf{4.5.5}\quad 计算
\[
\int_0^{\pi/2}\frac{\sin(2k-1)x}{\sin x}\,dx.
\]
\[
\left\langle\frac\pi2\right\rangle
\]

\noindent\textbf{提示}\quad
\[
\frac{\sin(2k-1)x}{\sin x}=1+2\sum_{i=1}^{k-1}\cos2ix.
\]

\noindent\textbf{4.5.6}\quad 证明：
\[
\int_0^{+\infty}f\left[\left(Ax-\frac Bx\right)^2\right]\,dx
=\frac1A\int_0^{+\infty}f(y^2)\,dy
\]
（假设其中的积分存在，且 $A,B>0$）.

\noindent\textbf{提示}\quad 在原式左端积分
\[
\int_0^{+\infty}f\left[\left(Ax-\frac Bx\right)^2\right]\,dx
\]
中令
\[
x=\frac{y+\sqrt{y^2+4AB}}{2A},
\quad\text{即}\quad y=Ax-\frac Bx
\]
（参看例 4.5.2），可得
\[
\text{左端}=\frac1{2A}\int_{-\infty}^{+\infty}f(y^2)\left(1+\frac y{\sqrt{y^2+4AB}}\right)dy
=\frac1{2A}\int_{-\infty}^{+\infty}f(y^2)\,dy
\]
（因为 $f(y^2)\dfrac y{\sqrt{y^2+4AB}}$ 为奇函数）
\[
=\text{右端}\quad(\text{因为 }f(y^2)\text{ 为偶函数}).
\]
另证：
\[
\text{左端}=\int_0^{+\infty}f\left[\left(Ax-\frac Bx\right)^2\right]\,dx
\]
\[
\stackrel{t=B/(Ax)}{=}\int_{+\infty}^0
f\left[\left(\frac Bt-At\right)^2\right]\left(-\frac B{At^2}\right)dt
=\int_0^{+\infty}f\left[\left(At-\frac Bt\right)^2\right]\frac B{At^2}\,dt\stackrel{\text{记}}{=}I
\]
\[
=\frac{\text{左端}+I}{2}
=\frac12\int_0^{+\infty}f\left[\left(At-\frac Bt\right)^2\right]\left(1+\frac B{At^2}\right)dt
\]
\[
=\frac1{2A}\int_0^{+\infty}f\left[\left(At-\frac Bt\right)^2\right]d\left(At-\frac Bt\right)
\stackrel{y=At-\frac Bt}{=}\frac1{2A}\int_{-\infty}^{+\infty}f(y^2)\,dy
=\text{右端}
\]
（因为 $f(y^2)$ 是偶函数）.

\noindent\textbf{注}\quad 能应用该等式的函数确实存在，例如 $f(x)=\dfrac1{x+1}, f(x)=e^{-x}$，等等.

\noindent\textbf{4.5.7}\quad 研究下列积分的敛散性：
\[
1)\ \int_{-\infty}^{+\infty}x^ne^{-\left(x^2+\frac1{x^2}\right)}\,dx\quad(n\text{ 为自然数});
\qquad
2)\ \int_0^{+\infty}\sin^2\left[\pi\left(x+\frac1x\right)\right]\,dx.
\]

\noindent\textbf{提示}\quad
$1)$（该积分收敛）因被积函数 $f(x)=x^ne^{-\left(x^2+\frac1{x^2}\right)}$，只要补充定义 $f(0)=0$，就在零点连续（故不是奇点）；在上、下限 $\pm\infty$ 处，只要注意到：当 $|x|$ 充分大时，有
\[
\left|x^{n+2}e^{-\left(x^2+\frac1{x^2}\right)}\right|\cdot\frac1{x^2}\le\frac1{x^2},
\]
用 $M$ 判别法，即知收敛；

$2)$（根据 Cauchy 准则：若对任意 $k$，有
\[
\int_{k+\frac\pi4}^{k+\frac{3\pi}4}f(x)\,dx\ge\varepsilon_0,
\]
则 $\displaystyle\int_a^{+\infty}f(x)\,dx$ 发散.）
\[
\sin^2\left[\pi\left(x+\frac1x\right)\right]\stackrel{t=x+1/x}{=}\sin^2(\pi t)\ge\frac12
\quad\left(t\in\left[k+\frac14,k+\frac34\right]\stackrel{\text{记}}{=}[t'_k,t''_k]\right).
\]
注意：$t\in[t'_k,t''_k]\Longleftrightarrow x\in[x'_k,x''_k]$，
\[
[x'_k,x''_k]=\left[\frac{t'_k+\sqrt{t_k'^2+4}}2,\frac{t''_k+\sqrt{t_k''^2+4}}2\right].
\]
\[
x''_k-x'_k
=\left.\left(\frac{t+\sqrt{t^2+4}}2\right)'\right|_{t=\xi}\cdot(t''_k-t'_k)
\ge\frac12(t''_k-t'_k)=\frac14.
\]
$\forall A>0$，$k$ 充分大时，$[x'_k,x''_k]$ 在 $A$ 的右边，而区间 $[x'_k,x''_k]$ 上积分
\[
\int_{x'_k}^{x''_k}\sin^2\left[\pi\left(x+\frac1x\right)\right]dx
\ge\frac12\int_{x'_k}^{x''_k}dx
=\frac12(x''_k-x'_k)\ge\frac18
\]
（不能任意小）. 根据 Cauchy 准则，原积分发散.

\noindent\textbf{4.5.8}\quad 设 $f(x)$ 在 $[a,+\infty)$ 上可微，且 $x\to+\infty$ 时 $f'(x)$ 单调递增趋于 $+\infty$，证明：
$\displaystyle\int_a^{+\infty}\sin(f(x))\,dx$ 和 $\displaystyle\int_a^{+\infty}\cos(f(x))\,dx$ 都收敛.

\noindent\textbf{提示}\quad $x\to+\infty$ 时，$f'(x)\nearrow+\infty$，因此 $\exists b>a$，使得 $x\in[b,+\infty)$ 时，有 $f'(x)>0$，$y=f(x)$ 严格递增，且有反函数 $x=g(y)$. 于是
\[
I\stackrel{\text{记}}{=}\int_a^{+\infty}\sin(f(x))\,dx
=\int_a^b\sin(f(x))\,dx+\int_b^{+\infty}\sin(f(x))\,dx=I_1+I_2,
\]
其中，$I_1$ 为正常积分（连续可积）；$I_2$：$x\ge b$，$f$ 有反函数 $f^{-1}=g$，
\[
I_2=\int_b^{+\infty}\sin(f(x))\,dx
\stackrel{x=g(y)}{=}\int_{f(b)}^{+\infty}\sin y\cdot g'(y)\,dy
=\int_{f(b)}^{+\infty}\sin y\cdot\frac1{f'(x)|_{x=g(y)}}\,dy.
\]
利用 Dirichlet 判别法，可知 $I_2$ 收敛，故 $I$ 收敛；$\displaystyle\int_a^{+\infty}\cos(f(x))\,dx$ 类似可证.

\noindent\textbf{☆4.5.9}\quad 设 $f(x)$ 为连续实值函数，对所有 $x$，有 $f(x)\ge0$，且
$\displaystyle\int_0^{+\infty}f(x)\,dx<+\infty$，求证：
\[
\frac1n\int_0^nxf(x)\,dx\longrightarrow0\quad(\text{当 }n\to\infty\text{ 时}).
\]
（中国科学院）

\noindent\textbf{提示}\quad $\forall\varepsilon>0$，$\exists A>0$，$0<\displaystyle\int_A^{+\infty}f(x)\,dx<\dfrac\varepsilon2$，再将 $A$ 固定，则
$0<\dfrac1n\displaystyle\int_0^Axf(x)\,dx\to0\ (n\to+\infty)$，因而 $\exists N>A$，$n>N$ 时，
$0<\dfrac1n\displaystyle\int_0^Axf(x)\,dx<\dfrac\varepsilon2$. 于是
\[
0<\frac1n\int_0^nxf(x)\,dx
=\frac1n\left(\int_0^A+\int_A^n\right)xf(x)\,dx
\le\frac1n\int_0^Axf(x)\,dx+\int_A^{+\infty}f(x)\,dx<\varepsilon.
\]
（因为 $\dfrac1n\displaystyle\int_A^nxf(x)\,dx=\displaystyle\int_A^n\dfrac xn f(x)\,dx
\le\displaystyle\int_A^n f(x)\,dx\le\displaystyle\int_A^{+\infty}f(x)\,dx<\dfrac\varepsilon2$）

\noindent\textbf{4.5.10}\quad 证明
\[
\lim_{x\to\infty}\int_0^{+\infty}\frac{e^{-tx}}{1+t^2}\,dt=0.
\]

\noindent\textbf{提示}\quad （原积分）
\[
I(x)=-\left.\frac{e^{-tx}}{x(1+t^2)}\right|_0^{+\infty}
+\frac1x\int_0^{+\infty}e^{-tx}\,d\frac1{1+t^2}
\]
\[
=\frac1x+\frac1x\int_0^{+\infty}e^{-tx}\,d\frac1{1+t^2},
\]
\[
|I(x)|\le\frac1x+\frac1x\int_0^{+\infty}\frac{2t}{(1+t^2)^2}\,dt=\frac2x\to0\quad(x\to+\infty).
\]

\noindent\textbf{4.5.11}\quad 设 $f(x)$ 是 $0\le x<+\infty$ 上的非负连续函数，并且满足
\[
1)\ \text{在 }0\le x<+\infty\text{ 上存在有界导数 }f'(x);\qquad
2)\ \int_0^{+\infty}f(x)\,dx<+\infty.
\]
求证 $\lim\limits_{x\to+\infty}f(x)=0$.（山东大学）

\noindent\textbf{提示}\quad 可参看例 4.5.24.（注意导数有界必一致连续）

\noindent\textbf{☆4.5.12}\quad $f(x)$ 在 $[a,+\infty)$ 上连续且 $\displaystyle\int_a^{+\infty}f(x)\,dx$ 收敛，问能否断定：
$\exists x_n\to+\infty$，使 $\lim\limits_{n\to+\infty}f(x_n)=0$？为什么？（南开大学）

\noindent\textbf{提示}\quad 能. $1)$ 若 $f(x)$ 无穷次变号或无穷次达到零，$\forall A>0$，在 $[A,+\infty)$ 内仍如此，则明显.

$2)$ 否则，不妨假设 $\exists A>0$，使 $x>A$ 时恒有 $f(x)>0$（恒有 $f(x)<0$ 可类似证明）. 于是，$\forall n\in\mathbb N$，
$\exists x_n>\max\{n,A\}$，使得 $f(x_n)<\dfrac1n$.

\noindent\textbf{4.5.13}\quad 设 $f(x)$ 在任一有限区间 $[0,a]\ (a>0)$ 上正常可积，在 $[0,+\infty)$ 上绝对可积，证明：
\[
\lim_{n\to\infty}\int_0^{+\infty}f(x)|\sin nx|\,dx
=\frac2\pi\int_0^{+\infty}f(x)\,dx.
\]
（南京大学）

\noindent\textbf{提示}\quad 可参看例 4.5.32.

\noindent\textbf{☆4.5.14}\quad 若函数 $p(t)$ 在 $[0,+\infty)$ 连续，且当 $t\to+\infty$ 时，
$p(t)=o(t^N)$（$N$ 为正整数）. 又 $\lambda<0$，证明：当 $t\to\infty$ 时
\[
\int_t^{+\infty}p(\tau)e^{\lambda\tau}\,d\tau=o(t^{N+1})e^{\lambda t}.
\]
（北京师范大学）

\noindent\textbf{提示}\quad 可参看例 4.5.30.

\noindent\textbf{4.5.15}\quad $\{C_n^k\}_{k=0}^n$ 为二项式系数，$A_n,G_n$ 分别表示它们的算术平均值与几何平均值，试证：
\[
\lim_{n\to\infty}\sqrt[n]{A_n}=2,\qquad
\lim_{n\to\infty}\sqrt[n]{G_n}=\sqrt e.
\]

\noindent\textbf{提示}\quad
\[
A_n=\frac1{n+1}(C_n^0+C_n^1+\cdots+C_n^{n-1}+C_n^n)
=\frac1{n+1}(1+1)^n=\frac1{n+1}2^n,
\]
\begin{equation}
G_n=\left(\prod_{k=0}^n C_n^k\right)^{\frac1{n+1}}
=e^{\frac1{n+1}\sum_{k=0}^n\ln C_n^k}. \tag{1}
\end{equation}
再应用 Stolz 公式.

\noindent\textbf{证}\quad
\[
\lim_{n\to\infty}\sqrt[n]{A_n}
=\lim_{n\to\infty}\sqrt[n]{\frac{2^n}{n+1}}=2.
\]
下证 $\displaystyle\lim_{n\to\infty}\sqrt[n]{G_n}=\sqrt e$.
\[
\begin{aligned}
\lim_{n\to\infty}\frac1{n(n+1)}\sum_{k=0}^n\ln C_n^k
&\stackrel{\text{Stolz公式}}{=}
\lim_{n\to\infty}\frac1{2n}\left(\sum_{k=0}^n\ln C_n^k-\sum_{k=0}^{n-1}\ln C_{n-1}^k\right)\\
&=\lim_{n\to\infty}\frac1{2n}\left(\sum_{k=1}^{n-1}\ln C_n^k-\sum_{k=1}^{n-2}\ln C_{n-1}^k\right)
\quad(\text{因为 }\ln C_n^0=\ln C_n^n=\ln C_{n-1}^0=0)\\
&=\lim_{n\to\infty}\frac1{2n}\left[\left(\sum_{k=1}^{n-2}\ln\frac{C_n^k}{C_{n-1}^k}\right)+\ln C_n^{n-1}\right]\\
&=\lim_{n\to\infty}\frac1{2n}\left[\sum_{k=1}^{n-2}\ln\left(\frac n{n-k}\right)+\ln n\right]
\quad\left(\text{因为 }C_n^k=\frac{n!}{k!(n-k)!}, C_n^{n-1}=C_n^1=n\right)\\
&=\lim_{n\to\infty}\frac12\left(-\frac1n\sum_{k=1}^{n-1}\ln\frac{n-k}n\right)
=-\frac12\int_0^1\ln(1-x)\,dx=\frac12.
\end{aligned}
\]
所以
\[
\lim_{n\to\infty}\sqrt[n]{G_n}
=\lim_{n\to\infty}\left(\prod_{k=0}^n C_n^k\right)^{\frac1{n(n+1)}}
=\lim_{n\to+\infty}e^{\frac1{n(n+1)}\sum_{k=0}^n\ln C_n^k}
=e^{1/2}=\sqrt e.
\]

\noindent\textbf{4.5.16}\quad 例 4.5.37 的逆命题不成立：即 $f(x)$ 在 $(0,1)$ 内单调，
$\displaystyle\lim_{n\to\infty}\dfrac1n\sum_{i=1}^{n-1}f\left(\dfrac in\right)$ 存在，
$\displaystyle\int_0^1f(x)\,dx$ 可以不收敛（考虑 $f(x)=\dfrac1x-\dfrac1{1-x}$）.

\noindent\textbf{4.5.17}\quad 已知积分
\[
\int_0^{+\infty}\frac{\sin\beta x}{x}\,dx=\frac\pi2\operatorname{sgn}\beta
\]
（见例 7.1.38），求积分
\[
\int_0^{+\infty}\frac{\sin x\cos xt}{x}\,dx.
\]
（华北电力大学）
\[
\left\langle \frac\pi2\ (|t|<1),\ \frac\pi4\ (t=\pm1),\ 0\ (|t|>1)\right\rangle
\]

\noindent\textbf{提示}\quad 可用积化和差公式.

\noindent\textbf{4.5.18}\quad 证明：
\[
\int_0^{+\infty}\frac{dx}{1+x^4}
=\int_0^{+\infty}\frac{x^2}{1+x^4}\,dx
=\frac\pi{2\sqrt2}.
\]
（北京航空航天大学）

\noindent\textbf{提示}\quad 见例 4.5.1 之证明.

\end{document}
